\input{preamble}

% OK, start here.
%
\begin{document}

\title{Diviseurs}


\maketitle

\phantomsection
\label{section-phantom}

\tableofcontents

\section{Introduction}
\label{section-introduction}

\noindent
Dans ce chapitre, nous étudions quelques questions très fondamentales liées à la
définition des diviseurs, etc. Une référence de base est  \cite{EGA}.



\section{Points associés}
\label{section-associated}

\noindent
Soit  $R$  un anneau et soit  $M$  un  $R$-module. Rappelons
qu'un  $\mathfrak p \subset R$  premier est  {\it  associé}  à
 $M$  s'il existe un élément de  $M$  dont l'annulateur est
 $\mathfrak p$. Voir Algèbre, Définition  \ref{algebra-definition-associated}. Voici la définition des
points associés pour les faisceaux quasi-cohérents sur les schémas comme donnée
dans  \cite[IV Definition 3.1.1]{EGA}.

\begin{definition}
\label{definition-associated}
Soit  $X$  un schéma. Soit  $\mathcal{F}$  un faisceau quasi-cohérent sur
 $X$.
\begin{enumerate}
\item On dit que  $x \in X$  est  {\it  associé}  à
 $\mathcal{F}$  si l'idéal maximal  $\mathfrak m_x$  est associé au  $\mathcal{O}_{X, x}$-module
 $\mathcal{F}_x$.
\item Nous désignons  $\text{Ass}(\mathcal{F})$  ou  $\text{Ass}_X(\mathcal{F})$  l'ensemble des points associés de
 $\mathcal{F}$.
\item Les points associés  {\it  de  $X$}  sont les
points associés de  $\mathcal{O}_X$.
\end{enumerate}
\end{definition}

\noindent
Ces définitions sont les plus utiles lorsque  $X$  est localement
noethérien et  $\mathcal{F}$  de type fini. Par exemple, il peut arriver qu'un point
générique d'une composante irréductible de  $X$  ne soit pas associé à
 $X$, voir l'exemple  \ref{example-no-associated-prime}. Dans le cas non noethérien, il peut
être plus pratique d'utiliser les points faiblement associés, voir la section
 \ref{section-weakly-associated}. Relions la notion schématique avec la notion algébrique sur les
ouverts affines; notez que cette correspondance fonctionne parfaitement uniquement
pour les schémas localement noethériens.

\begin{lemma}
\label{lemma-associated-affine-open}
Soit  $X$  un schéma. Soit  $\mathcal{F}$  un faisceau quasi-cohérent sur
 $X$. Soit  $\Spec(A) = U \subset X$  un ouvert affine, et posons  $M = \Gamma(U, \mathcal{F})$. Soit
 $x \in U$, et soit  $\mathfrak p \subset A$  le premier idéal correspondant.
\begin{enumerate}
\item Si  $\mathfrak p$  est associé à  $M$, alors  $x$  est associé à
 $\mathcal{F}$.
\item Si  $\mathfrak p$  est de type fini, alors la réciproque est également vraie.
\end{enumerate}
En particulier, si  $X$  est localement noethérien, alors l'équivalence
$$
\mathfrak p \in \text{Ass}(M) \Leftrightarrow x \in \text{Ass}(\mathcal{F})
$$
vaut pour toutes les paires  $(\mathfrak p, x)$  comme ci-dessus.
\end{lemma}

\begin{proof}
Cela découle de Algèbre, Lemme  \ref{algebra-lemma-associated-primes-localize}. Mais nous pouvons aussi raisonner
directement comme suit. Supposons que  $\mathfrak p$  soit associé à  $M$.
Alors il existe un  $m \in M$  dont l'annulateur est  $\mathfrak p$. Comme la
localisation est exacte, nous voyons que  $\mathfrak pA_{\mathfrak p}$  est l'annulateur de
 $m/1 \in M_{\mathfrak p}$. Comme  $M_{\mathfrak p} = \mathcal{F}_x$  (Schémas, Lemme  \ref{schemes-lemma-spec-sheaves}), nous concluons que
 $x$  est associé à  $\mathcal{F}$.

\medskip\noindent
Inversement, supposons que  $x$  soit associé à  $\mathcal{F}$, et que
 $\mathfrak p$  soit de type fini. Comme  $x$  est associé à  $\mathcal{F}$ , il
existe un élément  $m' \in M_{\mathfrak p}$  dont l'annulateur est  $\mathfrak pA_{\mathfrak p}$. Écrivons
 $m' = m/f$  pour quelque  $f \in A$,  $f \not \in \mathfrak p$. L'annulateur  $I$  de
 $m$  est un idéal de  $A$  tel que  $IA_{\mathfrak p} = \mathfrak pA_{\mathfrak p}$. Par conséquent
 $I \subset \mathfrak p$, et  $(\mathfrak p/I)_{\mathfrak p} = 0$. Puisque  $\mathfrak p$  est de type fini, il existe un
 $g \in A$,  $g \not \in \mathfrak p$  tel que  $g(\mathfrak p/I) = 0$. Ainsi, l'annulateur de
 $gm$  est  $\mathfrak p$  et nous avons gagné.

\medskip\noindent
Si  $X$  est localement noethérien, alors  $A$  est noethérien
(Propriétés, Lemme  \ref{properties-lemma-locally-Noetherian}) et  $\mathfrak p$  est toujours de type fini.
\end{proof}

\begin{lemma}
\label{lemma-ass-support}
Soit  $X$  un schéma. Soit  $\mathcal{F}$  un  $\mathcal{O}_X$-module
quasi-cohérent. Alors  $\text{Ass}(\mathcal{F}) \subset \text{Supp}(\mathcal{F})$.
\end{lemma}

\begin{proof}
Cela découle immédiatement des définitions.
\end{proof}

\begin{lemma}
\label{lemma-ses-ass}
Soit  $X$  un schéma. Soit  $0 \to \mathcal{F}_1 \to \mathcal{F}_2 \to \mathcal{F}_3 \to 0$  une suite exacte courte de
faisceaux quasi-cohérents sur  $X$. Alors  $\text{Ass}(\mathcal{F}_2) \subset
\text{Ass}(\mathcal{F}_1) \cup \text{Ass}(\mathcal{F}_3)$  et  $\text{Ass}(\mathcal{F}_1) \subset \text{Ass}(\mathcal{F}_2)$.
\end{lemma}

\begin{proof}
Pour chaque point  $x \in X$  la suite des germes  $0 \to \mathcal{F}_{1, x} \to \mathcal{F}_{2, x} \to \mathcal{F}_{3, x} \to 0$  est une suite
exacte courte de  $\mathcal{O}_{X, x}$-modules. Par conséquent, le lemme découle de
Algèbre, Lemme  \ref{algebra-lemma-ass}.
\end{proof}

\begin{lemma}
\label{lemma-finite-ass}
Soit  $X$  un schéma localement noethérien. Soit  $\mathcal{F}$  un
 $\mathcal{O}_X$-module cohérent. Alors  $\text{Ass}(\mathcal{F}) \cap U$  est fini pour tout ouvert
quasi-compact  $U \subset X$.
\end{lemma}

\begin{proof}
C'est vrai car l'ensemble des idéaux premiers associés d'un module fini sur un
anneau noethérien est fini, voir Algèbre, Lemme  \ref{algebra-lemma-finite-ass}. Pour passer des
schémas à l'algèbre, utilisez que  $U$  est une réunion finie d'ouverts
affines, chacun de ces ouverts est le spectre d'un anneau noethérien (Propriétés,
Lemme  \ref{properties-lemma-locally-Noetherian}),  $\mathcal{F}$  correspond à un module fini sur cet anneau
(Cohomologie des schémas, Lemme  \ref{coherent-lemma-coherent-Noetherian}), et enfin utilisez le Lemme
 \ref{lemma-associated-affine-open}.
\end{proof}

\begin{lemma}
\label{lemma-ass-zero}
Soit  $X$  un schéma localement noethérien. Soit  $\mathcal{F}$  un
 $\mathcal{O}_X$-module quasi-cohérent. Alors
$$
\mathcal{F} = 0 \Leftrightarrow \text{Ass}(\mathcal{F}) = \emptyset.
$$
\end{lemma}

\begin{proof}
Si  $\mathcal{F} = 0$, alors  $\text{Ass}(\mathcal{F}) = \emptyset$  par définition. Réciproquement, si
 $\text{Ass}(\mathcal{F}) = \emptyset$, alors  $\mathcal{F} = 0$  par Algèbre, Lemme  \ref{algebra-lemma-ass-zero}. Pour passer des
schémas à l'algèbre, restreignez-vous à tout ouvert affine et utilisez le Lemme
 \ref{lemma-associated-affine-open}.
\end{proof}

\begin{example}
\label{example-no-associated-prime}
Soit  $k$  un corps. L'anneau  $R = k[x_1, x_2, x_3, \ldots]/(x_i^2)$  est local avec un idéal maximal
localement nilpotent  $\mathfrak m$. Il n'existe aucun élément de  $R$  dont
l'annulateur soit  $\mathfrak m$. Donc  $\text{Ass}(R) = \emptyset$, et  $X = \Spec(R)$  est un exemple
d'un schéma qui n'a pas de points associés.
\end{example}

\begin{lemma}
\label{lemma-restriction-injective-open-contains-ass}
Soit  $X$  un schéma localement noethérien. Soit  $\mathcal{F}$  un
 $\mathcal{O}_X$-module quasi-cohérent. Si  $U \subset X$  est ouvert et  $\text{Ass}(\mathcal{F}) \subset U$,
alors  $\Gamma(X, \mathcal{F}) \to \Gamma(U, \mathcal{F})$  est injectif.
\end{lemma}

\begin{proof}
Soit  $s \in \Gamma(X, \mathcal{F})$  une section qui se restreint à zéro sur  $U$. Soit
 $\mathcal{F}' \subset \mathcal{F}$  l'image de l'application  $\mathcal{O}_X \to \mathcal{F}$  définie par  $s$. Alors
 $\text{Supp}(\mathcal{F}') \cap U = \emptyset$. D'autre part,  $\text{Ass}(\mathcal{F}') \subset \text{Ass}(\mathcal{F})$  d'après le Lemme  \ref{lemma-ses-ass}. Puisque
également  $\text{Ass}(\mathcal{F}') \subset \text{Supp}(\mathcal{F}')$  (Lemme  \ref{lemma-ass-support}), nous concluons  $\text{Ass}(\mathcal{F}') = \emptyset$. Par
conséquent,  $\mathcal{F}' = 0$  d'après le Lemme  \ref{lemma-ass-zero}.
\end{proof}

\begin{lemma}
\label{lemma-minimal-support-in-ass}
Soit  $X$  un schéma localement noethérien. Soit  $\mathcal{F}$  un
 $\mathcal{O}_X$-module quasi-cohérent. Soit  $x \in \text{Supp}(\mathcal{F})$  un point dans le support de
 $\mathcal{F}$  qui n’est pas une spécialisation d’un autre point de  $\text{Supp}(\mathcal{F})$.
Alors  $x \in \text{Ass}(\mathcal{F})$. En particulier, tout point générique d’une composante
irréductible de  $X$  est un point associé de  $X$.
\end{lemma}

\begin{proof}
Comme  $x \in \text{Supp}(\mathcal{F})$  le module  $\mathcal{F}_x$  n’est pas nul. Par conséquent
 $\text{Ass}(\mathcal{F}_x) \subset \Spec(\mathcal{O}_{X, x})$  n’est pas vide selon Algèbre, Lemme  \ref{algebra-lemma-ass-zero}. D’autre part, par
hypothèse  $\text{Supp}(\mathcal{F}_x) = \{\mathfrak m_x\}$. Comme  $\text{Ass}(\mathcal{F}_x) \subset \text{Supp}(\mathcal{F}_x)$  (Algèbre, Lemme  \ref{algebra-lemma-ass-support}) nous
voyons que  $\mathfrak m_x$  est associé à  $\mathcal{F}_x$  et nous avons gagné.
\end{proof}

\noindent
Le lemme suivant est l’analogue de Plus sur l’Algèbre, Lemme  \ref{more-algebra-lemma-check-injective-on-ass}.

\begin{lemma}
\label{lemma-check-injective-on-ass}
Laisse $X$ soit un schéma localement noethérien. Laisser $\varphi : \mathcal{F} \to \mathcal{G}$ être
une carte de quasi-cohérent $\mathcal{O}_X$-modules. Supposons que pour
chaque $x \in X$ au moins l'une des choses suivantes se produit
\begin{enumerate}
\item $\mathcal{F}_x \to \mathcal{G}_x$  est injectif, ou
\item $x \not \in \text{Ass}(\mathcal{F})$.
\end{enumerate}
Alors  $\varphi$  est injectif.
\end{lemma}

\begin{proof}
Les hypothèses impliquent que  $\text{Ass}(\Ker(\varphi)) = \emptyset$  et donc  $\Ker(\varphi) = 0$  par le Lemme
 \ref{lemma-ass-zero}.
\end{proof}

\begin{lemma}
\label{lemma-check-isomorphism-via-depth-and-ass}
Soit $X$ un schéma localement noethérien. Soit $\varphi : \mathcal{F} \to \mathcal{G}$ un
morphisme de $\mathcal{O}_X$-modules quasi-cohérents. Supposons $\mathcal{F}$ cohérent
et que, pour tout $x \in X$, l'une des situations suivantes se produise
\begin{enumerate}
\item $\mathcal{F}_x \to \mathcal{G}_x$  est un isomorphisme, ou
\item $\text{depth}(\mathcal{F}_x) \geq 2$  et  $x \not \in \text{Ass}(\mathcal{G})$.
\end{enumerate}
Alors  $\varphi$  est un isomorphisme.
\end{lemma}

\begin{proof}
Ceci est une traduction de Plus en Algèbre, Lemme  \ref{more-algebra-lemma-check-isomorphism-via-depth-and-ass}  dans le langage des
schémas.
\end{proof}




\section{Morphismes et points associés}
\label{section-morphisms-associated}

\noindent
Soit  $f : X \to S$  un morphisme de schémas. Soit  $\mathcal{F}$  un faisceau de
 $\mathcal{O}_X$-modules. Si  $s \in S$  est un point, il est souvent commode de
noter  $\mathcal{F}_s$  le  $\mathcal{O}_{X_s}$-module que l'on obtient en tirant en arrière
 $\mathcal{F}$  par le morphisme  $i_s : X_s \to X$. Ici,  $X_s$  est la fibre au sens
des schémas de  $f$  au-dessus de  $s$. Dans une formule
$$
\mathcal{F}_s = i_s^*\mathcal{F}
$$
Bien sûr, cette notation entre en conflit avec la notation déjà existante pour la
germe de  $\mathcal{F}$  en un point  $x \in X$  si  $f = \text{id}_X$. Cependant, la
notation est souvent pratique, comme dans la formulation du lemme suivant.

\begin{lemma}
\label{lemma-bourbaki}
Soit  $f : X \to S$  un morphisme de schémas. Soit  $\mathcal{F}$  un faisceau
quasi-cohérent sur  $X$  qui est plat sur  $S$. Soit  $\mathcal{G}$ 
un faisceau quasi-cohérent sur  $S$. Alors nous avons
$$
\text{Ass}_X(\mathcal{F} \otimes_{\mathcal{O}_X} f^*\mathcal{G})
\supset
\bigcup\nolimits_{s \in \text{Ass}_S(\mathcal{G})}
\text{Ass}_{X_s}(\mathcal{F}_s)
$$
et l'égalité tient si  $S$  est localement noethérien (pour la notation
 $\mathcal{F}_s$  voir ci-dessus).
\end{lemma}

\begin{proof}
Soit  $x \in X$  et soit  $s = f(x) \in S$. Posons  $B = \mathcal{O}_{X, x}$,  $A = \mathcal{O}_{S, s}$,
 $N = \mathcal{F}_x$ et  $M = \mathcal{G}_s$. Remarquez que la germe de  $\mathcal{F} \otimes_{\mathcal{O}_X} f^*\mathcal{G}$  en
 $x$  est égal au  $B$-module  $M \otimes_A N$. Ainsi  $x \in \text{Ass}_X(\mathcal{F} \otimes_{\mathcal{O}_X} f^*\mathcal{G})$  si
et seulement si  $\mathfrak m_B$  est dans  $\text{Ass}_B(M \otimes_A N)$. De même  $s \in \text{Ass}_S(\mathcal{G})$  et
 $x \in \text{Ass}_{X_s}(\mathcal{F}_s)$  si et seulement si  $\mathfrak m_A \in \text{Ass}_A(M)$  et  $\mathfrak m_B/\mathfrak m_A B \in
\text{Ass}_{B \otimes \kappa(\mathfrak m_A)}(N \otimes \kappa(\mathfrak m_A))$. Ainsi, le lemme
découle de Algèbre, Lemme  \ref{algebra-lemma-bourbaki-fibres}.
\end{proof}





\section{Points immergés}
\label{section-embedded}

\noindent
Soit  $R$  un anneau et soit  $M$  un  $R$-module. Rappelons
qu'un  $\mathfrak p \subset R$  premier est un  {\it  associé
immergé}  de  $M$  s'il est un idéal premier associé de
 $M$  qui n'est pas minimal parmi les idéaux premiers associés de
 $M$. Voir Algèbre, Définition  \ref{algebra-definition-embedded-primes}. Voici la définition des
points associés immergés pour les faisceaux quasi-cohérents sur les schémas telle
que donnée dans  \cite[IV Definition 3.1.1]{EGA}.

\begin{definition}
\label{definition-embedded}
Soit  $X$  un schéma. Soit  $\mathcal{F}$  un faisceau quasi-cohérent sur
 $X$.
\begin{enumerate}
\item Un  {\it  point associé immergé}  de  $\mathcal{F}$  est
un point associé qui n'est pas maximal parmi les points associés de  $\mathcal{F}$,
c'est-à-dire qu'il est la spécialisation d'un autre point associé de  $\mathcal{F}$.
\item Un point  $x$  de  $X$  est appelé un  {\it  point
incorporé}  si  $x$  est un point associé incorporé de
 $\mathcal{O}_X$.
\item Un  {\it  composant incorporé}  de  $X$  est un
sous-ensemble fermé irréductible  $Z = \overline{\{x\}}$  où  $x$  est un point
incorporé de  $X$.
\end{enumerate}
\end{definition}

\noindent
Dans le cas noethérien lorsque  $\mathcal{F}$  est cohérent, nous avons ce qui suit.

\begin{lemma}
\label{lemma-embedded}
Soit  $X$  un schéma localement noethérien. Soit  $\mathcal{F}$  un
 $\mathcal{O}_X$-module cohérent. Alors
\begin{enumerate}
\item les points génériques des composantes irréductibles de  $\text{Supp}(\mathcal{F})$  sont des
points associés de  $\mathcal{F}$, et
\item un point associé de  $\mathcal{F}$  est incorporé si et seulement si il n'est pas un
point générique d'une composante irréductible de  $\text{Supp}(\mathcal{F})$.
\end{enumerate}
En particulier un point incorporé de $X$  est un point associé de
 $X$ qui n’est pas un point générique d’une composante irréductible
de $X$.
\end{lemma}

\begin{proof}
Rappelez-vous que dans ce cas $Z = \text{Supp}(\mathcal{F})$ est fermé, voir Morphismes,
Lemme \ref{morphisms-lemma-support-finite-type}  et que les points génériques des composantes irréductibles de
 $Z$ sont des points associés de $\mathcal{F}$, voir le Lemme  \ref{lemma-minimal-support-in-ass}.
Enfin, nous avons  $\text{Ass}(\mathcal{F}) \subset Z$, par le Lemme  \ref{lemma-ass-support}. Ces résultats, combinés au
fait que  $Z$ est un espace topologique sobre et donc chaque point
de $Z$ est une spécialisation d'un point générique de $Z$, implique
(1) et (2).
\end{proof}

\begin{lemma}
\label{lemma-S1-no-embedded}
Laisse $X$ soit un schéma localement noethérien. Laisser $\mathcal{F}$ soit un
faisceau cohérent sur $X$. Alors les éléments suivants sont équivalents:
\begin{enumerate}
\item $\mathcal{F}$  n’a pas de points associés intégrés, et
\item $\mathcal{F}$  a la propriété  $(S_1)$.
\end{enumerate}
\end{lemma}

\begin{proof}
Ceci est l'algèbre, le lemme  \ref{algebra-lemma-criterion-no-embedded-primes}, combiné avec le lemme  \ref{lemma-associated-affine-open} 
ci-dessus.
\end{proof}

\begin{lemma}
\label{lemma-noetherian-dim-1-CM-no-embedded-points}
Soit  $X$  un schéma localement noethérien de dimension  $\leq 1$. Les
énoncés suivants sont équivalents
\begin{enumerate}
\item $X$  est de Cohen-Macaulay, et
\item $X$  n'a pas de points immergés.
\end{enumerate}
\end{lemma}

\begin{proof}
Cela découle du lemme  \ref{lemma-S1-no-embedded}  et des définitions.
\end{proof}

\begin{lemma}
\label{lemma-scheme-theoretically-dense-contain-embedded-points}
Soit  $X$  un schéma localement noethérien. Soit  $U \subset X$  un
sous-schéma ouvert. Les énoncés suivants sont équivalents
\begin{enumerate}
\item $U$  est schématiquement dense dans  $X$  (Morphismes, Définition
 \ref{morphisms-definition-scheme-theoretically-dense}),
\item $U$  est dense dans  $X$  et  $U$  contient tous les points
immergés de  $X$, et
\item $U$  contient tous les points associés de  $X$.
\end{enumerate}
\end{lemma}

\begin{proof}
La question est locale sur  $X$, par conséquent nous pouvons supposer que
 $X = \Spec(A)$  où  $A$  est un anneau noethérien. Alors  $U$  est
quasi-compact (Propriétés, Lemme  \ref{properties-lemma-immersion-into-noetherian}) donc  $U = D(f_1) \cup \ldots \cup D(f_n)$  (Algèbre, Lemme
 \ref{algebra-lemma-qc-open}). Dans cette situation,  $U$  est schématiquement dense dans
 $X$  si et seulement si  $A \to A_{f_1} \times \ldots \times A_{f_n}$  est injectif, voir Morphismes,
Exemple  \ref{morphisms-example-scheme-theoretic-closure}. La condition (3) traduite en algèbre signifie que pour
chaque idéal premier associé  $\mathfrak p$  de  $A$  il existe un
 $i$  avec  $f_i \not \in \mathfrak p$.

\medskip\noindent
Supposons (1), c'est-à-dire que  $A \to A_{f_1} \times \ldots \times A_{f_n}$  est injectif. Disons que  $a \in A$ 
a un annulateur un premier  $\mathfrak p$. Puisque  $a$  se projette dans un
élément non nul de  $A_{f_i}$  pour un certain  $i$  nous voyons que
 $f_i \not \in \mathfrak p$. Ainsi, (3) est vérifiée.

\medskip\noindent
Notez que  $U$  est dense dans  $\Spec(A)$  si et seulement si il contient
tous les idéaux premiers minimaux. Comme l'ensemble des premiers associés de
 $A$  est égal à l'union de l'ensemble des premiers minimaux et de
l'ensemble des premiers associés incorporés (voir Algèbre, Proposition
 \ref{algebra-proposition-minimal-primes-associated-primes}), nous voyons que (2) et (3) sont équivalents.

\medskip\noindent
Supposons (3). Cela signifie que chaque premier associé  $\mathfrak p$  de
 $A$  correspond à un premier de  $A_{f_i}$  pour un certain  $i$.
Alors  $A \to A_{f_1} \times \ldots \times A_{f_n}$  est injectif parce que  $A \to \prod_{\mathfrak p \in \text{Ass}(A)} A_\mathfrak p$  est injectif selon Algèbre,
Lemme  \ref{algebra-lemma-zero-at-ass-zero}. Ainsi nous voyons que (3) implique (1).
\end{proof}

\begin{lemma}
\label{lemma-remove-embedded-points}
Soit  $X$  un schéma localement noethérien. Soit  $\mathcal{F}$  un faisceau
cohérent sur  $X$. L'ensemble des sous-faisceaux cohérents
$$
\{
\mathcal{K} \subset \mathcal{F}
\mid
\text{Supp}(\mathcal{K})\text{ is nowhere dense in }\text{Supp}(\mathcal{F})
\}
$$
contient un élément maximal  $\mathcal{K}$. En posant  $\mathcal{F}' = \mathcal{F}/\mathcal{K}$ , nous avons ce
qui suit
\begin{enumerate}
\item $\text{Supp}(\mathcal{F}') = \text{Supp}(\mathcal{F})$,
\item $\mathcal{F}'$  n’a pas de points associés intégrés, et
\item il existe un ouvert dense  $U \subset X$  tel que  $U \cap \text{Supp}(\mathcal{F})$  est dense dans
 $\text{Supp}(\mathcal{F})$  et  $\mathcal{F}'|_U \cong \mathcal{F}|_U$.
\end{enumerate}
\end{lemma}

\begin{proof}
Cela découle de l’Algèbre, Lemme  \ref{algebra-lemma-remove-embedded-primes}  et  \ref{algebra-lemma-remove-embedded-primes-localize}. Notez que
 $U$  peut être pris comme le complément de l’adhérence de l’ensemble des
points associés intégrés de  $\mathcal{F}$.
\end{proof}

\begin{lemma}
\label{lemma-no-embedded-points-endos}
Soit  $X$  un schéma localement noethérien. Soit  $\mathcal{F}$  un
 $\mathcal{O}_X$-module cohérent sans points associés intégrés. Posons
$$
\mathcal{I}
=
\Ker(\mathcal{O}_X
\longrightarrow
\SheafHom_{\mathcal{O}_X}(\mathcal{F}, \mathcal{F})).
$$
Il s’agit d’un faisceau cohérent d’idéaux qui définit un sous-schéma fermé
 $Z \subset X$  sans points immergés. De plus, il existe un faisceau cohérent
 $\mathcal{G}$  sur  $Z$  tel que (a)  $\mathcal{F} = (Z \to X)_*\mathcal{G}$, (b)  $\mathcal{G}$  n’a pas
de points associés intégrés, et (c)  $\text{Supp}(\mathcal{G}) = Z$  (comme ensembles).
\end{lemma}

\begin{proof}
Certain des énoncés que nous avons vus dans la preuve de la cohomologie des
schémas, Lemme  \ref{coherent-lemma-coherent-support-closed}. Les autres découlent de Algèbre, Lemme  \ref{algebra-lemma-no-embedded-primes-endos}.
\end{proof}



\section{Points faiblement associés}
\label{section-weakly-associated}

\noindent
Soit  $R$  un anneau et soit  $M$  un  $R$-module. Rappelons
qu'un  $\mathfrak p \subset R$  premier est  {\it  faiblement
associé}  à  $M$  s'il existe un élément  $m$  de
 $M$  tel que  $\mathfrak p$  soit minimal parmi les premiers contenant
l'annulateur de  $m$. Voir Algèbre, Définition  \ref{algebra-definition-weakly-associated}. Si
 $R$  est un anneau local avec un idéal maximal  $\mathfrak m$, alors
 $\mathfrak m$  est faiblement associé à  $M$  si et seulement s'il existe un
élément  $m \in M$  dont l'annulateur a pour radical  $\mathfrak m$, voir Algèbre,
Lemme  \ref{algebra-lemma-weakly-ass-local}.

\begin{definition}
\label{definition-weakly-associated}
Soit  $X$  un schéma. Soit  $\mathcal{F}$  un faisceau quasi-cohérent sur
 $X$.
\begin{enumerate}
\item Nous disons  $x \in X$  est  {\it faiblement
associé} à $\mathcal{F}$ si l'idéal maximal $\mathfrak m_x$ est faiblement
associé à la $\mathcal{O}_{X, x}$-module $\mathcal{F}_x$.
\item Nous notons  $\text{WeakAss}(\mathcal{F})$  l'ensemble des points faiblement associés de
 $\mathcal{F}$.
\item Les points faiblement associés de  {\it  de
 $X$}  sont les points faiblement associés de  $\mathcal{O}_X$.
\end{enumerate}
\end{definition}

\noindent
Dans ce cas, sur tout ouvert affine, cela correspond exactement aux idéaux idéaux
premiers faiblement associés tels que définis ci-dessus. Voici l'énoncé précis.

\begin{lemma}
\label{lemma-weakly-associated-affine-open}
Soit  $X$  un schéma. Soit  $\mathcal{F}$  un faisceau quasi-cohérent sur
 $X$. Soit  $\Spec(A) = U \subset X$  un ouvert affine, et posons  $M = \Gamma(U, \mathcal{F})$. Soit
 $x \in U$, et soit  $\mathfrak p \subset A$  le premier correspondant. Les énoncés suivants
sont équivalents
\begin{enumerate}
\item $\mathfrak p$  est faiblement associé à  $M$, et
\item $x$  est faiblement associé à  $\mathcal{F}$.
\end{enumerate}
\end{lemma}

\begin{proof}
Cela découle de Algèbre, Lemme  \ref{algebra-lemma-weakly-ass-local}.
\end{proof}

\begin{lemma}
\label{lemma-weakly-ass-support}
Soit  $X$  un schéma. Soit  $\mathcal{F}$  un  $\mathcal{O}_X$-module
quasi-cohérent. Alors
$$
\text{Ass}(\mathcal{F}) \subset \text{WeakAss}(\mathcal{F}) \subset
\text{Supp}(\mathcal{F}).
$$
\end{lemma}

\begin{proof}
Cela découle immédiatement des définitions.
\end{proof}

\begin{lemma}
\label{lemma-ses-weakly-ass}
Soit  $X$  un schéma. Soit  $0 \to \mathcal{F}_1 \to \mathcal{F}_2 \to \mathcal{F}_3 \to 0$  une suite exacte courte de
faisceaux quasi-cohérents sur  $X$. Alors  $\text{WeakAss}(\mathcal{F}_2) \subset
\text{WeakAss}(\mathcal{F}_1) \cup \text{WeakAss}(\mathcal{F}_3)$  et  $\text{WeakAss}(\mathcal{F}_1) \subset \text{WeakAss}(\mathcal{F}_2)$.
\end{lemma}

\begin{proof}
Pour chaque point  $x \in X$  la suite des germes  $0 \to \mathcal{F}_{1, x} \to \mathcal{F}_{2, x} \to \mathcal{F}_{3, x} \to 0$  est une suite
exacte courte de  $\mathcal{O}_{X, x}$-modules. Par conséquent, le lemme découle de
Algèbre, Lemme  \ref{algebra-lemma-weakly-ass}.
\end{proof}

\begin{lemma}
\label{lemma-weakly-ass-zero}
Soit  $X$  un schéma. Soit  $\mathcal{F}$  un  $\mathcal{O}_X$-module
quasi-cohérent. Alors
$$
\mathcal{F} = (0) \Leftrightarrow \text{WeakAss}(\mathcal{F}) = \emptyset
$$
\end{lemma}

\begin{proof}
Découle du Lemme  \ref{lemma-weakly-associated-affine-open}  et de Algèbre, Lemme  \ref{algebra-lemma-weakly-ass-zero}
\end{proof}

\begin{lemma}
\label{lemma-restriction-injective-open-contains-weakly-ass}
Soit  $X$  un schéma. Soit  $\mathcal{F}$  un  $\mathcal{O}_X$-module
quasi-cohérent. Si  $U \subset X$  est ouvert et  $\text{WeakAss}(\mathcal{F}) \subset U$, alors  $\Gamma(X, \mathcal{F}) \to \Gamma(U, \mathcal{F})$  est
injectif.
\end{lemma}

\begin{proof}
Soit  $s \in \Gamma(X, \mathcal{F})$  une section qui se restreint à zéro sur  $U$. Soit
 $\mathcal{F}' \subset \mathcal{F}$  l'image de l'application  $\mathcal{O}_X \to \mathcal{F}$  définie par  $s$. Alors
 $\text{Supp}(\mathcal{F}') \cap U = \emptyset$. D'autre part,  $\text{WeakAss}(\mathcal{F}') \subset \text{WeakAss}(\mathcal{F})$  d'après le Lemme  \ref{lemma-ses-weakly-ass}. Puisque
également  $\text{WeakAss}(\mathcal{F}') \subset \text{Supp}(\mathcal{F}')$  (Lemme  \ref{lemma-weakly-ass-support}), nous concluons  $\text{WeakAss}(\mathcal{F}') = \emptyset$. Par
conséquent,  $\mathcal{F}' = 0$  d'après le Lemme  \ref{lemma-weakly-ass-zero}.
\end{proof}

\begin{lemma}
\label{lemma-minimal-support-in-weakly-ass}
Soit  $X$  un schéma. Soit  $\mathcal{F}$  un  $\mathcal{O}_X$-module
quasi-cohérent. Soit  $x \in \text{Supp}(\mathcal{F})$  un point dans le support de  $\mathcal{F}$  qui
n'est pas une spécialisation d'un autre point de  $\text{Supp}(\mathcal{F})$. Alors  $x \in \text{WeakAss}(\mathcal{F})$.
En particulier, tout point générique d'une composante irréductible de  $X$ 
est faiblement associé à  $\mathcal{O}_X$.
\end{lemma}

\begin{proof}
Puisque  $x \in \text{Supp}(\mathcal{F})$  le module  $\mathcal{F}_x$  n'est pas nul. Par conséquent,
 $\text{WeakAss}(\mathcal{F}_x) \subset \Spec(\mathcal{O}_{X, x})$  est non vide d'après Algèbre, Lemme  \ref{algebra-lemma-weakly-ass-zero}. D'autre part, par
hypothèse  $\text{Supp}(\mathcal{F}_x) = \{\mathfrak m_x\}$. Comme  $\text{WeakAss}(\mathcal{F}_x) \subset \text{Supp}(\mathcal{F}_x)$  (Algèbre, Lemme  \ref{algebra-lemma-weakly-ass-support}), nous
voyons que  $\mathfrak m_x$  est faiblement associé à  $\mathcal{F}_x$  et nous avons gagné.
\end{proof}

\begin{lemma}
\label{lemma-ass-weakly-ass}
Soit  $X$  un schéma. Soit  $\mathcal{F}$  un  $\mathcal{O}_X$-module
quasi-cohérent. Si  $\mathfrak m_x$  est un idéal de type fini de  $\mathcal{O}_{X, x}$, alors
$$
x \in \text{Ass}(\mathcal{F}) \Leftrightarrow
x \in \text{WeakAss}(\mathcal{F}).
$$
En particulier, si  $X$  est localement noethérien, alors  $\text{Ass}(\mathcal{F}) = \text{WeakAss}(\mathcal{F})$.
\end{lemma}

\begin{proof}
Voir Algèbre, Lemme  \ref{algebra-lemma-ass-weakly-ass}.
\end{proof}

\begin{lemma}
\label{lemma-weakass-pushforward}
Soit  $f : X \to S$  un morphisme quasi-compact et quasi-séparé de schémas. Soit
 $\mathcal{F}$  un  $\mathcal{O}_X$-module quasi-cohérent. Soit  $s \in S$  un point qui
n'est pas dans l'image de  $f$. Alors  $s$  n'est pas faiblement
associé à  $f_*\mathcal{F}$.
\end{lemma}

\begin{proof}
Considérons le changement de base  $f' : X' \to \Spec(\mathcal{O}_{S, s})$  de  $f$  par le morphisme
 $g : \Spec(\mathcal{O}_{S, s}) \to S$  et notons  $g' : X' \to X$  l'autre projection. Alors
$$
(f_*\mathcal{F})_s = (g^*f_*\mathcal{F})_s = (f'_*(g')^*\mathcal{F})_s
$$
La première égalité parce que  $g$  induit un isomorphisme sur les anneaux
locaux en  $s$  et la seconde par changement de base plat (Cohomology of
Schémas, Lemme  \ref{coherent-lemma-flat-base-change-cohomology}). Bien sûr,  $s \in \Spec(\mathcal{O}_{S, s})$  n'est pas dans l'image de
 $f'$. Ainsi, nous pouvons supposer que  $S$  est le spectre d'un
anneau local  $(A, \mathfrak m)$  et que  $s$  correspond à  $\mathfrak m$. D'après
Schémas, Lemme  \ref{schemes-lemma-push-forward-quasi-coherent} , le faisceau  $f_*\mathcal{F}$  est quasi-cohérent, disons
correspondant au  $A$-module  $M$. Comme  $s$  n'est pas
dans l'image de  $f$ , nous voyons que  $X = \bigcup_{a \in \mathfrak m} f^{-1}D(a)$  est un recouvrement
ouvert. Comme  $X$  est quasi-compact, nous pouvons trouver  $a_1, \ldots, a_n \in \mathfrak m$ 
tel que  $X = f^{-1}D(a_1) \cup \ldots \cup f^{-1}D(a_n)$. Il s'ensuit que
$$
M \to M_{a_1} \oplus \ldots \oplus M_{a_r}
$$
est injectif. Par conséquent, pour tout élément non nul  $m$  de la germe
 $M_\mathfrak p$ , il existe un  $i$  tel que  $a_i^n m$  est non nul pour tout
 $n \geq 0$. Ainsi,  $\mathfrak m$  n'est pas faiblement associé à  $M$.
\end{proof}

\begin{lemma}
\label{lemma-check-injective-on-weakass}
Soit  $X$  un schéma. Soit  $\varphi : \mathcal{F} \to \mathcal{G}$  une application de
 $\mathcal{O}_X$-modules quasi-cohérents. Supposons que pour chaque  $x \in X$  au
moins l'un des cas suivants se produise
\begin{enumerate}
\item $\mathcal{F}_x \to \mathcal{G}_x$  est injectif, ou
\item $x \not \in \text{WeakAss}(\mathcal{F})$.
\end{enumerate}
Alors  $\varphi$  est injectif.
\end{lemma}

\begin{proof}
Les hypothèses impliquent que  $\text{WeakAss}(\Ker(\varphi)) = \emptyset$  et donc  $\Ker(\varphi) = 0$  par le Lemme
 \ref{lemma-weakly-ass-zero}.
\end{proof}

\begin{lemma}
\label{lemma-depth-2-hartog}
Soit  $X$  un schéma localement noethérien. Soit  $\mathcal{F}$  un
 $\mathcal{O}_X$-module cohérent. Soit  $j : U \to X$  un sous-schéma ouvert tel que pour
 $x \in X \setminus U$  nous avons  $\text{depth}(\mathcal{F}_x) \geq 2$. Alors
$$
\mathcal{F} \longrightarrow j_*(\mathcal{F}|_U)
$$
est un isomorphisme et par conséquent  $\Gamma(X, \mathcal{F}) \to \Gamma(U, \mathcal{F})$  est également un isomorphisme.
\end{lemma}

\begin{proof}
Nous affirmons que le Lemme  \ref{lemma-check-isomorphism-via-depth-and-ass}  s'applique à la carte affichée dans le
lemme. Soit  $x \in X$. Si  $x \in U$, alors la carte est un isomorphisme sur
les germes comme  $j_*(\mathcal{F}|_U)|_U = \mathcal{F}|_U$. Si  $x \in X \setminus U$, alors  $x \not \in \text{Ass}(j_*(\mathcal{F}|_U))$  (Lemmes
 \ref{lemma-weakass-pushforward}  et  \ref{lemma-weakly-ass-support}). Puisque nous avons supposé  $\text{depth}(\mathcal{F}_x) \geq 2$ , cela
termine la preuve.
\end{proof}

\begin{lemma}
\label{lemma-weakass-reduced}
Soit  $X$  un schéma réduit. Alors les points faiblement associés de
 $X$  sont exactement les points génériques des composantes irréductibles
de  $X$.
\end{lemma}

\begin{proof}
Suivant de Algèbre, Lemme  \ref{algebra-lemma-reduced-weakly-ass-minimal}.
\end{proof}



\section{Morphismes et points faiblement associés}
\label{section-morphisms-weakly-associated}

\begin{lemma}
\label{lemma-weakly-ass-reverse-functorial}
Soit  $f : X \to S$  un morphisme affine de schémas. Soit  $\mathcal{F}$  un
 $\mathcal{O}_X$-module quasi-cohérent. Alors nous avons
$$
\text{WeakAss}_S(f_*\mathcal{F}) \subset f(\text{WeakAss}_X(\mathcal{F}))
$$
\end{lemma}

\begin{proof}
On peut supposer que  $X$  et  $S$  sont affines, donc  $X \to S$ 
provient d'une application d'anneaux  $A \to B$. Alors  $\mathcal{F} = \widetilde M$  pour un
certain  $B$-module  $M$. D'après le Lemme  \ref{lemma-weakly-associated-affine-open} , les
points faiblement associés de  $\mathcal{F}$  correspondent exactement aux idéaux
idéaux premiers faiblement associés de  $M$. De même, les points
faiblement associés de  $f_*\mathcal{F}$  correspondent exactement aux idéaux idéaux
premiers faiblement associés de  $M$  en tant que  $A$-module. Par
conséquent, le lemme découle de Algèbre, Lemme  \ref{algebra-lemma-weakly-ass-reverse-functorial}.
\end{proof}

\begin{lemma}
\label{lemma-ass-functorial-equal}
Soit  $f : X \to S$  un morphisme affine de schémas. Soit  $\mathcal{F}$  un
 $\mathcal{O}_X$-module quasi-cohérent. Si  $X$  est localement noethérien,
alors nous avons
$$
f(\text{Ass}_X(\mathcal{F})) =
\text{Ass}_S(f_*\mathcal{F}) =
\text{WeakAss}_S(f_*\mathcal{F}) =
f(\text{WeakAss}_X(\mathcal{F}))
$$
\end{lemma}

\begin{proof}
Nous pouvons supposer que  $X$  et  $S$  sont affines, donc
 $X \to S$  provient d'un morphisme d'anneaux  $A \to B$. Comme  $X$ 
est localement noethérien, l'anneau  $B$  est noethérien, voir Propriétés,
Lemme  \ref{properties-lemma-locally-Noetherian}. Écrivons  $\mathcal{F} = \widetilde M$  pour un certain  $B$-module
 $M$. D'après le Lemme  \ref{lemma-associated-affine-open} , les points associés de  $\mathcal{F}$ 
correspondent exactement aux idéaux premiers associés de  $M$, et tout
idéal premier associé de  $M$  en tant que  $A$-module est un point
associé de  $f_*\mathcal{F}$. Ainsi, l'inclusion
$$
f(\text{Ass}_X(\mathcal{F})) \subset \text{Ass}_S(f_*\mathcal{F})
$$
découle du Lemme  \ref{algebra-lemma-ass-functorial-Noetherian} d'Algèbre. Nous avons l'inclusion
$$
\text{Ass}_S(f_*\mathcal{F}) \subset \text{WeakAss}_S(f_*\mathcal{F})
$$
d'après le Lemme  \ref{lemma-weakly-ass-support}. Nous avons l'inclusion
$$
\text{WeakAss}_S(f_*\mathcal{F}) \subset f(\text{WeakAss}_X(\mathcal{F}))
$$
d'après le Lemme  \ref{lemma-weakly-ass-reverse-functorial}. Les ensembles externes sont égaux d'après le Lemme
 \ref{lemma-ass-weakly-ass} , donc nous avons l'égalité partout.
\end{proof}

\begin{lemma}
\label{lemma-weakly-associated-finite}
Soit  $f : X \to S$  un morphisme fini de schémas. Soit  $\mathcal{F}$  un
 $\mathcal{O}_X$-module quasi-cohérent. Alors  $\text{WeakAss}(f_*\mathcal{F}) = f(\text{WeakAss}(\mathcal{F}))$.
\end{lemma}

\begin{proof}
Nous pouvons supposer que  $X$  et  $S$  sont affines, donc
 $X \to S$  provient d'un morphisme fini d'anneaux  $A \to B$. Notons
 $\mathcal{F} = \widetilde M$  pour un certain  $B$-module  $M$. D'après le Lemme
 \ref{lemma-weakly-associated-affine-open} , les points faiblement associés de  $\mathcal{F}$  correspondent
exactement aux idéaux idéaux premiers faiblement associés de  $M$. De
même, les points faiblement associés de  $f_*\mathcal{F}$  correspondent exactement aux
idéaux idéaux premiers faiblement associés de  $M$  en tant que
 $A$-module. Ainsi, le lemme découle du Lemme d'Algèbre  \ref{algebra-lemma-weakly-ass-finite-ring-map}.
\end{proof}

\begin{lemma}
\label{lemma-weakly-ass-pullback}
Soit  $f : X \to S$  un morphisme de schémas. Soit  $\mathcal{G}$  un
 $\mathcal{O}_S$-module quasi-cohérent. Soit  $x \in X$  avec  $s = f(x)$. Si
 $f$  est plat en  $x$, le point  $x$  est un point
générique de la fibre  $X_s$, et  $s \in \text{WeakAss}_S(\mathcal{G})$, alors  $x \in \text{WeakAss}(f^*\mathcal{G})$.
\end{lemma}

\begin{proof}
Soient  $A = \mathcal{O}_{S, s}$,  $B = \mathcal{O}_{X, x}$, et  $M = \mathcal{G}_s$. Soit  $m \in M$  un élément
dont l’annulateur  $I = \{a \in A \mid am = 0\}$  a un radical  $\mathfrak m_A$. Alors  $m \otimes 1$  a
pour annulateur  $I B$  car  $A \to B$  est fidèlement plat. Ainsi, il
suffit de voir que  $\sqrt{I B} = \mathfrak m_B$. Cela découle du fait que l'idéal maximal de
 $B/\mathfrak m_AB$  est localement nilpotent (voir Algèbre, Lemme  \ref{algebra-lemma-minimal-prime-reduced-ring}) et de
l’hypothèse que  $\sqrt{I} = \mathfrak m_A$. Certains détails sont omis.
\end{proof}

\begin{lemma}
\label{lemma-weakly-ass-change-fields}
Soit  $K/k$  une extension de corps. Soit  $X$  un schéma sur
 $k$. Soit  $\mathcal{F}$  un  $\mathcal{O}_X$-module quasi-cohérent. Soit
 $y \in X_K$  avec image  $x \in X$. Si  $y$  est un point faiblement
associé du tiré en arrière  $\mathcal{F}_K$, alors  $x$  est un point
faiblement associé de  $\mathcal{F}$.
\end{lemma}

\begin{proof}
Ceci est la traduction de Algèbre, Lemme  \ref{algebra-lemma-weakly-ass-change-fields}  dans le langage des schémas.
\end{proof}

\noindent
Voici un lemme simple où nous constatons que les images directes ont souvent une
profondeur d'au moins 2.

\begin{lemma}
\label{lemma-depth-pushforward}
Soit  $f : X \to S$  un morphisme de schémas quasi-compact et quasi-séparé. Soit
 $\mathcal{F}$  un  $\mathcal{O}_X$-module quasi-cohérent. Soit  $s \in S$.
\begin{enumerate}
\item Si  $s \not \in f(X)$, alors  $s$  n'est pas faiblement associé à  $f_*\mathcal{F}$.
\item Si  $s \not \in f(X)$  et  $\mathcal{O}_{S, s}$  est noethérien, alors  $s$  n'est pas
associé à  $f_*\mathcal{F}$.
\item Si  $s \not \in f(X)$,  $(f_*\mathcal{F})_s$  est un  $\mathcal{O}_{S, s}$-module fini, et  $\mathcal{O}_{S, s}$  est
noethérien, alors  $\text{depth}((f_*\mathcal{F})_s) \geq 2$.
\item Si  $\mathcal{F}$  est plat sur  $S$  et  $a \in \mathfrak m_s$  est un non-diviseur de
zéro, alors  $a$  est un non-diviseur de zéro sur  $(f_*\mathcal{F})_s$.
\item Si  $\mathcal{F}$  est plat sur  $S$  et  $a, b \in \mathfrak m_s$  est une suite régulière,
alors  $a$  est un non-diviseur de zéro sur  $(f_*\mathcal{F})_s$  et  $b$ 
est un non-diviseur de zéro sur  $(f_*\mathcal{F})_s/a(f_*\mathcal{F})_s$.
\item Si  $\mathcal{F}$  est plat sur  $S$  et  $(f_*\mathcal{F})_s$  est un
 $\mathcal{O}_{S, s}$-module fini, alors  $\text{depth}((f_*\mathcal{F})_s) \geq
\min(2, \text{depth}(\mathcal{O}_{S, s}))$.
\end{enumerate}
\end{lemma}

\begin{proof}
La partie (1) est le Lemme  \ref{lemma-weakass-pushforward}. La partie (2) découle de (1) et du Lemme
 \ref{lemma-ass-weakly-ass}.

\medskip\noindent
Preuve de la partie (3). Pour montrer que la profondeur est  $\geq 2$ , il
suffit de montrer que  $\Hom_{\mathcal{O}_{S, s}}(\kappa(s), (f_*\mathcal{F})_s) = 0$  et  $\Ext^1_{\mathcal{O}_{S, s}}(\kappa(s), (f_*\mathcal{F})_s) = 0$, voir Algèbre, Lemme
 \ref{algebra-lemma-depth-ext}. En utilisant la suite exacte  $0 \to \mathfrak m_s \to \mathcal{O}_{S, s} \to \kappa(s) \to 0$ , il suffit de prouver que
l'application
$$
\Hom_{\mathcal{O}_{S, s}}(\mathcal{O}_{S, s}, (f_*\mathcal{F})_s)
\to
\Hom_{\mathcal{O}_{S, s}}(\mathfrak m_s, (f_*\mathcal{F})_s)
$$
est un isomorphisme. Par changement de base plat (Cohomologie des schémas, Lemme
 \ref{coherent-lemma-flat-base-change-cohomology}), nous pouvons remplacer  $S$  par  $\Spec(\mathcal{O}_{S, s})$  et
 $X$  par  $\Spec(\mathcal{O}_{S, s}) \times_S X$. Notons  $\mathfrak m \subset \mathcal{O}_S$  le faisceau d'idéaux de
 $s$. Alors nous constatons que
$$
\Hom_{\mathcal{O}_{S, s}}(\mathfrak m_s, (f_*\mathcal{F})_s) =
\Hom_{\mathcal{O}_S}(\mathfrak m, f_*\mathcal{F}) =
\Hom_{\mathcal{O}_X}(f^*\mathfrak m, \mathcal{F})
$$
la première égalité car  $S$  est local avec le point fermé  $s$  et
la deuxième égalité par adjonction pour  $f^*, f_*$  sur des modules
quasi-cohérents. Cependant, puisque  $s \not \in f(X)$  nous voyons que  $f^*\mathfrak m = \mathcal{O}_X$. En
revenant sur les arguments, nous obtenons l'égalité désirée.

\medskip\noindent
Pour la preuve de (4), (5) et (6), nous utilisons le changement de base plat
(Cohomologie des schémas, Lemme  \ref{coherent-lemma-flat-base-change-cohomology}) pour réduire au cas où  $S$ 
est le spectre de  $\mathcal{O}_{S, s}$. Alors un non-diviseur  $a \in \mathcal{O}_{S, s}$  détermine une
suite exacte courte
$$
0 \to \mathcal{O}_S \xrightarrow{a} \mathcal{O}_S \to
\mathcal{O}_S/a \mathcal{O}_S \to 0
$$
Puisque  $\mathcal{F}$  est plat sur  $S$, nous obtenons une suite exacte
$$
0 \to \mathcal{F} \xrightarrow{a} \mathcal{F} \to
\mathcal{F}/a\mathcal{F} \to 0
$$
En poussant vers l'avant, nous obtenons une suite exacte
$$
0 \to f_*\mathcal{F} \xrightarrow{a} f_*\mathcal{F} \to
f_*(\mathcal{F}/a\mathcal{F})
$$
. Cela prouve (4) et montre que  $f_*\mathcal{F}/ af_*\mathcal{F} \subset f_*(\mathcal{F}/a\mathcal{F})$. Si  $b$  est un diviseur de
zéro non nul sur  $\mathcal{O}_{S, s}/a\mathcal{O}_{S, s}$, alors le même argument exact montre que
 $b : \mathcal{F}/a\mathcal{F} \to \mathcal{F}/a\mathcal{F}$  est injectif. En poussant vers l'avant, nous concluons que
$$
b : f_*(\mathcal{F}/a\mathcal{F}) \to f_*(\mathcal{F}/a\mathcal{F})
$$
est injectif et donc aussi  $b : f_*\mathcal{F}/ af_*\mathcal{F} \to f_*\mathcal{F}/ af_*\mathcal{F}$  est injectif. Cela prouve (5). La partie
(6) découle de (4) et (5) et des définitions.
\end{proof}











\section{Assassin relatif}
\label{section-relative-assassin}

\noindent
Soit  $A \to B$  un morphisme de anneaux. Soit  $N$  un
 $B$-module. Rappelons qu'un idéal premier  $\mathfrak q \subset B$  est dit appartenir
à l'assassin relatif de  $N$  sur  $B/A$  si  $\mathfrak q$  est un idéal
premier associé de  $N \otimes_A \kappa(\mathfrak p)$. Ici  $\mathfrak p = A \cap \mathfrak q$. Voir Algèbre, Définition
 \ref{algebra-definition-relative-assassin}. Voici la définition de l'assassin relatif pour les faisceaux
quasi-cohérents sur un morphisme de schémas.

\begin{definition}
\label{definition-relative-assassin}
Laisse $f : X \to S$ soit un morphisme de schémas. Soit $\mathcal{F}$ être
quasi-cohérent $\mathcal{O}_X$-module. Le  {\it assassin relatif
de $\mathcal{F}$ dans $X$ au-dessus $S$} est l'ensemble
$$
\text{Ass}_{X/S}(\mathcal{F}) =
\bigcup\nolimits_{s \in S} \text{Ass}_{X_s}(\mathcal{F}_s)
$$
où  $\mathcal{F}_s = (X_s \to X)^*\mathcal{F}$  est la restriction de  $\mathcal{F}$  à la fibre de  $f$  en
 $s$.
\end{definition}

\noindent
Il y a encore un avertissement selon lequel cela est préférable lorsque les fibres
de  $f$  sont localement noethériennes et que  $\mathcal{F}$  est de type
fini. Dans le cas général, nous devrions probablement utiliser l'assassin relatif
faible (défini dans la section suivante). Relions la notion schématique avec la
notion algébrique sur les ouverts affines; notez que cette correspondance
fonctionne parfaitement uniquement pour les morphismes de schémas dont les fibres
sont localement noethériennes.

\begin{lemma}
\label{lemma-relative-assassin-affine-open}
Soit  $f : X \to S$  un morphisme de schémas. Soit  $\mathcal{F}$  un faisceau
quasi-cohérent sur  $X$. Soit  $U \subset X$  et  $V \subset S$  des ouverts
affines avec  $f(U) \subset V$. Écrire  $U = \Spec(A)$,  $V = \Spec(R)$, et définir
 $M = \Gamma(U, \mathcal{F})$. Soit  $x \in U$, et soit  $\mathfrak p \subset A$  le premier idéal
correspondant. Alors
$$
\mathfrak p \in \text{Ass}_{A/R}(M) \Rightarrow
x \in \text{Ass}_{X/S}(\mathcal{F})
$$
Si toutes les fibres  $X_s$  de  $f$  sont localement noethériennes,
alors  $\mathfrak p \in \text{Ass}_{A/R}(M) \Leftrightarrow
x \in \text{Ass}_{X/S}(\mathcal{F})$  pour toutes les paires  $(\mathfrak p, x)$  comme ci-dessus.
\end{lemma}

\begin{proof}
L'ensemble  $\text{Ass}_{A/R}(M)$  est défini en Algèbre, Définition  \ref{algebra-definition-relative-assassin}. Choisissez
une paire  $(\mathfrak p, x)$. Soit  $s = f(x)$. Soit  $\mathfrak r \subset R$  le premier situé sous
 $\mathfrak p$, c'est-à-dire le premier correspondant à  $s$. Soit
 $\mathfrak p' \subset A \otimes_R \kappa(\mathfrak r)$  le premier dont l'image inverse est  $\mathfrak p$, c'est-à-dire le
premier correspondant à  $x$  considéré comme un point de sa fibre
 $X_s$. Alors  $\mathfrak p \in \text{Ass}_{A/R}(M)$  si et seulement si  $\mathfrak p'$  est un premier
associé de  $M \otimes_R \kappa(\mathfrak r)$, voir Algèbre, Lemme  \ref{algebra-lemma-compare-relative-assassins}. Notez que l'anneau
 $A \otimes_R \kappa(\mathfrak r)$  correspond à  $U_s$  et que le module  $M \otimes_R \kappa(\mathfrak r)$  correspond au
faisceau quasi-cohérent  $\mathcal{F}_s|_{U_s}$. Ainsi  $x$  est un point associé de
 $\mathcal{F}_s$  par le Lemme  \ref{lemma-associated-affine-open}. L'implication inverse est valable si
 $\mathfrak p'$  est de type fini, ce qui permet de voir que la dernière phrase est
vraie.
\end{proof}

\begin{lemma}
\label{lemma-base-change-relative-assassin}
Soit  $f : X \to S$  un morphisme de schémas. Soit  $\mathcal{F}$  un
 $\mathcal{O}_X$-module quasi-cohérent. Soit  $g : S' \to S$  un morphisme de schémas.
Considérons le diagramme de changement de base
$$
\xymatrix{
X' \ar[d] \ar[r]_{g'} & X \ar[d] \\
S' \ar[r]^g & S
}
$$
et posons  $\mathcal{F}' = (g')^*\mathcal{F}$. Soit  $x' \in X'$  un point avec les images  $x \in X$,
 $s' \in S'$  et  $s \in S$. Supposons  $f$  localement de type fini.
Alors  $x' \in \text{Ass}_{X'/S'}(\mathcal{F}')$  si et seulement si  $x \in \text{Ass}_{X/S}(\mathcal{F})$  et  $x'$  correspond à un
point générique d'une composante irréductible de  $\Spec(\kappa(s') \otimes_{\kappa(s)} \kappa(x))$.
\end{lemma}

\begin{proof}
Considérons le morphisme  $X'_{s'} \to X_s$  des fibres. Comme  $X_{s'} = X_s \times_{\Spec(\kappa(s))} \Spec(\kappa(s'))$ , il s'agit
d'un morphisme plat. De plus,  $\mathcal{F}'_{s'}$  est le tiré en arrière de  $\mathcal{F}_s$ 
via ce morphisme. Comme  $X_s$  est localement de type fini sur le schéma
noethérien  $\Spec(\kappa(s))$ , nous avons que  $X_s$  est localement noethérien,
voir Morphismes, Lemme  \ref{morphisms-lemma-finite-type-noetherian}. Ainsi, nous pouvons appliquer le Lemme
 \ref{lemma-bourbaki}  et nous voyons que
$$
\text{Ass}_{X'_{s'}}(\mathcal{F}'_{s'}) =
\bigcup\nolimits_{x \in \text{Ass}(\mathcal{F}_s)} \text{Ass}((X'_{s'})_x).
$$
Ainsi, pour prouver le lemme, il suffit de montrer que les points associés de la
fibre  $(X'_{s'})_x$  du morphisme  $X'_{s'} \to X_s$  au-dessus de  $x$  sont ses
points génériques. Notez que  $(X'_{s'})_x = \Spec(\kappa(s') \otimes_{\kappa(s)} \kappa(x))$  en tant que schémas. D'après Algèbre,
Lemme  \ref{algebra-lemma-tensor-fields-CM} , l'anneau  $\kappa(s') \otimes_{\kappa(s)} \kappa(x)$  est un anneau de Cohen-Macaulay
noethérien. Par conséquent, ses idéaux premiers associés sont ses idéaux premiers
minimaux, voir Algèbre, Proposition  \ref{algebra-proposition-minimal-primes-associated-primes}  (les idéaux premiers minimaux
sont associés) et Algèbre, Lemme  \ref{algebra-lemma-criterion-no-embedded-primes}  (pas d'idéaux premiers incorporés).
\end{proof}

\begin{remark}
\label{remark-base-change-relative-assassin}
Avec les notations et hypothèses comme dans le Lemme  \ref{lemma-base-change-relative-assassin} , nous voyons
qu'il est toujours le cas que  $(g')^{-1}(\text{Ass}_{X/S}(\mathcal{F})) \supset
\text{Ass}_{X'/S'}(\mathcal{F}')$. Si le morphisme  $S' \to S$  est
quasi-fini localement, alors nous avons en fait
$$
(g')^{-1}(\text{Ass}_{X/S}(\mathcal{F}))
=
\text{Ass}_{X'/S'}(\mathcal{F}')
$$
car dans ce cas les extensions de corps  $\kappa(s')/\kappa(s)$  sont toujours finies. En
fait, cela est vrai de manière plus générale pour tout morphisme  $g : S' \to S$  tel
que toutes les extensions de corps  $\kappa(s')/\kappa(s)$  soient algébriques, car dans ce
cas tous les idéaux premiers de  $\kappa(s') \otimes_{\kappa(s)} \kappa(x)$  sont des premiers maximaux (et
minimaux), voir Algèbre, Lemme  \ref{algebra-lemma-integral-over-field}.
\end{remark}




\section{Assassin faible relatif}
\label{section-relative-weak-assassin}

\begin{definition}
\label{definition-relative-weak-assassin}
Soit  $f : X \to S$  un morphisme de schémas. Soit  $\mathcal{F}$  un
 $\mathcal{O}_X$-module quasi-cohérent. L'assassin faible relatif
 {\it  de  $\mathcal{F}$  dans  $X$  sur
 $S$}  est l'ensemble
$$
\text{WeakAss}_{X/S}(\mathcal{F}) =
\bigcup\nolimits_{s \in S} \text{WeakAss}(\mathcal{F}_s)
$$
où  $\mathcal{F}_s = (X_s \to X)^*\mathcal{F}$  est la restriction de  $\mathcal{F}$  à la fibre de  $f$  en
 $s$.
\end{definition}

\begin{lemma}
\label{lemma-relative-weak-assassin-assassin-finite-type}
Soit  $f : X \to S$  un morphisme de schémas qui est localement de type fini. Soit
 $\mathcal{F}$  un  $\mathcal{O}_X$-module quasi-cohérent. Alors  $\text{WeakAss}_{X/S}(\mathcal{F}) = \text{Ass}_{X/S}(\mathcal{F})$.
\end{lemma}

\begin{proof}
C'est vrai parce que les fibres de  $f$  sont des schémas localement
noethériens, et les points associés et faiblement associés coïncident sur les
schémas localement noethériens, voir le Lemme  \ref{lemma-ass-weakly-ass}.
\end{proof}

\begin{lemma}
\label{lemma-relative-weak-assassin-finite}
Soit  $f : X \to S$  un morphisme de schémas. Soit  $i : Z \to X$  un morphisme fini.
Soit  $\mathcal{F}$  un  $\mathcal{O}_Z$-module quasi-cohérent. Alors  $\text{WeakAss}_{X/S}(i_*\mathcal{F}) =
i(\text{WeakAss}_{Z/S}(\mathcal{F}))$.
\end{lemma}

\begin{proof}
Soit  $i_s : Z_s \to X_s$  le morphisme induit entre les fibres. Alors  $(i_*\mathcal{F})_s = i_{s, *}(\mathcal{F}_s)$  d'après
Cohomologie des schémas, Lemme  \ref{coherent-lemma-affine-base-change}  et le fait que  $i$  est
affine. Par conséquent, nous pouvons appliquer le Lemme  \ref{lemma-weakly-associated-finite}  pour
conclure.
\end{proof}





\section{Idéaux de Fitting}
\label{section-fitting-ideals}

\noindent
Cette section est la continuation de la discussion dans Plus sur l'Algèbre,
Section  \ref{more-algebra-section-fitting-ideals}. Soit  $S$  un schéma. Soit  $\mathcal{F}$  un
 $\mathcal{O}_S$-module quasi-cohérent de type fini. Dans cette situation, nous
pouvons construire les idéaux de Fitting
$$
0 = \text{Fit}_{-1}(\mathcal{F}) \subset \text{Fit}_0(\mathcal{F}) \subset
\text{Fit}_1(\mathcal{F}) \subset \ldots \subset \mathcal{O}_S
$$
comme la suite des idéaux quasi-cohérents caractérisés par la propriété suivante:
pour tout ouvert affine  $U = \Spec(A)$  de  $S$ , si  $\mathcal{F}|_U$  correspond
au  $A$-module  $M$, alors  $\text{Fit}_i(\mathcal{F})|_U$  correspond à l'idéal
 $\text{Fit}_i(M) \subset A$. Cela est bien défini et constitue un faisceau quasi-cohérent
d'idéaux parce que si  $f \in A$, alors le  $i$-ième idéal de Fitting de
 $M_f$  sur  $A_f$  est égal à  $\text{Fit}_i(M) A_f$  selon More on Algebra, Lemme
 \ref{more-algebra-lemma-fitting-ideal-basics}.

\medskip\noindent
Alternativement, nous pouvons construire les idéaux de Fitting en termes de
présentations locales de  $\mathcal{F}$. À savoir, si  $U \subset X$  est ouvert, et
$$
\bigoplus\nolimits_{i \in I} \mathcal{O}_U \to
\mathcal{O}_U^{\oplus n} \to \mathcal{F}|_U \to 0
$$
est une présentation de  $\mathcal{F}$  sur  $U$, alors  $\text{Fit}_r(\mathcal{F})|_U$  est
généré par les  $(n - r) \times (n - r)$ mineurs de la matrice définissant la première flèche
de la présentation. Cela est compatible avec la construction ci-dessus car c'est
ainsi que l'idéal de Fitting d'un module sur un anneau est réellement défini.
Quelques détails omis.

\begin{lemma}
\label{lemma-base-change-fitting-ideal}
Soit  $f : T \to S$  un morphisme de schémas. Soit  $\mathcal{F}$  un
 $\mathcal{O}_S$-module quasi-cohérent de type fini. Alors  $f^{-1}\text{Fit}_i(\mathcal{F}) \cdot \mathcal{O}_T =
\text{Fit}_i(f^*\mathcal{F})$.
\end{lemma}

\begin{proof}
Suit immédiatement de Plus sur Algèbre, Lemme  \ref{more-algebra-lemma-fitting-ideal-basics}  partie (3).
\end{proof}

\begin{lemma}
\label{lemma-fitting-ideal-of-finitely-presented}
Soit  $S$  un schéma. Soit  $\mathcal{F}$  un  $\mathcal{O}_S$-module de
présentation finie. Alors  $\text{Fit}_r(\mathcal{F})$  est un idéal quasi-cohérent de type fini.
\end{lemma}

\begin{proof}
Suit immédiatement de Plus sur Algèbre, Lemme  \ref{more-algebra-lemma-fitting-ideal-basics}  partie (4).
\end{proof}

\begin{lemma}
\label{lemma-on-subscheme-cut-out-by-Fit-0}
Soit  $S$  un schéma. Soit  $\mathcal{F}$  un  $\mathcal{O}_S$-module
quasi-cohérent de type fini. Soit  $Z_0 \subset S$  le sous-schéma fermé défini par
 $\text{Fit}_0(\mathcal{F})$. Soit  $Z \subset S$  le support schématique de  $\mathcal{F}$. Alors
\begin{enumerate}
\item $Z \subset Z_0 \subset S$  en tant que sous-schémas fermés,
\item $Z = Z_0 = \text{Supp}(\mathcal{F})$  en tant que sous-ensembles fermés,
\item il existe un  $\mathcal{O}_{Z_0}$-module quasi-cohérent de type fini  $\mathcal{G}_0$  avec
$$
(Z_0 \to X)_*\mathcal{G}_0 = \mathcal{F}.
$$
\end{enumerate}
\end{lemma}

\begin{proof}
Rappelons que  $Z$  est localement défini par l'annulateur de  $\mathcal{F}$,
voir Morphismes, Définition  \ref{morphisms-definition-scheme-theoretic-support}  (qui utilise Morphismes, Lemme
 \ref{morphisms-lemma-scheme-theoretic-support}  pour définir  $Z$). Par conséquent, nous voyons que
 $Z \subset Z_0$  schématiquement par More on Algebra, Lemme  \ref{more-algebra-lemma-fitting-ideal-basics}  partie (6).
D'autre part, nous avons  $Z = \text{Supp}(\mathcal{F})$  schématiquement par Morphismes, Lemme
 \ref{morphisms-lemma-scheme-theoretic-support}  et nous avons  $Z_0 = Z$  schématiquement par More on Algebra, Lemme
 \ref{more-algebra-lemma-fitting-ideal-basics}  partie (7). Enfin, pour obtenir  $\mathcal{G}_0$  comme dans la partie
(3), nous pouvons soit utiliser le fait que nous avons  $\mathcal{G}$  sur
 $Z$  comme dans Morphismes, Lemme  \ref{morphisms-lemma-scheme-theoretic-support}  et définir  $\mathcal{G}_0 = (Z \to Z_0)_*\mathcal{G}$ ,
soit utiliser Morphismes, Lemme  \ref{morphisms-lemma-i-star-equivalence}  et le fait que  $\text{Fit}_0(\mathcal{F})$  annule
 $\mathcal{F}$  selon Plus sur Algèbre, Lemme  \ref{more-algebra-lemma-fitting-ideal-basics}  partie (6).
\end{proof}

\begin{lemma}
\label{lemma-fitting-ideal-generate-locally}
Soit  $S$  un schéma. Soit  $\mathcal{F}$  un  $\mathcal{O}_S$-module
quasi-cohérent de type fini. Soit  $s \in S$. Alors  $\mathcal{F}$  peut être
généré par  $r$  éléments dans un voisinage de  $s$  si et seulement
si  $\text{Fit}_r(\mathcal{F})_s = \mathcal{O}_{S, s}$.
\end{lemma}

\begin{proof}
Suit immédiatement de Plus sur l’Algèbre, Lemme  \ref{more-algebra-lemma-fitting-ideal-generate-locally}.
\end{proof}

\begin{lemma}
\label{lemma-fitting-ideal-finite-locally-free}
Soit  $S$  un schéma. Soit  $\mathcal{F}$  un  $\mathcal{O}_S$-module
quasi-cohérent de type fini. Soit  $r \geq 0$. Les éléments suivants sont
équivalents
\begin{enumerate}
\item $\mathcal{F}$  est localement libre de rang fini  $r$
\item $\text{Fit}_{r - 1}(\mathcal{F}) = 0$  et  $\text{Fit}_r(\mathcal{F}) = \mathcal{O}_S$, et
\item $\text{Fit}_k(\mathcal{F}) = 0$  pour  $k < r$  et  $\text{Fit}_k(\mathcal{F}) = \mathcal{O}_S$  pour  $k \geq r$.
\end{enumerate}
\end{lemma}

\begin{proof}
Suit immédiatement de Plus sur l’Algèbre, Lemme  \ref{more-algebra-lemma-fitting-ideal-finite-locally-free}.
\end{proof}

\begin{lemma}
\label{lemma-locally-free-rank-r-pullback}
Soit  $S$  un schéma. Soit  $\mathcal{F}$  un  $\mathcal{O}_S$-module
quasi-cohérent de type fini. Les sous-schéma fermés
$$
S = Z_{-1} \supset Z_0 \supset Z_1 \supset Z_2 \ldots
$$
définis par les idéaux de Fitting de  $\mathcal{F}$  ont les propriétés suivantes
\begin{enumerate}
\item L’intersection  $\bigcap Z_r$  est vide.
\item Le foncteur  $(\Sch/S)^{opp} \to \textit{Sets}$  défini par la règle
$$
T \longmapsto
\left\{
\begin{matrix}
\{*\} & \text{if }\mathcal{F}_T\text{ is locally generated by }
\leq r\text{ sections} \\
\emptyset & \text{otherwise}
\end{matrix}
\right.
$$
est représentable par le sous-schéma ouvert  $S \setminus Z_r$.
\item Le foncteur  $F_r : (\Sch/S)^{opp} \to \textit{Sets}$  défini par la règle
$$
T \longmapsto
\left\{
\begin{matrix}
\{*\} & \text{if }\mathcal{F}_T\text{ locally free rank }r\\
\emptyset & \text{otherwise}
\end{matrix}
\right.
$$
est représentable par le sous-schéma localement fermé  $Z_{r - 1} \setminus Z_r$  de
 $S$.
\end{enumerate}
Si  $\mathcal{F}$  est de présentation finie, alors  $Z_r \to S$,  $S \setminus Z_r \to S$ et
 $Z_{r - 1} \setminus Z_r \to S$  sont de présentation finie.
\end{lemma}

\begin{proof}
La partie (1) est vraie car sur chaque ouvert affine  $U$ , il existe un
entier  $n$  tel que  $\text{Fit}_n(\mathcal{F})|_U = \mathcal{O}_U$. À savoir, nous pouvons prendre
 $n$  comme étant le nombre de générateurs de  $\mathcal{F}$  sur
 $U$, voir More on Algèbre, Section  \ref{more-algebra-section-fitting-ideals}.

\medskip\noindent
Pour tout morphisme  $g : T \to S$ , nous voyons d'après les lemmes  \ref{lemma-base-change-fitting-ideal}  et
 \ref{lemma-fitting-ideal-generate-locally}  que  $\mathcal{F}_T$  est localement engendré par  $\leq r$  sections si
et seulement si  $\text{Fit}_r(\mathcal{F}) \cdot \mathcal{O}_T = \mathcal{O}_T$. Cela prouve (2).

\medskip\noindent
Pour tout morphisme  $g : T \to S$ , nous voyons d'après les lemmes  \ref{lemma-base-change-fitting-ideal}  et
 \ref{lemma-fitting-ideal-finite-locally-free}  que  $\mathcal{F}_T$  est libre de rang  $r$  si et seulement si
 $\text{Fit}_r(\mathcal{F}) \cdot \mathcal{O}_T = \mathcal{O}_T$  et  $\text{Fit}_{r - 1}(\mathcal{F}) \cdot \mathcal{O}_T = 0$. Cela prouve (3).

\medskip\noindent
Supposons que  $\mathcal{F}$  soit de présentation finie. Alors chacun des morphismes
 $Z_r \to S$  est de présentation finie puisque  $\text{Fit}_r(\mathcal{F})$  est de type fini
(Lemme  \ref{lemma-fitting-ideal-of-finitely-presented}  et Morphismes, Lemme  \ref{morphisms-lemma-closed-immersion-finite-presentation}). Cela implique que
 $Z_{r - 1} \setminus Z_r$  est un ouvert rétrocompact dans  $Z_r$  (Propriétés, Lemme
 \ref{properties-lemma-quasi-coherent-finite-type-ideals}) et donc le morphisme  $Z_{r - 1} \setminus Z_r \to Z_r$  est également de présentation
finie.
\end{proof}

\noindent
Indépendamment du Lemme  \ref{lemma-locally-free-rank-r-pullback} , le lemme suivant ne tient pas si
 $\mathcal{F}$  est un module quasi-cohérent de type fini. À savoir, la
stratification existe toujours mais il n'est pas vrai qu'elle représente le
foncteur  $F_{flat}$  en général.

\begin{lemma}
\label{lemma-finite-presentation-module}
Soit  $S$  un schéma. Soit  $\mathcal{F}$  un  $\mathcal{O}_S$-module de
présentation finie. Soit  $S = Z_{-1} \supset Z_0 \supset Z_1 \supset \ldots$  comme dans le Lemme  \ref{lemma-locally-free-rank-r-pullback}. Posons
 $S_r = Z_{r - 1} \setminus Z_r$. Alors  $S' = \coprod_{r \geq 0} S_r$  représente le foncteur
$$
F_{flat} : \Sch/S \longrightarrow \textit{Sets},\quad\quad
T \longmapsto
\left\{
\begin{matrix}
\{*\} & \text{if }\mathcal{F}_T\text{ flat over }T\\
\emptyset & \text{otherwise}
\end{matrix}
\right.
$$
De plus,  $\mathcal{F}|_{S_r}$  est localement libre de rang  $r$  et les morphismes
 $S_r \to S$  et  $S' \to S$  sont de présentation finie.
\end{lemma}

\begin{proof}
Supposons que  $g : T \to S$  soit un morphisme de schémas tel que le tiré en arrière
 $\mathcal{F}_T = g^*\mathcal{F}$  soit plat. Alors  $\mathcal{F}_T$  est un  $\mathcal{O}_T$-module plat de
présentation finie. Par conséquent,  $\mathcal{F}_T$  est fini et localement libre,
voir Propriétés, Lemme  \ref{properties-lemma-finite-locally-free}. Ainsi  $T = \coprod_{r \geq 0} T_r$, où  $\mathcal{F}_T|_{T_r}$  est
localement libre de rang  $r$. Cela implique que
$$
F_{flat} = \coprod\nolimits_{r \geq 0} F_r
$$
dans la catégorie des faisceaux de Zariski sur  $\Sch/S$  où  $F_r$  est
comme dans le Lemme  \ref{lemma-locally-free-rank-r-pullback}. Il s'ensuit que  $F_{flat}$  est représenté par
 $\coprod_{r \geq 0} (Z_{r - 1} \setminus Z_r)$  où  $Z_r$  est comme dans le Lemme  \ref{lemma-locally-free-rank-r-pullback}. Les autres
énoncés découlent également du lemme.
\end{proof}

\begin{example}
\label{example-not-fp-MB}
Soit  $R = \prod_{n \in \mathbf{N}} \mathbf{F}_2$. Soit  $I \subset R$  l'idéal des éléments  $a = (a_n)_{n \in \mathbf{N}}$  dont
presque toutes les composantes sont nulles. Soit  $\mathfrak m$  un idéal maximal
contenant  $I$. Alors  $M = R/\mathfrak m$  est un  $R$-module fini et
plat, car  $R$  est absolument plat (Plus sur Algèbre, Lemme  \ref{more-algebra-lemma-product-fields-absolutely-flat}).
Posons  $S = \Spec(R)$  et  $\mathcal{F} = \widetilde{M}$. Les sous-schémas fermés du Lemme  \ref{lemma-locally-free-rank-r-pullback} 
sont  $S = Z_{-1}$,  $Z_0 = \Spec(R/\mathfrak m)$ et  $Z_i = \emptyset$  pour  $i > 0$. Mais
 $\text{id} : S \to S$  ne se factorise pas par  $(S \setminus Z_0) \amalg Z_0$  car  $\mathfrak m$  est un point
non isolé de  $S$. Ainsi, le Lemme  \ref{lemma-finite-presentation-module}  ne tient pas pour les
modules de type fini.
\end{example}







\section{Le lieu singulier d'un morphisme}
\label{section-singular-locus-morphism}

\noindent
Soit  $f : X \to S$  un morphisme de type fini de schémas. L'ensemble  $U$ 
des points où  $f$  est lisse est un ouvert de  $X$  (d'après
Morphismes, Définition  \ref{morphisms-definition-smooth}). Dans de nombreuses situations, il est utile
d'avoir un sous-schéma fermé canonique  $\text{Sing}(f) \subset X$  dont le complémentaire est
 $U$  et dont la formation commute avec tout changement de base.

\medskip\noindent
Si  $f$  est de présentation finie, alors un choix serait de considérer le
sous-schéma fermé  $Z$  défini par des fonctions qui sont affinement
localement « strictement standard » au sens de Lissage des morphismes d'anneaux, Définition
 \ref{smoothing-definition-strictly-standard}. Il découle de Lissage des morphismes d'anneaux, Lemme  \ref{smoothing-lemma-strictly-standard-base-change}  que si
 $f' : X' \to S'$  est le changement de base de  $f$  par un morphisme
 $S' \to S$, alors  $Z' \subset S' \times_S Z$  où  $Z'$  est le sous-schéma fermé de
 $X'$  défini par des fonctions qui sont affinement localement strictement
standard. Cependant, l'égalité n'est pas claire. La notion d'élément strictement
standard a été utile dans le chapitre sur le théorème de Popescu. Le sous-schéma
fermé défini par ces éléments n'est (autant que nous le sachions) pas utilisé dans
la littérature\footnote{Si  $f$  est un morphisme
d'intersection complète locale (Plus sur les morphismes, Définition  \ref{more-morphisms-definition-lci}),
alors le sous-schéma fermé délimité par les éléments strictement standard locaux
est la chose correcte à examiner.}.

\medskip\noindent
Si  $f$  est plat, de présentation finie, et que les fibres de  $f$ 
sont toutes équidimensionnelles de dimension  $d$, alors l'idéal de
fitting  $d$ de  $\Omega_{X/S}$  est utilisé pour obtenir un bon sous-schéma
fermé. Pour tout morphisme de type fini, les sous-schémas fermés de  $X$ 
définis par les idéaux de fitting de  $\Omega_{X/S}$  définissent une stratification
de  $X$  en termes de rang de  $\Omega_{X/S}$  dont la formation commute avec
le changement de base. Cela peut être utile; c'est lié aux dimensions d'immersion
des fibres, voir Variétés, Section  \ref{varieties-section-embedding-dimension}.

\begin{lemma}
\label{lemma-base-change-and-fitting-ideal-omega}
Soit  $f : X \to S$  un morphisme de schémas qui est localement de type fini. Soit
 $X = Z_{-1} \supset Z_0 \supset Z_1 \supset \ldots$  les sous-schémas fermés définis par les idéaux de fitting de
 $\Omega_{X/S}$. Alors la formation de  $Z_i$  commute avec tout changement de
base.
\end{lemma}

\begin{proof}
Remarquez que  $\Omega_{X/S}$  est un  $\mathcal{O}_X$-module quasi-cohérent de type fini
(Morphismes, Lemme  \ref{morphisms-lemma-finite-type-differentials}) donc les idéaux de Fitting sont définis. Si
 $f' : X' \to S'$  est le changement de base de  $f$  par  $g : S' \to S$, alors
 $\Omega_{X'/S'} = (g')^*\Omega_{X/S}$  où  $g' : X' \to X$  est la projection (Morphismes, Lemme  \ref{morphisms-lemma-base-change-differentials}).
D'où  $(g')^{-1}\text{Fit}_i(\Omega_{X/S}) \cdot \mathcal{O}_{X'} =
\text{Fit}_i(\Omega_{X'/S'})$. Cela signifie que
$$
Z'_i = (g')^{-1}(Z_i) = Z_i \times_X X'
$$
schématiquement et c'est le sens de l'énoncé du lemme.
\end{proof}

\noindent
Le  $0$-ième idéal de Fitting de  $\Omega$  définit le « lieu de
ramification » du morphisme.

\begin{lemma}
\label{lemma-zero-fitting-ideal-omega-unramified}
Soit  $f : X \to S$  un morphisme de schémas qui est localement de type fini. Le
sous-schéma fermé  $Z \subset X$  défini par le  $0$-ième idéal de Fitting de
 $\Omega_{X/S}$  est exactement l'ensemble des points où  $f$  n'est pas non
ramifié.
\end{lemma}

\begin{proof}
D'après le Lemme  \ref{lemma-on-subscheme-cut-out-by-Fit-0} , le complément de  $Z$  est exactement le
lieu où  $\Omega_{X/S}$  est nul. Il s'agit exactement de l'ensemble des points où
 $f$  est non ramifié par les Morphismes, Lemme  \ref{morphisms-lemma-unramified-omega-zero}.
\end{proof}

\begin{lemma}
\label{lemma-d-fitting-ideal-omega-smooth}
Soit  $f : X \to S$  un morphisme de schémas. Soit  $d \geq 0$  un entier. Supposons
\begin{enumerate}
\item que $f$  est plat,
\item que $f$  est localement de présentation finie, et
\item que toute fibre non vide de  $f$  est équidimensionnelle de dimension
 $d$.
\end{enumerate}
Soit  $Z \subset X$  le sous-schéma fermé défini par le nième idéal de Fitting
 $d$ de  $\Omega_{X/S}$. Alors  $Z$  est exactement l'ensemble des
points où  $f$  n'est pas lisse.
\end{lemma}

\begin{proof}
D'après le Lemme  \ref{lemma-locally-free-rank-r-pullback} , le complément de  $Z$  est exactement le
lieu où  $\Omega_{X/S}$  peut être engendré par au plus  $d$  éléments. Par
conséquent, le lemme découle de Morphismes, Lemme  \ref{morphisms-lemma-smooth-at-point}.
\end{proof}







\section{Modules sans torsion}
\label{section-torsion-free}

\noindent
Cette section est l'analogue de Plus sur l'algèbre, Section  \ref{more-algebra-section-torsion-flat}  pour les
modules quasi-cohérents.

\begin{lemma}
\label{lemma-torsion-sections}
Soit  $X$  un schéma intègre avec point générique  $\eta$. Soit
 $\mathcal{F}$  un  $\mathcal{O}_X$-module quasi-cohérent. Soit  $U \subset X$  ouvert non
vide et  $s \in \mathcal{F}(U)$. Les énoncés suivants sont équivalents
\begin{enumerate}
\item pour certains  $x \in U$  l'image de  $s$  dans  $\mathcal{F}_x$  est de
torsion,
\item pour tous  $x \in U$  l'image de  $s$  dans  $\mathcal{F}_x$  est de torsion,
\item l'image de  $s$  dans  $\mathcal{F}_\eta$  est nulle,
\item l'image de  $s$  dans  $j_*\mathcal{F}_\eta$  est nulle, où  $j : \eta \to X$  est le
morphisme d'inclusion.
\end{enumerate}
\end{lemma}

\begin{proof}
Omis.
\end{proof}

\begin{definition}
\label{definition-torsion}
Soit  $X$  un schéma intègre. Soit  $\mathcal{F}$  un  $\mathcal{O}_X$-module
quasi-cohérent.
\begin{enumerate}
\item On dit qu'une section locale de  $\mathcal{F}$  est  {\it 
torsion}  si elle satisfait les conditions équivalentes du Lemme
 \ref{lemma-torsion-sections}.
\item On dit que  $\mathcal{F}$  est  {\it  sans torsion}  si
chaque section à torsion de  $\mathcal{F}$  est  $0$.
\end{enumerate}
\end{definition}

\noindent
Voici le lemme obligatoire comparant cela à la notion algébrique habituelle.

\begin{lemma}
\label{lemma-check-torsion-on-affines}
Soit  $X$  un schéma intègre. Soit  $\mathcal{F}$  un  $\mathcal{O}_X$-module
quasi-cohérent. Les conditions suivantes sont équivalentes
\begin{enumerate}
\item $\mathcal{F}$  est sans torsion,
\item pour  $U \subset X$  ouvert affine,  $\mathcal{F}(U)$  est un  $\mathcal{O}(U)$-module sans
torsion.
\end{enumerate}
\end{lemma}

\begin{proof}
Omis.
\end{proof}

\begin{lemma}
\label{lemma-torsion}
Soit  $X$  un schéma intègre. Soit  $\mathcal{F}$  un  $\mathcal{O}_X$-module
quasi-cohérent. Les sections de torsion de  $\mathcal{F}$  forment un
 $\mathcal{O}_X$-sous-module quasi-cohérent  $\mathcal{F}_{tors} \subset \mathcal{F}$. Le module quotient
 $\mathcal{F}/\mathcal{F}_{tors}$  est sans torsion.
\end{lemma}

\begin{proof}
Omis. Voir Plus sur Algèbre, Lemme  \ref{more-algebra-lemma-torsion}  pour l'analogue algébrique.
\end{proof}

\begin{lemma}
\label{lemma-flat-torsion-free}
Soit  $X$  un schéma intègre. Tout  $\mathcal{O}_X$-module quasi-cohérent plat
est sans torsion.
\end{lemma}

\begin{proof}
Omis. Voir Plus sur Algèbre, Lemme  \ref{more-algebra-lemma-flat-torsion-free}.
\end{proof}

\begin{lemma}
\label{lemma-flat-pullback-torsion}
Soit  $f : X \to Y$  un morphisme plat de schémas intègres. Soit  $\mathcal{G}$  un
 $\mathcal{O}_Y$-module quasi-cohérent sans torsion. Alors  $f^*\mathcal{G}$  est un
 $\mathcal{O}_X$-module sans torsion.
\end{lemma}

\begin{proof}
Omis. Voir Plus sur Algèbre, Lemme  \ref{more-algebra-lemma-flat-pullback-torsion}  pour l'analogue algébrique.
\end{proof}

\begin{lemma}
\label{lemma-flat-over-integral-integral-fibre}
Soit  $f : X \to Y$  un morphisme plat de schémas. Si  $Y$  est intègre et la
fibre générique de  $f$  est intègre, alors  $X$  est intègre.
\end{lemma}

\begin{proof}
L'analogue algébrique est le suivant: soit  $A$  un domaine avec le corps
des fractions  $K$  et soit  $B$  une  $A$-algèbre plate
telle que  $B \otimes_A K$  soit un domaine. Alors  $B$  est un domaine. Ceci
est vrai parce que  $B$  est sans torsion selon Plus sur l’Algèbre, Lemme
 \ref{more-algebra-lemma-flat-torsion-free}  et donc  $B \subset B \otimes_A K$.
\end{proof}

\begin{lemma}
\label{lemma-check-torsion}
Soit  $X$  un schéma intègre. Soit  $\mathcal{F}$  un  $\mathcal{O}_X$-module
quasi-cohérent. Alors  $\mathcal{F}$  est sans torsion si et seulement si
 $\mathcal{F}_x$  est un  $\mathcal{O}_{X, x}$-module sans torsion pour tout  $x \in X$.
\end{lemma}

\begin{proof}
Omis. Voir Plus sur Algèbre, Lemme  \ref{more-algebra-lemma-check-torsion}.
\end{proof}

\begin{lemma}
\label{lemma-extension-torsion-free}
Soit  $X$  un schéma intègre. Soit  $0 \to \mathcal{F} \to \mathcal{F}' \to \mathcal{F}'' \to 0$  une suite exacte courte de
 $\mathcal{O}_X$-modules quasi-cohérents. Si  $\mathcal{F}$  et  $\mathcal{F}''$  sont sans
torsion, alors  $\mathcal{F}'$  est sans torsion.
\end{lemma}

\begin{proof}
Omis. Voir Plus sur Algèbre, Lemme  \ref{more-algebra-lemma-extension-torsion-free}  pour l'analogue algébrique.
\end{proof}

\begin{lemma}
\label{lemma-torsion-free-finite-noetherian-domain}
Soit  $X$  un schéma intégral localement noethérien avec le point générique
 $\eta$. Soit  $\mathcal{F}$  un  $\mathcal{O}_X$-module cohérent non nul. Les
énoncés suivants sont équivalents
\begin{enumerate}
\item $\mathcal{F}$  est sans torsion,
\item $\eta$  est le seul idéal premier associé de  $\mathcal{F}$,
\item $\eta$  est dans le support de  $\mathcal{F}$  et  $\mathcal{F}$  possède la
propriété  $(S_1)$, et
\item $\eta$  est dans le support de  $\mathcal{F}$  et  $\mathcal{F}$  n'a aucun idéal
premier associé incorporé.
\end{enumerate}
\end{lemma}

\begin{proof}
Ceci est une traduction de Plus sur Algèbre, Lemme  \ref{more-algebra-lemma-torsion-free-finite-noetherian-domain}  dans le langage
des schémas. Nous omettons la traduction.
\end{proof}

\begin{lemma}
\label{lemma-torsion-free-over-regular-dim-1}
Soit  $X$  un schéma régulier intègre de dimension  $\leq 1$. Soit
 $\mathcal{F}$  un  $\mathcal{O}_X$-module cohérent. Les énoncés suivants sont
équivalents
\begin{enumerate}
\item $\mathcal{F}$  est sans torsion,
\item $\mathcal{F}$  est localement libre de type fini.
\end{enumerate}
\end{lemma}

\begin{proof}
Il est clair qu'un module localement libre de type fini est sans torsion. Pour la
réciproque, nous montrerons que si  $\mathcal{F}$  est sans torsion, alors
 $\mathcal{F}_x$  est un  $\mathcal{O}_{X, x}$-module libre pour tout  $x \in X$. Cela suffit
d'après Algèbre, Lemme  \ref{algebra-lemma-finite-projective}  et le fait que  $\mathcal{F}$  est cohérent. Si
 $\dim(\mathcal{O}_{X, x}) = 0$, alors  $\mathcal{O}_{X, x}$  est un corps et l'énoncé est clair. Si
 $\dim(\mathcal{O}_{X, x}) = 1$, alors  $\mathcal{O}_{X, x}$  est un anneau à valuation discrète (Algèbre,
Lemme  \ref{algebra-lemma-characterize-dvr}) et  $\mathcal{F}_x$  est sans torsion. Par conséquent,  $\mathcal{F}_x$ 
est libre d'après Plus sur Algèbre, Lemme  \ref{more-algebra-lemma-dedekind-torsion-free-flat}.
\end{proof}

\begin{lemma}
\label{lemma-hom-into-torsion-free}
Soit  $X$  un schéma intègre. Soient  $\mathcal{F}$,  $\mathcal{G}$  des
 $\mathcal{O}_X$-modules quasi-cohérents. Si  $\mathcal{G}$  est sans torsion et
 $\mathcal{F}$  est de présentation finie, alors  $\SheafHom_{\mathcal{O}_X}(\mathcal{F}, \mathcal{G})$  est sans torsion.
\end{lemma}

\begin{proof}
L'énoncé a du sens car  $\SheafHom_{\mathcal{O}_X}(\mathcal{F}, \mathcal{G})$  est quasi-cohérent selon Schémas, Section
 \ref{schemes-section-quasi-coherent}. Pour voir que l'énoncé est vrai, voir Plus en Algèbre, Lemme
 \ref{more-algebra-lemma-hom-into-torsion-free}. Certains détails sont omis.
\end{proof}

\begin{lemma}
\label{lemma-isom-depth-2-torsion-free}
Soit  $X$  un schéma intègre localement noethérien. Soit  $\varphi : \mathcal{F} \to \mathcal{G}$  une
application de  $\mathcal{O}_X$-modules quasi-cohérents. Supposons que  $\mathcal{F}$ 
soit cohérent, que  $\mathcal{G}$  soit sans torsion, et que pour chaque
 $x \in X$ , l'une des situations suivantes se produise
\begin{enumerate}
\item $\mathcal{F}_x \to \mathcal{G}_x$  est un isomorphisme, ou
\item $\text{depth}(\mathcal{F}_x) \geq 2$.
\end{enumerate}
Alors  $\varphi$  est un isomorphisme.
\end{lemma}

\begin{proof}
Ceci est une traduction de Plus en Algèbre, Lemme  \ref{more-algebra-lemma-isom-depth-2-torsion-free}  dans le langage des
schémas.
\end{proof}








\section{Rangs des modules}
\label{section-rank}

\noindent
Cette section est l'analogue de Plus sur l'Algèbre, Section  \ref{more-algebra-section-rank}  pour les
modules quasi-cohérents. Rappelons que si  $X$  est un schéma intègre avec
le point générique  $\eta$, alors l'anneau local  $\mathcal{O}_{X, \eta}$  est égal au
corps résiduel  $\kappa(\eta)$, voir Morphismes, Lemme  \ref{morphisms-lemma-integral-scheme-rational-functions}.

\begin{definition}
\label{definition-rank}
Laisse $X$ soit un schéma intègre avec un point générique $\eta$.
Laisser  $\mathcal{F}$ être un quasi-cohérent $\mathcal{O}_X$-module. Le
 {\it rang}  de  $\mathcal{F}$  est  $\dim_{\kappa(\eta)} \mathcal{F}_\eta$.
\end{definition}

\noindent
Cela ne contredit pas la définition du rang d'un module localement libre (dans le
cas où les deux définitions s'appliquent). Il ne semble pas judicieux de définir
le rang d'un module quasi-cohérent si le schéma  $X$  n'est pas
irréductible ou si l'anneau local  $\mathcal{O}_{X, \eta}$  au point générique n'est pas un
corps\footnote{S'il  $X$  est irréductible et  $\mathcal{F}_\eta$ 
est libre sur  $\mathcal{O}_{X, \eta}$ , il semblerait naturel d'utiliser le rang de ce module
libre. Attention, certaines références définissent le rang en utilisant la
longueur de  $\mathcal{F}_\eta$.}

\begin{lemma}
\label{lemma-check-rank-on-affines}
Soit  $X$  un schéma intègre. Soit  $\mathcal{F}$  un  $\mathcal{O}_X$-module
quasi-cohérent. Soit  $r$  un cardinal. Les énoncés suivants sont
équivalents
\begin{enumerate}
\item $\mathcal{F}$  a pour rang  $r$,
\item pour tout ouvert affine non vide  $U \subset X$ ,  $\mathcal{F}(U)$  est un
 $\mathcal{O}_X(U)$-module de rang  $r$, et
\item pour un certain  $U \subset X$  ouvert affine non vide  $\mathcal{F}(U)$  est un
 $\mathcal{O}_X(U)$-module de rang  $r$.
\end{enumerate}
\end{lemma}

\begin{proof}
Omis.
\end{proof}

\begin{lemma}
\label{lemma-dominant-pullback-rank}
Soit  $f : X \to Y$  un morphisme dominant de schémas intègres. Soit  $\mathcal{G}$  un
 $\mathcal{O}_Y$-module quasi-cohérent. Alors le rang de  $\mathcal{G}$  sur  $Y$ 
est égal au rang de  $f^*\mathcal{G}$  sur  $X$.
\end{lemma}

\begin{proof}
Utilisez le Lemme  \ref{lemma-check-rank-on-affines}, Morphismes, Lemme  \ref{morphisms-lemma-dominant-between-integral}, et Plus sur
Algèbre, Lemme  \ref{more-algebra-lemma-dominant-pullback-rank}.
\end{proof}

\begin{lemma}
\label{lemma-rank-additive}
Soit  $X$  un schéma intègre. Soit  $0 \to \mathcal{F} \to \mathcal{F}' \to \mathcal{F}'' \to 0$  une suite exacte courte de
 $\mathcal{O}_X$-modules quasi-cohérents. Alors le rang de  $\mathcal{F}'$  est la somme
des rangs de  $\mathcal{F}$  et  $\mathcal{F}''$.
\end{lemma}

\begin{proof}
Omis. Voir Plus sur Algèbre, Lemme  \ref{more-algebra-lemma-rank-additive}.
\end{proof}

\begin{lemma}
\label{lemma-rank-tensor}
Soit  $X$  un schéma intègre. Soient  $\mathcal{F}$  et  $\mathcal{G}$  des
 $\mathcal{O}_X$-modules quasi-cohérents.
\begin{enumerate}
\item Le rang de  $\mathcal{F} \otimes_{\mathcal{O}_X} \mathcal{G}$  est le produit des rangs de  $\mathcal{F}$  et  $\mathcal{G}$.
\item Si  $\mathcal{F}$  est de présentation finie, alors le rang de  $\SheafHom_{\mathcal{O}_X}(\mathcal{F}, \mathcal{G})$  est le
produit des rangs de  $\mathcal{F}$  et  $\mathcal{G}$.
\end{enumerate}
\end{lemma}

\begin{proof}
Omis. Notez que la partie (2) a du sens car  $\SheafHom_{\mathcal{O}_X}(\mathcal{F}, \mathcal{G})$  est quasi-cohérent selon
l'article (10) dans Schémas, Section  \ref{schemes-section-quasi-coherent}. La version algébrique de ce
lemme est Plus sur l’algèbre, Lemme  \ref{more-algebra-lemma-rank-tensor}.
\end{proof}

\begin{lemma}
\label{lemma-finite-module-rank}
Soit  $X$  un schéma intègre. Soit  $\mathcal{F}$  un  $\mathcal{O}_X$-module
quasi-cohérent de type fini. Alors
\begin{enumerate}
\item $\mathcal{F}$  a un rang fini  $r \geq 0$, et
\item il existe un ouvert non vide  $U \subset X$  tel que  $\mathcal{F}|_U$  soit libre de rang
 $r$.
\end{enumerate}
\end{lemma}

\begin{proof}
Voir Plus sur l’algèbre, Lemme  \ref{more-algebra-lemma-finite-module-rank}  pour la version algébrique.
\end{proof}











\section{Modules réflexifs}
\label{section-reflexive}

\noindent
Cette section est l'analogue de Plus sur l'Algèbre, Section  \ref{more-algebra-section-reflexive}  pour les
modules cohérents sur les schémas localement noethériens. La raison de travailler
avec des modules cohérents est que  $\SheafHom_{\mathcal{O}_X}(\mathcal{F}, \mathcal{G})$  est cohérent pour chaque paire de
 $\mathcal{O}_X$-modules cohérents  $\mathcal{F}, \mathcal{G}$, voir Modules, Lemme  \ref{modules-lemma-internal-hom-locally-kernel-direct-sum}.

\begin{definition}
\label{definition-reflexive}
Soit  $X$  un schéma intègre localement noethérien. Soit  $\mathcal{F}$  un
 $\mathcal{O}_X$-module cohérent. Le  {\it  enveloppe
réflexive}  de  $\mathcal{F}$  est le  $\mathcal{O}_X$-module
$$
\mathcal{F}^{**} = \SheafHom_{\mathcal{O}_X}(
\SheafHom_{\mathcal{O}_X}(\mathcal{F}, \mathcal{O}_X), \mathcal{O}_X)
$$
Nous disons que  $\mathcal{F}$  est  {\it  réflexif}  si
l'application naturelle  $j : \mathcal{F} \longrightarrow \mathcal{F}^{**}$  est un isomorphisme.
\end{definition}

\noindent
Il découle du Lemme  \ref{lemma-dual-reflexive}  que l'adhérence réflexive est un
 $\mathcal{O}_X$-module réflexif. Vous pouvez utiliser la même définition pour définir
des modules réflexifs dans des situations plus générales, mais cela ne semble pas
très utile. Voici le lemme obligatoire comparant cela à la notion algébrique
habituelle.

\begin{lemma}
\label{lemma-check-reflexive-on-affines}
Soit  $X$  un schéma intègre localement noethérien. Soit  $\mathcal{F}$  un
 $\mathcal{O}_X$-module cohérent. Les énoncés suivants sont équivalents
\begin{enumerate}
\item $\mathcal{F}$  est réflexif,
\item pour  $U \subset X$  ouvert affine,  $\mathcal{F}(U)$  est un  $\mathcal{O}(U)$-module réflexif.
\end{enumerate}
\end{lemma}

\begin{proof}
Omis.
\end{proof}

\begin{remark}
\label{remark-different-reflexive}
Si  $X$  est un schéma de type fini sur un corps, alors parfois une notion
différente de modules réflexifs est utilisée (voir par exemple  \cite[bottom of page 5 and Definition 1.1.9]{HL}).
Cette autre notion utilise  $R\SheafHom$  dans un complexe dualisant  $\omega_X^\bullet$  au
lieu de dans  $\mathcal{O}_X$  et devrait probablement avoir un nom différent car elle
peut être différente lorsque  $X$  n'est pas Gorenstein. Par exemple, si
 $X = \Spec(k[t^3, t^4, t^5])$, alors un calcul montre que le fibré dualisant  $\omega_X$  n'est pas
réflexif selon notre définition, mais il est réflexif selon l'autre définition car
 $\omega_X \to \SheafHom(\SheafHom(\omega_X, \omega_X), \omega_X)$  est un isomorphisme.
\end{remark}

\begin{lemma}
\label{lemma-reflexive-torsion-free}
Soit  $X$  un schéma intègre localement noethérien. Soit  $\mathcal{F}$  un
 $\mathcal{O}_X$-module cohérent.
\begin{enumerate}
\item Si  $\mathcal{F}$  est réflexif, alors  $\mathcal{F}$  est sans torsion.
\item L'application  $j : \mathcal{F} \longrightarrow \mathcal{F}^{**}$  est injective si et seulement si  $\mathcal{F}$  est sans
torsion.
\end{enumerate}
\end{lemma}

\begin{proof}
Omis. Voir Plus sur Algèbre, Lemme  \ref{more-algebra-lemma-reflexive-torsion-free}.
\end{proof}

\begin{lemma}
\label{lemma-check-reflexive}
Soit  $X$  un schéma intègre localement noethérien. Soit  $\mathcal{F}$  un
 $\mathcal{O}_X$-module cohérent. Les énoncés suivants sont équivalents
\begin{enumerate}
\item $\mathcal{F}$  est réflexif,
\item $\mathcal{F}_x$  est un  $\mathcal{O}_{X, x}$-module réflexif pour tout  $x \in X$,
\item $\mathcal{F}_x$  est un  $\mathcal{O}_{X, x}$-module réflexif pour tous les points fermés
 $x \in X$.
\end{enumerate}
\end{lemma}

\begin{proof}
Par modules, le lemme  \ref{modules-lemma-stalk-internal-hom}  on voit que (1) et (2) sont équivalents.
Puisque chaque point de  $X$  se spécialise vers un point fermé
(Propriétés, lemme  \ref{properties-lemma-locally-Noetherian-closed-point}), on voit que (2) et (3) sont équivalents.
\end{proof}

\begin{lemma}
\label{lemma-flat-pullback-reflexive}
Soit  $f : X \to Y$  un morphisme plat de schémas localement noethériens intégraux.
Soit  $\mathcal{G}$  un  $\mathcal{O}_Y$-module réflexif cohérent. Alors  $f^*\mathcal{G}$  est
un  $\mathcal{O}_X$-module réflexif cohérent.
\end{lemma}

\begin{proof}
Omis. Voir Plus sur Algèbre, Lemme  \ref{more-algebra-lemma-flat-pullback-reflexive}  pour l'analogue algébrique.
\end{proof}

\begin{lemma}
\label{lemma-sequence-reflexive}
Soit  $X$  un schéma intègre localement noethérien. Soit  $0 \to \mathcal{F} \to \mathcal{F}' \to \mathcal{F}''$  une
suite exacte de  $\mathcal{O}_X$-modules cohérents. Si  $\mathcal{F}'$  est réflexif et
 $\mathcal{F}''$  est sans torsion, alors  $\mathcal{F}$  est réflexif.
\end{lemma}

\begin{proof}
Omis. Voir Plus sur Algèbre, Lemme  \ref{more-algebra-lemma-sequence-reflexive}.
\end{proof}

\begin{lemma}
\label{lemma-dual-reflexive}
Soit  $X$  un schéma intègre localement noethérien. Soit  $\mathcal{F}$,
 $\mathcal{G}$  des  $\mathcal{O}_X$-modules cohérents. Si  $\mathcal{G}$  est réflexif,
alors  $\SheafHom_{\mathcal{O}_X}(\mathcal{F}, \mathcal{G})$  est réflexif.
\end{lemma}

\begin{proof}
L'énoncé a un sens car  $\SheafHom_{\mathcal{O}_X}(\mathcal{F}, \mathcal{G})$  est cohérent selon Cohomologie des Schémas,
Lemme  \ref{coherent-lemma-tensor-hom-coherent}. Pour voir que l'énoncé est vrai, voir Plus sur Algèbre, Lemme
 \ref{more-algebra-lemma-dual-reflexive}. Certains détails sont omis.
\end{proof}

\begin{remark}
\label{remark-tensor}
Soit  $X$  un schéma intègre localement noethérien. Grâce au Lemme
 \ref{lemma-dual-reflexive} , nous savons que le enveloppe réflexive  $\mathcal{F}^{**}$  d’un
 $\mathcal{O}_X$-module cohérent est cohérent réflexif. Considérons la catégorie
 $\mathcal{C}$  des  $\mathcal{O}_X$-modules cohérents réflexifs. La prise des enveloppes
réflexives donne un adjoint à gauche du foncteur d’inclusion  $\mathcal{C} \to \textit{Coh}(\mathcal{O}_X)$.
Remarquons que  $\mathcal{C}$  est une catégorie additive avec des noyaux et des
conoyaux. En effet, étant donné  $\varphi : \mathcal{F} \to \mathcal{G}$  dans  $\mathcal{C}$, le noyau usuel
 $\Ker(\varphi)$  est réflexif (Lemme  \ref{lemma-sequence-reflexive}) et l’enveloppe réflexive
 $\Coker(\varphi)^{**}$  du conoyau usuel est le conoyau dans  $\mathcal{C}$. De plus,
 $\mathcal{C}$  hérite d’un produit tensoriel
$$
\mathcal{F} \otimes_\mathcal{C} \mathcal{G} =
(\mathcal{F} \otimes_{\mathcal{O}_X} \mathcal{G})^{**}
$$
qui est associatif et symétrique. Il existe un Hom interne au sens où pour trois
objets  $\mathcal{F}, \mathcal{G}, \mathcal{H}$  de  $\mathcal{C}$ , nous avons l’identité
$$
\Hom_\mathcal{C}(\mathcal{F} \otimes_\mathcal{C} \mathcal{G}, \mathcal{H}) =
\Hom_\mathcal{C}(\mathcal{F},
\SheafHom_{\mathcal{O}_X}(\mathcal{G}, \mathcal{H}))
$$
voir Modules, Lemme  \ref{modules-lemma-internal-hom}. Dans  $\mathcal{C}$ , chaque objet  $\mathcal{F}$  a
un objet  {\it  dual }   $\SheafHom_{\mathcal{O}_X}(\mathcal{F}, \mathcal{O}_X)$. Sans conditions
supplémentaires sur  $X$ , il peut arriver que
$$
\SheafHom_{\mathcal{O}_X}(\mathcal{F}, \mathcal{G}) \not \cong
\SheafHom_{\mathcal{O}_X}(\mathcal{F}, \mathcal{O}_X)
\otimes_\mathcal{C} \mathcal{G}
\quad\text{and}\quad
\mathcal{F} \otimes_\mathcal{C}
\SheafHom_{\mathcal{O}_X}(\mathcal{F}, \mathcal{O}_X)
\not \cong \mathcal{O}_X
$$
pour  $\mathcal{F}, \mathcal{G}$  de rang  $1$  dans  $\mathcal{C}$. Pour donner un exemple,
soient  $X = \Spec(R)$  où  $R$  est comme dans Plus sur l'Algèbre, Exemple
 \ref{more-algebra-example-ring-not-S2}  et soient  $\mathcal{F}, \mathcal{G}$  les modules correspondant à  $M$.
Calcul omis.
\end{remark}

\begin{lemma}
\label{lemma-reflexive-depth-2}
Soit  $X$  un schéma intègre localement noethérien. Soit  $\mathcal{F}$  un
 $\mathcal{O}_X$-module cohérent. Les énoncés suivants sont équivalents
\begin{enumerate}
\item $\mathcal{F}$  est réflexif,
\item pour chaque  $x \in X$  l'un des cas suivants se produit
\begin{enumerate}
\item $\mathcal{F}_x$  est un  $\mathcal{O}_{X, x}$-module réflexif, ou
\item $\text{depth}(\mathcal{F}_x) \geq 2$.
\end{enumerate}
\end{enumerate}
\end{lemma}

\begin{proof}
Omis. Voir Plus sur Algèbre, Lemme  \ref{more-algebra-lemma-reflexive-depth-2}.
\end{proof}

\begin{lemma}
\label{lemma-reflexive-S2}
Soit  $X$  un schéma intègre localement noethérien. Soit  $\mathcal{F}$  un
 $\mathcal{O}_X$-module cohérent réflexif. Soit  $x \in X$.
\begin{enumerate}
\item Si  $\text{depth}(\mathcal{O}_{X, x}) \geq 2$, alors  $\text{depth}(\mathcal{F}_x) \geq 2$.
\item Si  $X$  est  $(S_2)$, alors  $\mathcal{F}$  est  $(S_2)$.
\end{enumerate}
\end{lemma}

\begin{proof}
Omis. Voir Plus sur Algèbre, Lemme  \ref{more-algebra-lemma-reflexive-S2}.
\end{proof}

\begin{lemma}
\label{lemma-reflexive-S2-extend}
Soit  $X$  un schéma intègre localement noethérien. Soit  $j : U \to X$  un
sous-schéma ouvert avec complémentaire  $Z$. Supposons que  $\mathcal{O}_{X, z}$ 
ait une profondeur  $\geq 2$  pour tous  $z \in Z$. Alors  $j^*$  et
 $j_*$  définissent une équivalence de catégories entre la catégorie des
 $\mathcal{O}_X$-modules cohérents réflexifs et la catégorie des  $\mathcal{O}_U$-modules
cohérents réflexifs.
\end{lemma}

\begin{proof}
Soit  $\mathcal{F}$  un  $\mathcal{O}_X$-module cohérent réflexif. Pour  $z \in Z$ , la
germe  $\mathcal{F}_z$  a une profondeur  $\geq 2$  d'après le Lemme  \ref{lemma-reflexive-S2}.
Ainsi,  $\mathcal{F} \to j_*j^*\mathcal{F}$  est un isomorphisme d'après le Lemme  \ref{lemma-depth-2-hartog}.
Inversement, soit  $\mathcal{G}$  un  $\mathcal{O}_U$-module cohérent réflexif. Il suffit
de montrer que  $j_*\mathcal{G}$  est un  $\mathcal{O}_X$-module cohérent réflexif. Pour le
prouver, nous pouvons supposer que  $X$  est affine. D'après les
Propriétés, Lemme  \ref{properties-lemma-lift-finite-presentation}  il existe un  $\mathcal{O}_X$-module cohérent
 $\mathcal{F}$  avec  $\mathcal{G} = j^*\mathcal{F}$. Après avoir remplacé  $\mathcal{F}$  par son
enveloppe réflexive, nous pouvons supposer que  $\mathcal{F}$  est réflexif (voir la
discussion ci-dessus et en particulier le Lemme  \ref{lemma-dual-reflexive}). D'après ce qui
précède,  $j_*\mathcal{G} = j_*j^*\mathcal{F} = \mathcal{F}$  comme désiré.
\end{proof}

\noindent
Si le schéma est normal, alors réflexif est la même chose que sans torsion et
 $(S_2)$.

\begin{lemma}
\label{lemma-reflexive-over-normal}
Soit  $X$  un schéma normal intègre localement noethérien. Soit
 $\mathcal{F}$  un  $\mathcal{O}_X$-modèle cohérent. Les énoncés suivants sont
équivalents:
\begin{enumerate}
\item $\mathcal{F}$  est réflexif,
\item $\mathcal{F}$  est sans torsion et possède la propriété  $(S_2)$, et
\item il existe un sous-schéma ouvert  $j : U \to X$  tel que
\begin{enumerate}
\item chaque composante irréductible de  $X \setminus U$  a une codimension  $\geq 2$  dans
 $X$.
\item $j^*\mathcal{F}$  est de type fini localement libre, et
\item $\mathcal{F} = j_*j^*\mathcal{F}$.
\end{enumerate}
\end{enumerate}
\end{lemma}

\begin{proof}
En utilisant le Lemme  \ref{lemma-check-reflexive-on-affines} , l'équivalence de (1) et (2) découle de More on
Algebra, Lemme  \ref{more-algebra-lemma-reflexive-over-normal}. Soit  $U \subset X$  comme en (3). Par Propriétés, Lemme
 \ref{properties-lemma-criterion-normal}  nous voyons que  $\text{depth}(\mathcal{O}_{X, x}) \geq 2$  pour  $x \not \in U$. Comme un module de
type fini localement libre est réflexif, nous en concluons que (3) implique (1)
par le Lemme  \ref{lemma-reflexive-S2-extend}.

\medskip\noindent
Supposons (1). Soit  $U \subset X$  le sous-schéma ouvert maximal tel que
 $j^*\mathcal{F} = \mathcal{F}|_U$  soit localement libre de type fini. Donc (3)(b) est satisfait. Soit
 $x \in X$  un point. Si  $\mathcal{F}_x$  est un  $\mathcal{O}_{X, x}$-module libre, alors
 $x \in U$, voir Modules, Lemme  \ref{modules-lemma-finite-presentation-stalk-free}. Si  $\dim(\mathcal{O}_{X, x}) \leq 1$, alors
 $\mathcal{O}_{X, x}$  est soit un corps, soit un anneau de valuation discrète (Propriétés,
Lemme  \ref{properties-lemma-criterion-normal}) et donc  $\mathcal{F}_x$  est libre (Plus sur Algèbre, Lemme
 \ref{more-algebra-lemma-dedekind-torsion-free-flat}). Ainsi  $x \not \in U \Rightarrow \dim(\mathcal{O}_{X, x}) \geq 2$. Ensuite, Propriétés, Lemme  \ref{properties-lemma-codimension-local-ring}  montre
que (3)(a) est satisfait. Par le Lemme déjà utilisé des Propriétés,  \ref{properties-lemma-criterion-normal} ,
nous voyons également que  $\text{depth}(\mathcal{O}_{X, x}) \geq 2$  pour  $x \not \in U$  et donc (3)(c) découle du
Lemme  \ref{lemma-reflexive-S2-extend}.
\end{proof}

\begin{lemma}
\label{lemma-describe-reflexive-hull}
Soit  $X$  un schéma normal localement noethérien intègre de point
générique  $\eta$. Soient  $\mathcal{F}$,  $\mathcal{G}$  des  $\mathcal{O}_X$-modules
cohérents. Soit  $T : \mathcal{G}_\eta \to \mathcal{F}_\eta$  une application linéaire. Alors  $T$  se
prolonge en une application  $\mathcal{G} \to \mathcal{F}^{**}$  de  $\mathcal{O}_X$-modules si et seulement
si
\begin{itemize}
\item[(*)] pour tout  $x \in X$  avec  $\dim(\mathcal{O}_{X, x}) = 1$  nous avons
$$
T\left(\Im(\mathcal{G}_x \to \mathcal{G}_\eta)\right) \subset
\Im(\mathcal{F}_x \to \mathcal{F}_\eta).
$$
\end{itemize}
\end{lemma}

\begin{proof}
. Comme  $\mathcal{F}^{**}$  est sans torsion et que  $\mathcal{F}_\eta = \mathcal{F}^{**}_\eta$  un prolongement, s'il
existe, est unique. Ainsi, il suffit de prouver le lemme sur les éléments d'un
recouvrement ouvert de  $X$, c'est-à-dire que nous pouvons supposer que
 $X$  est affine. Dans ce cas, nous posons la question algébrique suivante:
Soit  $R$  un anneau normal noethérien de corps des fractions  $K$,
soient  $M$,  $N$  des  $R$-modules finis, soit
 $T : M \otimes_R K \to N \otimes_R K$  une application  $K$-linéaire. Quand  $T$  se
prolonge-t-elle en une application  $N \to M^{**}$? Selon More sur Algèbre, Lemme
 \ref{more-algebra-lemma-describe-reflexive-hull}  cela se produit si et seulement si  $N_\mathfrak p$  se transforme en
 $(M/M_{tors})_\mathfrak p$  pour chaque hauteur  $1$  premier  $\mathfrak p$  de
 $R$. C'est exactement la condition  $(*)$  du lemme.
\end{proof}

\begin{lemma}
\label{lemma-reflexive-over-regular-dim-2}
Soit  $X$  un schéma régulier de dimension  $\leq 2$. Soit  $\mathcal{F}$ 
un  $\mathcal{O}_X$-module cohérent. Les affirmations suivantes sont équivalentes
\begin{enumerate}
\item $\mathcal{F}$  est réflexif,
\item $\mathcal{F}$  est localement libre de type fini.
\end{enumerate}
\end{lemma}

\begin{proof}
Il est clair qu'un module fini localement libre est réflexif. Pour la réciproque,
nous montrerons que si  $\mathcal{F}$  est réflexif, alors  $\mathcal{F}_x$  est un
 $\mathcal{O}_{X, x}$-module libre pour tout  $x \in X$. Cela suffit d'après Algèbre,
Lemme  \ref{algebra-lemma-finite-projective}  et le fait que  $\mathcal{F}$  est cohérent. Si  $\dim(\mathcal{O}_{X, x}) = 0$,
alors  $\mathcal{O}_{X, x}$  est un corps et l'énoncé est clair. Si  $\dim(\mathcal{O}_{X, x}) = 1$, alors
 $\mathcal{O}_{X, x}$  est un anneau de valuation discrète (Algèbre, Lemme  \ref{algebra-lemma-characterize-dvr}) et
 $\mathcal{F}_x$  est sans torsion. Par conséquent,  $\mathcal{F}_x$  est libre d'après
Encore sur Algèbre, Lemme  \ref{more-algebra-lemma-dedekind-torsion-free-flat}. Si  $\dim(\mathcal{O}_{X, x}) = 2$, alors  $\mathcal{O}_{X, x}$  est un
anneau local régulier de dimension  $2$. D'après Encore sur Algèbre, Lemme
 \ref{more-algebra-lemma-reflexive-over-normal}  nous voyons que  $\mathcal{F}_x$  a une profondeur  $\geq 2$. Par
conséquent,  $\mathcal{F}_x$  est libre d'après Algèbre, Lemme  \ref{algebra-lemma-regular-mcm-free}.
\end{proof}







\section{Diviseurs de Cartier effectifs}
\label{section-effective-Cartier-divisors}

\noindent
Nous définissons la notion de diviseur de Cartier effectif avant tout autre type
de diviseur.

\begin{definition}
\label{definition-effective-Cartier-divisor}
Soit  $S$  un schéma.
\begin{enumerate}
\item une {\it sous-schéma fermé localement principal}  de
 $S$ est un sous-schéma fermé dont le faisceau d'idéaux est localement
engendré par un seul élément.
\item une {\it diviseur de Cartier effectif} sur $S$ 
est un sous-schéma fermé  $D \subset S$ dont le faisceau idéal $\mathcal{I}_D \subset \mathcal{O}_S$  est
inversible  $\mathcal{O}_S$-module.
\end{enumerate}
\end{definition}

\noindent
Ainsi, un diviseur de Cartier effectif est un sous-schéma fermé localement
principal, mais la réciproque n'est pas toujours vraie. Les diviseurs de Cartier
effectifs sont des sous-schémas fermés de codimension pure  $1$  dans le
sens le plus fort possible. En effet, ils sont localement définis par un seul
élément qui est un non-diviseur de zéro. En particulier, ils sont partout denses
nulle part.

\begin{lemma}
\label{lemma-characterize-effective-Cartier-divisor}
Soit  $S$  un schéma. Soit  $D \subset S$  un sous-schéma fermé. Les
assertions suivantes sont équivalentes:
\begin{enumerate}
\item Le sous-schéma  $D$  est un diviseur de Cartier effectif sur  $S$.
\item Pour chaque  $x \in D$ , il existe un voisinage ouvert affine  $\Spec(A) = U \subset S$  de
 $x$  tel que  $U \cap D = \Spec(A/(f))$  avec  $f \in A$  un non-diviseur de zéro.
\end{enumerate}
\end{lemma}

\begin{proof}
Supposons (1). Pour chaque  $x \in D$ , il existe un voisinage ouvert affine
 $\Spec(A) = U \subset S$  de  $x$  tel que  $\mathcal{I}_D|_U \cong \mathcal{O}_U$. En d'autres termes, il existe
une section  $f \in \Gamma(U, \mathcal{I}_D)$  qui engendre librement la restriction  $\mathcal{I}_D|_U$. Ainsi
 $f \in A$, et l'application de multiplication  $f : A \to A$  est injective. De
plus, comme  $\mathcal{I}_D$  est quasi-cohérent, nous voyons que  $D \cap U = \Spec(A/(f))$.

\medskip\noindent
Supposons (2). Soit  $x \in D$. D'après l'hypothèse, il existe un voisinage
ouvert affine  $\Spec(A) = U \subset S$  de  $x$  tel que  $U \cap D = \Spec(A/(f))$  avec  $f \in A$ 
un non-diviseur de zéro. Alors  $\mathcal{I}_D|_U \cong \mathcal{O}_U$  puisque c'est égal à  $\widetilde{(f)} \cong \widetilde{A} \cong \mathcal{O}_U$.
Bien sûr,  $\mathcal{I}_D$  restreint au sous-schéma ouvert  $S \setminus D$  est isomorphe
à  $\mathcal{O}_{S \setminus D}$. Ainsi  $\mathcal{I}_D$  est un  $\mathcal{O}_S$-module inversible.
\end{proof}

\begin{lemma}
\label{lemma-complement-locally-principal-closed-subscheme}
Soit  $S$  un schéma. Soit  $Z \subset S$  un sous-schéma fermé localement
principal. Soit  $U = S \setminus Z$. Alors  $U \to S$  est un morphisme affine.
\end{lemma}

\begin{proof}
La question est locale sur  $S$, voir Morphismes, Lemme  \ref{morphisms-lemma-characterize-affine}.
Ainsi nous pouvons supposer  $S = \Spec(A)$  et  $Z = V(f)$  pour un certain
 $f \in A$. Dans ce cas,  $U = D(f) = \Spec(A_f)$  est affine donc  $U \to S$  est affine.
\end{proof}

\begin{lemma}
\label{lemma-complement-effective-Cartier-divisor}
Soit  $S$  un schéma. Soit  $D \subset S$  un diviseur de Cartier effectif.
Soit  $U = S \setminus D$. Alors  $U \to S$  est un morphisme affine et  $U$  est
schématiquement dense dans  $S$.
\end{lemma}

\begin{proof}
L'affinité est le Lemme  \ref{lemma-complement-locally-principal-closed-subscheme}. La question de la densité est locale sur
 $S$, voir Morphismes, Lemme  \ref{morphisms-lemma-characterize-scheme-theoretically-dense}. Ainsi nous pouvons supposer
 $S = \Spec(A)$  et  $D$  correspondant au non-diviseur de zéro  $f \in A$,
voir Lemme  \ref{lemma-characterize-effective-Cartier-divisor}. Ainsi  $A \subset A_f$  ce qui implique que  $U \subset S$  est
schématiquement dense, voir Morphismes, Exemple  \ref{morphisms-example-scheme-theoretic-closure}.
\end{proof}

\begin{lemma}
\label{lemma-effective-Cartier-makes-dimension-drop}
Soit  $S$  un schéma. Soit  $D \subset S$  un diviseur de Cartier effectif.
Soit  $s \in D$. Si  $\dim_s(S) < \infty$, alors  $\dim_s(D) < \dim_s(S)$.
\end{lemma}

\begin{proof}
Supposons  $\dim_s(S) < \infty$. Soit  $U = \Spec(A) \subset S$  un voisinage ouvert affine de
 $s$  tel que  $\dim(U) = \dim_s(S)$  et tel que  $D = V(f)$  pour un certain non-zéro
diviseur  $f \in A$  (voir le Lemme  \ref{lemma-characterize-effective-Cartier-divisor}). Rappelons que  $\dim(U)$  est
la dimension de Krull de l'anneau  $A$  et que  $\dim(U \cap D)$  est la
dimension de Krull de l'anneau  $A/(f)$. Alors  $f$  n'est contenu dans
aucun idéal premier minimal de  $A$. Ainsi, toute chaîne maximale de
premiers dans  $A/(f)$, considérée comme une chaîne de premiers dans
 $A$, peut être étendue en ajoutant un idéal premier minimal.
\end{proof}

\begin{definition}
\label{definition-sum-effective-Cartier-divisors}
Soit  $S$  un schéma. Étant donné des diviseurs de Cartier effectifs
 $D_1$,  $D_2$  sur  $S$ , nous posons  $D = D_1 + D_2$  égal au
sous-schéma fermé de  $S$  correspondant au faisceau quasi-cohérent
d'idéaux  $\mathcal{I}_{D_1}\mathcal{I}_{D_2} \subset \mathcal{O}_S$. Nous appelons ceci la somme  {\it  des
diviseurs de Cartier effectifs  $D_1$  et  $D_2$}.
\end{definition}

\noindent
Il est clair que nous pouvons définir la somme  $\sum n_iD_i$  donnée par un nombre
fini de diviseurs de Cartier effectifs  $D_i$  sur  $X$  et des
entiers non négatifs  $n_i$.

\begin{lemma}
\label{lemma-sum-effective-Cartier-divisors}
La somme de deux diviseurs de Cartier effectifs est un diviseur de Cartier
effectif.
\end{lemma}

\begin{proof}
Omit. Localement  $f_1, f_2 \in A$  sont des non-zéros-diviseurs, alors aussi
 $f_1f_2 \in A$  est un non-zéros-diviseur.
\end{proof}

\begin{lemma}
\label{lemma-difference-effective-Cartier-divisors}
Soit  $X$  un schéma. Soient  $D, D'$  deux diviseurs de Cartier
effectifs sur  $X$. Si  $D \subset D'$  (en tant que sous-schémas fermés de
 $X$), alors il existe un diviseur de Cartier effectif  $D''$  tel
que  $D' = D + D''$.
\end{lemma}

\begin{proof}
Omis.
\end{proof}

\begin{lemma}
\label{lemma-sum-closed-subschemes-effective-Cartier}
Soit  $X$  un schéma. Soient  $Z, Y$  deux sous-schémas fermés de
 $X$  avec des faisceaux d'idéaux  $\mathcal{I}$  et  $\mathcal{J}$. Si
 $\mathcal{I}\mathcal{J}$  définit un diviseur de Cartier effectif  $D \subset X$, alors
 $Z$  et  $Y$  sont des diviseurs de Cartier effectifs et
 $D = Z + Y$.
\end{lemma}

\begin{proof}
En appliquant le Lemme  \ref{lemma-characterize-effective-Cartier-divisor} , nous obtenons la situation algébrique
suivante:  $A$  est un anneau,  $I, J \subset A$  des idéaux et  $f \in A$  un
élément non-diviseur de zéro tel que  $IJ = (f)$. Ainsi, le résultat découle de
Algèbre, Lemme  \ref{algebra-lemma-product-ideals-principal}.
\end{proof}

\begin{lemma}
\label{lemma-sum-effective-Cartier-divisors-union}
Soit  $X$  un schéma. Soient  $D, D' \subset X$  des diviseurs de Cartier
effectifs tels que l'intersection schématique  $D \cap D'$  soit un diviseur de
Cartier effectif sur  $D'$. Alors  $D + D'$  est l'union schématique de
 $D$  et  $D'$.
\end{lemma}

\begin{proof}
Voir Morphismes, Définition  \ref{morphisms-definition-scheme-theoretic-intersection-union}  pour la définition de l'intersection et
de l'union schématiques. Pour prouver le lemme en travaillant localement (en
utilisant le Lemme  \ref{lemma-characterize-effective-Cartier-divisor}) nous obtenons le problème algébrique suivant:
Étant donné un anneau  $A$  et des non-diviseurs de zéro  $f_1, f_2 \in A$  tels
que  $f_1$  s'envoie sur un non-diviseur de zéro dans  $A/f_2A$, montrer
que  $f_1A \cap f_2A = f_1f_2A$. Nous omettons l'argument direct.
\end{proof}

\noindent
Rappelons que nous avons défini l'image réciproque d'un sous-schéma fermé sous
tout morphisme de schémas dans Schémas, Définition  \ref{schemes-definition-inverse-image-closed-subscheme}.

\begin{lemma}
\label{lemma-pullback-locally-principal}
Soit  $f : S' \to S$  un morphisme de schémas. Soit  $Z \subset S$  un sous-schéma fermé
localement principal. Alors l'image inverse  $f^{-1}(Z)$  est un sous-schéma fermé
localement principal de  $S'$.
\end{lemma}

\begin{proof}
Omis.
\end{proof}

\begin{definition}
\label{definition-pullback-effective-Cartier-divisor}
Soit  $f : S' \to S$  un morphisme de schémas. Soit  $D \subset S$  un diviseur de
Cartier effectif. Nous disons que le tiré en arrière  {\it  de
 $D$  par  $f$  est défini }  si le sous-schéma fermé
 $f^{-1}(D) \subset S'$  est un diviseur de Cartier effectif. Dans ce cas, nous le notons soit
 $f^*D$  soit  $f^{-1}(D)$  et nous l'appelons le tiré en arrière
 {\it  du diviseur de Cartier effectif }.
\end{definition}

\noindent
La condition selon laquelle  $f^{-1}(D)$  est un diviseur de Cartier effectif est
souvent satisfaite en pratique. Voici un lemme exemple.

\begin{lemma}
\label{lemma-pullback-effective-Cartier-defined}
Soit  $f : X \to Y$  un morphisme de schémas. Soit  $D \subset Y$  un diviseur de
Cartier effectif. Le tiré en arrière de  $D$  par  $f$  est défini
dans chacun des cas suivants:
\begin{enumerate}
\item $f(x) \not \in D$  pour tout point faiblement associé  $x$  de  $X$,
\item $X$,  $Y$  intègre et  $f$  dominant,
\item $X$  réduit et  $f(\xi) \not \in D$  pour tout point générique  $\xi$  de
toute composante irréductible de  $X$,
\item $X$  est localement noethérien et  $f(x) \not \in D$  pour tout point associé
 $x$  de  $X$,
\item $X$  est localement noethérien, n'a pas de points incorporés, et
 $f(\xi) \not \in D$  pour tout point générique  $\xi$  d'une composante irréductible
de  $X$,
\item $f$  est plat, et
\item ajouter plus ici si nécessaire.
\end{enumerate}
\end{lemma}

\begin{proof}
La question est locale sur  $X$, et nous réduisons donc au cas où
 $X = \Spec(A)$,  $Y = \Spec(R)$,  $f$  est donné par  $\varphi : R \to A$  et
 $D = \Spec(R/(t))$  où  $t \in R$  est un diviseur de zéro non nul. L'objectif dans
chaque cas est de montrer que  $\varphi(t) \in A$  est un diviseur de zéro non nul.

\medskip\noindent
Dans le cas (1), cela découle du Lemme d'Algèbre  \ref{algebra-lemma-weakly-ass-zero-divisors}. Le cas (4) est un
cas particulier du (1) par le Lemme  \ref{lemma-ass-weakly-ass}. Le cas (5) découle du (4) et des
définitions. Le cas (3) est un cas particulier du (1) par le Lemme  \ref{lemma-weakass-reduced}.
Le cas (2) est un cas particulier du (3). Si  $R \to A$  est plat, alors
 $t : R \to R$  étant injectif montre que  $t : A \to A$  est injectif. Cela prouve (6).
\end{proof}

\begin{lemma}
\label{lemma-pullback-effective-Cartier-divisors-additive}
Soit  $f : S' \to S$  un morphisme de schémas. Soient  $D_1$,  $D_2$  des
diviseurs de Cartier effectifs sur  $S$. Si les pullbacks de  $D_1$ 
et  $D_2$  sont définis, alors le pullback de  $D = D_1 + D_2$  est défini et
 $f^*D = f^*D_1 + f^*D_2$.
\end{lemma}

\begin{proof}
Omis.
\end{proof}





\section{Diviseurs de Cartier effectifs et faisceaux inversibles}
\label{section-effective-Cartier-invertible}

\noindent
Puisqu'un diviseur de Cartier effectif possède un faisceau d'idéaux inversible
(Définition  \ref{definition-effective-Cartier-divisor}), la définition suivante a un sens.

\begin{definition}
\label{definition-invertible-sheaf-effective-Cartier-divisor}
Soit  $S$  un schéma. Soit  $D \subset S$  un diviseur de Cartier effectif
avec le faisceau d'idéaux  $\mathcal{I}_D$.
\begin{enumerate}
\item Le  {\it faisceau inversible $\mathcal{O}_S(D)$ associé
à $D$} est défini par
$$
\mathcal{O}_S(D) =
\SheafHom_{\mathcal{O}_S}(\mathcal{I}_D, \mathcal{O}_S) =
\mathcal{I}_D^{\otimes -1}.
$$
\item Le  {\it section canonique}, généralement noté
 $1$  ou  $1_D$, est la section globale de  $\mathcal{O}_S(D)$ correspondant
à l'application d'inclusion $\mathcal{I}_D \to \mathcal{O}_S$.
\item Nous écrivons  $\mathcal{O}_S(-D) = \mathcal{O}_S(D)^{\otimes -1} = \mathcal{I}_D$.
\item Étant donné un second diviseur de Cartier effectif  $D' \subset S$ , nous définissons
 $\mathcal{O}_S(D - D') =
\mathcal{O}_S(D) \otimes_{\mathcal{O}_S} \mathcal{O}_S(-D')$.
\end{enumerate}
\end{definition}

\noindent
Quelques commentaires. Nous verrons ci-dessous que l'assignation  $D \mapsto \mathcal{O}_S(D)$ 
transforme l'addition des diviseurs de Cartier effectifs (Définition  \ref{definition-sum-effective-Cartier-divisors})
en addition dans le groupe de Picard de  $S$  (Lemme  \ref{lemma-invertible-sheaf-sum-effective-Cartier-divisors}).
Cependant, l'expression  $D - D'$  dans la définition ci-dessus n'a aucun sens
géométrique. Plus précisément, nous pouvons considérer l'ensemble des diviseurs de
Cartier effectifs sur  $S$  comme un monoïde commutatif  $\text{EffCart}(S)$  dont
l'élément zéro est le diviseur de Cartier effectif vide. Alors l'assignation
 $(D, D') \mapsto \mathcal{O}_S(D - D')$  définit un homomorphisme de groupes
$$
\text{EffCart}(S)^{gp} \longrightarrow \Pic(S)
$$
où le côté gauche est l'achèvement en groupe de  $\text{EffCart}(S)$. En d'autres termes,
lorsque nous écrivons  $\mathcal{O}_S(D - D')$ , nous pouvons penser à  $D - D'$  comme un
élément de  $\text{EffCart}(S)^{gp}$.

\begin{lemma}
\label{lemma-conormal-effective-Cartier-divisor}
Soit  $S$  un schéma et soit  $D \subset S$  un diviseur de Cartier effectif.
Alors le faisceau conormal est  $\mathcal{C}_{D/S} = \mathcal{I}_D|_D =
\mathcal{O}_S(-D)|_D$  et le faisceau normal est
 $\mathcal{N}_{D/S} = \mathcal{O}_S(D)|_D$.
\end{lemma}

\begin{proof}
Cela découle du Morphismes, Lemme,  \ref{morphisms-lemma-affine-conormal}.
\end{proof}

\begin{lemma}
\label{lemma-ses-add-divisor}
Soit  $X$  un schéma. Soient  $D, D' \subset X$  des diviseurs de Cartier
effectifs avec  $D \subset D'$  et soit  $C \subset X$  un diviseur de Cartier effectif
tel que  $D' = D + C$  (qui existe selon le Lemme  \ref{lemma-difference-effective-Cartier-divisors}). Alors il existe une
suite exacte courte
$$
0 \to \mathcal{O}_X(-D)|_C \to \mathcal{O}_{D'} \to \mathcal{O}_D \to 0
$$
de  $\mathcal{O}_X$-modules.
\end{lemma}

\begin{proof}
Dans l'énoncé du lemme et dans la démonstration, nous utilisons l'équivalence des
morphismes, Lemme  \ref{morphisms-lemma-i-star-equivalence} , pour penser aux modules quasi-cohérents sur les
sous-schémas fermés de  $X$  comme des modules quasi-cohérents sur
 $X$. Soit  $\mathcal{I}$  le faisceau d'idéaux de  $D$  dans
 $D'$  et soit  $\mathcal{I}_D = \mathcal{O}_X(-D)$  le faisceau d'idéaux de  $D$  dans
 $X$. Alors il existe une surjection canonique  $\mathcal{I}_D \to \mathcal{I}$. Ainsi, il
suffit de montrer que le noyau de cette application est égal à  $\mathcal{I}_C\mathcal{I}_D$  (avec
la notation évidente). Cela peut être vérifié localement, donc nous pouvons
supposer  $X = \Spec(A)$,  $D = \Spec(A/fA)$  et  $C = \Spec(A/gA)$  où  $f, g \in A$  sont des
non-diviseurs de zéro (Lemme  \ref{lemma-characterize-effective-Cartier-divisor}). Alors  $\mathcal{I}_D$  correspond à
 $fA$, le non-diviseur de zéro  $fg \in A$  génère l'idéal de  $D'$,
et  $\mathcal{I}$  correspond à  $I = fA/fgA$. Ainsi, le résultat est clair.
\end{proof}

\begin{lemma}
\label{lemma-invertible-sheaf-sum-effective-Cartier-divisors}
Soit  $S$  un schéma. Soient  $D_1$,  $D_2$  des diviseurs de
Cartier effectifs sur  $S$. Soit  $D = D_1 + D_2$. Alors il existe un
isomorphisme unique
$$
\mathcal{O}_S(D_1) \otimes_{\mathcal{O}_S} \mathcal{O}_S(D_2)
\longrightarrow
\mathcal{O}_S(D)
$$
qui envoie  $1_{D_1} \otimes 1_{D_2}$  sur  $1_D$.
\end{lemma}

\begin{proof}
Omis.
\end{proof}

\begin{lemma}
\label{lemma-pullback-effective-Cartier-divisors}
Soit  $f : S' \to S$  un morphisme de schémas. Soit  $D$  un diviseur de
Cartier effectif sur  $S$. Si le tiré en arrière de  $D$  est
défini alors  $f^*\mathcal{O}_S(D) = \mathcal{O}_{S'}(f^*D)$  et la section canonique  $1_D$  se tirent en
arrière sur la section canonique  $1_{f^*D}$.
\end{lemma}

\begin{proof}
Omis.
\end{proof}

\begin{definition}
\label{definition-regular-section}
Soit  $(X, \mathcal{O}_X)$  un espace localement annelé. Soit  $\mathcal{L}$  un faisceau
inversible sur  $X$. Une section globale  $s \in \Gamma(X, \mathcal{L})$  est appelée
 {\it  section régulière}  si l'application
 $\mathcal{O}_X \to \mathcal{L}$,  $f \mapsto fs$  est injective.
\end{definition}

\begin{lemma}
\label{lemma-regular-section-structure-sheaf}
Soit  $X$  un espace localement annelé. Soit  $f \in \Gamma(X, \mathcal{O}_X)$. Les énoncés
suivants sont équivalents:
\begin{enumerate}
\item $f$  est une section régulière, et
\item pour tout  $x \in X$  l'image  $f \in \mathcal{O}_{X, x}$  est un non-diviseur de zéro.
\end{enumerate}
Si  $X$  est un schéma, cela équivaut également à
\begin{enumerate}
\item[(3)] pour tout ouvert affine  $\Spec(A) = U \subset X$  l'image  $f \in A$  est un non-diviseur de
zéro,
\item[(4)] il existe un recouvrement par ouverts affines  $X = \bigcup \Spec(A_i)$  tel que l'image de
 $f$  dans  $A_i$  soit un non-diviseur de zéro pour tout
 $i$.
\end{enumerate}
\end{lemma}

\begin{proof}
Omis.
\end{proof}

\noindent
Notez qu'une section globale  $s$  d'un  $\mathcal{O}_X$-module inversible
 $\mathcal{L}$  peut être vue comme une application de  $\mathcal{O}_X$-module
 $s : \mathcal{O}_X \to \mathcal{L}$. Son dual est donc une application  $s : \mathcal{L}^{\otimes -1} \to \mathcal{O}_X$. (Voir Modules,
Définition  \ref{modules-definition-powers}  pour la définition du faisceau inversible dual.)

\begin{definition}
\label{definition-zero-scheme-s}
Soit  $X$  un schéma. Soit  $\mathcal{L}$  un faisceau inversible. Soit
 $s \in \Gamma(X, \mathcal{L})$  une section globale. Le  {\it  schéma des
zéros}  de  $s$  est le sous-schéma fermé  $Z(s) \subset X$  défini par
le faisceau quasi-cohérent d'idéaux  $\mathcal{I} \subset \mathcal{O}_X$  qui est l'image de l'application
 $s : \mathcal{L}^{\otimes -1} \to \mathcal{O}_X$.
\end{definition}

\begin{lemma}
\label{lemma-zero-scheme}
Soit  $X$  un schéma. Soit  $\mathcal{L}$  un faisceau inversible. Soit
 $s \in \Gamma(X, \mathcal{L})$.
\begin{enumerate}
\item Considérons des immersions fermées  $i : Z \to X$  telles que  $i^*s \in \Gamma(Z, i^*\mathcal{L})$  soit
ordonné zéro par inclusion. Le schéma des zéros  $Z(s)$  est l'élément maximal
de cet ensemble ordonné.
\item Pour tout morphisme de schémas  $f : Y \to X$ , nous avons  $f^*s = 0$  dans
 $\Gamma(Y, f^*\mathcal{L})$  si et seulement si  $f$  se factorise par  $Z(s)$.
\item Le schéma des zéros  $Z(s)$  est un sous-schéma fermé localement principal.
\item Le schéma zéro  $Z(s)$  est un diviseur de Cartier effectif si et seulement si
 $s$  est une section régulière de  $\mathcal{L}$.
\end{enumerate}
\end{lemma}

\begin{proof}
Omis.
\end{proof}

\begin{lemma}
\label{lemma-characterize-OD}
\begin{slogan}
Les diviseurs de Cartier effectifs sur un schéma sont les mêmes que les faisceaux
inversibles avec une section globale régulière fixée.
\end{slogan}
Soit  $X$  un schéma.
\begin{enumerate}
\item Si  $D \subset X$  est un diviseur de Cartier effectif, alors la section canonique
 $1_D$  de  $\mathcal{O}_X(D)$  est régulière.
\item Inversement, si  $s$  est une section régulière du faisceau inversible
 $\mathcal{L}$, alors il existe un unique diviseur de Cartier effectif  $D = Z(s) \subset X$ 
et un unique isomorphisme  $\mathcal{O}_X(D) \to \mathcal{L}$  qui envoie  $1_D$  sur  $s$.
\end{enumerate}
Les constructions  $D \mapsto (\mathcal{O}_X(D), 1_D)$  et  $(\mathcal{L}, s) \mapsto Z(s)$  fournissent des applications
mutuellement inverses.
$$
\left\{
\begin{matrix}
\text{effective Cartier divisors on }X
\end{matrix}
\right\}
\leftrightarrow
\left\{
\begin{matrix}
\text{isomorphism classes of pairs }(\mathcal{L}, s)\\
\text{consisting of an invertible }
\mathcal{O}_X\text{-module}\\
\mathcal{L}\text{ and a regular global section }s
\end{matrix}
\right\}
$$
\end{lemma}

\begin{proof}
Omis.
\end{proof}

\begin{remark}
\label{remark-ses-regular-section}
Soit  $X$  un schéma,  $\mathcal{L}$  un  $\mathcal{O}_X$-module inversible, et
 $s$  une section régulière de  $\mathcal{L}$. Alors le schéma des zéros
 $D = Z(s)$  est un diviseur de Cartier effectif sur  $X$  et il existe des
suites exactes courtes
$$
0 \to \mathcal{O}_X \to \mathcal{L} \to i_*(\mathcal{L}|_D) \to 0
\quad\text{and}\quad
0 \to \mathcal{L}^{\otimes -1} \to \mathcal{O}_X \to i_*\mathcal{O}_D \to 0.
$$
Étant donné un diviseur de Cartier effectif  $D \subset X$  en utilisant les lemmes
 \ref{lemma-characterize-OD}  et  \ref{lemma-conormal-effective-Cartier-divisor} , nous obtenons
$$
0 \to \mathcal{O}_X \to \mathcal{O}_X(D) \to i_*(\mathcal{N}_{D/X}) \to 0
\quad\text{and}\quad
0 \to \mathcal{O}_X(-D) \to \mathcal{O}_X \to i_*(\mathcal{O}_D) \to 0
$$
\end{remark}






\section{Diviseurs de Cartier effectifs sur les schémas noethériens}
\label{section-Noetherian-effective-Cartier}

\noindent
Dans le cadre localement noethérien, la plupart des discussions sur les diviseurs
de Cartier effectifs et les sections régulières se simplifient quelque peu.

\begin{lemma}
\label{lemma-regular-section-associated-points}
Soit  $X$  un schéma localement noethérien. Soit  $\mathcal{L}$  un
 $\mathcal{O}_X$-module inversible. Soit  $s \in \Gamma(X, \mathcal{L})$. Alors  $s$  est une
section régulière si et seulement si  $s$  ne s'annule pas aux points
associés de  $X$.
\end{lemma}

\begin{proof}
Omission. Indice: réduire au cas affine et  $\mathcal{L}$  trivial, puis utiliser le
lemme  \ref{lemma-regular-section-structure-sheaf}  et Algèbre, lemme  \ref{algebra-lemma-ass-zero-divisors}.
\end{proof}

\begin{lemma}
\label{lemma-effective-Cartier-in-points}
Soit  $X$  un schéma localement noethérien. Soit  $D \subset X$  un
sous-schéma fermé correspondant au faisceau d'idéaux quasi-cohérent  $\mathcal{I} \subset \mathcal{O}_X$.
\begin{enumerate}
\item Si pour chaque  $x \in D$  l'idéal  $\mathcal{I}_x \subset \mathcal{O}_{X, x}$  peut être engendré par un seul
élément, alors  $D$  est localement principal.
\item Si pour chaque  $x \in D$  l'idéal  $\mathcal{I}_x \subset \mathcal{O}_{X, x}$  peut être engendré par un seul
élément non diviseur de zéro, alors  $D$  est un diviseur de Cartier
effectif.
\end{enumerate}
\end{lemma}

\begin{proof}
Soit  $\Spec(A)$  un voisinage affine d'un point  $x \in D$. Soit  $\mathfrak p \subset A$ 
le premier idéal correspondant à  $x$. Soit  $I \subset A$  l'idéal
définissant la trace de  $D$  sur  $\Spec(A)$. Puisque  $A$  est
noethérien (car  $X$  est localement noethérien), l'idéal  $I$  est
engendré par un nombre fini d'éléments, disons  $I = (f_1, \ldots, f_r)$. Sous l'hypothèse de
(1), nous avons  $I_\mathfrak p = (f)$  pour un certain  $f \in A_\mathfrak p$. Alors  $f_i = g_i f$  pour
un certain  $g_i \in A_\mathfrak p$. Écrivons  $g_i = a_i/h_i$  et  $f = f'/h$  pour certains
 $a_i, h_i, f', h \in A$,  $h_i, h \not \in \mathfrak p$. Alors  $I_{h_1 \ldots h_r h} \subset A_{h_1 \ldots h_r h}$  est principal, car il est engendré
par  $f'$. Cela prouve (1). Pour (2) nous pouvons supposer  $I = (f)$.
L'hypothèse implique que l'image de  $f$  dans  $A_\mathfrak p$  est un élément
non diviseur de zéro. Alors  $f$  est un élément non diviseur de zéro dans
un voisinage de  $x$  par Algèbre, Lemme  \ref{algebra-lemma-regular-sequence-in-neighbourhood}. Cela prouve (2).
\end{proof}

\begin{lemma}
\label{lemma-effective-Cartier-codimension-1}
Soit  $X$  un schéma localement noethérien.
\begin{enumerate}
\item Soit  $D \subset X$  un sous-schéma fermé localement principal. Soit  $\xi \in D$  un
point générique d'une composante irréductible de  $D$. Alors  $\dim(\mathcal{O}_{X, \xi}) \leq 1$.
\item Soit  $D \subset X$  un diviseur de Cartier effectif. Soit  $\xi \in D$  un point
générique d'une composante irréductible de  $D$. Alors  $\dim(\mathcal{O}_{X, \xi}) = 1$.
\end{enumerate}
\end{lemma}

\begin{proof}
Preuve de (1). Par hypothèse, nous pouvons supposer  $X = \Spec(A)$  et  $D = \Spec(A/(f))$ 
où  $A$  est un anneau noethérien et  $f \in A$. Soit  $\xi$ 
correspondant à l'idéal premier  $\mathfrak p \subset A$. L'hypothèse que  $\xi$  est un
point générique d'une composante irréductible de  $D$  signifie que
 $\mathfrak p$  est minimal sur  $(f)$. Ainsi  $\dim(A_\mathfrak p) \leq 1$  par Algèbre, Lemme
 \ref{algebra-lemma-minimal-over-1}.

\medskip\noindent
Preuve de (2). D'après la partie (1), nous voyons que  $\dim(\mathcal{O}_{X, \xi}) \leq 1$. D'autre part,
l'équation locale  $f$  est un non-diviseur dans  $A_\mathfrak p$  par le Lemme
 \ref{lemma-characterize-effective-Cartier-divisor}  ce qui implique que la dimension est au moins  $1$  (parce
qu'il doit y avoir un premier dans  $A_\mathfrak p$  ne contenant pas  $f$ 
selon le Lemme élémentaire de l'Algèbre  \ref{algebra-lemma-Zariski-topology}).
\end{proof}

\begin{lemma}
\label{lemma-integral-effective-Cartier-divisor-dvr}
Soit  $X$  un schéma noethérien. Soit  $D \subset X$  un sous-schéma fermé
intégral qui est aussi un diviseur de Cartier effectif. Alors l'anneau local de
 $X$  au point générique de  $D$  est un anneau de valuation
discrète.
\end{lemma}

\begin{proof}
D'après le Lemme  \ref{lemma-characterize-effective-Cartier-divisor} , nous pouvons supposer que  $X = \Spec(A)$  et
 $D = \Spec(A/(f))$  où  $A$  est un anneau noethérien et  $f \in A$  est un
non-diviseur de zéro. L'hypothèse que  $D$  est intégral signifie que
 $(f)$  est premier. Par conséquent, l'anneau local de  $X$  au point
générique est  $A_{(f)}$ , qui est un anneau local noethérien dont l'idéal
maximal est engendré par un non-diviseur de zéro. Ainsi, c'est un anneau de
valuation discrète d'après Algèbre, Lemme  \ref{algebra-lemma-characterize-dvr}.
\end{proof}

\begin{lemma}
\label{lemma-effective-Cartier-divisor-Sk}
Soit  $X$  un schéma localement noethérien. Soit  $D \subset X$  un diviseur
de Cartier effectif. Si  $X$  est  $(S_k)$, alors  $D$  est
 $(S_{k - 1})$.
\end{lemma}

\begin{proof}
Soit  $x \in D$. Alors  $\mathcal{O}_{D, x} = \mathcal{O}_{X, x}/(f)$  où  $f \in \mathcal{O}_{X, x}$  est un non-diviseur de zéro.
Par hypothèse, nous avons  $\text{depth}(\mathcal{O}_{X, x}) \geq \min(\dim(\mathcal{O}_{X, x}), k)$. Par Algèbre, Lemme  \ref{algebra-lemma-depth-drops-by-one}  nous avons
 $\text{depth}(\mathcal{O}_{D, x}) = \text{depth}(\mathcal{O}_{X, x}) - 1$  et par Algèbre, Lemme  \ref{algebra-lemma-one-equation}   $\dim(\mathcal{O}_{D, x}) = \dim(\mathcal{O}_{X, x}) - 1$. Il s'ensuit que
 $\text{depth}(\mathcal{O}_{D, x}) \geq \min(\dim(\mathcal{O}_{D, x}), k - 1)$  comme désiré.
\end{proof}

\begin{lemma}
\label{lemma-normal-effective-Cartier-divisor-S1}
Soit  $X$  un schéma normal localement noethérien. Soit  $D \subset X$  un
diviseur de Cartier effectif. Alors  $D$  est  $(S_1)$.
\end{lemma}

\begin{proof}
Par Propriétés, Lemme  \ref{properties-lemma-criterion-normal}  nous voyons que  $X$  est  $(S_2)$.
Ainsi, nous concluons par le Lemme  \ref{lemma-effective-Cartier-divisor-Sk}.
\end{proof}

\begin{lemma}
\label{lemma-weil-divisor-is-cartier-UFD}
Soit  $X$  un schéma localement noethérien. Soit  $D \subset X$  un
sous-schéma fermé intègre. Supposons que
\begin{enumerate}
\item $D$  ait une codimension  $1$  dans  $X$, et
\item $\mathcal{O}_{X, x}$  est un UFD pour tout  $x \in D$.
\end{enumerate}
Alors  $D$  est un diviseur de Cartier effectif.
\end{lemma}

\begin{proof}
Soit  $x \in D$  et soit  $A = \mathcal{O}_{X, x}$. Soit  $\mathfrak p \subset A$  correspondant au point
générique de  $D$. Alors  $A_\mathfrak p$  a une dimension  $1$  par
l'assumption (1). Ainsi  $\mathfrak p$  est un idéal premier de hauteur  $1$.
Puisque  $A$  est un UFD, cela implique que  $\mathfrak p = (f)$  pour un certain
 $f \in A$. Bien sûr,  $f$  est un non-diviseur de zéro et nous concluons
par le Lemme  \ref{lemma-effective-Cartier-in-points}.
\end{proof}

\begin{lemma}
\label{lemma-codim-1-part}
Soit  $X$  un schéma noethérien. Soit  $Z \subset X$  un sous-schéma fermé.
Supposons qu'il existe
\begin{enumerate}
\item une collection de diviseurs de Cartier effectifs intègres  $D_i \subset X$,
 $i \in I$
\item un sous-ensemble fermé  $Z' \subset X$  dont toutes les composantes irréductibles ont
une codimension  $\geq 2$  dans  $X$  (Topologie, Définition
 \ref{topology-definition-codimension})
\end{enumerate}
de sorte que  $Z \subset Z' \cup \bigcup_{i \in I} D_i$  est d'un point de vue théorique des ensembles. Alors il
existe des entiers  $a_i \geq 0$  avec  $a_i = 0$  pour presque tous les
 $i \in I$  de sorte que le diviseur de Cartier effectif
$$
D = \sum a_i D_i
$$
est contenu dans  $Z$  et que le morphisme d'inclusion  $D \to Z$  est un
isomorphisme hors de la codimension  $2$  dans  $X$  (dans le sens
qu'il existe un ouvert  $U \subset Z$  tel que  $D \cap U \to Z \cap U$  est un isomorphisme et
que chaque composante irréductible de  $Z \setminus U$  a une codimension  $\geq 2$ 
dans  $X$). Lorsque  $Z$  est rare dans  $X$ , l'existence
de  $D_i$,  $i \in I$  et  $Z'$  est garantie si  $\mathcal{O}_{X, x}$  est un
UFD pour tous les  $x \in Z$  ou si  $X$  est régulier.
\end{lemma}

\begin{proof}
Soit  $\xi_i \in D_i$  le point générique et soit  $\mathcal{O}_i = \mathcal{O}_{X, \xi_i}$  l'anneau local qui est
un anneau de valuation discrète d'après le Lemme  \ref{lemma-integral-effective-Cartier-divisor-dvr}. Soit  $a_i \geq 0$ 
la valuation minimale d'un élément de  $\mathcal{I}_{Z, \xi_i} \subset \mathcal{O}_i$. Bien sûr  $a_i > 0$ 
seulement si  $D_i$  est une composante irréductible de  $Z$  et donc
 $a_i > 0$  seulement pour un nombre fini de  $i \in I$. Nous affirmons que le
diviseur de Cartier effectif  $D = \sum a_i D_i$  fonctionne.

\medskip\noindent
À savoir, supposons que  $x \in X$. Soit  $A = \mathcal{O}_{X, x}$. Soient  $D_1, \ldots, D_n$  les
diviseurs mutuellement distincts  $D_i$  tels que  $x \in D_i$  et
 $a_i > 0$. Pour  $1 \leq i \leq n$ , soit  $f_i \in A$  une équation locale pour
 $D_i$. Alors  $f_i$  est un élément premier de  $A$  et
 $\mathcal{O}_i = A_{(f_i)}$. Soit  $I = \mathcal{I}_{Z, x} \subset A$  la germe du faisceau d'idéaux de  $Z$. Par
notre choix de  $a_i$ , nous avons  $I A_{(f_i)} = f_i^{a_i}A_{(f_i)}$. Nous affirmons que
 $I \subset (\prod_{i = 1, \ldots, n} f_i^{a_i})$.

\medskip\noindent
Preuve de l'assertion. La carte de localisation  $\varphi : A/(f_i) \to A_{(f_i)}/f_iA_{(f_i)}$  est injective puisque
l'idéal premier  $(f_i)$  est l'image réciproque de l'idéal maximal
 $f_iA_{(f_i)}$. Par induction sur  $n$ , nous déduisons que  $\varphi_n : A/(f_i^n)\to A_{(f_i)}/f_i^nA_{(f_i)}$  est
également injective. Puisque  $\varphi_{a_i}(I) = 0$, nous avons  $I \subset (f_i^{a_i})$. Ainsi, pour
tout  $x \in I$, nous pouvons écrire  $x = f_1^{a_1}x_1$  pour un certain  $x_1 \in A$.
Puisque  $D_1, \ldots, D_n$  sont deux à deux distincts,  $f_i$  est une unité dans
 $A_{(f_j)}$  pour  $i \not = j$. En comparant  $x$  et  $x_1$  en
 $A_{(f_i)}$  pour  $n \geq i > 1$, nous avons encore  $x_1 \in (f_i^{a_i})$. En répétant le
processus précédent, nous écrivons par induction  $x_i = f_{i + 1}^{a_{i + 1}}x_{i + 1}$  pour tout
 $n > i \geq 1$. En conclusion,  $x \in (\prod_{i = 1, \ldots n} f_i^{a_i})$  pour tout  $x \in I$  comme souhaité.

\medskip\noindent
La réclamation montre que  $\mathcal{I}_Z \subset \mathcal{I}_D$, c'est-à-dire que  $D \subset Z$. De plus,
nous voyons également que  $D$  et  $Z$  concordent au  $\xi_i$,
ce qui prouve que  $D \to Z$  est un isomorphisme hors de la codimension
 $2$  sur  $X$.

\medskip\noindent
Pour voir les dernières déclarations, nous argumentons comme suit. Un anneau local
régulier est un UFD (Plus sur Algèbre, Lemme  \ref{more-algebra-lemma-regular-local-UFD}), donc il suffit de
raisonner dans le cas du UFD. Dans ce cas, soient  $D_i$  les composantes
irréductibles de  $Z$  qui ont une codimension  $1$  dans
 $X$. Par le Lemme  \ref{lemma-weil-divisor-is-cartier-UFD} , chaque  $D_i$  est un diviseur de
Cartier effectif.
\end{proof}

\begin{lemma}
\label{lemma-codimension-1-is-effective-Cartier}
Soit  $Z \subset X$  un sous-schéma fermé d'un schéma noethérien. Supposons
\begin{enumerate}
\item $Z$  n'a pas de points incorporés,
\item chaque composante irréductible de  $Z$  a une codimension  $1$  dans
 $X$.
\item chaque anneau local  $\mathcal{O}_{X, x}$,  $x \in Z$  est un UFD ou  $X$  est
régulier.
\end{enumerate}
Alors  $Z$  est un diviseur de Cartier effectif.
\end{lemma}

\begin{proof}
Soit  $D = \sum a_i D_i$  comme dans le Lemme  \ref{lemma-codim-1-part}  où  $D_i \subset Z$  sont les
composantes irréductibles de  $Z$. Si  $D \to Z$  n'est pas un
isomorphisme, alors  $\mathcal{O}_Z \to \mathcal{O}_D$  a un noyau non nul situé en codimension
 $\geq 2$. Cela signifierait que  $Z$  a des points incorporés, ce qui
est interdit par l'hypothèse (1). Ainsi  $D \cong Z$  comme désiré.
\end{proof}

\begin{lemma}
\label{lemma-UFD-one-equation-CM}
Soit  $R$  un UFD noethérien. Soit  $I \subset R$  un idéal tel que
 $R/I$  n'a pas de premiers incorporés et tel que chaque premier minimal
au-dessus de  $I$  a une hauteur  $1$. Alors  $I = (f)$  pour un
certain  $f \in R$.
\end{lemma}

\begin{proof}
D'après le Lemme  \ref{lemma-codimension-1-is-effective-Cartier} , le faisceau d'idéaux  $\tilde I$  est inversible
sur  $\Spec(R)$. D'après Plus sur Algèbre, Lemme  \ref{more-algebra-lemma-UFD-Pic-trivial} , il est engendré
par un seul élément.
\end{proof}

\begin{lemma}
\label{lemma-effective-Cartier-divisor-is-a-sum}
Soit  $X$  un schéma noethérien. Soit  $D \subset X$  un diviseur de Cartier
effectif. Supposons qu'il existe des diviseurs de Cartier effectifs intègres
 $D_i \subset X$  tels que  $D \subset \bigcup D_i$ , théoriquement par ensembles. Alors
 $D = \sum a_i D_i$  pour un certain  $a_i \geq 0$. L'existence de  $D_i$  est
garantie si  $\mathcal{O}_{X, x}$  est un UFD pour tout  $x \in D$  ou si  $X$  est
régulier.
\end{lemma}

\begin{proof}
Choisissez  $a_i$  comme dans le Lemme  \ref{lemma-codim-1-part}  et définissez
 $D' = \sum a_i D_i$. Alors  $D' \to D$  est une inclusion de diviseurs de Cartier
effectifs qui est un isomorphisme en dehors de la codimension  $2$  sur
 $X$. Choisissez  $x \in X$. Définissez  $A = \mathcal{O}_{X, x}$  et laissez
 $f, f' \in A$  être le non-diviseur de zéro générant l'idéal de  $D, D'$  dans
 $A$. Alors  $f = gf'$  pour quelque  $g \in A$. De plus, pour chaque
idéal premier  $\mathfrak p$  de hauteur  $\leq 1$  de  $A$ , nous voyons
que  $g$  se mappe sur une unité de  $A_\mathfrak p$. Cela implique que
 $g$  est une unité parce que les idéaux premiers minimaux au-dessus de
 $(g)$  ont une hauteur  $1$  (Algèbre, Lemme  \ref{algebra-lemma-minimal-over-1}).
\end{proof}

\begin{lemma}
\label{lemma-quasi-projective-Noetherian-pic-effective-Cartier}
\begin{slogan}
Sur un schéma projectif, chaque fibré en droite a une section méromorphe
régulière.
\end{slogan}
Soit  $X$  un schéma noethérien qui possède un faisceau inversible ample.
Alors chaque  $\mathcal{O}_X$-module inversible est isomorphe à
$$
\mathcal{O}_X(D - D') =
\mathcal{O}_X(D) \otimes_{\mathcal{O}_X} \mathcal{O}_X(D')^{\otimes -1}
$$
pour certains diviseurs de Cartier effectifs  $D, D'$  dans  $X$. De
plus, étant donné un sous-ensemble fini  $E \subset X$ , nous pouvons choisir
 $D, D'$  tel que  $E \cap D = \emptyset$  et  $E \cap D' = \emptyset$. Si  $X$  est
quasi-affine, alors nous pouvons choisir  $D' = \emptyset$.
\end{lemma}

\begin{proof}
Soit  $x_1, \ldots, x_n$  les points associés de  $X$  (Lemme  \ref{lemma-finite-ass}).

\medskip\noindent
Si  $X$  est quasi-affine et  $\mathcal{N}$  est un module  $\mathcal{O}_X$
inversible quelconque, alors nous pouvons choisir une section  $t$  de
 $\mathcal{N}$  qui ne s'annule en aucun des points de  $E \cup \{x_1, \ldots, x_n\}$, voir Propriétés,
Lemme  \ref{properties-lemma-quasi-affine-invertible-nonvanishing-section}. Ensuite,  $t$  est une section régulière de
 $\mathcal{N}$  selon le Lemme  \ref{lemma-regular-section-associated-points}. Par conséquent,  $\mathcal{N} \cong \mathcal{O}_X(D)$  où
 $D = Z(t)$  est le diviseur de Cartier effectif correspondant à  $t$,
voir Lemme  \ref{lemma-characterize-OD}. Puisque  $E \cap D = \emptyset$  par construction, nous avons terminé
dans ce cas.

\medskip\noindent
En revenant au cas général, soit  $\mathcal{L}$  un faisceau inversible ample sur
 $X$. Il existe un  $m > 0$  et une section  $s \in \Gamma(X, \mathcal{L}^{\otimes m})$  tels que
 $X_s$  soit affine et tels que  $E \cup \{x_1, \ldots, x_n\} \subset X_s$  (Propriétés, Lemme
 \ref{properties-lemma-ample-finite-set-in-principal-affine}).

\medskip\noindent
Soit  $\mathcal{N}$  un  $\mathcal{O}_X$-module inversible arbitraire. Dans le cas
quasi-affine, nous pouvons trouver une section  $t \in \mathcal{N}(X_s)$  qui ne s'annule en
aucun point de  $E \cup \{x_1, \ldots, x_n\}$. D'après les Propriétés, Lemme  \ref{properties-lemma-invert-s-sections} , nous
voyons que pour un certain  $e \geq 0$ , la section  $s^e|_{X_s} t$  s'étend à une
section globale  $\tau$  de  $\mathcal{L}^{\otimes me} \otimes \mathcal{N}$. Ainsi,  $\mathcal{L}^{\otimes me} \otimes \mathcal{N}$  et  $\mathcal{L}^{\otimes me}$ 
sont toutes deux des faisceaux inversibles qui possèdent des sections globales ne
s'annulant en aucun point de  $E \cup \{x_1, \ldots, x_n\}$. Par conséquent, ce sont des sections
régulières d'après le Lemme  \ref{lemma-regular-section-associated-points}. Donc  $\mathcal{L}^{\otimes me} \otimes \mathcal{N} \cong \mathcal{O}_X(D)$  et  $\mathcal{L}^{\otimes me} \cong \mathcal{O}_X(D')$  pour
certains diviseurs de Cartier effectifs  $D$  et  $D'$, voir le
Lemme  \ref{lemma-characterize-OD}. Par construction  $E \cap D = \emptyset$  et  $E \cap D' = \emptyset$  et la
démonstration est terminée.
\end{proof}

\begin{lemma}
\label{lemma-wedge-product-ses}
Soit  $X$  un schéma régulier intègre de dimension  $2$. Soit
 $i : D \to X$  l'immersion d'un diviseur de Cartier effectif. Soit  $\mathcal{F} \to \mathcal{F}' \to i_*\mathcal{G} \to 0$  une
suite exacte de  $\mathcal{O}_X$-modules cohérents. Supposons que
\begin{enumerate}
\item $\mathcal{F}, \mathcal{F}'$  soient localement libres de rang  $r$  sur un ouvert non vide
de  $X$,
\item $D$  est un schéma intègre,
\item $\mathcal{G}$  est un  $\mathcal{O}_D$-module localement libre de rang fini
 $s$.
\end{enumerate}
Alors  $\mathcal{L} = (\wedge^r\mathcal{F})^{**}$  et  $\mathcal{L}' = (\wedge^r \mathcal{F}')^{**}$  sont des  $\mathcal{O}_X$-modules inversibles et
 $\mathcal{L}' \cong \mathcal{L}(k D)$  pour un certain  $k \in \{0, \ldots, \min(s, r)\}$.
\end{lemma}

\begin{proof}
La première assertion découle du Lemme  \ref{lemma-reflexive-over-regular-dim-2}  car l'hypothèse (1) implique
que  $\mathcal{L}$  et  $\mathcal{L}'$  ont un rang  $1$. Comme prendre
 $\wedge^r$  et les doubles duaux sont des foncteurs, nous obtenons donc une
application canonique  $\sigma : \mathcal{L} \to \mathcal{L}'$  qui est un isomorphisme sur l'ouvert non vide
de (1), et donc non nulle. Pour terminer la démonstration, il suffit de voir que
 $\sigma$  considéré comme une section globale de  $\mathcal{L}' \otimes \mathcal{L}^{\otimes -1}$  ne s'annule en
aucun point de codimension de  $X$, sauf au point générique de
 $D$  et là avec un ordre de disparition au plus égal à  $\min(s, r)$.

\medskip\noindent
Traduit en algèbre, nous arrivons au problème suivant: Soit  $(A, \mathfrak m, \kappa)$  un anneau
de valuation discrète avec corps des fractions  $K$. Soit  $M \to M' \to N \to 0$  une
suite exacte de  $A$-modules finis avec  $\dim_K(M \otimes K) = \dim_K(M' \otimes K) = r$  et avec  $N \cong \kappa^{\oplus s}$.
Montrer que l'application induite  $L = \wedge^r(M)^{**} \to L' = \wedge^r(M')^{**}$  s'annule d'ordre au plus
 $\min(s, r)$. Nous utiliserons le théorème de structure pour les modules sur
 $A$, voir Plus sur Algèbre, Lemme  \ref{more-algebra-lemma-generalized-valuation-ring-modules}  ou  \ref{more-algebra-lemma-modules-PID}. Diviser
un  $A$-module fini par un sous-module de torsion ne change pas le double
dual. Ainsi, nous pouvons remplacer  $M$  par  $M/M_{tors}$  et  $M'$ 
par  $M'/\Im(M_{tors} \to M')$  et supposer que  $M$  est sans torsion. Alors  $M \to M'$ 
est injectif et  $M'_{tors} \to N$  est injectif. Par conséquent, nous pouvons remplacer
 $M'$  par  $M'/M'_{tors}$  et  $N$  par  $N/M'_{tors}$. Ainsi, nous
réduisons au cas où  $M$  et  $M'$  sont libres de rang  $r$ 
et  $N \cong \kappa^{\oplus s}$. Dans ce cas,  $\sigma$  est le déterminant de  $M \to M'$  et
s'annule d'ordre  $s$ , par exemple d'après Algèbre, Lemme  \ref{algebra-lemma-order-vanishing-determinant}.
\end{proof}









\section{Complémentaires d'ouverts affines}
\label{section-complement-affine-open}

\noindent
Dans cette section, nous discutons du résultat selon lequel le complémentaire d'un
ouvert affine dans une variété a une codimension pure  $1$.

\begin{lemma}
\label{lemma-affine-punctured-spec}
Soit  $(A, \mathfrak m)$  un anneau local noethérien. Le spectre ponctué  $U = \Spec(A) \setminus \{\mathfrak m\}$  de
 $A$  est affine si et seulement si  $\dim(A) \leq 1$.
\end{lemma}

\begin{proof}
Si  $\dim(A) = 0$, alors  $U$  est vide et donc affine (égale au spectre de
l'anneau  $0$ ). Si  $\dim(A) = 1$, alors nous pouvons choisir un élément
 $f \in \mathfrak m$  qui n'est contenu dans aucun des nombres finis de premiers minimaux
de  $A$  (Algèbre, Lemme  \ref{algebra-lemma-Noetherian-irreducible-components}  et  \ref{algebra-lemma-silly}). Alors  $U = \Spec(A_f)$ 
est affine.

\medskip\noindent
Le cas inverse est plus intéressant. Nous allons donner une preuve quelque peu non
standard et discuter de l'argument standard dans une remarque ci-dessous.
Supposons que  $U = \Spec(B)$  soit affine. Puisque l'affinité et la dimension ne sont
pas affectées en passant à la réduction, nous pouvons remplacer  $A$  par
le quotient par son idéal d'éléments nilpotents et supposer que  $A$  est
réduit. Posons  $Q = B/A$  considéré comme un  $A$-module. Le support de
 $Q$  est  $\{\mathfrak m\}$  comme  $A_\mathfrak p = B_\mathfrak p$  pour tous les premiers non
maximaux  $\mathfrak p$  de  $A$. Nous pouvons supposer  $\dim(A) \geq 1$, donc
comme ci-dessus nous pouvons choisir  $f \in \mathfrak m$  qui n'est contenu dans aucun des
idéaux minimaux de  $A$. Puisque  $A$  est réduit, cela implique
que  $f$  est un diviseur de zéro non nul dans  $A$. Le même
argument montre que  $f : B \to B$  est injectif car  $\Spec(B)$  est ouvert dans
 $\Spec(A)$. Notez que  $\dim(A/fA) = \dim(A) - 1$  d'après Algèbre, Lemme  \ref{algebra-lemma-one-equation}. En
appliquant le lemme du serpent à la multiplication par  $f$  sur la suite
exacte courte  $0 \to A \to B \to Q \to 0$ , nous obtenons
$$
0 \to Q[f] \to A/fA \to B/fB \to Q/fQ \to 0
$$
où  $Q[f] = \Ker(f : Q \to Q)$. Cela implique que  $Q[f]$  est un  $A$-module fini.
Puisque le support de  $Q[f]$  est  $\{\mathfrak m\}$ , nous voyons  $l = \text{length}_A(Q[f]) < \infty$ 
(Algèbre, Lemme  \ref{algebra-lemma-support-point}). Posons  $l_n = \text{length}_A(Q[f^n])$. La suite exacte
$$
0 \to Q[f^n] \to Q[f^{n + 1}] \xrightarrow{f^n} Q[f]
$$
montre par induction que  $l_n < \infty$  et que  $l_n \leq l_{n + 1}$. En considérant la suite
exacte
$$
0 \to Q[f] \to Q[f^{n + 1}] \xrightarrow{f} Q[f^n] \to Q/fQ
$$
et nous voyons que l'image de  $Q[f^n]$  dans  $Q/fQ$  a une longueur de
 $l_n - l_{n + 1} + l \leq l$. Puisque  $Q = \bigcup Q[f^n]$  nous trouvons que la longueur de  $Q/fQ$ 
est au plus  $l$, c'est-à-dire bornée. Ainsi  $Q/fQ$  est un
 $A$-module fini. Par conséquent,  $A/fA \to B/fB$  est un morphisme d'anneaux
fini, en particulier il induit une application fermée sur les spectres (Algèbre,
lemmes  \ref{algebra-lemma-integral-going-up}  et  \ref{algebra-lemma-going-up-closed}). D'autre part,  $\Spec(B/fB)$  est le spectre
percé de  $\Spec(A/fA)$. Ceci est une contradiction à moins de  $\Spec(B/fB) = \emptyset$ , ce qui
signifie que  $\dim(A/fA) = 0$  comme souhaité.
\end{proof}

\begin{remark}
\label{remark-affine-punctured-spectrum-standard-proof}
Si  $(A, \mathfrak m)$  est un domaine local noethérien normal de dimension  $\geq 2$ 
et  $U$  est le spectre puncturé de  $A$, alors  $\Gamma(U, \mathcal{O}_U) = A$. Cette
version algébrique du théorème de Hartogs découle du fait que  $A = \bigcap_{\text{height}(\mathfrak p) = 1} A_\mathfrak p$  que
nous avons vu en Algèbre, Lemme  \ref{algebra-lemma-normal-domain-intersection-localizations-height-1}. Ainsi, dans ce cas,  $U$  ne
peut pas être affine (car cela obligerait  $\mathfrak m$  à être un point de
 $U$). Cela est souvent utilisé comme point de départ de la preuve du
Lemme  \ref{lemma-affine-punctured-spec}. Pour réduire le cas d'un anneau local noethérien général à ce
cas, nous complétons d'abord (pour obtenir un anneau local de Nagata), puis
remplaçons  $A$  par  $A/\mathfrak q$  pour un idéal premier minimal approprié,
puis normalisons. Chacune de ces étapes ne change pas la dimension et nous
obtenons une contradiction. Vous pouvez sauter l'étape de la complétion, mais
alors la normalisation n'est généralement pas un domaine noethérien. Cependant,
c'est toujours un domaine de Krull de la même dimension (cela se prouve en
utilisant Krull-Akizuki) et on peut appliquer le même argument.
\end{remark}

\begin{remark}
\label{remark-affine-puctured-spectrum-general}
Il n'est pas clair comment caractériser les anneaux locaux non noethériens
 $(A, \mathfrak m)$  dont le spectre ponctué est affine. Un tel anneau possède un idéal
fini  $I$  avec  $\mathfrak m = \sqrt{I}$. Bien sûr, si nous pouvons prendre
 $I$  généré par  $1$  élément, alors  $A$  a un spectre
ponctué affine; cela donne de nombreux exemples non noethériens. Inversement, il
découle de l'argument dans la démonstration du Lemme  \ref{lemma-affine-punctured-spec}  qu'un tel anneau
ne peut pas posséder un non-diviseur  $f \in \mathfrak m$  avec  $H^0_I(A/fA) = 0$  (donc
 $A$  ne peut pas avoir une suite régulière de longueur  $2$). De
plus, il en va de même pour tout anneau  $A'$  qui est la cible d'un
homomorphisme local d'anneaux locaux  $A \to A'$  tel que  $\mathfrak m_{A'} = \sqrt{\mathfrak mA'}$.
\end{remark}

\begin{lemma}
\label{lemma-complement-affine-open-immersion}
\begin{reference}
\cite[EGA IV, Corollaire 21.12.7]{EGA4}
\end{reference}
Soit  $X$  un schéma localement noethérien. Soit  $U \subset X$  un
sous-schéma ouvert tel que le morphisme d'inclusion  $U \to X$  soit affine. Pour
tout point générique  $\xi$  d'une composante irréductible de  $X \setminus U$ ,
l'anneau local  $\mathcal{O}_{X, \xi}$  a une dimension  $\leq 1$. Si  $U$  est
dense ou si  $\xi$  se trouve dans l'adhérence de  $U$, alors
 $\dim(\mathcal{O}_{X, \xi}) = 1$.
\end{lemma}

\begin{proof}
Puisque  $\xi$  est un point générique de  $X \setminus U$, nous voyons que
$$
U_\xi = U \times_X \Spec(\mathcal{O}_{X, \xi}) \subset
\Spec(\mathcal{O}_{X, \xi})
$$
est le spectre poncturé de  $\mathcal{O}_{X, \xi}$  (indice: utilisez Schémas, Lemme
 \ref{schemes-lemma-specialize-points}). Comme  $U \to X$  est affine, nous voyons que  $U_\xi \to \Spec(\mathcal{O}_{X, \xi})$  est
affine (Morphismes, Lemme  \ref{morphisms-lemma-base-change-affine}) et nous concluons que  $U_\xi$  est
affine. Donc  $\dim(\mathcal{O}_{X, \xi}) \leq 1$  d'après le Lemme  \ref{lemma-affine-punctured-spec}. Si  $\xi \in \overline{U}$, alors il
existe une spécialisation  $\eta \to \xi$  où  $\eta \in U$  (il suffit de prendre
 $\eta$  un point générique d'une composante irréductible de  $\overline{U}$  qui
contient  $\xi$; puisque  $\overline{U}$  est localement noethérien, donc
localement a un nombre fini de composantes irréductibles, nous voyons que
 $\eta \in U$). Alors  $\eta \in \Spec(\mathcal{O}_{X, \xi})$  et nous voyons que la dimension ne peut pas être
 $0$.
\end{proof}

\begin{lemma}
\label{lemma-complement-affine-open}
Soit  $X$  un schéma séparé et localement noethérien. Soit  $U \subset X$  un
ouvert affine. Pour chaque point générique  $\xi$  d'une composante
irréductible de  $X \setminus U$ , l'anneau local  $\mathcal{O}_{X, \xi}$  a pour dimension
 $\leq 1$. Si  $U$  est dense ou si  $\xi$  est dans l'adhérence de
 $U$, alors  $\dim(\mathcal{O}_{X, \xi}) = 1$.
\end{lemma}

\begin{proof}
Ceci découle du Lemme  \ref{lemma-complement-affine-open-immersion}  car le morphisme  $U \to X$  est affine par
Morphismes, Lemme  \ref{morphisms-lemma-affine-permanence}.
\end{proof}

\noindent
Le lemme suivant peut parfois être utilisé pour produire des diviseurs de Cartier
effectifs.

\begin{lemma}
\label{lemma-complement-open-affine-effective-cartier-divisor}
Soit  $X$  un schéma séparé noethérien. Soit  $U \subset X$  un ouvert affine
dense. Si  $\mathcal{O}_{X, x}$  est un UFD pour tout  $x \in X \setminus U$, alors il existe un
diviseur de Cartier effectif  $D \subset X$  avec  $U = X \setminus D$.
\end{lemma}

\begin{proof}
Puisque  $X$  est noethérien, le complément  $X \setminus U$  a un nombre fini
de composantes irréductibles  $D_1, \ldots, D_r$  (Propriétés, Lemme  \ref{properties-lemma-Noetherian-irreducible-components}  appliqué
à la structure de sous-schéma réduit induite sur  $X \setminus U$). Chaque
 $D_i \subset X$  a une codimension  $1$  d'après le Lemme  \ref{lemma-complement-affine-open}  (et
Propriétés, Lemme  \ref{properties-lemma-codimension-local-ring}). Ainsi,  $D_i$  est un diviseur de Cartier
effectif d'après le Lemme  \ref{lemma-weil-divisor-is-cartier-UFD}. Par conséquent, nous pouvons prendre
 $D = D_1 + \ldots + D_r$.
\end{proof}

\begin{lemma}
\label{lemma-complement-open-affine-effective-cartier-divisor-bis}
Soit  $X$  un schéma noethérien avec diagonale affine. Soit  $U \subset X$  un
ouvert affine dense. Si  $\mathcal{O}_{X, x}$  est un UFD pour tout  $x \in X \setminus U$, alors il
existe un diviseur de Cartier effectif  $D \subset X$  avec  $U = X \setminus D$.
\end{lemma}

\begin{proof}
Puisque  $X$  est noethérien, le complément  $X \setminus U$  a un nombre fini
de composantes irréductibles  $D_1, \ldots, D_r$  (Propriétés, Lemme  \ref{properties-lemma-Noetherian-irreducible-components}  appliqué
à la structure de sous-schéma réduit induite sur  $X \setminus U$). Nous considérons
 $D_i$  comme un sous-schéma fermé réduit de  $X$. Soit  $X = \bigcup_{j \in J} X_j$ 
un recouvrement ouvert affine de  $X$. Pour tout  $j$  dans
 $J$, posons  $U_j = U \cap X_j$. Puisque  $X$  a une diagonale affine, le
schéma
$$
U_j = X \times_{(X \times X)} (U \times X_j)
$$
est affine. Par conséquent, comme  $X_j$  est séparé, il s'ensuit du Lemme
 \ref{lemma-complement-open-affine-effective-cartier-divisor}  et de sa démonstration que pour tous  $j \in J$  et  $1 \leq i \leq r$ ,
l'intersection  $D_i \cap X_j$  est soit vide, soit un diviseur de Cartier effectif
dans  $X_j$. Ainsi  $D_i \subset X$  est un diviseur de Cartier effectif (car il
s'agit d'une propriété locale). Par conséquent, nous pouvons prendre  $D = D_1 + \ldots + D_r$.
\end{proof}

\begin{lemma}
\label{lemma-regular-ample-family}
Soit  $X$  un schéma quasi-compact, régulier avec diagonale affine. Alors
 $X$  possède une famille abondante de modules inversibles (Morphismes,
Définition  \ref{morphisms-definition-family-ample-invertible-modules}).
\end{lemma}

\begin{proof}
Remarquons que  $X$  est une réunion finie disjointe de schémas intègres
(Propriétés, Lemmata  \ref{properties-lemma-regular-normal}  et  \ref{properties-lemma-normal-Noetherian}). Ainsi, on peut supposer que
 $X$  est intégral ainsi que Noethérien, régulier et ayant une diagonale
affine. Soit  $x \in X$. Choisissons un voisinage ouvert affine  $U \subset X$  de
 $x$. Puisque  $X$  est intégral,  $U$  est dense dans
 $X$. D'après More on Algebra, Lemme  \ref{more-algebra-lemma-regular-local-UFD} , les anneaux locaux de
 $X$  sont des UFD. Par conséquent, d'après le Lemme  \ref{lemma-complement-open-affine-effective-cartier-divisor-bis} , nous
pouvons trouver un diviseur de Cartier effectif  $D \subset X$  dont le complément
est  $U$. Alors la section canonique  $s = 1_D$  de  $\mathcal{L} = \mathcal{O}_X(D)$, voir la
Définition  \ref{definition-invertible-sheaf-effective-Cartier-divisor}, s'annule exactement le long de  $D$  donc
 $U = X_s$. Ainsi, les deux conditions dans Morphismes, Définition  \ref{morphisms-definition-family-ample-invertible-modules} 
sont satisfaites et nous avons terminé.
\end{proof}




\section{Normes}
\label{section-norms}

\noindent
Soit  $\pi : X \to Y$  un morphisme fini de schémas et soit  $d \geq 1$  un entier.
Disons qu'il existe un  {\it  norme de degré  $d$  pour
 $\pi$}\footnote{Ceci est une notation non
standard.}  s'il existe une application multiplicative
$$
\text{Norm}_\pi : \pi_*\mathcal{O}_X \to \mathcal{O}_Y
$$
de faisceaux telle que
\begin{enumerate}
\item la composition  $\mathcal{O}_Y \xrightarrow{\pi^\sharp} \pi_*\mathcal{O}_X
\xrightarrow{\text{Norm}_\pi} \mathcal{O}_Y$  égale  $g \mapsto g^d$, et
\item pour  $V \subset Y$  ouvert si  $f \in \mathcal{O}_X(\pi^{-1}V)$  est zéro en  $x \in \pi^{-1}(V)$, alors
 $\text{Norm}_\pi(f)$  est zéro en  $\pi(x)$.
\end{enumerate}
Nous observons que la condition (1) oblige  $\pi$  à être surjective. Puisque
 $\text{Norm}_\pi$  est multiplicatif, il envoie les unités sur des unités; par
conséquent, étant donné  $y \in Y$, si  $f$  est une fonction régulière
sur  $X$  définie mais non nulle en tout  $x \in X$  avec  $\pi(x) = y$,
alors  $\text{Norm}_\pi(f)$  est définie et ne s'annule pas en  $y$. Cela tient sans
nécessiter (2); en fait, les constructions dans cette section ne nécessiteront que
la condition (1) et seules certaines propriétés d'annulation (qui sont utilisées
en particulier dans la démonstration du Lemme  \ref{lemma-norm-ample}) nécessiteront la
propriété (2).

\begin{lemma}
\label{lemma-finite-trivialize-invertible-upstairs}
Soit  $\pi : X \to Y$  un morphisme fini de schémas. Soit  $\mathcal{L}$  un
 $\mathcal{O}_X$-module inversible. Soit  $y \in Y$. Il existe un voisinage ouvert
 $V \subset Y$  de  $y$  tel que  $\mathcal{L}|_{\pi^{-1}(V)}$  soit trivial.
\end{lemma}

\begin{proof}
De toute évidence, nous pouvons supposer que  $Y$  et donc  $X$ 
sont affines. Puisque  $\pi$  est fini, la fibre  $\pi^{-1}(\{y\})$  au-dessus de
 $y$  est finie. Puisque  $X$  est affine, nous pouvons choisir
 $s \in \Gamma(X, \mathcal{L})$  ne s'annulant en aucun point de  $\pi^{-1}(\{y\})$. Cela découle des
Propriétés, Lemme  \ref{properties-lemma-quasi-affine-invertible-nonvanishing-section} , mais nous donnons également un argument direct. À
savoir, nous pouvons choisir un ensemble fini  $E \subset X$  de points fermés tel
que chaque  $x \in \pi^{-1}(\{y\})$  se spécialise en un certain point de  $E$. Pour
 $x \in E$ , notons  $i_x : x \to X$  l'immersion fermée. Alors  $\mathcal{L} \to \bigoplus_{x \in E} i_{x, *}i_x^*\mathcal{L}$  est une
application surjective de  $\mathcal{O}_X$-modules quasi-cohérents, et donc
l'application
$$
\Gamma(X, \mathcal{L}) \to
\bigoplus\nolimits_{x \in E} \mathcal{L}_x/\mathfrak m_x\mathcal{L}_x
$$
est surjectif (puisque la prise de sections globales est un foncteur exact sur la
catégorie des  $\mathcal{O}_X$-modules quasi-cohérents, voir Schémas, Lemme
 \ref{schemes-lemma-equivalence-quasi-coherent}). Ainsi, nous pouvons trouver un  $s \in \Gamma(X, \mathcal{L})$  ne s'annulant en aucun
point se spécialisant en un point de  $E$. Alors  $X_s \subset X$  est un
voisinage ouvert de  $\pi^{-1}(\{y\})$. Puisque  $\pi$  est fini, donc fermé, nous
en concluons qu'il existe un voisinage ouvert  $V \subset Y$  de  $y$  dont
l'image réciproque est contenue dans  $X_s$  comme souhaité.
\end{proof}

\begin{lemma}
\label{lemma-norm-invertible}
Soit  $\pi : X \to Y$  un morphisme fini de schémas. S'il existe une norme de degré
 $d$  pour  $\pi$, alors il existe un homomorphisme de groupes
abéliens
$$
\text{Norm}_\pi : \Pic(X) \to \Pic(Y)
$$
tel que  $\text{Norm}_\pi(\pi^*\mathcal{N}) \cong \mathcal{N}^{\otimes d}$  pour tous les  $\mathcal{O}_Y$-modules inversibles  $\mathcal{N}$.
\end{lemma}

\begin{proof}
Nous utiliserons la correspondance entre les classes d'isomorphisme de modules
 $\mathcal{O}_X$ inversibles et les éléments de  $H^1(X, \mathcal{O}_X^*)$  donnée dans Cohomology,
Lemme  \ref{cohomology-lemma-h1-invertible}  sans autre mention. Nous expliquons comment prendre la norme
d'un module  $\mathcal{O}_X$ inversible  $\mathcal{L}$. À savoir, d'après le Lemme
 \ref{lemma-finite-trivialize-invertible-upstairs} , il existe un recouvrement ouvert  $Y = \bigcup V_j$  tel que  $\mathcal{L}|_{\pi^{-1}V_j}$ 
soit trivial. Choisissez une section génératrice  $s_j \in \mathcal{L}(\pi^{-1}V_j)$  pour chaque
 $j$. Sur les intersections  $\pi^{-1}V_j \cap \pi^{-1}V_{j'}$ , nous pouvons écrire
$$
s_j = u_{jj'} s_{j'}
$$
pour un  $u_{jj'} \in \mathcal{O}^*_X(\pi^{-1}V_j \cap \pi^{-1}V_{j'})$ unique. Ainsi, nous pouvons considérer les éléments
$$
v_{jj'} = \text{Norm}_\pi(u_{jj'}) \in \mathcal{O}_Y^*(V_j \cap V_{j'})
$$
Ces éléments satisfont la condition de cocycle (car  $u_{jj'}$  le fait et
 $\text{Norm}_\pi$  est multiplicatif) et définissent donc un  $\mathcal{O}_Y$-module
inversible. Nous omettons la vérification que: cela est bien défini, additif sur
les groupes de Picard, et satisfait la propriété  $\text{Norm}_\pi(\pi^*\mathcal{N}) \cong \mathcal{N}^{\otimes d}$  pour tous les
 $\mathcal{O}_Y$-modules inversibles  $\mathcal{N}$.
\end{proof}

\begin{lemma}
\label{lemma-norm-map-invertible}
Soit  $\pi : X \to Y$  un morphisme fini de schémas. Supposons qu'il existe une norme
de degré  $d$  pour  $\pi$. Pour toute application linéaire
 $\mathcal{O}_X$  $\varphi : \mathcal{L} \to \mathcal{L}'$  de  $\mathcal{O}_X$-modules inversibles, il existe une
application linéaire  $\mathcal{O}_Y$
$$
\text{Norm}_\pi(\varphi) :
\text{Norm}_\pi(\mathcal{L})
\longrightarrow
\text{Norm}_\pi(\mathcal{L}')
$$
avec  $\text{Norm}_\pi(\mathcal{L})$,  $\text{Norm}_\pi(\mathcal{L}')$  comme dans le Lemme  \ref{lemma-norm-invertible}. De plus, pour
 $y \in Y$ , les éléments suivants sont équivalents
\begin{enumerate}
\item $\varphi$  est nul en un point de  $x \in X$  avec  $\pi(x) = y$, et
\item $\text{Norm}_\pi(\varphi)$  est nul à  $y$.
\end{enumerate}
\end{lemma}

\begin{proof}
Nous choisissons un recouvrement ouvert  $Y = \bigcup V_j$  tel que  $\mathcal{L}$  et
 $\mathcal{L}'$  soient triviaux sur les ouverts  $\pi^{-1}V_j$. Cela est possible
d'après le Lemme  \ref{lemma-finite-trivialize-invertible-upstairs}. Choisissez des sections génératrices  $s_j$ 
et  $s'_j$  de  $\mathcal{L}$  et  $\mathcal{L}'$  sur les ouverts  $\pi^{-1}V_j$.
Alors  $\varphi(s_j) = f_js'_j$  pour certains  $f_j \in \mathcal{O}_X(\pi^{-1}V_j)$. Définissez  $\text{Norm}_\pi(\varphi)$  comme étant
la multiplication par  $\text{Norm}_\pi(f_j)$  sur  $V_j$. Un calcul simple impliquant
les cocycles utilisés pour construire  $\text{Norm}_\pi(\mathcal{L})$,  $\text{Norm}_\pi(\mathcal{L}')$  dans la
démonstration du Lemme  \ref{lemma-norm-invertible}  montre que cela définit une application comme
indiqué dans le lemme. L'énoncé final découle de la condition (2) dans la
définition d'une application norme de degré  $d$. Certains détails sont
omis.
\end{proof}

\begin{lemma}
\label{lemma-norm-ample}
Soit  $\pi : X \to Y$  un morphisme fini de schémas. Supposons que  $X$  possède
un faisceau inversible ample et qu'il existe une norme de degré  $d$  pour
 $\pi$. Alors  $Y$  possède un faisceau inversible ample.
\end{lemma}

\begin{proof}
Soit  $\mathcal{L}$  le faisceau inversible ample sur  $X$  qui nous est donné
par hypothèse. Nous prouverons que  $\mathcal{N} = \text{Norm}_\pi(\mathcal{L})$  est ample sur  $Y$.

\medskip\noindent
Puisque  $X$  est quasi-compact (Propriétés, Définition  \ref{properties-definition-ample}) et
que  $X \to Y$  est surjectif (par l'existence de  $\text{Norm}_\pi$), nous voyons que
 $Y$  est quasi-compact. Soit  $y \in Y$  un point. Pour terminer la
preuve, nous montrerons qu'il existe une section  $t$  d'une certaine
puissance tensorielle positive de  $\mathcal{N}$  qui ne s'annule pas en  $y$ 
de sorte que  $Y_t$  soit affine. Pour ce faire, choisissez un voisinage
ouvert affine  $V \subset Y$  de  $y$. Choisissez  $n \gg 0$  et une section
 $s \in \Gamma(X, \mathcal{L}^{\otimes n})$  telle que
$$
\pi^{-1}(\{y\}) \subset X_s \subset \pi^{-1}V
$$
par Propriétés, Lemme  \ref{properties-lemma-ample-finite-set-in-principal-affine}. Ensuite,  $t = \text{Norm}_\pi(s)$  est une section de
 $\mathcal{N}^{\otimes n}$  qui ne s'annule pas en  $x$  et avec  $Y_t \subset V$, voir Lemme
 \ref{lemma-norm-map-invertible}. Alors  $Y_t$  est affine par Propriétés, Lemme  \ref{properties-lemma-affine-cap-s-open}.
\end{proof}

\begin{lemma}
\label{lemma-norm-quasi-affine}
Soit  $\pi : X \to Y$  un morphisme fini de schémas. Supposons que  $X$  soit
quasi-affine et qu'il existe une norme de degré  $d$  pour  $\pi$.
Alors  $Y$  est quasi-affine.
\end{lemma}

\begin{proof}
D'après les Propriétés, Lemme  \ref{properties-lemma-quasi-affine-O-ample} , nous voyons que  $\mathcal{O}_X$  est un
faisceau inversible ample sur  $X$. La démonstration du Lemme  \ref{lemma-norm-ample} 
montre que  $\text{Norm}_\pi(\mathcal{O}_X) = \mathcal{O}_Y$  est un  $\mathcal{O}_Y$-module inversible ample. Par
conséquent, les Propriétés, Lemme  \ref{properties-lemma-quasi-affine-O-ample}  montrent que  $Y$  est
quasi-affine.
\end{proof}

\begin{lemma}
\label{lemma-finite-locally-free-has-norm}
Soit  $\pi : X \to Y$  un morphisme fini localement libre de degré  $d \geq 1$. Alors
il existe une norme canonique de degré  $d$  dont la formation commute avec
tout changement de base.
\end{lemma}

\begin{proof}
Soit  $V \subset Y$  un ouvert affine tel que  $(\pi_*\mathcal{O}_X)|_V$  soit fini libre de rang
 $d$. En choisissant une base, nous obtenons un isomorphisme
$$
\mathcal{O}_V^{\oplus d} \cong (\pi_*\mathcal{O}_X)|_V
$$
Pour chaque multiplication  $f \in \pi_*\mathcal{O}_X(V) = \mathcal{O}_X(\pi^{-1}(V))$  par  $f$  définit un endomorphisme
 $\mathcal{O}_V$-linéaire  $m_f$  du fibré vectoriel libre affiché. Ainsi, nous
obtenons une matrice  $d \times d$   $M_f \in \text{Mat}(d \times d, \mathcal{O}_Y(V))$  et nous pouvons définir
$$
\text{Norm}_\pi(f) = \det(M_f)
$$
. Comme le déterminant d'une matrice est indépendant du choix de la base choisie,
nous voyons que cela est bien défini, ce qui signifie également que cette
construction se collera pour donner une application globale comme souhaité. La
compatibilité avec le changement de base découle directement de la construction.

\medskip\noindent
La propriété (1) découle du fait que le déterminant d'une matrice diagonale
 $d \times d$  avec des entrées  $g, g, \ldots, g$  est  $g^d$. Pour voir la
propriété (2), nous pouvons changer de base et supposer que  $Y$  est le
spectre d'un corps  $k$. Puis  $X = \Spec(A)$  avec  $A$  une
 $k$-algèbre avec  $\dim_k(A) = d$. S'il existe un  $x \in X$  tel que
 $f \in A$  s'annule en  $x$, alors il existe une application
 $A \to \kappa$  dans un corps telle que  $f$  se mappe à zéro dans
 $\kappa$. Alors  $f : A \to A$  ne peut pas être surjectif, donc  $\det(f : A \to A) = 0$ 
comme souhaité.
\end{proof}

\begin{lemma}
\label{lemma-norm-in-normal-case}
Soit  $\pi : X \to Y$  un morphisme fini et surjectif avec  $X$  et  $Y$ 
intégraux et  $Y$  normal. Alors il existe une norme de degré  $[R(X) : R(Y)]$ 
pour  $\pi$.
\end{lemma}

\begin{proof}
Soit  $\Spec(B) \subset Y$  un ouvert affine et soit  $\Spec(A) \subset X$  son image réciproque.
Alors  $A$  et  $B$  sont des domaines. Soit  $K$  le corps
des fractions de  $A$  et  $L$  le corps des fractions de
 $B$. Image:
$$
\xymatrix{
L \ar[r] & K \\
B \ar[u] \ar[r] & A \ar[u]
}
$$
Étant donné que  $K/L$  est une extension finie, il existe une application
norme  $\text{Norm}_{K/L} : K^* \to L^*$  de degré  $d = [K : L]$; ceci est donné en mappant  $f \in K$ 
vers  $\det_L(f : K \to K)$  comme dans la démonstration du Lemme  \ref{lemma-finite-locally-free-has-norm}. Observe que
le polynôme caractéristique de  $f : K \to K$  est une puissance du polynôme minimal
de  $f$  sur  $L$; en particulier  $\text{Norm}_{K/L}(f)$  est une puissance du
coefficient constant du polynôme minimal de  $f$  sur  $L$. Par
conséquent, d’après l’Algèbre, Lemme  \ref{algebra-lemma-minimal-polynomial-normal-domain} ,  $\text{Norm}_{K/L}$  mappe  $A$ 
dans  $B$. Cela détermine un système compatible de fonctions sur les
sections sur des affines et donc une application norme globale  $\text{Norm}_\pi$  de
degré  $d$.

\medskip\noindent
La propriété (1) découle immédiatement de la construction. Pour voir la propriété
(2), soit  $f \in A$  contenu dans l'idéal premier  $\mathfrak p \subset A$. Soit
 $f^m + b_1 f^{m - 1} + \ldots + b_m$  le polynôme minimal de  $f$  sur  $L$. Par Algèbre,
Lemme  \ref{algebra-lemma-minimal-polynomial-normal-domain} , nous avons  $b_i \in B$. Donc  $b_0 \in B \cap \mathfrak p$. Puisque
 $\text{Norm}_{K/L}(f) = b_0^{d/m}$  (voir ci-dessus), nous concluons que la norme s'annule au point image
de  $\mathfrak p$.
\end{proof}

\begin{lemma}
\label{lemma-Frobenius-gives-norm-for-reduction}
Soit  $X$  un schéma noethérien. Soit  $p$  un nombre premier tel
que  $p\mathcal{O}_X = 0$. Alors pour un certain  $e > 0$ , il existe une norme de degré
 $p^e$  pour  $X_{red} \to X$  où  $X_{red}$  est la réduction de  $X$.
\end{lemma}

\begin{proof}
Soit  $A$  un anneau noethérien avec  $pA = 0$. Soit  $I \subset A$ 
l'idéal des éléments nilpotents. Alors  $I^n = 0$  pour un certain  $n$ 
(Algèbre, Lemme  \ref{algebra-lemma-Noetherian-power}). Choisissez  $e$  tel que  $p^e \geq n$. Alors
$$
A/I \longrightarrow A,\quad
f \bmod I \longmapsto f^{p^e}
$$
est bien défini. Cela produit une norme de degré  $p^e$  pour  $\Spec(A/I) \to \Spec(A)$.
Maintenant, si  $X$  est obtenu en recollant certains schémas affines
 $\Spec(A_i)$  alors pour un certain  $e \gg 0$  ces applications se recollent en
une application de norme pour  $X_{red} \to X$. Détails omis.
\end{proof}

\begin{proposition}
\label{proposition-push-down-ample}
Soit  $\pi : X \to Y$  un morphisme fini surjectif de schémas. Supposons que
 $X$  ait un  $\mathcal{O}_X$-module inversible ample. Si
\begin{enumerate}
\item $\pi$  est fini et localement libre, soit
\item $Y$  est un schéma normal intègre, ou
\item $Y$  est noethérien,  $p\mathcal{O}_Y = 0$, et  $X = Y_{red}$,
\end{enumerate}
alors  $Y$  a un  $\mathcal{O}_Y$-module inversible ample.
\end{proposition}

\begin{proof}
Le cas (1) découle d'une combinaison des lemmes  \ref{lemma-finite-locally-free-has-norm}  et  \ref{lemma-norm-ample}. Le
cas (3) découle d'une combinaison des lemmes  \ref{lemma-Frobenius-gives-norm-for-reduction}  et  \ref{lemma-norm-ample}. Dans le
cas (2), nous remplaçons d'abord  $X$  par une composante irréductible de
 $X$  qui domine  $Y$  (considéré comme un sous-schéma fermé réduit
de  $X$). Ensuite, nous pouvons appliquer le lemme  \ref{lemma-norm-in-normal-case}.
\end{proof}

\begin{lemma}
\label{lemma-push-down-quasi-affine}
Soit  $\pi : X \to Y$  un morphisme fini surjectif de schémas. Supposons que
 $X$  soit quasi-affine. Si soit
\begin{enumerate}
\item $\pi$  est fini et localement libre, soit
\item $Y$  est un schéma intègre normal
\end{enumerate}
alors  $Y$  est quasi-affine.
\end{lemma}

\begin{proof}
Le cas (1) découle d'une combinaison des lemmes  \ref{lemma-finite-locally-free-has-norm}  et  \ref{lemma-norm-quasi-affine}. Dans
le cas (2), nous remplaçons d'abord  $X$  par une composante irréductible
de  $X$  qui domine  $Y$  (considéré comme un sous-schéma fermé
réduit de  $X$). Ensuite, nous pouvons appliquer le lemme  \ref{lemma-norm-in-normal-case}.
\end{proof}





\section{Diviseurs de Cartier effectifs relatifs}
\label{section-effective-Cartier-morphisms}

\noindent
Le lemme suivant montre qu'un diviseur de Cartier effectif qui est plat sur la
base est réellement une « famille de diviseurs de Cartier effectifs » sur la base.
Par exemple, la restriction à toute fibre est un diviseur de Cartier effectif.

\begin{lemma}
\label{lemma-relative-Cartier}
Soit  $f : X \to S$  un morphisme de schémas. Soit  $D \subset X$  un sous-schéma fermé.
Supposons
\begin{enumerate}
\item $D$  est un diviseur de Cartier effectif, et
\item $D \to S$  est un morphisme plat.
\end{enumerate}
Puis, pour tout morphisme de schémas  $g : S' \to S$ , le tiré en arrière  $(g')^{-1}D$ 
est un diviseur de Cartier effectif sur  $X' = S' \times_S X$  où  $g' : X' \to X$  est la
projection.
\end{lemma}

\begin{proof}
En utilisant le Lemme  \ref{lemma-characterize-effective-Cartier-divisor} , nous traduisons ceci comme suit en algèbre.
Soit  $A \to B$  une application d'anneaux et  $h \in B$. Supposons que
 $h$  soit un non-diviseur de zéro et que  $B/hB$  soit plat sur
 $A$. Alors
$$
0 \to B \xrightarrow{h} B \to B/hB \to 0
$$
est une suite exacte courte de  $A$-modules avec  $B/hB$  plat sur
 $A$. D'après le Lemme d'Algèbre  \ref{algebra-lemma-flat-tor-zero} , cette suite reste exacte
lors du produit tensoriel sur  $A$  avec n'importe quel module, en
particulier avec n'importe quelle  $A$-algèbre  $A'$.
\end{proof}

\noindent
Ce lemme est la motivation pour la définition suivante.

\begin{definition}
\label{definition-relative-effective-Cartier-divisor}
Soit  $f : X \to S$  un morphisme de schémas. Un  {\it  diviseur de
Cartier effectif relatif}  sur  $X/S$  est un diviseur de Cartier
effectif  $D \subset X$  tel que  $D \to S$  soit un morphisme plat de schémas.
\end{definition}

\noindent
Nous avertissons le lecteur que cela peut être une notation non standard. En
particulier, dans  \cite[IV, Section 21.15]{EGA}  la notion de diviseur relatif est discutée
uniquement lorsque  $X \to S$  est plat et localement de présentation finie.
Notre définition est un peu plus générale. Cependant, il s'avère que si
 $x \in D$  alors  $X \to S$  est plat à  $x$  dans de nombreux cas (mais
pas toujours).

\begin{lemma}
\label{lemma-sum-relative-effective-Cartier-divisor}
Soit  $f : X \to S$  un morphisme de schémas. Si  $D_1, D_2 \subset X$  sont des diviseurs de
Cartier effectifs relatifs sur  $X/S$  alors il en est de même de
 $D_1 + D_2$  (Définition  \ref{definition-sum-effective-Cartier-divisors}).
\end{lemma}

\begin{proof}
Cela se traduit par le fait algébrique suivant: Soit  $A \to B$  une application
d'anneaux et  $h_1, h_2 \in B$. Supposons que les  $h_i$  sont des diviseurs de
zéro et que  $B/h_iB$  est plat sur  $A$. Alors  $h_1h_2$  est un
non-diviseur de zéro et  $B/h_1h_2B$  est plat sur  $A$. La raison est que
nous avons une suite exacte courte
$$
0 \to B/h_1B \to B/h_1h_2B \to B/h_2B \to 0
$$
où la première flèche est donnée par la multiplication par  $h_2$. Comme les
deux modules extérieurs sont plats sur  $A$, il en est de même pour le
module du milieu, voir Algèbre, Lemme  \ref{algebra-lemma-flat-ses}.
\end{proof}

\begin{lemma}
\label{lemma-difference-relative-effective-Cartier-divisor}
Soit  $f : X \to S$  un morphisme de schémas. Si  $D_1, D_2 \subset X$  sont des diviseurs de
Cartier effectifs relatifs sur  $X/S$  et  $D_1 \subset D_2$  comme sous-schémas
fermés, alors le diviseur de Cartier effectif  $D$  tel que  $D_2 = D_1 + D$ 
(Lemme  \ref{lemma-difference-effective-Cartier-divisors}) est un diviseur de Cartier relatif effectif sur  $X/S$.
\end{lemma}

\begin{proof}
Cela se traduit par le fait algébrique suivant: Soit  $A \to B$  un homomorphisme
d'anneaux et  $h_1, h_2 \in B$. Supposons que les  $h_i$  soient des diviseurs de
zéro non nuls, que  $B/h_iB$  soit plat sur  $A$, et que  $(h_2) \subset (h_1)$.
Alors nous pouvons écrire  $h_2 = h h_1$  où  $h \in B$  est un diviseur de zéro non
nul. Nous obtenons une suite exacte courte
$$
0 \to B/hB \to B/h_2B \to B/h_1B \to 0
$$
où la première flèche est donnée par la multiplication par  $h_1$. Puisque
les deux modules de droite sont plats sur  $A$, il en est de même pour
celui du milieu, voir Algèbre, Lemme  \ref{algebra-lemma-flat-ses}.
\end{proof}

\begin{lemma}
\label{lemma-flat-at-x}
Soit  $f : X \to S$  un morphisme de schémas. Soit  $D \subset X$  un diviseur de
Cartier relatif effectif sur  $X/S$. Si  $x \in D$  et  $\mathcal{O}_{X, x}$  sont
noethériens, alors  $f$  est plat en  $x$.
\end{lemma}

\begin{proof}
Posons  $A = \mathcal{O}_{S, f(x)}$  et  $B = \mathcal{O}_{X, x}$. Soit  $h \in B$  un élément qui génère
l'idéal de  $D$. Alors  $h$  est un non-diviseur de zéro dans
 $B$  tel que  $B/hB$  soit une  $A$-algèbre locale plate. Soit
 $I \subset A$  un idéal de type fini. Considérons le diagramme commutatif
$$
\xymatrix{
0 \ar[r] &
B \ar[r]_h &
B \ar[r] &
B/hB \ar[r] & 0 \\
0 \ar[r] &
B \otimes_A I \ar[r]^h \ar[u] &
B \otimes_A I \ar[r] \ar[u] &
B/hB \otimes_A I \ar[r] \ar[u] & 0
}
$$
. La suite inférieure est exacte courte car  $B/hB$  est plat sur
 $A$, voir Algèbre, Lemme  \ref{algebra-lemma-flat-tor-zero}. La flèche verticale de droite est
injective car  $B/hB$  est plat sur  $A$, voir Algèbre, Lemme
 \ref{algebra-lemma-flat}. Par conséquent, la multiplication par  $h$  est surjective
sur le noyau  $K$  de la flèche verticale du milieu. Par le lemme de
Nakayama, voir Algèbre, Lemme  \ref{algebra-lemma-NAK} , nous concluons que  $K= 0$. Par
conséquent,  $B$  est plat sur  $A$, voir Algèbre, Lemme
 \ref{algebra-lemma-flat}.
\end{proof}

\noindent
Le lemme suivant repose sur la version algébrique de l'ouverture du lieu plat. La
version schématique peut être trouvée dans More on Morphismes, Section
 \ref{more-morphisms-section-open-flat}.

\begin{lemma}
\label{lemma-flat-relative-Cartier-divisor}
Soit  $f : X \to S$  un morphisme de schémas. Soit  $D \subset X$  un diviseur de
Cartier relatif effectif. Si  $f$  est localement de présentation finie,
alors il existe un sous-schéma ouvert  $U \subset X$  tel que  $D \subset U$  et tel que
 $f|_U : U \to S$  soit plat.
\end{lemma}

\begin{proof}
Choisissez  $x \in D$. Il suffit de trouver un voisinage ouvert  $U \subset X$  de
 $x$  tel que  $f|_U$  soit plat. Ainsi, le lemme se réduit au cas où
 $X = \Spec(B)$  et  $S = \Spec(A)$  sont affines et où  $D$  est donné par un
élément non diviseur de zéro  $h \in B$. Par hypothèse,  $B$  est une
 $A$-algèbre de présentation finie et  $B/hB$  est une
 $A$-algèbre plate. Nous allons utiliser l'approximation absolue de
Noether.

\medskip\noindent
Écrire  $B = A[x_1, \ldots, x_n]/(g_1, \ldots, g_m)$. Supposons que  $h$  est l'image de  $h' \in A[x_1, \ldots, x_n]$.
Choisir une  $\mathbf{Z}$-sous-algèbre de type fini  $A_0 \subset A$  telle que tous les
coefficients des polynômes  $h', g_1, \ldots, g_m$  soient dans  $A_0$. Alors nous
pouvons définir  $B_0 = A_0[x_1, \ldots, x_n]/(g_1, \ldots, g_m)$  et  $h_0$  l'image de  $h'$  dans
 $B_0$. Alors  $B = B_0 \otimes_{A_0} A$  et  $B/hB = B_0/h_0B_0 \otimes_{A_0} A$. D'après le Lemme  \ref{algebra-lemma-flat-finite-presentation-limit-flat} 
d'Algèbre, nous pouvons, après avoir élargi  $A_0$, supposer que
 $B_0/h_0B_0$  est plat sur  $A_0$. Soit  $K_0 = \Ker(h_0 : B_0 \to B_0)$. Comme  $B_0$  est
de type fini sur  $\mathbf{Z}$ , nous voyons que  $K_0$  est un idéal de type
fini. Soit  $A_1 \subset A$  une  $\mathbf{Z}$-sous-algèbre de type fini contenant
 $A_0$  et notons  $B_1$,  $h_1$,  $K_1$  les objets
correspondants sur  $A_1$. D'après More on Algèbre, Lemme  \ref{more-algebra-lemma-base-change-H1-regular} ,
l'application  $K_0 \otimes_{A_0} A_1 \to K_1$  est surjective. D'autre part, le noyau de  $h : B \to B$ 
est nul par hypothèse. Ainsi, chaque élément de  $K_0$  se mappe sur zéro
dans  $K_1$  pour des sous-anneaux  $A_1 \subset A$ suffisamment grands. Comme
 $K_0$  est de type fini, nous en concluons que  $K_1 = 0$  pour un choix
approprié de  $A_1$.

\medskip\noindent
Mettez  $f_1 : X_1 \to S_1$  égal à  $\Spec$  de l'anneau carte  $A_1 \to B_1$. Mettez
 $D_1 = \Spec(B_1/h_1B_1)$. Puisque  $B = B_1 \otimes_{A_1} A$, c'est-à-dire  $X = X_1 \times_{S_1} S$, il suffit
maintenant de prouver le lemme pour  $X_1 \to S_1$  et le diviseur de Cartier
effectif relatif  $D_1$, voir Morphismes, Lemme  \ref{morphisms-lemma-base-change-module-flat}. Ainsi, nous
avons réduit au cas où  $A$  est un anneau noethérien. Dans ce cas, nous
savons que le morphisme d'anneaux  $A \to B$  est plat en chaque premier
 $\mathfrak q$  de  $V(h)$  d'après le Lemme  \ref{lemma-flat-at-x}. Combiné avec le fait
que le lieu de platitude est ouvert dans ce cas, voir Algèbre, Théorème
 \ref{algebra-theorem-openness-flatness} , nous réussissons.
\end{proof}

\noindent
Il y a aussi le lemme suivant (dont l'idée est apparemment due à Michael Artin,
voir  \cite{Nobile}) qui n'a besoin d'aucune hypothèse de finitude.

\begin{lemma}
\label{lemma-michael-artin}
Soit  $f : X \to S$  un morphisme de schémas. Soit  $D \subset X$  un diviseur de
Cartier effectif relatif sur  $X/S$. Si  $f$  est plat en tous les
points de  $X \setminus D$, alors  $f$  est plat.
\end{lemma}

\begin{proof}
Cela se traduit par le fait algébrique suivant: Soit  $A \to B$  un morphisme
d’anneaux et  $h \in B$. Supposons que  $h$  soit un non-diviseur de
zéro, que  $B/hB$  soit plat sur  $A$, et que la localisation
 $B_h$  soit plate sur  $A$. Alors  $B$  est plat sur
 $A$. La raison est que nous avons une suite exacte courte
$$
0 \to B \to B_h \to \colim_n (1/h^n)B/B \to 0
$$
et que le deuxième et le troisième termes sont plats sur  $A$, ce qui
implique que  $B$  est plat sur  $A$  (voir Algèbre, Lemme
 \ref{algebra-lemma-flat-ses}). Notez qu'une colimite filtrée de modules plats est plate (voir
Algèbre, Lemme  \ref{algebra-lemma-colimit-flat}) et que par induction sur  $n$  chaque
 $(1/h^n)B/B \cong B/h^nB$  est plat sur  $A$  puisqu'il s'insère dans la suite exacte
courte
$$
0 \to B/h^{n - 1}B \xrightarrow{h} B/h^nB \to B/hB \to 0
$$
Quelques détails omis.
\end{proof}

\begin{example}
\label{example-relative-cartier-ambient-space-not-flat}
Voici un exemple de diviseur de Cartier relatif effectif  $D$  où le schéma
ambiant n'est pas plat dans un voisinage de  $D$. À savoir, soit
 $A = k[t]$  et
$$
B = k[t, x, y, x^{-1}y, x^{-2}y, \ldots]/(ty, tx^{-1}y, tx^{-2}y, \ldots)
$$
Alors  $B$  n'est pas plat sur  $A$  mais  $B/xB \cong A$  est plat sur
 $A$. De plus,  $x$  est un non-diviseur de zéro et définit donc un
diviseur de Cartier relatif effectif dans  $\Spec(B)$  sur  $\Spec(A)$.
\end{example}

\noindent
Si le schéma ambiant est plat et localement de présentation finie sur la base,
alors nous pouvons caractériser un diviseur de Cartier relatif effectif en termes
de ses fibres. Voir aussi Compléments sur les morphismes, Lemme  \ref{more-morphisms-lemma-slice-given-element}  pour une approche
légèrement différente de ce lemme.

\begin{lemma}
\label{lemma-fibre-Cartier}
Soit  $\varphi : X \to S$  un morphisme plat qui est localement de présentation finie. Soit
 $Z \subset X$  un sous-schéma fermé. Soit  $x \in Z$  avec une image  $s \in S$.
\begin{enumerate}
\item Si  $Z_s \subset X_s$  est un diviseur de Cartier dans un voisinage de  $x$,
alors il existe un ouvert  $U \subset X$  et un diviseur de Cartier relatif effectif
 $D \subset U$  tels que  $Z \cap U \subset D$  et  $Z_s \cap U = D_s$.
\item Si  $Z_s \subset X_s$  est un diviseur de Cartier dans un voisinage de  $x$, que
le morphisme  $Z \to X$  est de présentation finie, et que  $Z \to S$  est plat
en  $x$, alors nous pouvons choisir  $U$  et  $D$  tels que
 $Z \cap U = D$.
\item Si  $Z_s \subset X_s$  est un diviseur de Cartier dans un voisinage de  $x$  et
 $Z$  est un sous-schéma fermé localement principal de  $X$  dans un
voisinage de  $x$, alors nous pouvons choisir  $U$  et  $D$ 
tels que  $Z \cap U = D$.
\end{enumerate}
En particulier, si  $Z \to S$  est localement de présentation finie et plat et
que toutes les fibres  $Z_s \subset X_s$  sont des diviseurs de Cartier effectifs, alors
 $Z$  est un diviseur de Cartier effectif relatif. De même, si  $Z$ 
est un sous-schéma fermé localement principal de  $X$  tel que toutes les
fibres  $Z_s \subset X_s$  sont des diviseurs de Cartier effectifs, alors  $Z$ 
est un diviseur de Cartier effectif relatif.
\end{lemma}

\begin{proof}
Choisissez des voisinages ouverts affines  $\Spec(A)$  de  $s$  et
 $\Spec(B)$  de  $x$  tels que  $\varphi(\Spec(B)) \subset \Spec(A)$. Soit  $\mathfrak p \subset A$  l'idéal
premier correspondant à  $s$. Soit  $\mathfrak q \subset B$  l'idéal premier
correspondant à  $x$. Soit  $I \subset B$  l'idéal correspondant à
 $Z$. D'après l'hypothèse initiale du lemme, nous savons que  $A \to B$ 
est plat et de présentation finie. L'hypothèse en (1) signifie que, après avoir
réduit  $\Spec(B)$, nous pouvons supposer que  $I(B \otimes_A \kappa(\mathfrak p))$  est engendré par un
seul élément qui est un diviseur de zéro dans  $B \otimes_A \kappa(\mathfrak p)$. Disons que
 $f \in I$  se mappe sur ce générateur. Nous affirmons qu'après avoir inversé un
élément  $g \in B$,  $g \not \in \mathfrak q$  le sous-schéma fermé  $D = V(f) \subset \Spec(B_g)$  est un
diviseur de Cartier effectif relatif.

\medskip\noindent
D'après Algèbre, Lemme  \ref{algebra-lemma-flat-finite-presentation-limit-flat} , nous pouvons trouver un morphisme d'anneaux
de type fini plat  $A_0 \to B_0$  d'anneaux noethériens, un élément  $f_0 \in B_0$, un
morphisme d'anneaux  $A_0 \to A$  et un isomorphisme  $A \otimes_{A_0} B_0 \cong B$. Si  $\mathfrak p_0 = A_0 \cap \mathfrak p$ 
alors nous voyons que
$$
B \otimes_A \kappa(\mathfrak p) =
\left(B_0 \otimes_{A_0} \kappa(\mathfrak p_0)\right)
\otimes_{\kappa(\mathfrak p_0))} \kappa(\mathfrak p)
$$
, d'où  $f_0$  est un non-diviseur de zéro dans  $B_0 \otimes_{A_0} \kappa(\mathfrak p_0)$. D'après
Algèbre, Lemme  \ref{algebra-lemma-grothendieck} , nous voyons que  $f_0$  est un non-diviseur de
zéro dans  $(B_0)_{\mathfrak q_0}$  où  $\mathfrak q_0 = B_0 \cap \mathfrak q$  et que  $(B_0/f_0B_0)_{\mathfrak q_0}$  est plat sur
 $A_0$. Donc d'après Algèbre, Lemme  \ref{algebra-lemma-regular-sequence-in-neighbourhood}  et Théorème d'Algèbre
 \ref{algebra-theorem-openness-flatness} , il existe un  $g_0 \in B_0$,  $g_0 \not \in \mathfrak q_0$  tel que  $f_0$  est un
non-diviseur de zéro dans  $(B_0)_{g_0}$  et tel que  $(B_0/f_0B_0)_{g_0}$  est plat sur
 $A_0$. Par conséquent, nous voyons que  $D_0 = V(f_0) \subset \Spec((B_0)_{g_0})$  est un diviseur de
Cartier relatif effectif. Puisque nous savons que cette propriété est préservée
par changement de base, voir le Lemme  \ref{lemma-relative-Cartier}, nous obtenons l'affirmation
mentionnée ci-dessus avec  $g$  égal à l'image de  $g_0$  dans
 $B$.

\medskip\noindent
À ce stade, nous avons démontré (1). Pour voir (2), considérons l'immersion fermée
 $Z \to D$. L'application d'anneaux surjective  $u : \mathcal{O}_{D, x} \to \mathcal{O}_{Z, x}$  est une application
d'algèbres locales plates sur  $\mathcal{O}_{S, s}$, essentiellement de présentation finie,
et qui devient un isomorphisme après division par  $\mathfrak m_s$. Par conséquent,
c'est un isomorphisme, voir Algèbre, Lemme  \ref{algebra-lemma-mod-injective-general}. Il s'ensuit que
 $Z \to D$  est un isomorphisme dans un voisinage de  $x$, voir Algèbre,
Lemme  \ref{algebra-lemma-local-isomorphism}. Pour voir (3), après avoir éventuellement rétréci
 $U$ , nous pouvons supposer que l'idéal de  $D$  est engendré par
un seul non-diviseur de zéro  $f$  et que l'idéal de  $Z$  est
engendré par un élément  $g$. Alors  $f = gh$. Mais  $g|_{U_s}$  et
 $f|_{U_s}$  définissent le même diviseur de Cartier effectif dans un voisinage de
 $x$. Par conséquent,  $h|_{X_s}$  est une unité dans  $\mathcal{O}_{X_s, x}$, donc
 $h$  est une unité dans  $\mathcal{O}_{X, x}$  et donc  $h$  est une unité
dans un voisinage ouvert de  $x$. C.-à-d.,  $Z \cap U = D$  après avoir
rétréci  $U$.

\medskip\noindent
Les affirmations finales du lemme suivent immédiatement des parties (2) et (3),
combinées avec le fait que  $Z \to S$  est localement de présentation finie si et
seulement si  $Z \to X$  est de présentation finie, voir Morphismes, Lemme
 \ref{morphisms-lemma-composition-finite-presentation}  et  \ref{morphisms-lemma-finite-presentation-permanence}.
\end{proof}



\section{Le cône normal d'une immersion}
\label{section-normal-cone}

\noindent
Laisse $i : Z \to X$ être une immersion fermée. Soit $\mathcal{I} \subset \mathcal{O}_X$ soit le faisceau
quasi-cohérent d'idéaux correspondant. Considérons le faisceau quasi-cohérent
gradué $\mathcal{O}_X$-algèbres $\bigoplus_{n \geq 0} \mathcal{I}^n/\mathcal{I}^{n + 1}$. Comme les faisceaux $\mathcal{I}^n/\mathcal{I}^{n + 1}$ sont tous
annulés par $\mathcal{I}$ cette algèbre graduée correspond à un faisceau
quasi-cohérent gradué $\mathcal{O}_Z$-algèbres par morphismes, Lemme \ref{morphisms-lemma-i-star-equivalence}. Ce
gradué quasi-cohérent $\mathcal{O}_Z$-algèbre est appelée la
 {\it algèbre conormale
de $Z$ dans $X$}  et est souvent simplement désigné
 $\bigoplus_{n \geq 0} \mathcal{I}^n/\mathcal{I}^{n + 1}$ par l'abus de notation mentionné dans Morphismes, Section \ref{morphisms-section-closed-immersions-quasi-coherent}.

\medskip\noindent
Soit  $f : Z \to X$  une immersion. Nous définissons l'algèbre conormale de
 $f$  comme l'algèbre conormale de l'immersion fermée  $i : Z \to X \setminus \partial Z$, où
 $\partial Z = \overline{Z} \setminus Z$. Elle est souvent notée  $\bigoplus_{n \geq 0} \mathcal{I}^n/\mathcal{I}^{n + 1}$  où  $\mathcal{I}$  est le faisceau
d'idéaux de l'immersion fermée  $i : Z \to X \setminus \partial Z$.

\begin{definition}
\label{definition-conormal-sheaf}
Laisse $f : Z \to X$ être une immersion. Le {\it algèbre
conormale $\mathcal{C}_{Z/X, *}$  de  $Z$ dans $X$}  ou le
 {\it algèbre conormale de $f$} est le
faisceau quasi-cohérent gradué $\mathcal{O}_Z$-algèbres $\bigoplus_{n \geq 0} \mathcal{I}^n/\mathcal{I}^{n + 1}$  décrit ci-dessus.
\end{definition}

\noindent
Ainsi,  $\mathcal{C}_{Z/X, 1} = \mathcal{C}_{Z/X}$  est le faisceau conormal de l'immersion. De plus,  $\mathcal{C}_{Z/X, 0} = \mathcal{O}_Z$ 
et  $\mathcal{C}_{Z/X, n}$  est un  $\mathcal{O}_Z$-module quasi-cohérent caractérisé par la
propriété
\begin{equation}
\label{equation-conormal-in-degree-n}
i_*\mathcal{C}_{Z/X, n} = \mathcal{I}^n/\mathcal{I}^{n + 1}
\end{equation}
où  $i : Z \to X \setminus \partial Z$  et  $\mathcal{I}$  est le faisceau d'idéaux de  $i$  comme
ci-dessus. Enfin, noter qu'il existe un morphisme surjectif canonique
\begin{equation}
\label{equation-conormal-algebra-quotient}
\text{Sym}^*(\mathcal{C}_{Z/X}) \longrightarrow \mathcal{C}_{Z/X, *}
\end{equation}
d'algèbres  $\mathcal{O}_Z$ graduées quasi-cohérentes qui est un isomorphisme dans les
degrés  $0$  et  $1$.

\begin{lemma}
\label{lemma-affine-conormal-sheaf}
Soit  $i : Z \to X$  une immersion. L'algèbre conormale de  $i$  possède les
propriétés suivantes:
\begin{enumerate}
\item Soit  $U \subset X$  un ouvert quelconque tel que  $i(Z)$  est un sous-ensemble
fermé de  $U$. Soit  $\mathcal{I} \subset \mathcal{O}_U$  le faisceau d'idéaux correspondant au
sous-schéma fermé  $i(Z) \subset U$. Alors
$$
\mathcal{C}_{Z/X, *} =
i^*\left(\bigoplus\nolimits_{n \geq 0} \mathcal{I}^n\right) =
i^{-1}\left(
\bigoplus\nolimits_{n \geq 0} \mathcal{I}^n/\mathcal{I}^{n + 1}
\right)
$$
\item Pour tout ouvert affine  $\Spec(R) = U \subset X$  tel que  $Z \cap U = \Spec(R/I)$  il existe un
isomorphisme canonique  $\Gamma(Z \cap U, \mathcal{C}_{Z/X, *}) = \bigoplus_{n \geq 0} I^n/I^{n + 1}$.
\end{enumerate}
\end{lemma}

\begin{proof}
Principalement clair d'après les définitions. Noter que donné un anneau
 $R$  et un idéal  $I$  de  $R$  nous avons  $I^n/I^{n + 1} = I^n \otimes_R R/I$.
Détails omis.
\end{proof}

\begin{lemma}
\label{lemma-conormal-algebra-functorial}
Soit
$$
\xymatrix{
Z \ar[r]_i \ar[d]_f & X \ar[d]^g \\
Z' \ar[r]^{i'} & X'
}
$$
soit un diagramme commutatif dans la catégorie des schémas. Supposons
 $i$,  $i'$  immersions. Il existe un morphisme canonique d'algèbres
 $\mathcal{O}_Z$ graduées
$$
f^*\mathcal{C}_{Z'/X', *}
\longrightarrow
\mathcal{C}_{Z/X, *}
$$
caractérisé par la propriété suivante: Pour chaque paire d'ouverts affines
 $(\Spec(R) = U \subset X, \Spec(R') = U' \subset X')$  avec  $g(U) \subset U'$  tels que  $Z \cap U = \Spec(R/I)$  et  $Z' \cap U' = \Spec(R'/I')$  la flèche
induite
$$
\Gamma(Z' \cap U', \mathcal{C}_{Z'/X', *}) =
\bigoplus\nolimits (I')^n/(I')^{n + 1}
\longrightarrow
\bigoplus\nolimits_{n \geq 0} I^n/I^{n + 1} =
\Gamma(Z \cap U, \mathcal{C}_{Z/X, *})
$$
est celle induite par le morphisme d'anneaux  $f^\sharp : R' \to R$  qui a la propriété
 $f^\sharp(I') \subset I$.
\end{lemma}

\begin{proof}
Soient  $\partial Z' = \overline{Z'} \setminus Z'$  et  $\partial Z = \overline{Z} \setminus Z$. Ce sont des sous-ensembles fermés de
 $X'$  et de  $X$. En remplaçant  $X'$  par  $X' \setminus \partial Z'$  et
 $X$  par  $X \setminus \big(g^{-1}(\partial Z') \cup \partial Z\big)$ , nous voyons que nous pouvons supposer que
 $i$  et  $i'$  sont des immersions fermées.

\medskip\noindent
Le fait que  $g \circ i$  se factorise à travers  $i'$  implique que
 $g^*\mathcal{I}'$  se projette dans  $\mathcal{I}$  sous l'application canonique
 $g^*\mathcal{I}' \to \mathcal{O}_X$, voir Schémas, Lemmata  \ref{schemes-lemma-characterize-closed-subspace}  et  \ref{schemes-lemma-restrict-map-to-closed}. Par conséquent,
nous obtenons une application induite de faisceaux quasi-cohérents  $g^*((\mathcal{I}')^n/(\mathcal{I}')^{n + 1}) \to
\mathcal{I}^n/\mathcal{I}^{n + 1}$.
Le tiré en arrière par  $i$  donne  $i^*g^*((\mathcal{I}')^n/(\mathcal{I}')^{n + 1}) \to
i^*(\mathcal{I}^n/\mathcal{I}^{n + 1})$. Notez que  $i^*(\mathcal{I}^n/\mathcal{I}^{n + 1}) = \mathcal{C}_{Z/X, n}$.
D'autre part,  $i^*g^*((\mathcal{I}')^n/(\mathcal{I}')^{n + 1}) =
f^*(i')^*((\mathcal{I}')^n/(\mathcal{I}')^{n + 1}) =
f^*\mathcal{C}_{Z'/X', n}$. Cela donne l'application désirée.

\medskip\noindent
Vérifier que l'application est localement décrite comme l'application donnée
 $(I')^n/(I')^{n + 1} \to I^n/I^{n + 1}$  est une question de déroulement des définitions et est omis. Une
autre observation est que pour tout  $x \in i(Z)$  donné, il existe effectivement
des voisinages ouverts affines  $U$,  $U'$  avec  $f(U) \subset U'$  et
 $Z \cap U$  ainsi que  $U' \cap Z'$  fermés tels que  $x \in U$. Preuve omise.
Par conséquent, l'exigence du lemme caractérise effectivement l'application (et
aurait pu être utilisée pour la définir).
\end{proof}

\begin{lemma}
\label{lemma-conormal-algebra-functorial-flat}
Soit
$$
\xymatrix{
Z \ar[r]_i \ar[d]_f & X \ar[d]^g \\
Z' \ar[r]^{i'} & X'
}
$$
soit un diagramme de produits fibrés dans la catégorie des schémas avec
 $i$,  $i'$  immersions. Alors l'application canonique  $f^*\mathcal{C}_{Z'/X', *} \to \mathcal{C}_{Z/X, *}$ 
du Lemme  \ref{lemma-conormal-algebra-functorial}  est surjective. Si  $g$  est plat, alors c'est un
isomorphisme.
\end{lemma}

\begin{proof}
Soit  $R' \to R$  un morphisme d'anneaux, et  $I' \subset R'$  un idéal. Posons
 $I = I'R$. Alors  $(I')^n/(I')^{n + 1} \otimes_{R'} R \to I^n/I^{n + 1}$  est surjective. Si  $R' \to R$  est plat, alors
 $I^n = (I')^n \otimes_{R'} R$  et nous voyons que la carte est un isomorphisme.
\end{proof}

\begin{definition}
\label{definition-normal-cone}
Soit  $i : Z \to X$  une immersion de schémas. Le cône normal
 {\it   $C_ZX$}  de  $Z$  dans
 $X$  est
$$
C_ZX = \underline{\Spec}_Z(\mathcal{C}_{Z/X, *})
$$
voir Constructions, Définitions  \ref{constructions-definition-cone}  et  \ref{constructions-definition-abstract-cone}. Le fibré normal
 {\it }  de  $Z$  dans  $X$  est le
fibré vectoriel
$$
N_ZX = \underline{\Spec}_Z(\text{Sym}(\mathcal{C}_{Z/X}))
$$
voir Constructions, Définitions  \ref{constructions-definition-vector-bundle}  et  \ref{constructions-definition-abstract-vector-bundle}.
\end{definition}

\noindent
Ainsi,  $C_ZX \to Z$  est un cône sur  $Z$  et  $N_ZX \to Z$  est un fibré
vectoriel sur  $Z$  (rappelons que dans notre terminologie, cela n'implique
pas que le faisceau conormal soit un faisceau localement libre de type fini). De
plus, la surjection canonique (\ref{equation-conormal-algebra-quotient}) des algèbres graduées définit une
immersion fermée canonique
\begin{equation}
\label{equation-normal-cone-in-normal-bundle}
C_ZX \longrightarrow N_ZX
\end{equation}
des cônes sur  $Z$.





\section{Faisceaux d'idéaux réguliers}
\label{section-regular-ideal-sheaves}

\noindent
Dans cette section, nous généralisons la notion de diviseur de Cartier effectif à
des codimensions supérieures. Rappelons qu'une suite d'éléments  $f_1, \ldots, f_r$  d'un
anneau  $R$  est une  {\it  suite régulière} 
si pour chaque  $i = 1, \ldots, r$  l'élément  $f_i$  est un non-diviseur de zéro sur
 $R/(f_1, \ldots, f_{i - 1})$  et  $R/(f_1, \ldots, f_r) \not = 0$, voir Algèbre, Définition  \ref{algebra-definition-regular-sequence}. Il existe
trois conditions plus faibles et étroitement liées que nous pouvons imposer. La
première consiste à supposer que  $f_1, \ldots, f_r$  est une  {\it 
suite Koszul-régulière}, c’est-à-dire que  $H_i(K_\bullet(f_1, \ldots, f_r)) = 0$  pour
 $i > 0$, voir Plus sur l'Algèbre, Définition  \ref{more-algebra-definition-koszul-regular-sequence}. La suite est
appelée une  {\it  suite  $H_1$-régulière}  si
 $H_1(K_\bullet(f_1, \ldots, f_r)) = 0$. Une autre condition que nous pouvons imposer est qu'avec
 $J = (f_1, \ldots, f_r)$, l'application
$$
R/J[T_1, \ldots, T_r]
\longrightarrow
\bigoplus\nolimits_{n \geq 0}
J^n/J^{n + 1}
$$
qui associe  $T_i$  à  $f_i \bmod J^2$  est un isomorphisme. Dans ce cas, nous
disons que  $f_1, \ldots, f_r$  est une  {\it  suite
quasi-régulière}, voir Algèbre, Définition  \ref{algebra-definition-quasi-regular-sequence}. Étant donné un
 $R$-module  $M$ , il existe également une notion de
 $M$-séquence régulière et de  $M$-séquence quasi-régulière.

\medskip\noindent
Nous pouvons généraliser cela au cas des espaces annelés comme suit. Soit
 $X$  un espace annelé et soit  $f_1, \ldots, f_r \in \Gamma(X, \mathcal{O}_X)$. Nous disons que  $f_1, \ldots, f_r$ 
est une  {\it  séquence régulière}  si pour chaque
 $i = 1, \ldots, r$ , l'application
\begin{equation}
\label{equation-map-regular}
f_i :
\mathcal{O}_X/(f_1, \ldots, f_{i - 1})
\longrightarrow
\mathcal{O}_X/(f_1, \ldots, f_{i - 1})
\end{equation}
est une application injective de faisceaux. Nous disons que  $f_1, \ldots, f_r$  est une
 {\it  séquence Koszul-régulière}  si le complexe de
Koszul
\begin{equation}
\label{equation-koszul}
K_\bullet(\mathcal{O}_X, f_\bullet),
\end{equation}
voir Modules, Définition  \ref{modules-definition-koszul-complex}, est acyclique dans les degrés  $> 0$.
Nous disons que  $f_1, \ldots, f_r$  est une séquence  {\it 
 $H_1$-régulière}  si le complexe de Koszul  $K_\bullet(\mathcal{O}_X, f_\bullet)$  est exact
dans le degré  $1$. Enfin, nous disons que  $f_1, \ldots, f_r$  est une séquence
 {\it  quasi-régulière}  si l'application
\begin{equation}
\label{equation-map-quasi-regular}
\mathcal{O}_X/\mathcal{J}[T_1, \ldots, T_r]
\longrightarrow
\bigoplus\nolimits_{d \geq 0} \mathcal{J}^d/\mathcal{J}^{d + 1}
\end{equation}
est un isomorphisme de faisceaux où  $\mathcal{J} \subset \mathcal{O}_X$  est le faisceau d'idéaux engendré
par  $f_1, \ldots, f_r$. (Il existe également une notion de séquence
 $\mathcal{F}$-régulière et  $\mathcal{F}$-quasi-régulière pour un  $\mathcal{O}_X$-module
 $\mathcal{F}$  donné que nous introduirons ici si nous en avons besoin.)

\begin{lemma}
\label{lemma-types-regular-sequences-implications}
Soit  $X$  un espace annelé. Soit  $f_1, \ldots, f_r \in \Gamma(X, \mathcal{O}_X)$. Nous avons les implications
suivantes  $f_1, \ldots, f_r$  est une suite régulière  $\Rightarrow$   $f_1, \ldots, f_r$  est une
suite Koszul-régulière  $\Rightarrow$   $f_1, \ldots, f_r$  est une suite
 $H_1$-régulière  $\Rightarrow$   $f_1, \ldots, f_r$  est une suite quasi-régulière.
\end{lemma}

\begin{proof}
Puisque nous pouvons vérifier l'exactitude sur les fibres, une suite  $f_1, \ldots, f_r$ 
est une suite régulière si et seulement si les applications
$$
f_i :
\mathcal{O}_{X, x}/(f_1, \ldots, f_{i - 1})
\longrightarrow
\mathcal{O}_{X, x}/(f_1, \ldots, f_{i - 1})
$$
sont injectives pour tout  $x \in X$. Les autres types de régularité peuvent
également être vérifiés fibre par fibre (détails omis). Par conséquent, les
implications découlent de Plus sur Algèbre, Lemme  \ref{more-algebra-lemma-regular-koszul-regular},  \ref{more-algebra-lemma-koszul-regular-H1-regular}, et
 \ref{more-algebra-lemma-H1-regular-quasi-regular}.
\end{proof}

\begin{definition}
\label{definition-regular-ideal-sheaf}
\begin{reference}
Le concept de faisceau d'idéaux Koszul-régulier a été introduit dans  \cite[Expose VII, Definition 1.4]{SGA6} 
où il était appelé un faisceau d'idéaux régulier.
\end{reference}
Soit  $X$  un espace annelé. Soit  $\mathcal{J} \subset \mathcal{O}_X$  un faisceau d'idéaux.
\begin{enumerate}
\item On dit que  $\mathcal{J}$  est  {\it  régulier}  si pour
chaque  $x \in \text{Supp}(\mathcal{O}_X/\mathcal{J})$  il existe un voisinage ouvert  $x \in U \subset X$  et une suite
régulière  $f_1, \ldots, f_r \in \mathcal{O}_X(U)$  telle que  $\mathcal{J}|_U$  est engendrée par  $f_1, \ldots, f_r$.
\item On dit que  $\mathcal{J}$  est  {\it  Koszul-régulier} 
si pour chaque  $x \in \text{Supp}(\mathcal{O}_X/\mathcal{J})$  il existe un voisinage ouvert  $x \in U \subset X$  et une
suite Koszul-régulière  $f_1, \ldots, f_r \in \mathcal{O}_X(U)$  telle que  $\mathcal{J}|_U$  est engendrée par
 $f_1, \ldots, f_r$.
\item On dit que  $\mathcal{J}$  est  {\it 
 $H_1$-régulier}  si pour chaque  $x \in \text{Supp}(\mathcal{O}_X/\mathcal{J})$  il existe un
voisinage ouvert  $x \in U \subset X$  et une suite  $H_1$-régulière  $f_1, \ldots, f_r \in \mathcal{O}_X(U)$ 
telle que  $\mathcal{J}|_U$  est engendrée par  $f_1, \ldots, f_r$.
\item Nous disons que  $\mathcal{J}$  est  {\it 
quasi-régulier}  si pour chaque  $x \in \text{Supp}(\mathcal{O}_X/\mathcal{J})$  il existe un voisinage
ouvert  $x \in U \subset X$  et une suite quasi-régulière  $f_1, \ldots, f_r \in \mathcal{O}_X(U)$  telle que
 $\mathcal{J}|_U$  est engendrée par  $f_1, \ldots, f_r$.
\end{enumerate}
\end{definition}

\noindent
De nombreuses propriétés de cette notion découlent immédiatement des notions
correspondantes pour les suites régulières et quasi-régulières dans les anneaux.

\begin{lemma}
\label{lemma-regular-quasi-regular-scheme}
Soit  $X$  un espace annelé. Soit  $\mathcal{J}$  un faisceau d'idéaux. Nous
avons les implications suivantes:  $\mathcal{J}$  est régulier  $\Rightarrow$ 
 $\mathcal{J}$  est Koszul-régulier  $\Rightarrow$   $\mathcal{J}$  est
 $H_1$-régulier  $\Rightarrow$   $\mathcal{J}$  est quasi-régulier.
\end{lemma}

\begin{proof}
Le lemme se réduit immédiatement au Lemme  \ref{lemma-types-regular-sequences-implications}.
\end{proof}

\begin{lemma}
\label{lemma-quasi-regular-ideal}
Soit  $X$  un espace localement annelé. Soit  $\mathcal{J} \subset \mathcal{O}_X$  un faisceau
d'idéaux. Alors  $\mathcal{J}$  est quasi-régulier si et seulement si les conditions
suivantes sont satisfaites:
\begin{enumerate}
\item $\mathcal{J}$  est un  $\mathcal{O}_X$-module de type fini,
\item $\mathcal{J}/\mathcal{J}^2$  est un  $\mathcal{O}_X/\mathcal{J}$-module localement libre de type fini, et
\item les applications canoniques
$$
\text{Sym}^n_{\mathcal{O}_X/\mathcal{J}}(\mathcal{J}/\mathcal{J}^2)
\longrightarrow
\mathcal{J}^n/\mathcal{J}^{n + 1}
$$
sont des isomorphismes pour tout  $n \geq 0$.
\end{enumerate}
\end{lemma}

\begin{proof}
Il est clair que si  $U \subset X$  est un ouvert tel que  $\mathcal{J}|_U$  soit engendré
par une suite quasi-régulière  $f_1, \ldots, f_r \in \mathcal{O}_X(U)$  alors  $\mathcal{J}|_U$  est de type fini,
 $\mathcal{J}|_U/\mathcal{J}^2|_U$  est libre avec pour base  $f_1, \ldots, f_r$, et les applications de (3)
sont des isomorphismes car elles sont une formulation indépendante des coordonnées
de la partie de degré  $n$  de (\ref{equation-map-quasi-regular}). Il en découle qu'il est
clair qu'être quasi-régulier implique les conditions (1), (2) et (3).

\medskip\noindent
Inversement, supposons que (1), (2) et (3) soient vérifiés. Choisissons un point
 $x \in \text{Supp}(\mathcal{O}_X/\mathcal{J})$. Il existe alors un voisinage  $U \subset X$  de  $x$  tel que
 $\mathcal{J}|_U/\mathcal{J}^2|_U$  soit libre de rang  $r$  sur  $\mathcal{O}_U/\mathcal{J}|_U$. Après avoir
éventuellement restreint  $U$ , nous pouvons supposer qu'il existe
 $f_1, \ldots, f_r \in \mathcal{J}(U)$  qui se projettent sur une base de  $\mathcal{J}|_U/\mathcal{J}^2|_U$  comme
 $\mathcal{O}_U/\mathcal{J}|_U$-module. En particulier, nous voyons que les images de  $f_1, \ldots, f_r$ 
dans  $\mathcal{J}_x/\mathcal{J}^2_x$  engendrent. Par conséquent, d'après le lemme de Nakayama
(Algèbre, Lemme  \ref{algebra-lemma-NAK}), nous voyons que  $f_1, \ldots, f_r$  engendrent la germe
 $\mathcal{J}_x$. Ainsi, puisque  $\mathcal{J}$  est de type fini, d'après Modules, Lemme
 \ref{modules-lemma-finite-type-surjective-on-stalk} , après avoir restreint  $U$ , nous pouvons supposer que
 $f_1, \ldots, f_r$  engendrent  $\mathcal{J}$. Enfin, à partir de (3) et de l'isomorphisme
 $\mathcal{J}|_U/\mathcal{J}^2|_U = \bigoplus \mathcal{O}_U/\mathcal{J}|_U f_i$ , il est clair que  $f_1, \ldots, f_r \in \mathcal{O}_X(U)$  est une suite quasi-régulière.
\end{proof}

\begin{lemma}
\label{lemma-generate-regular-ideal}
Soit  $(X, \mathcal{O}_X)$  un espace localement annelé. Soit  $\mathcal{J} \subset \mathcal{O}_X$  un faisceau
d'idéaux. Soient  $x \in X$  et  $f_1, \ldots, f_r \in \mathcal{J}_x$  dont les images fournissent une base
de l'espace vectoriel  $\kappa(x)$  $\mathcal{J}_x/\mathfrak m_x\mathcal{J}_x$.
\begin{enumerate}
\item Si  $\mathcal{J}$  est quasi-régulier, alors il existe un voisinage ouvert tel que
 $f_1, \ldots, f_r \in \mathcal{O}_X(U)$  forme une suite quasi-régulière générant  $\mathcal{J}|_U$.
\item Si  $\mathcal{J}$  est  $H_1$-régulier, alors il existe un voisinage ouvert tel
que  $f_1, \ldots, f_r \in \mathcal{O}_X(U)$  forme une suite  $H_1$-régulière générant  $\mathcal{J}|_U$.
\item Si  $\mathcal{J}$  est Koszul-régulier, alors il existe un voisinage ouvert tel que
 $f_1, \ldots, f_r \in \mathcal{O}_X(U)$  forme une suite Koszul-régulière générant  $\mathcal{J}|_U$.
\end{enumerate}
\end{lemma}

\begin{proof}
Tout d'abord, supposons que  $\mathcal{J}$  est quasi-régulier. Nous pouvons choisir
un voisinage ouvert  $U \subset X$  de  $x$  et une suite quasi-régulière
 $g_1, \ldots, g_s \in \mathcal{O}_X(U)$  qui engendre  $\mathcal{J}|_U$. Notez que cela implique que  $\mathcal{J}/\mathcal{J}^2$ 
est libre de rang  $s$  sur  $\mathcal{O}_U/\mathcal{J}|_U$  (voir le Lemme  \ref{lemma-quasi-regular-ideal}  et sa
démonstration) et donc  $r = s$. Nous pouvons réduire  $U$  et supposer
 $f_1, \ldots, f_r \in \mathcal{J}(U)$. Ainsi nous pouvons écrire
$$
f_i = \sum a_{ij} g_j
$$
pour certains  $a_{ij} \in \mathcal{O}_X(U)$. Par hypothèse, la matrice  $A = (a_{ij})$  se mappe sur
une matrice inversible sur  $\kappa(x)$. Par conséquent, après avoir encore réduit
 $U$ , nous pouvons supposer que  $(a_{ij})$  est inversible. Ainsi nous
voyons que  $f_1, \ldots, f_r$  donnent une base pour  $(\mathcal{J}/\mathcal{J}^2)|_U$  ce qui prouve que
 $f_1, \ldots, f_r$  est une suite quasi-régulière sur  $U$.

\medskip\noindent
Notez que pour prouver (2) et (3), nous pouvons, puisque les hypothèses de (2) et
(3) sont plus fortes que l'hypothèse de (1), déjà supposer que  $f_1, \ldots, f_r \in \mathcal{J}(U)$  et
 $f_i = \sum a_{ij}g_j$  avec  $(a_{ij})$  inversible comme ci-dessus, où maintenant
 $g_1, \ldots, g_r$  est une suite  $H_1$-régulière ou Koszul-régulière. Puisque le
complexe de Koszul sur  $f_1, \ldots, f_r$  est isomorphe au complexe de Koszul sur
 $g_1, \ldots, g_r$  via la matrice  $(a_{ij})$  (voir Plus sur Algèbre, Lemme
 \ref{more-algebra-lemma-change-basis}), nous concluons que  $f_1, \ldots, f_r$  est  $H_1$-régulier ou
Koszul-régulier comme souhaité.
\end{proof}

\begin{lemma}
\label{lemma-regular-ideal-sheaf-quasi-coherent}
Tout faisceau d'idéaux régulier, Koszul-régulier,  $H_1$-régulier ou
quasi-régulier sur un schéma est un faisceau d'idéaux quasi-cohérent de type fini.
\end{lemma}

\begin{proof}
Cela découle du fait qu'un tel faisceau d'idéaux est localement engendré par un
nombre fini de sections. Et tout faisceau d'idéaux localement engendré par des
sections sur un schéma est quasi-cohérent, voir Schémas, Lemme  \ref{schemes-lemma-closed-subspace-scheme}.
\end{proof}

\begin{lemma}
\label{lemma-regular-ideal-sheaf-scheme}
Soit  $X$  un schéma. Soit  $\mathcal{J}$  un faisceau d'idéaux. Alors
 $\mathcal{J}$  est régulier (resp. \ Koszul-régulier,
 $H_1$-régulier, quasi-régulier) si et seulement si pour chaque  $x \in \text{Supp}(\mathcal{O}_X/\mathcal{J})$ 
il existe un voisinage affine ouvert  $x \in U \subset X$,  $U = \Spec(A)$  tel que
 $\mathcal{J}|_U = \widetilde{I}$  et tel que  $I$  soit engendré par une suite régulière (resp.
\ Koszul-régulière,  $H_1$-régulière, quasi-régulière)
 $f_1, \ldots, f_r \in A$.
\end{lemma}

\begin{proof}
Par hypothèse, nous pouvons trouver un voisinage ouvert  $U$  de
 $x$  au-dessus duquel  $\mathcal{J}$  est engendré par une suite régulière
(resp. \ Koszul-régulière,  $H_1$-régulière, quasi-régulière)
 $f_1, \ldots, f_r \in \mathcal{O}_X(U)$. Après avoir restreint  $U$ , nous pouvons supposer que
 $U$  est affine, disons  $U = \Spec(A)$. Puisque  $\mathcal{J}$  est
quasi-cohérent d'après le Lemme  \ref{lemma-regular-ideal-sheaf-quasi-coherent} , nous voyons que  $\mathcal{J}|_U = \widetilde{I}$  pour un
certain idéal  $I \subset A$. Nous pouvons maintenant utiliser le fait que
$$
\widetilde{\ } : \text{Mod}_A \longrightarrow \QCoh(\mathcal{O}_U)
$$
est une équivalence de catégories qui préserve l'exactitude. Par exemple, le fait
que les fonctions  $f_i$  génèrent  $\mathcal{J}$  signifie que les
 $f_i$, vus comme des éléments de  $A$ , génèrent  $I$. Le
fait que (\ref{equation-map-regular}) soit injectif (resp. \ (\ref{equation-koszul}) soit
exact, (\ref{equation-koszul}) soit exact en degré  $1$, (\ref{equation-map-quasi-regular}) soit un
isomorphisme) implique la propriété correspondante de l'application  $A/(f_1, \ldots, f_{i - 1}) \to A/(f_1, \ldots, f_{i - 1})$ 
(resp. \  le complexe  $K_\bullet(A, f_1, \ldots, f_r)$, l'application  $A/I[T_1, \ldots, T_r] \to \bigoplus I^n/I^{n + 1}$). Ainsi,
 $f_1, \ldots, f_r \in A$  est une suite régulière (resp. \ Koszul-régulière,
 $H_1$-régulière, quasi-régulière) de l’anneau  $A$.
\end{proof}

\begin{lemma}
\label{lemma-Noetherian-scheme-regular-ideal}
Soit  $X$  un schéma localement noethérien. Soit  $\mathcal{J} \subset \mathcal{O}_X$  un faisceau
quasi-cohérent d’idéaux. Soit  $x$  un point du support de  $\mathcal{O}_X/\mathcal{J}$.
Les assertions suivantes sont équivalentes
\begin{enumerate}
\item $\mathcal{J}_x$  est généré par une suite régulière dans  $\mathcal{O}_{X, x}$,
\item $\mathcal{J}_x$  est généré par une suite Koszul-régulière dans  $\mathcal{O}_{X, x}$,
\item $\mathcal{J}_x$  est généré par une suite  $H_1$-régulière dans  $\mathcal{O}_{X, x}$,
\item $\mathcal{J}_x$  est généré par une suite quasi-régulière dans  $\mathcal{O}_{X, x}$,
\item il existe un voisinage affine  $U = \Spec(A)$  de  $x$  tel que  $\mathcal{J}|_U = \widetilde{I}$  et
 $I$  sont générés par une suite régulière dans  $A$, et
\item il existe un voisinage affine  $U = \Spec(A)$  de  $x$  tel que  $\mathcal{J}|_U = \widetilde{I}$  et
 $I$  sont générés par une suite Koszul-régulière dans  $A$, et
\item il existe un voisinage affine  $U = \Spec(A)$  de  $x$  tel que  $\mathcal{J}|_U = \widetilde{I}$  et
 $I$  sont générés par une suite  $H_1$-régulière dans  $A$,
et
\item il existe un voisinage affine  $U = \Spec(A)$  de  $x$  tel que  $\mathcal{J}|_U = \widetilde{I}$  et
 $I$  soit généré par une suite quasi-régulière dans  $A$,
\item il existe un voisinage  $U$  de  $x$  tel que  $\mathcal{J}|_U$  soit
régulier, et
\item il existe un voisinage  $U$  de  $x$  tel que  $\mathcal{J}|_U$  soit
Koszul-régulier, et
\item il existe un voisinage  $U$  de  $x$  tel que  $\mathcal{J}|_U$  soit
 $H_1$-régulier, et
\item il existe un voisinage  $U$  de  $x$  tel que  $\mathcal{J}|_U$  soit
quasi-régulier.
\end{enumerate}
En particulier, sur un schéma localement noethérien, les notions de faisceau
d'idéaux régulier, Koszul-régulier,  $H_1$-régulier, ou quasi-régulier
coïncident toutes.
\end{lemma}

\begin{proof}
Il résulte du Lemme  \ref{lemma-regular-ideal-sheaf-scheme}  que (5)  $\Leftrightarrow$  (9), (6)  $\Leftrightarrow$  (10),
(7)  $\Leftrightarrow$  (11), et (8)  $\Leftrightarrow$  (12). Il est clair que (5)  $\Rightarrow$ 
(1), (6)  $\Rightarrow$  (2), (7)  $\Rightarrow$  (3), et (8)  $\Rightarrow$  (4). Nous
avons (1)  $\Rightarrow$  (5) par Algèbre, Lemme  \ref{algebra-lemma-regular-sequence-in-neighbourhood}. Nous avons (9)
 $\Rightarrow$  (10)  $\Rightarrow$  (11)  $\Rightarrow$  (12) par le Lemme  \ref{lemma-regular-quasi-regular-scheme}.
Enfin, (4)  $\Rightarrow$  (1) par Algèbre, Lemme  \ref{algebra-lemma-quasi-regular-regular}. Maintenant, tous les
12 énoncés sont équivalents.
\end{proof}













\section{Immersions régulières}
\label{section-regular-immersions}

\noindent
Soit  $i : Z \to X$  une immersion de schémas. Par définition, cela signifie qu'il
existe un sous-schéma ouvert  $U \subset X$  tel que  $Z$  soit identifié à un
sous-schéma fermé de  $U$. Soit  $\mathcal{I} \subset \mathcal{O}_U$  le faisceau quasi-cohérent
d'idéaux correspondant. Supposons que  $U' \subset X$  soit un second sous-schéma
ouvert de ce type, et notons  $\mathcal{I}' \subset \mathcal{O}_{U'}$  le faisceau quasi-cohérent d'idéaux
correspondant. Alors  $\mathcal{I}|_{U \cap U'} = \mathcal{I}'|_{U \cap U'}$. De plus, le support de  $\mathcal{O}_U/\mathcal{I}$  est
 $Z$  qui est contenu dans  $U \cap U'$  et est également le support de
 $\mathcal{O}_{U'}/\mathcal{I}'$. Il s'ensuit donc d'après la Définition  \ref{definition-regular-ideal-sheaf}  que  $\mathcal{I}$ 
est un idéal régulier si et seulement si  $\mathcal{I}'$  est un idéal régulier. De
même pour être Koszul-régulier,  $H_1$-régulier, ou quasi-régulier.

\begin{definition}
\label{definition-regular-immersion}
\begin{reference}
Le concept d'une immersion Koszul-régulière a été introduit dans  \cite[Expose VII, Definition 1.4]{SGA6}  où
il était appelé une immersion régulière.
\end{reference}
Soit  $i : Z \to X$  une immersion de schémas. Choisissez un sous-schéma ouvert
 $U \subset X$  tel que  $i$  identifie  $Z$  avec un sous-schéma fermé
de  $U$  et désignez  $\mathcal{I} \subset \mathcal{O}_U$  le faisceau quasi-cohérent d'idéaux
correspondant.
\begin{enumerate}
\item Nous disons que  $i$  est une  {\it  immersion
régulière}  si  $\mathcal{I}$  est régulière.
\item Nous disons que  $i$  est une  {\it  immersion
Koszul-régulière}  si  $\mathcal{I}$  est Koszul-régulière.
\item Nous disons que  $i$  est une  {\it  immersion
 $H_1$-régulière}  si  $\mathcal{I}$  est  $H_1$-régulier.
\item Nous disons que  $i$  est une  {\it  immersion
quasi-régulière}  si  $\mathcal{I}$  est quasi-régulière.
\end{enumerate}
\end{definition}

\noindent
La discussion ci-dessus montre que cela est indépendant du choix de  $U$.
Les conditions sont énumérées dans l'ordre décroissant de force, voir le Lemme
 \ref{lemma-regular-quasi-regular-immersion}. Une immersion fermée de Koszul-régulière est lisse localement une
immersion régulière, voir le Lemme  \ref{lemma-koszul-regular-smooth-locally-regular}. Dans le cas localement
noethérien, les quatre notions coïncident, voir le Lemme  \ref{lemma-Noetherian-scheme-regular-ideal}.

\begin{lemma}
\label{lemma-regular-quasi-regular-immersion}
Soit  $i : Z \to X$  une immersion de schémas. Nous avons les implications suivantes:
 $i$  est régulière  $\Rightarrow$   $i$  est Koszul-régulière
 $\Rightarrow$   $i$  est  $H_1$-régulière  $\Rightarrow$   $i$  est
quasi-régulière.
\end{lemma}

\begin{proof}
Le lemme se réduit immédiatement au Lemme  \ref{lemma-regular-quasi-regular-scheme}.
\end{proof}

\begin{lemma}
\label{lemma-regular-immersion-noetherian}
Soit  $i : Z \to X$  une immersion de schémas. Supposons que  $X$  soit
localement noethérien. Alors  $i$  est régulier  $\Leftrightarrow$   $i$ 
est de Koszul-régulier  $\Leftrightarrow$   $i$  est  $H_1$-régulier
 $\Leftrightarrow$   $i$  est quasi-régulier.
\end{lemma}

\begin{proof}
Suit immédiatement du Lemme  \ref{lemma-regular-quasi-regular-immersion}  et du Lemme  \ref{lemma-Noetherian-scheme-regular-ideal}.
\end{proof}

\begin{lemma}
\label{lemma-flat-base-change-regular-immersion}
Soit  $i : Z \to X$  une immersion régulière (resp. \  de
Koszul-régularité,  $H_1$-régulière, quasi-régulière). Soit  $X' \to X$  un
morphisme plat. Alors le changement de base  $i' : Z \times_X X' \to X'$  est une immersion
régulière (resp. \  de Koszul-régularité,  $H_1$-régulière,
quasi-régulière).
\end{lemma}

\begin{proof}
Via le Lemme  \ref{lemma-regular-ideal-sheaf-scheme}  cela se traduit en les énoncés algébriques dans Algèbre,
Lemmata  \ref{algebra-lemma-flat-increases-depth}  et  \ref{algebra-lemma-flat-base-change-quasi-regular}  et Plus sur Algèbre, Lemme  \ref{more-algebra-lemma-koszul-regular-flat-base-change}.
\end{proof}

\begin{lemma}
\label{lemma-quasi-regular-immersion}
Soit  $i : Z \to X$  une immersion de schémas. Alors  $i$  est une immersion
quasi-régulière si et seulement si les conditions suivantes sont satisfaites
\begin{enumerate}
\item $i$  est localement de présentation finie,
\item le faisceau conormal  $\mathcal{C}_{Z/X}$  est localement libre de type fini, et
\item l'application (\ref{equation-conormal-algebra-quotient}) est un isomorphisme.
\end{enumerate}
\end{lemma}

\begin{proof}
Une immersion ouverte est localement de présentation finie. Par conséquent, nous
pouvons remplacer  $X$  par un sous-schéma ouvert  $U \subset X$  tel que
 $i$  identifie  $Z$  avec un sous-schéma fermé de  $U$,
c’est-à-dire que nous pouvons supposer que  $i$  est une immersion fermée.
Soit  $\mathcal{I} \subset \mathcal{O}_X$  le faisceau d'idéaux quasi-cohérents correspondant. Rappelons,
voir Morphismes, Lemme  \ref{morphisms-lemma-closed-immersion-finite-presentation} , que  $\mathcal{I}$  est de type fini si et
seulement si  $i$  est localement de présentation finie. Par conséquent,
l'équivalence découle du Lemme  \ref{lemma-quasi-regular-ideal}  et du dépliement des définitions.
\end{proof}

\begin{lemma}
\label{lemma-transitivity-conormal-quasi-regular}
Soient  $Z \to Y \to X$  des immersions de schémas. Supposons que  $Z \to Y$  est
 $H_1$-régulier. Alors la suite canonique de Morphismes, Lemme  \ref{morphisms-lemma-transitivity-conormal}
$$
0 \to i^*\mathcal{C}_{Y/X} \to
\mathcal{C}_{Z/X} \to
\mathcal{C}_{Z/Y} \to 0
$$
est exacte et localement scindée.
\end{lemma}

\begin{proof}
Puisque  $\mathcal{C}_{Z/Y}$  est localement libre de type fini (voir le Lemme  \ref{lemma-quasi-regular-immersion} 
et le Lemme  \ref{lemma-regular-quasi-regular-scheme}), il suffit de prouver que la suite est exacte. D'après
ce qui a été prouvé dans Morphismes, Lemme  \ref{morphisms-lemma-transitivity-conormal} , il suffit de montrer que
la première application est injective. En travaillant localement de manière
affine, cela se réduit à la question suivante: Supposons que nous ayons un anneau
 $A$  et des idéaux  $I \subset J \subset A$. Supposons que  $J/I \subset A/I$  soit engendré
par une suite  $H_1$-régulière. Cela implique-t-il que  $I/I^2 \otimes_A A/J \to J/J^2$  est
injective? Notez que  $I/I^2 \otimes_A A/J = I/IJ$. Par conséquent, nous essayons de prouver que
 $I \cap J^2 = IJ$. Ceci est le résultat de Plus sur Algèbre, Lemme  \ref{more-algebra-lemma-conormal-sequence-H1-regular}.
\end{proof}

\noindent
Une composition d'immersions presque régulières peut ne pas être presque
régulière, voir Algèbre, Remarque  \ref{algebra-remark-join-quasi-regular-sequences}. Les autres types d'immersions
régulières sont préservés par composition.

\begin{lemma}
\label{lemma-composition-regular-immersion}
Soient  $i : Z \to Y$  et  $j : Y \to X$  des immersions de schémas.
\begin{enumerate}
\item Si  $i$  et  $j$  sont des immersions régulières, il en est de même
pour  $j \circ i$.
\item Si  $i$  et  $j$  sont des immersions Koszul-régulières, il en est
de même pour  $j \circ i$.
\item Si  $i$  et  $j$  sont des immersions  $H_1$-régulières, il en
est de même pour  $j \circ i$.
\item Si  $i$  est une immersion  $H_1$-régulière et  $j$  est une
immersion quasi-régulière, alors  $j \circ i$  est une immersion quasi-régulière.
\end{enumerate}
\end{lemma}

\begin{proof}
La version algébrique de (1) est Algèbre, Lemme  \ref{algebra-lemma-join-regular-sequences}. La version
algébrique de (2) est Plus sur Algèbre, Lemme  \ref{more-algebra-lemma-join-koszul-regular-sequences}. La version algébrique
de (3) est Plus sur Algèbre, Lemme  \ref{more-algebra-lemma-join-H1-regular-sequences}. La version algébrique de (4) est
Plus sur Algèbre, Lemme  \ref{more-algebra-lemma-join-quasi-regular-H1-regular}.
\end{proof}

\begin{lemma}
\label{lemma-permanence-regular-immersion}
Soient  $i : Z \to Y$  et  $j : Y \to X$  des immersions de schémas. Supposons que
 $j$  soit localement de présentation finie et que la séquence
$$
0 \to i^*\mathcal{C}_{Y/X} \to
\mathcal{C}_{Z/X} \to
\mathcal{C}_{Z/Y} \to 0
$$
des morphismes, le Lemme  \ref{morphisms-lemma-transitivity-conormal}  est exact et localement scindé.
\begin{enumerate}
\item Si  $j \circ i$  est une immersion quasi-régulière, il en est de même pour
 $i$.
\item Si  $j \circ i$  est une immersion  $H_1$-régulière, il en est de même pour
 $i$.
\item Si  $j$  et  $j \circ i$  sont toutes deux des immersions de
Koszul-régulières, il en est de même pour  $i$.
\end{enumerate}
\end{lemma}

\begin{proof}
Après avoir restreint  $Y$  et  $X$ , nous pouvons supposer que
 $i$  et  $j$  sont des immersions fermées. Notons  $\mathcal{I} \subset \mathcal{O}_X$  le
faisceau d'idéaux de  $Y$  et  $\mathcal{J} \subset \mathcal{O}_X$  le faisceau d'idéaux de
 $Z$. La suite conormale est  $0 \to \mathcal{I}/\mathcal{I}\mathcal{J}
\to \mathcal{J}/\mathcal{J}^2 \to
\mathcal{J}/(\mathcal{I} + \mathcal{J}^2) \to 0$. Soient  $z \in Z$  et posons
 $y = i(z)$,  $x = j(y) = j(i(z))$. Choisissons  $f_1, \ldots, f_n \in \mathcal{I}_x$  qui se projettent sur une
base de  $\mathcal{I}_x/\mathfrak m_z\mathcal{I}_x$. Étendons cela à  $f_1, \ldots, f_n, g_1, \ldots, g_m \in \mathcal{J}_x$  qui se projettent sur une base
de  $\mathcal{J}_x/\mathfrak m_z\mathcal{J}_x$. Cela est possible car nous avons supposé que la suite des
faisceaux conormaux est scindée dans un voisinage de  $z$, donc
 $\mathcal{I}_x/\mathfrak m_x\mathcal{I}_x \to
\mathcal{J}_x/\mathfrak m_x\mathcal{J}_x$  est injective.

\medskip\noindent
Preuve de (1). D'après le Lemme  \ref{lemma-generate-regular-ideal} , nous pouvons trouver un voisinage
affine ouvert  $U$  de  $x$  tel que  $f_1, \ldots, f_n, g_1, \ldots, g_m$  forme une suite
quasi-régulière générant  $\mathcal{J}$. Puisque  $i$  est localement de
présentation finie, nous pouvons en outre supposer que  $\mathcal{I}$  est engendré
par  $f_1, \ldots, f_n$  (utiliser Morphismes, Lemme  \ref{morphisms-lemma-closed-immersion-finite-presentation}, le lemme de Nakayama,
et Modules, Lemme  \ref{modules-lemma-finite-type-surjective-on-stalk}). Ainsi, d'après Algèbre, Lemme  \ref{algebra-lemma-truncate-quasi-regular} , nous
voyons que  $g_1, \ldots, g_m$  induit une suite quasi-régulière sur  $Y \cap U$ 
définissant  $Z$.

\medskip\noindent
Preuve de (2). Exactement la même que la preuve de (1) sauf en utilisant Plus sur
Algèbre, Lemme  \ref{more-algebra-lemma-truncate-H1-regular}.

\medskip\noindent
Preuve de (3). D'après le Lemme  \ref{lemma-generate-regular-ideal}  (appliqué deux fois), nous pouvons
trouver un voisinage affine ouvert  $U$  de  $x$  tel que
 $f_1, \ldots, f_n$  forme une suite Koszul-régulière générant  $\mathcal{I}$  et
 $f_1, \ldots, f_n, g_1, \ldots, g_m$  forme une suite Koszul-régulière générant  $\mathcal{J}$. Par
conséquent, d'après Plus sur Algèbre, Lemme  \ref{more-algebra-lemma-truncate-koszul-regular} , nous voyons que
 $g_1, \ldots, g_m$  induit une suite Koszul-régulière sur  $Y \cap U$  coupant
 $Z$.
\end{proof}

\begin{lemma}
\label{lemma-extra-permanence-regular-immersion-noetherian}
Soient  $i : Z \to Y$  et  $j : Y \to X$  des immersions de schémas. Choisissez
 $z \in Z$  et désignez  $y \in Y$,  $x \in X$  les points correspondants.
Supposons que  $X$  soit localement noethérien. Les conditions suivantes
sont équivalentes
\begin{enumerate}
\item $i$  est une immersion régulière dans un voisinage de  $z$  et
 $j$  est une immersion régulière dans un voisinage de  $y$.
\item $i$  et  $j \circ i$  sont des immersions régulières dans un voisinage de
 $z$,
\item . $j \circ i$  est une immersion régulière dans un voisinage de  $z$  et la
suite conormale
$$
0 \to i^*\mathcal{C}_{Y/X} \to
\mathcal{C}_{Z/X} \to
\mathcal{C}_{Z/Y} \to 0
$$
est exacte scindée dans un voisinage de  $z$.
\end{enumerate}
\end{lemma}

\begin{proof}
. Puisque  $X$  (et donc  $Y$) est localement noethérien, les 4
types d’immersions régulières coïncident, et de plus, nous pouvons vérifier si un
morphisme est une immersion régulière au niveau des anneaux locaux, voir le Lemme
 \ref{lemma-Noetherian-scheme-regular-ideal}. L’implication (1)  $\Rightarrow$  (2) est le Lemme  \ref{lemma-composition-regular-immersion}.
L’implication (2)  $\Rightarrow$  (3) est le Lemme  \ref{lemma-transitivity-conormal-quasi-regular}. Ainsi, il suffit de
prouver que (3) implique (1).

\medskip\noindent
Supposons (3). Posons  $A = \mathcal{O}_{X, x}$. Désignons par  $I \subset A$  le noyau de
l'application surjective  $\mathcal{O}_{X, x} \to \mathcal{O}_{Y, y}$  et par  $J \subset A$  le noyau de
l'application surjective  $\mathcal{O}_{X, x} \to \mathcal{O}_{Z, z}$. Notez que toute suite minimale d'éléments
générant  $J$  dans  $A$  est une suite quasi-régulière donc
régulière, voir le Lemme  \ref{lemma-generate-regular-ideal}. Par hypothèse, la suite conormale
$$
0 \to I/IJ \to J/J^2 \to J/(I + J^2) \to 0
$$
est exactement scindé en une suite de  $A/J$-modules. Par conséquent, nous
pouvons choisir un système minimal de générateurs  $f_1, \ldots, f_n, g_1, \ldots, g_m$  de  $J$ 
avec  $f_1, \ldots, f_n \in I$  un système minimal de générateurs de  $I$. Comme indiqué
ci-dessus,  $f_1, \ldots, f_n, g_1, \ldots, g_m$  est une suite régulière dans  $A$. Il ressort
directement de la définition d'une suite régulière que  $f_1, \ldots, f_n$  est une suite
régulière dans  $A$  et  $\overline{g}_1, \ldots, \overline{g}_m$  est une suite régulière dans
 $A/I$. Ainsi,  $j$  est une immersion régulière en  $y$  et
 $i$  est une immersion régulière en  $z$.
\end{proof}

\begin{remark}
\label{remark-not-always-extra-permanence}
Dans la situation du Lemme  \ref{lemma-extra-permanence-regular-immersion-noetherian} , les parties (1), (2), (3) sont
 {\bf  non }  équivalentes à « $j \circ i$  et
 $j$  sont des immersions régulières en  $z$  et  $y$ ». Un
exemple est  $X = \mathbf{A}^1_k = \Spec(k[x])$,  $Y = \Spec(k[x]/(x^2))$  et  $Z = \Spec(k[x]/(x))$.
\end{remark}

\begin{lemma}
\label{lemma-koszul-regular-smooth-locally-regular}
Soit  $i : Z \to X$  une immersion fermée koszul-régulière. Alors il existe un
morphisme lisse surjectif  $X' \to X$  tel que le changement de base  $i' : Z \times_X X' \to X'$ 
de  $i$  soit une immersion régulière.
\end{lemma}

\begin{proof}
Nous pouvons supposer que  $X$  est affine et que l'idéal de  $Z$ 
est engendré par une suite koszul-régulière en remplaçant  $X$  par les
éléments d'un recouvrement ouvert affine approprié (ouverts affines comme dans le
Lemme  \ref{lemma-regular-ideal-sheaf-scheme}). Le cas affine est Développement sur Algèbre, Lemme
 \ref{more-algebra-lemma-Koszul-regular-flat-locally-regular}.
\end{proof}

\begin{lemma}
\label{lemma-immersion-regular-regular-immersion}
Soit  $i : Z \to X$  une immersion. Si  $Z$  et  $X$  sont des schémas
réguliers, alors  $i$  est une immersion régulière.
\end{lemma}

\begin{proof}
Soit  $z \in Z$. D'après le Lemme  \ref{lemma-Noetherian-scheme-regular-ideal} , il suffit de montrer que le
noyau de  $\mathcal{O}_{X, z} \to \mathcal{O}_{Z, z}$  est engendré par une suite régulière. Cela découle des Lemme
d'Algèbre  \ref{algebra-lemma-regular-quotient-regular}  et  \ref{algebra-lemma-regular-ring-CM}.
\end{proof}





\section{Immersions régulières relatives}
\label{section-relative-regular-immersion}

\noindent
Dans cette section, nous considérons la propriété de changement de base pour les
immersions régulières. Le lemme suivant ne tient pas pour les immersions
régulières ou pour les immersions de Koszul, voir Exemples, Lemme  \ref{examples-lemma-base-change-regular-sequence}.

\begin{lemma}
\label{lemma-relative-regular-immersion}
Soit  $f : X \to S$  un morphisme de schémas. Soit  $i : Z \subset X$  une immersion.
Supposons
\begin{enumerate}
\item $i$  est un  $H_1$-régulier (resp.\ immersion
quasi-régulière, et
\item $Z \to S$  est un morphisme plat.
\end{enumerate}
Alors, pour tout morphisme de schémas  $g : S' \to S$ , le changement de base
 $Z' = S' \times_S Z \to X' = S' \times_S X$  est une immersion  $H_1$-régulière (resp. quasi-régulière)
\ .
\end{lemma}

\begin{proof}
En déroulant les définitions et en utilisant le Lemme  \ref{lemma-regular-ideal-sheaf-scheme} , cela se
traduit par More on Algebra, Lemme  \ref{more-algebra-lemma-relative-regular-immersion-algebra}.
\end{proof}

\noindent
Ce lemme est la motivation pour la définition suivante.

\begin{definition}
\label{definition-relative-H1-regular-immersion}
Soit  $f : X \to S$  un morphisme de schémas. Soit  $i : Z \to X$  une immersion.
\begin{enumerate}
\item On dit que  $i$  est une  {\it  immersion quasi-régulière
relative }  si  $Z \to S$  est plate et  $i$  est une immersion
quasi-régulière.
\item On dit que  $i$  est une  {\it  immersion
 $H_1$-régulière relative }  si  $Z \to S$  est plate et
 $i$  est une immersion  $H_1$-régulière.
\end{enumerate}
\end{definition}

\noindent
Nous avertissons le lecteur que cela peut être une notation non standard. Le lemme
 \ref{lemma-relative-regular-immersion}  garantit que les immersions quasi-régulières relatives (resp.
\ $H_1$-régulières) sont préservées sous tout changement de base.
Une immersion relative  $H_1$-régulière est une immersion quasi-régulière
relative, voir le lemme  \ref{lemma-regular-quasi-regular-immersion}. Veuillez consulter le lemme  \ref{lemma-flat-relative-H1-regular}  (ou
le lemme  \ref{lemma-relative-regular-immersion-flat-in-neighbourhood}) qui montre que si  $Z \to X$  est une immersion relative
 $H_1$-régulière (ou quasi-régulière) et que le schéma ambiant est (plat et)
localement de présentation finie sur  $S$, alors  $Z \to X$  est en fait
une immersion régulière et il en reste de même après tout changement de base.

\begin{lemma}
\label{lemma-quasi-regular-immersion-flat-at-x}
Soit  $f : X \to S$  un morphisme de schémas. Soit  $Z \to X$  une immersion
quasi-régulière relative. Si  $x \in Z$  et  $\mathcal{O}_{X, x}$  sont noethériens, alors
 $f$  est plat en  $x$.
\end{lemma}

\begin{proof}
Soit  $f_1, \ldots, f_r \in \mathcal{O}_{X, x}$  une suite quasi-régulière définissant l'idéal de  $Z$  en
 $x$. D'après le Lemme d'Algèbre  \ref{algebra-lemma-quasi-regular-regular} , nous savons que
 $f_1, \ldots, f_r$  est une suite régulière. Par conséquent,  $f_r$  est un élément
non-diviseur de zéro sur  $\mathcal{O}_{X, x}/(f_1, \ldots, f_{r - 1})$  tel que le quotient est un
 $\mathcal{O}_{S, f(x)}$-module plat. D'après le Lemme  \ref{lemma-flat-at-x} , nous en concluons que
 $\mathcal{O}_{X, x}/(f_1, \ldots, f_{r - 1})$  est un  $\mathcal{O}_{S, f(x)}$-module plat. En continuant par induction, nous
trouvons que  $\mathcal{O}_{X, x}$  est un  $\mathcal{O}_{S, s}$-module plat.
\end{proof}

\begin{lemma}
\label{lemma-relative-regular-immersion-flat-in-neighbourhood}
Soit  $X \to S$  un morphisme de schémas. Soit  $Z \to X$  une immersion.
Supposons
\begin{enumerate}
\item que $X \to S$  est plat et localement de présentation finie,
\item que $Z \to X$  est une immersion quasi-régulière relative.
\end{enumerate}
Alors  $Z \to X$  est une immersion régulière et la même chose reste vraie après
tout changement de base.
\end{lemma}

\begin{proof}
Choisissez  $x \in Z$  avec l'image  $s \in S$. Pour le prouver, il suffit de
trouver un voisinage affine de  $x$  contenu dans  $U$  tel que le
résultat soit valide sur cet ouvert affine. Ainsi, nous pouvons supposer que
 $X$  est affine et qu'il existe une suite quasi-régulière  $f_1, \ldots, f_r \in \Gamma(X, \mathcal{O}_X)$ 
telle que  $Z = V(f_1, \ldots, f_r)$. D'après More on Algebra, Lemme  \ref{more-algebra-lemma-relative-regular-immersion-algebra} , la suite
 $f_1|_{X_s}, \ldots, f_r|_{X_s}$  est une suite quasi-régulière dans  $\Gamma(X_s, \mathcal{O}_{X_s})$. Comme  $X_s$ 
est noethérien, cela implique, éventuellement après avoir légèrement réduit
 $X$ , que  $f_1|_{X_s}, \ldots, f_r|_{X_s}$  est une suite régulière, voir Algebra, Lemmata
 \ref{algebra-lemma-quasi-regular-regular}  et  \ref{algebra-lemma-regular-sequence-in-neighbourhood}. D'après le Lemme  \ref{lemma-fibre-Cartier} , il découle que
 $Z_1 = V(f_1) \subset X$  est un diviseur de Cartier effectif relatif, encore après avoir
éventuellement un peu rétréci  $X$ . En appliquant à nouveau le même lemme,
mais cette fois à  $Z_2 = V(f_1, f_2) \subset Z_1$ , nous voyons que  $Z_2 \subset Z_1$  est un diviseur de
Cartier effectif relatif. Et ainsi de suite jusqu'à ce que l'on atteigne
 $Z = Z_n = V(f_1, \ldots, f_n)$. Comme être un diviseur de Cartier effectif relatif se conserve par
tout changement de base arbitraire, voir le Lemme  \ref{lemma-relative-Cartier}, nous voyons
également que l'énoncé final du lemme est vrai.
\end{proof}

\begin{remark}
\label{remark-relative-regular-immersion-elements}
La codimension d'une immersion quasi-régulière relative, si elle est constante, ne
change pas après un changement de base. En fait, si nous avons un morphisme
d'anneaux  $A \to B$  et une suite quasi-régulière  $f_1, \ldots, f_r \in B$  telle que
 $B/(f_1, \ldots, f_r)$  est plate sur  $A$, alors pour tout morphisme d'anneaux
 $A \to A'$  nous avons une suite quasi-régulière  $f_1 \otimes 1, \ldots, f_r \otimes 1$  dans  $B' = B \otimes_A A'$ 
selon More on Algebra, Lemme  \ref{more-algebra-lemma-relative-regular-immersion-algebra}  (qui a été utilisé dans la démonstration
du Lemme  \ref{lemma-relative-regular-immersion}  ci-dessus). Maintenant, la démonstration du Lemme
 \ref{lemma-relative-regular-immersion-flat-in-neighbourhood}  montre que si  $A \to B$  est plate et localement de présentation
finie, alors pour tout idéal premier  $\mathfrak q' \subset B'$  la suite  $f_1 \otimes 1, \ldots, f_r \otimes 1$  est même
une suite régulière dans l'anneau local  $B'_{\mathfrak q'}$.
\end{remark}

\begin{lemma}
\label{lemma-flat-relative-H1-regular}
Soit  $X \to S$  un morphisme de schémas. Soit  $Z \to X$  une immersion
 $H_1$-régulière relative. Supposons que  $X \to S$  soit localement de
présentation finie. Alors
\begin{enumerate}
\item il existe un sous-schéma ouvert  $U \subset X$  tel que  $Z \subset U$  et tel que
 $U \to S$  soit plat, et
\item $Z \to X$  est une immersion régulière et il en reste de même après tout
changement de base.
\end{enumerate}
\end{lemma}

\begin{proof}
Choisir  $x \in Z$. Pour prouver (1), il suffit de trouver un voisinage ouvert
 $U \subset X$  de  $x$  tel que  $U \to S$  soit plat. Ainsi, le lemme se
réduit au cas où  $X = \Spec(B)$  et  $S = \Spec(A)$  sont affines et où  $Z$  est
donné par une suite  $H_1$-régulière  $f_1, \ldots, f_r \in B$. Par hypothèse,
 $B$  est une  $A$-algèbre de présentation finie et  $B/(f_1, \ldots, f_r)B$ 
est une  $A$-algèbre plate. Nous allons utiliser l'approximation
noethérienne absolue.

\medskip\noindent
Écrivez  $B = A[x_1, \ldots, x_n]/(g_1, \ldots, g_m)$. Supposons que  $f_i$  soit l'image de  $f_i' \in A[x_1, \ldots, x_n]$.
Choisissez une  $\mathbf{Z}$-sous-algèbre de type fini  $A_0 \subset A$  telle que tous
les coefficients des polynômes  $f_1', \ldots, f_r', g_1, \ldots, g_m$  soient dans  $A_0$. Nous
définissons  $B_0 = A_0[x_1, \ldots, x_n]/(g_1, \ldots, g_m)$  et nous notons  $f_{i, 0}$  l'image de  $f_i'$  dans
 $B_0$. Ensuite  $B = B_0 \otimes_{A_0} A$  et
$$
B/(f_1, \ldots, f_r) =
B_0/(f_{0, 1}, \ldots, f_{0, r}) \otimes_{A_0} A.
$$
. Par Algèbre, Lemme  \ref{algebra-lemma-flat-finite-presentation-limit-flat} , nous pouvons, après avoir agrandi  $A_0$,
supposer que  $B_0/(f_{0, 1}, \ldots, f_{0, r})$  est plat sur  $A_0$. Il se peut qu'à ce moment le
groupe de cohomologie de Koszul  $H_1(K_\bullet(B_0, f_{0, 1}, \ldots, f_{0, r}))$  ne soit pas nul. D'un autre côté,
comme  $B_0$  est Noethérien, c'est un  $B_0$-module de type fini.
Soient  $\xi_1, \ldots, \xi_n \in H_1(K_\bullet(B_0, f_{0, 1}, \ldots, f_{0, r}))$  des générateurs. Soit  $A_0 \subset A_1 \subset A$  une
 $\mathbf{Z}$-sous-algèbre de type fini plus grande de  $A$. Notons
 $f_{1, i}$  l'image de  $f_{0, i}$  dans  $B_1 = B_0 \otimes_{A_0} A_1$. Par Plus sur Algèbre,
Lemme  \ref{more-algebra-lemma-base-change-H1-regular} , l'application
$$
H_1(K_\bullet(B_0, f_{0, 1}, \ldots, f_{0, r})) \otimes_{A_0} A_1
\longrightarrow
H_1(K_\bullet(B_1, f_{1, 1}, \ldots, f_{1, r}))
$$
est surjective. De plus, il est clair que le colimite (sur toutes les choix de
 $A_1$  comme ci-dessus) des complexes  $K_\bullet(B_1, f_{1, 1}, \ldots, f_{1, r})$  est le complexe
 $K_\bullet(B, f_1, \ldots, f_r)$  qui est acyclique en degré  $1$. Par conséquent,
$$
\colim_{A_0 \subset A_1 \subset A}
H_1(K_\bullet(B_1, f_{1, 1}, \ldots, f_{1, r}))
= 0
$$
d'après Algèbre, Lemme  \ref{algebra-lemma-directed-colimit-exact}. Ainsi, nous pouvons trouver un choix de
 $A_1$  tel que  $\xi_1, \ldots, \xi_n$  se projette tous vers zéro dans  $H_1(K_\bullet(B_1, f_{1, 1}, \ldots, f_{1, r}))$. En
d'autres termes, le groupe de cohomologie de Koszul  $H_1(K_\bullet(B_1, f_{1, 1}, \ldots, f_{1, r}))$  est nul.

\medskip\noindent
Considérons le morphisme de schémas affines  $X_1 \to S_1$  égal à  $\Spec$  du
morphisme d'anneaux  $A_1 \to B_1$  et  $Z_1 = \Spec(B_1/(f_{1, 1}, \ldots, f_{1, r}))$. Puisque  $B = B_1 \otimes_{A_1} A$,
c'est-à-dire  $X = X_1 \times_{S_1} S$, et de même  $Z = Z_1 \times_S S_1$, il suffit maintenant de
prouver (1) pour  $X_1 \to S_1$  et l'immersion  $H_1$-relative
 $Z_1 \to X_1$-régulière, voir Morphismes, Lemme  \ref{morphisms-lemma-base-change-module-flat}. Nous avons donc
réduit au cas où  $X \to S$  est un morphisme de type fini de schémas de Noether.
Dans ce cas, nous savons que  $X \to S$  est plat en tout point de  $Z$ 
d'après le Lemme  \ref{lemma-quasi-regular-immersion-flat-at-x}. Combiné avec le fait que le lieu plat est ouvert
dans ce cas, voir Algèbre, Théorème  \ref{algebra-theorem-openness-flatness} , nous voyons que (1) est vrai. La
partie (2) découle alors d'une application du Lemme  \ref{lemma-relative-regular-immersion-flat-in-neighbourhood}.
\end{proof}

\noindent
Si le schéma ambiant est plat et localement de présentation finie sur la base,
alors nous pouvons caractériser une immersion quasi-régulière relative en termes
de ses fibres.

\begin{lemma}
\label{lemma-fibre-quasi-regular}
Soit  $\varphi : X \to S$  un morphisme plat qui est localement de présentation finie. Soit
 $T \subset X$  un sous-schéma fermé. Soit  $x \in T$  avec une image  $s \in S$.
\begin{enumerate}
\item Si  $T_s \subset X_s$  est une immersion quasi-régulière dans un voisinage de
 $x$, alors il existe un ouvert  $U \subset X$  et une immersion
quasi-régulière relative  $Z \subset U$  tels que  $Z_s = T_s \cap U_s$  et  $T \cap U \subset Z$.
\item Si  $T_s \subset X_s$  est une immersion quasi-régulière dans un voisinage de
 $x$, le morphisme  $T \to X$  est de présentation finie, et  $T \to S$ 
est plat en  $x$, alors nous pouvons choisir  $U$  et  $Z$ 
comme dans (1) tels que  $T \cap U = Z$.
\item Si  $T_s \subset X_s$  est une immersion quasi-régulière dans un voisinage de
 $x$, et  $T$  est définie par  $c$  équations dans un
voisinage de  $x$, où  $c = \dim_x(X_s) - \dim_x(T_s)$, alors nous pouvons choisir
 $U$  et  $Z$  comme dans (1) de telle sorte que  $T \cap U = Z$.
\end{enumerate}
Dans chaque cas,  $Z \to U$  est une immersion régulière selon le Lemme
 \ref{lemma-relative-regular-immersion-flat-in-neighbourhood}. En particulier, si  $T \to S$  est localement de présentation finie
et plate et que toutes les fibres  $T_s \subset X_s$  sont des immersions
quasi-régulières, alors  $T \to X$  est une immersion quasi-régulière relative.
\end{lemma}

\begin{proof}
Choisissez des voisinages ouverts affines  $\Spec(A)$  de  $s$  et
 $\Spec(B)$  de  $x$  tels que  $\varphi(\Spec(B)) \subset \Spec(A)$. Soit  $\mathfrak p \subset A$  l'idéal
premier correspondant à  $s$. Soit  $\mathfrak q \subset B$  l'idéal premier
correspondant à  $x$. Soit  $I \subset B$  l'idéal correspondant à
 $T$. D'après l'hypothèse initiale du lemme, nous savons que  $A \to B$ 
est plat et de présentation finie. L'hypothèse en (1) signifie que, après avoir
rétréci  $\Spec(B)$, nous pouvons supposer que  $I(B \otimes_A \kappa(\mathfrak p))$  est engendré par une
suite quasi-régulière d'éléments. Après avoir éventuellement localisé  $B$ 
en certains  $g \in B$,  $g \not \in \mathfrak q$ , nous pouvons supposer qu'il existe
 $f_1, \ldots, f_r \in I$  qui se projettent sur une suite quasi-régulière dans  $B \otimes_A \kappa(\mathfrak p)$  qui
génère  $I(B \otimes_A \kappa(\mathfrak p))$. Par l'Algèbre, les lemmes  \ref{algebra-lemma-quasi-regular-regular}  et  \ref{algebra-lemma-regular-sequence-in-neighbourhood} , nous
pouvons supposer qu'après une autre localisation,  $f_1, \ldots, f_r \in I$  forment une suite
régulière dans  $B \otimes_A \kappa(\mathfrak p)$. Selon le lemme  \ref{lemma-fibre-Cartier} , il s'ensuit que
 $Z_1 = V(f_1) \subset \Spec(B)$  est un diviseur de Cartier effectif relatif, encore une fois après
avoir éventuellement localisé  $B$. En appliquant à nouveau le même lemme,
mais maintenant à  $Z_2 = V(f_1, f_2) \subset Z_1$ , nous voyons que  $Z_2 \subset Z_1$  est un diviseur de
Cartier effectif relatif. Et ainsi de suite jusqu'à atteindre  $Z = Z_n = V(f_1, \ldots, f_n)$.
Ensuite,  $Z \to \Spec(B)$  est une immersion régulière et  $Z$  est plat sur
 $S$, en particulier  $Z \to \Spec(B)$  est une immersion quasi-régulière
relative sur  $\Spec(A)$. Cela prouve (1).

\medskip\noindent
Pour voir (2), considérons l'immersion fermée  $Z \to D$. L'application
surjective de anneaux  $u : \mathcal{O}_{D, x} \to \mathcal{O}_{Z, x}$  est une application d'algèbres locales plates
 $\mathcal{O}_{S, s}$ qui sont essentiellement de présentation finie, et qui devient un
isomorphisme après division par  $\mathfrak m_s$. Par conséquent, c'est un
isomorphisme, voir Algèbre, Lemme  \ref{algebra-lemma-mod-injective-general}. Il s'ensuit que  $Z \to D$  est
un isomorphisme dans un voisinage de  $x$, voir Algèbre, Lemme
 \ref{algebra-lemma-local-isomorphism}.

\medskip\noindent
Pour voir (3), après avoir éventuellement restreint  $U$ , nous pouvons
supposer que l'idéal de  $Z$  est engendré par une suite régulière
 $f_1, \ldots, f_r$  (voir notre construction de  $Z$  ci-dessus) et que l'idéal de
 $T$  est engendré par  $g_1, \ldots, g_c$. Nous affirmons que  $c = r$. À
savoir,
\begin{align*}
\dim_x(X_s) & = \dim(\mathcal{O}_{X_s, x}) +
\text{trdeg}_{\kappa(s)}(\kappa(x)), \\
\dim_x(T_s) & = \dim(\mathcal{O}_{T_s, x}) +
\text{trdeg}_{\kappa(s)}(\kappa(x)), \\
\dim(\mathcal{O}_{X_s, x}) & = \dim(\mathcal{O}_{T_s, x}) + r
\end{align*}
les deux premières égalités par l’Algèbre, Lemme  \ref{algebra-lemma-dimension-at-a-point-finite-type-field}  et la seconde par
 $r$  fois en appliquant l’Algèbre, Lemme  \ref{algebra-lemma-one-equation}. Comme  $T \subset Z$ 
nous voyons que  $f_i = \sum b_{ij} g_j$. Mais les idéaux de  $Z$  et  $T$ 
définissent le même sous-schéma fermé quasi-régulier de  $X_s$  dans un
voisinage de  $x$. Par conséquent, la matrice  $(b_{ij}) \bmod \mathfrak m_x$  est inversible
(certains détails omis). Ainsi  $(b_{ij})$  est inversible dans un voisinage
ouvert de  $x$. En d’autres termes,  $T \cap U = Z$  après avoir restreint
 $U$.

\medskip\noindent
Les déclarations finales du lemme suivent immédiatement de la partie (2),
combinées avec le fait que  $Z \to S$  est localement de présentation finie si et
seulement si  $Z \to X$  est de présentation finie, voir Morphismes, Lemme
 \ref{morphisms-lemma-composition-finite-presentation}  et  \ref{morphisms-lemma-finite-presentation-permanence}.
\end{proof}

\noindent
Le lemme suivant est une amélioration de Morphismes, Lemme  \ref{morphisms-lemma-section-smooth-morphism}.

\begin{lemma}
\label{lemma-section-smooth-regular-immersion}
Soit  $f : X \to S$  un morphisme lisse de schémas. Soit  $\sigma : S \to X$  une section de
 $f$. Alors  $\sigma$  est une immersion régulière.
\end{lemma}

\begin{proof}
D'après Schémas, Lemme  \ref{schemes-lemma-semi-diagonal}  le morphisme  $\sigma$  est une immersion.
Après avoir remplacé  $X$  par un voisinage ouvert de  $\sigma(S)$ , nous
pouvons supposer que  $\sigma$  est une immersion fermée. Soit  $T = \sigma(S)$  le
sous-schéma fermé correspondant de  $X$. Puisque  $T \to S$  est un
isomorphisme, il est plat et de présentation finie. De plus, un morphisme lisse
est plat et localement de présentation finie, voir Morphismes, Lemme  \ref{morphisms-lemma-smooth-flat} 
et  \ref{morphisms-lemma-smooth-locally-finite-presentation}. Ainsi, selon le Lemme  \ref{lemma-fibre-quasi-regular}, il suffit de montrer que
 $T_s \subset X_s$  est un sous-schéma fermé quasi-régulier. Cela découle immédiatement
des morphismes, Lemme  \ref{morphisms-lemma-section-smooth-morphism} , mais nous pouvons aussi le voir directement
comme suit. Soit  $k$  un corps et soit  $A$  une
 $k$-algèbre lisse. Soit  $\mathfrak m \subset A$  un idéal maximal dont le corps
résiduel est  $k$. Alors  $\mathfrak m$  est engendré par une suite
quasi-régulière, éventuellement après avoir remplacé  $A$  par  $A_g$ 
pour certains  $g \in A$,  $g \not \in \mathfrak m$. Dans Algèbre, Lemme  \ref{algebra-lemma-characterize-smooth-over-field}  nous
avons prouvé que  $A_{\mathfrak m}$  est un anneau local régulier, donc  $\mathfrak mA_{\mathfrak m}$  est
engendré par une suite régulière. Cela implique en effet que  $\mathfrak m$  est
engendré par une suite régulière (après avoir remplacé  $A$  par
 $A_g$  pour certains  $g \in A$,  $g \not \in \mathfrak m$), voir Algèbre, Lemme
 \ref{algebra-lemma-regular-sequence-in-neighbourhood}.
\end{proof}

\noindent
Le lemme suivant a une sorte de réciproque, voir Lemme  \ref{lemma-push-regular-immersion-thru-smooth}.

\begin{lemma}
\label{lemma-lift-regular-immersion-to-smooth}
Soit
$$
\xymatrix{
Y \ar[rd]_j \ar[rr]_i & & X \ar[ld] \\
& S
}
$$
soit un diagramme commutatif de morphismes de schémas. Supposons que  $X \to S$ 
soit lisse, et que  $i$,  $j$  soient des immersions. Si
 $j$  est une immersion régulière (resp. \ Koszul-régulière,
 $H_1$-régulière, quasi-régulière), alors il en est de même pour
 $i$.
\end{lemma}

\begin{proof}
Nous pouvons écrire  $i$  comme la composition
$$
Y \to Y \times_S X \to X
$$
. D'après le Lemme  \ref{lemma-section-smooth-regular-immersion} , la première flèche est une immersion régulière.
La deuxième flèche est un changement de base plat de  $Y \to S$, par conséquent
c'est une immersion régulière (resp. \ Koszul-régulière,
 $H_1$-régulière, quasi-régulière), voir Lemme  \ref{lemma-flat-base-change-regular-immersion}. Nous concluons
par une application du Lemme  \ref{lemma-composition-regular-immersion}.
\end{proof}

\begin{lemma}
\label{lemma-immersion-lci-into-smooth-regular-immersion}
Soit
$$
\xymatrix{
Y \ar[rd] \ar[rr]_i & & X \ar[ld] \\
& S
}
$$
un diagramme commutatif de morphismes de schémas. Supposons que  $Y \to S$  soit
syntomique,  $X \to S$  lisse, et  $i$  une immersion. Alors  $i$ 
est une immersion régulière.
\end{lemma}

\begin{proof}
Après avoir remplacé  $X$  par un voisinage ouvert de  $i(Y)$ , nous
pouvons supposer que  $i$  est une immersion fermée. Soit  $T = i(Y)$  le
sous-schéma fermé correspondant de  $X$. Comme  $T \cong Y$ , le morphisme
 $T \to S$  est plat et de présentation finie (Morphismes, Lemmata  \ref{morphisms-lemma-syntomic-locally-finite-presentation} 
et  \ref{morphisms-lemma-syntomic-flat}). De plus, un morphisme lisse est plat et localement de
présentation finie (Morphismes, Lemmata  \ref{morphisms-lemma-smooth-flat}  et  \ref{morphisms-lemma-smooth-locally-finite-presentation}). Ainsi,
selon le Lemme  \ref{lemma-fibre-quasi-regular}, il suffit de montrer que  $T_s \subset X_s$  est un
sous-schéma fermé quasi-régulier. Comme  $X_s$  est localement de type fini
sur un corps, il est noethérien (Morphismes, Lemme  \ref{morphisms-lemma-finite-type-noetherian}). Ainsi, nous
pouvons vérifier que  $T_s \subset X_s$  est une immersion quasi-régulière en les points,
voir le Lemme  \ref{lemma-Noetherian-scheme-regular-ideal}. Prenons  $t \in T_s$. Par Morphismes, Lemme  \ref{morphisms-lemma-local-complete-intersection} ,
l'anneau local  $\mathcal{O}_{T_s, t}$  est une intersection complète locale sur  $\kappa(s)$.
L'anneau local  $\mathcal{O}_{X_s, t}$  est régulier, voir Algèbre, Lemme  \ref{algebra-lemma-characterize-smooth-over-field}. Par
Algèbre, Lemme  \ref{algebra-lemma-lci-local} , nous voyons que le noyau de la surjection
 $\mathcal{O}_{X_s, t} \to \mathcal{O}_{T_s, t}$  est engendré par une suite régulière, ce que nous devions montrer.
\end{proof}

\begin{lemma}
\label{lemma-immersion-smooth-into-smooth-regular-immersion}
Soit
$$
\xymatrix{
Y \ar[rd] \ar[rr]_i & & X \ar[ld] \\
& S
}
$$
un diagramme commutatif de morphismes de schémas. Supposons que  $Y \to S$  est
lisse,  $X \to S$  lisse, et  $i$  une immersion. Alors  $i$  est
une immersion régulière.
\end{lemma}

\begin{proof}
Il s'agit d'un cas particulier du Lemme  \ref{lemma-immersion-lci-into-smooth-regular-immersion}  car un morphisme lisse est
syntomique, voir Morphismes, Lemme  \ref{morphisms-lemma-smooth-syntomic}.
\end{proof}

\begin{lemma}
\label{lemma-push-regular-immersion-thru-smooth}
Soit
$$
\xymatrix{
Y \ar[rd]_j \ar[rr]_i & & X \ar[ld] \\
& S
}
$$
un diagramme commutatif de morphismes de schémas. Supposons  $X \to S$  lisse et
 $i$  et  $j$  des immersions. Si  $i$  est une immersion de
Koszul-régulière (resp. \ $H_1$-régulière, quasi-régulière),
alors il en est de même pour  $j$.
\end{lemma}

\begin{proof}
Nous utiliserons le Lemme  \ref{lemma-regular-quasi-regular-immersion}  sans autre mention. Soit  $y \in Y$  un
point quelconque. Posons  $x = i(y)$  et posons  $s = j(y)$. Il suffit de prouver
le résultat après avoir remplacé  $X$  et  $S$  par des voisinages
ouverts  $U$  et  $V$  de  $x$  et  $s$  et
 $Y$  par un voisinage ouvert de  $y$  dans  $i^{-1}(U) \cap j^{-1}(V)$.

\medskip\noindent
Nous prouvons d'abord le résultat pour  $X = \mathbf{A}^n_S$. Après avoir remplacé
 $S$  par un ouvert affine  $V$  et remplacé  $Y$  par
 $j^{-1}(V)$ , nous pouvons supposer que  $j$  est une immersion fermée et
que  $S$  est affine. Écrivons  $S = \Spec(A)$. Alors  $j : Y \to S$  définit un
isomorphisme de  $Y$  vers le sous-schéma fermé  $\Spec(A/I)$  pour un
certain idéal  $I \subset A$. La carte  $i : Y = \Spec(A/I) \to
\mathbf{A}^n_S = \Spec(A[x_1, \ldots, x_n])$  correspond à un homomorphisme
d'algèbres  $A$-algèbre  $i^\sharp : A[x_1, \ldots, x_n] \to A/I$. Choisir  $a_i \in A$  qui se mappe à
 $i^\sharp(x_i)$  dans  $A/I$. Observez que l'idéal de l'immersion fermée
 $i$  est
$$
J = (x_1 - a_1, \ldots, x_n - a_n) + IA[x_1, \ldots, x_n].
$$
. Posons  $K = (x_1 - a_1, \ldots, x_n - a_n)$. Nous affirmons que la suite
$$
0 \to K/KJ \to J/J^2 \to J/(K + J^2) \to 0
$$
est exacte scindée. Pour voir cela, notez que  $K/K^2$  est libre avec pour
base  $x_i - a_i$  sur l'anneau  $A[x_1, \ldots, x_n]/K \cong A$. Par conséquent,  $K/KJ$  est
libre avec la même base sur l'anneau  $A[x_1, \ldots, x_n]/J \cong A/I$. D'autre part, prendre les
dérivées donne une carte
$$
\text{d}_{A[x_1, \ldots, x_n]/A} :
J/J^2
\longrightarrow
\Omega_{A[x_1, \ldots, x_n]/A} \otimes_{A[x_1, \ldots, x_n]}
A[x_1, \ldots, x_n]/J
$$
qui associe les générateurs  $x_i - a_i$  aux éléments de base  $\text{d}x_i$  du
module libre sur la droite. La revendication s'ensuit. De plus, notez que
 $x_1 - a_1, \ldots, x_n - a_n$  est une suite régulière dans  $A[x_1, \ldots, x_n]$  avec un anneau quotient
 $A[x_1, \ldots, x_n]/(x_1 - a_1, \ldots, x_n - a_n) \cong A$. Ainsi, nous avons une factorisation
$$
Y \to V(x_1 - a_1, \ldots, x_n - a_n) \to \mathbf{A}^n_S
$$
de notre immersion fermée  $i$  où la composition est Koszul-régulière
(resp. \ $H_1$-régulière, quasi-régulière), la deuxième flèche
est une immersion régulière, et la suite conormale associée est scindée.
Maintenant, le résultat s'ensuit du Lemme  \ref{lemma-permanence-regular-immersion}.

\medskip\noindent
Ensuite, nous prouvons que le résultat vaut si  $i$  est
 $H_1$-régulier. À savoir, en rétrécissant comme dans le premier paragraphe
de la preuve, nous pouvons supposer que  $Y$,  $X$ et  $S$ 
sont affines. Dans ce cas, nous pouvons choisir une immersion fermée  $h : X \to \mathbf{A}^n_S$ 
au-dessus de  $S$  pour un certain  $n$. Notez que  $h$  est
une immersion régulière d'après le Lemme  \ref{lemma-immersion-smooth-into-smooth-regular-immersion}. Par conséquent,
 $h \circ i$  est une immersion  $H_1$-régulière, voir le Lemme  \ref{lemma-composition-regular-immersion} 
(notez que cette étape ne fonctionne pas dans le « cas quasi-régulier »). Ainsi,
nous réduisons au cas affine de  $X = \mathbf{A}^n_S$  et  $S$  que nous avons montré
ci-dessus.

\medskip\noindent
Enfin, supposons que  $i$  soit quasi-régulier. Après rétrécissement comme
dans le premier paragraphe de la démonstration, nous pouvons utiliser les
morphismes, le lemme  \ref{morphisms-lemma-smooth-etale-over-affine-space}  pour factoriser  $f$  en  $X \to \mathbf{A}^n_S \to S$  où
le premier morphisme  $X \to \mathbf{A}^n_S$  est  \'étale. Cela réduit le problème
aux deux cas (a)  $X = \mathbf{A}^n_S$  et (b)  $f$  est  \'étale. Le cas
(a) a été traité dans le deuxième paragraphe de la démonstration. Le cas (b) est
traité dans le paragraphe suivant.

\medskip\noindent
Supposons que  $f$  soit  \'étale. Après réduction, nous pouvons
supposer que  $X$,  $Y$ et  $S$  sont des immersions
affines  $i$  et  $j$  fermées (petit détail omis). Disons
 $S = \Spec(A)$,  $X = \Spec(B)$  et  $Y = \Spec(B/J) = \Spec(A/I)$. En rétrécissant encore, nous pouvons
supposer que  $J$  est engendré par une suite quasi-régulière.
L'application d'anneaux  $A \to B$  est  \'étale, donc formellement
 \'étale (Algèbre, Lemme  \ref{algebra-lemma-formally-etale-etale}). Ainsi  $\bigoplus I^n/I^{n + 1} \cong \bigoplus J^n/J^{n + 1}$  par Algèbre,
Lemme  \ref{algebra-lemma-formally-etale-lift-infinitesimal}. Comme  $J$  est engendré par une suite quasi-régulière,
il en est de même pour  $I$. Cela termine la démonstration.
\end{proof}











\section{Fonctions et sections méromorphes}
\label{section-meromorphic-functions}

\noindent
Cette section contient uniquement les définitions générales et quelques résultats
élémentaires. Voir  \cite{misconceptions}  pour certains pièges possibles \footnote
{Danger, Will Robinson! }.

\medskip\noindent
Soit  $(X, \mathcal{O}_X)$  un espace localement annelé. Pour tout ouvert  $U \subset X$ , nous
avons défini l'ensemble  $\mathcal{S}(U) \subset \mathcal{O}_X(U)$  des sections régulières de  $\mathcal{O}_X$  sur
 $U$, voir la Définition  \ref{definition-regular-section}. La restriction d'une section
régulière à un ouvert plus petit est régulière. Ainsi,  $\mathcal{S} : U \mapsto \mathcal{S}(U)$  est un
sous-faisceau (d'ensembles) de  $\mathcal{O}_X$. Nous notons parfois  $\mathcal{S} = \mathcal{S}_X$  si
nous voulons indiquer la dépendance à  $X$. De plus,  $\mathcal{S}(U)$  est un
sous-ensemble multiplicatif de l'anneau  $\mathcal{O}_X(U)$  pour chaque  $U$. Par
conséquent, nous pouvons considérer le préfaisceau d'anneaux
$$
U \longmapsto \mathcal{S}(U)^{-1} \mathcal{O}_X(U),
$$
, voir Modules, Lemme  \ref{modules-lemma-simple-invert}.

\begin{definition}
\label{definition-sheaf-meromorphic-functions}
Soit  $(X, \mathcal{O}_X)$  un espace localement annelé. Le faisceau
 {\it  des fonctions méromorphes sur  $X$} 
est le faisceau  {\it   $\mathcal{K}_X$}  associé au
préfaisceau affiché ci-dessus. Une  {\it  fonction méromorphe
}  sur  $X$  est une section globale de  $\mathcal{K}_X$.
\end{definition}

\noindent
Comme chaque élément de chaque  $\mathcal{S}(U)$  est un non-diviseur de zéro sur
 $\mathcal{O}_X(U)$ , nous voyons que l'application naturelle de faisceaux d'anneaux
 $\mathcal{O}_X \to \mathcal{K}_X$  est injective.

\begin{example}
\label{example-no-change}
Soit  $A = \mathbf{C}[x, \{y_\alpha\}_{\alpha \in \mathbf{C}}]/
((x - \alpha)y_\alpha, y_\alpha y_\beta)$. Tout élément de  $A$  peut être écrit de manière unique
comme  $f(x) + \sum \lambda_\alpha y_\alpha$  avec  $f(x) \in \mathbf{C}[x]$  et  $\lambda_\alpha \in \mathbf{C}$. Soit  $X = \Spec(A)$. Dans ce
cas,  $\mathcal{O}_X = \mathcal{K}_X$, puisque sur tout ouvert affine  $D(f)$ , l'anneau
 $A_f$  tout non-diviseur de zéro est une unité (preuve omise).
\end{example}

\noindent
Soit  $(X, \mathcal{O}_X)$  un espace localement annelé. Soit  $\mathcal{F}$  un faisceau de
 $\mathcal{O}_X$-modules. Considérons le préfaisceau  $U \mapsto \mathcal{S}(U)^{-1}\mathcal{F}(U)$. Sa faisceautisation
est le faisceau  $\mathcal{F} \otimes_{\mathcal{O}_X} \mathcal{K}_X$, voir Modules, Lemme  \ref{modules-lemma-simple-invert-module}.

\begin{definition}
\label{definition-meromorphic-section}
Soit  $X$  un espace localement annelé. Soit  $\mathcal{F}$  un faisceau de
 $\mathcal{O}_X$-modules.
\begin{enumerate}
\item Nous notons  $\mathcal{K}_X(\mathcal{F})$  le faisceau de  $\mathcal{K}_X$-modules qui est la
faisceautisation du préfaisceau  $U \mapsto \mathcal{S}(U)^{-1}\mathcal{F}(U)$. Équivalemment  $\mathcal{K}_X(\mathcal{F}) =
\mathcal{F} \otimes_{\mathcal{O}_X} \mathcal{K}_X$  (voir
ci-dessus).
\item Une section méromorphe  {\it  de  $\mathcal{F}$}  est
une section globale de  $\mathcal{K}_X(\mathcal{F})$.
\end{enumerate}
\end{definition}

\noindent
En particulier, nous avons
$$
\mathcal{K}_X(\mathcal{F})_x
=
\mathcal{F}_x \otimes_{\mathcal{O}_{X, x}} \mathcal{K}_{X, x}
=
\mathcal{S}_x^{-1}\mathcal{F}_x
$$
pour tout point  $x \in X$. Cependant, il faut faire attention car il se peut
que  $\mathcal{S}_x$  ne soit pas l'ensemble des diviseurs de zéro dans l'anneau local
 $\mathcal{O}_{X, x}$. À savoir, il existe toujours une application injective
$$
\mathcal{K}_{X, x} \longrightarrow Q(\mathcal{O}_{X, x})
$$
au quotient total de l'anneau. Il est également surjectif si et seulement si
 $\mathcal{S}_x$  est l'ensemble des diviseurs de zéro non nuls dans  $\mathcal{O}_{X, x}$. Les
faisceaux de sections méromorphes ne sont pas en général des modules
quasi-cohérents, mais ils possèdent certaines propriétés en commun avec les
modules quasi-cohérents.

\begin{definition}
\label{definition-pullback-meromorphic-sections}
Soit  $f : (X, \mathcal{O}_X) \to (Y, \mathcal{O}_Y)$  un morphisme d'espaces localement annelés. Nous disons que
 {\it  les rétrodéductions de fonctions méromorphes sont
définies pour  $f$}  si pour chaque paire d'ouverts
 $U \subset X$,  $V \subset Y$  telle que  $f(U) \subset V$, et pour toute section
 $s \in \Gamma(V, \mathcal{S}_Y)$  le rétrodéduit  $f^\sharp(s) \in \Gamma(U, \mathcal{O}_X)$  est un élément de  $\Gamma(U, \mathcal{S}_X)$.
\end{definition}

\noindent
Dans ce cas, il existe une application induite  $f^\sharp : f^{-1}\mathcal{K}_Y \to \mathcal{K}_X$, en d'autres termes,
nous obtenons un diagramme commutatif de morphismes d'espaces annelés
$$
\xymatrix{
(X, \mathcal{K}_X) \ar[r] \ar[d]^f &
(X, \mathcal{O}_X) \ar[d]^f \\
(Y, \mathcal{K}_Y) \ar[r] &
(Y, \mathcal{O}_Y)
}
$$
Nous notons parfois  $f^*(s) = f^\sharp(s)$  pour une section  $s \in \Gamma(Y, \mathcal{K}_Y)$.

\begin{lemma}
\label{lemma-pullback-meromorphic-sections-defined}
Soit  $f : X \to Y$  un morphisme de schémas. Dans chacun des cas suivants, les tirés
en arrière des fonctions méromorphes sont définis.
\begin{enumerate}
\item chaque point faiblement associé de  $X$  se projette sur un point générique
d'une composante irréductible de  $Y$,
\item $X$,  $Y$  sont intègres et  $f$  est dominant,
\item $X$  est intègre et le point générique de  $X$  se projette sur le
point générique d'une composante irréductible de  $Y$,
\item $X$  est réduit et chaque point générique de chaque composante
irréductible de  $X$  se projette sur le point générique d'une composante
irréductible de  $Y$,
\item $X$  est localement noethérien et tout point associé de  $X$  se
projette sur un point générique d'une composante irréductible de  $Y$.
\item $X$  est localement noethérien, n'a pas de points immergés et tout point
générique d'une composante irréductible de  $X$  se projette sur le point
générique d'une composante irréductible de  $Y$, et
\item $f$  est plat.
\end{enumerate}
\end{lemma}

\begin{proof}
La question est locale sur  $X$  et  $Y$. Nous nous réduisons donc
au cas où  $X = \Spec(A)$,  $Y = \Spec(R)$  et  $f$  sont donnés par un morphisme
d'anneaux  $\varphi : R \to A$. D'après la caractérisation des sections régulières du
faisceau structural dans le Lemme  \ref{lemma-regular-section-structure-sheaf} , nous devons montrer que
 $R \to A$  envoie les non-diviseurs de zéro sur des non-diviseurs de zéro. Soit
 $t \in R$  un non-diviseur de zéro.

\medskip\noindent
Si  $R \to A$  est plat, alors  $t : R \to R$  étant injectif montre que
 $t : A \to A$  est injectif. Cela prouve (7).

\medskip\noindent
Dans les autres cas, nous notons que  $t$  n'est contenu dans aucun des
idéaux premiers minimaux de  $R$  (car chaque élément d'un idéal premier
minimal dans un anneau est un diviseur de zéro). Ainsi, dans le cas (1), nous
voyons que  $\varphi(t)$  n'est contenu dans aucun idéal idéal premier faiblement
associé de  $A$. Par conséquent, ce cas découle de Algèbre, Lemme
 \ref{algebra-lemma-weakly-ass-zero-divisors}. Le cas (5) est un cas particulier de (1) selon le Lemme
 \ref{lemma-ass-weakly-ass}. Le cas (6) découle de (5) et des définitions. Le cas (4) est un cas
particulier de (1) selon le Lemme  \ref{lemma-weakass-reduced}. Les cas (2) et (3) sont des cas
particuliers du (4).
\end{proof}


\begin{lemma}
\label{lemma-meromorphic-weakass-finite}
Soit  $X$  un schéma tel que
\begin{enumerate}
\item[(a)] chaque point faiblement associé de  $X$  est un point générique d'une
composante irréductible de  $X$, et
\item[(b)] tout ouvert quasi-compact a un nombre fini de composantes irréductibles.
\end{enumerate}
Soit  $X^0$  l'ensemble des points génériques des composantes irréductibles
de  $X$. Alors nous avons
$$
\mathcal{K}_X =
\bigoplus\nolimits_{\eta \in X^0} j_{\eta, *}\mathcal{O}_{X, \eta} =
\prod\nolimits_{\eta \in X^0} j_{\eta, *}\mathcal{O}_{X, \eta}
$$
où  $j_\eta : \Spec(\mathcal{O}_{X, \eta}) \to X$  est l'application canonique des Schémas, Section  \ref{schemes-section-points}. De
plus,
\begin{enumerate}
\item $\mathcal{K}_X$  est un faisceau quasi-cohérent d'algèbres  $\mathcal{O}_X$,
\item pour tout  $\mathcal{O}_X$-module quasi-cohérent  $\mathcal{F}$  le faisceau
$$
\mathcal{K}_X(\mathcal{F}) =
\bigoplus\nolimits_{\eta \in X^0} j_{\eta, *}\mathcal{F}_\eta =
\prod\nolimits_{\eta \in X^0} j_{\eta, *}\mathcal{F}_\eta
$$
de sections méromorphes de  $\mathcal{F}$  est quasi-cohérent,
\item $\mathcal{S}_x \subset \mathcal{O}_{X, x}$  est l'ensemble des non-diviseurs de zéro pour tout  $x \in X$,
\item $\mathcal{K}_{X, x}$  est l'anneau total des fractions de  $\mathcal{O}_{X, x}$  pour tout
 $x \in X$,
\item $\mathcal{K}_X(U)$  est égal à l'anneau total des fractions de  $\mathcal{O}_X(U)$  pour tout
ouvert affine  $U \subset X$,
\item l'anneau des fonctions rationnelles de  $X$  (Morphismes, Définition
 \ref{morphisms-definition-rational-function}) est l'anneau des fonctions méromorphes sur  $X$, en formule:
 $R(X) = \Gamma(X, \mathcal{K}_X)$.
\end{enumerate}
\end{lemma}

\begin{proof}
Remarquez qu'une somme directe localement finie de faisceaux de modules est égale
au produit puisque vous pouvez le vérifier sur les fibres, par exemple. Ensuite,
puisque  $\mathcal{K}_X(\mathcal{F}) =
\mathcal{F} \otimes_{\mathcal{O}_X} \mathcal{K}_X$ , nous voyons que (2) découle des autres énoncés. De plus,
remarquez que la partie (6) découle de l'énoncé initial du lemme et du Morphismes,
Lemme  \ref{morphisms-lemma-integral-scheme-rational-functions}  lorsque  $X^0$  est fini; le cas général de (6) découle de
cela par recollement (argument omis).

\medskip\noindent
Soit  $j : Y = \coprod\nolimits_{\eta \in X^0} \Spec(\mathcal{O}_{X, \eta}) \to X$  le coproduit des morphismes  $j_\eta$. Nous devons montrer
que  $\mathcal{K}_X = j_*\mathcal{O}_Y$. Notez d'abord que  $\mathcal{K}_Y = \mathcal{O}_Y$  puisque  $Y$  est une
réunion disjointe de spectres d'anneaux locaux de dimension  $0$: dans un
anneau local de dimension zéro, tout élément qui n'est pas un diviseur de zéro est
une unité. Ensuite, notez que les images réciproques de fonctions méromorphes sont
définies pour  $j$  d'après le Lemme  \ref{lemma-pullback-meromorphic-sections-defined}. Cela donne une
application
$$
\mathcal{K}_X \longrightarrow j_*\mathcal{O}_Y.
$$
Soit  $\Spec(A) = U \subset X$  un ouvert affine. Alors  $A$  est un anneau avec un
nombre fini de premiers minimaux  $\mathfrak q_1, \ldots, \mathfrak q_t$  et chaque idéal premier faiblement
associé de  $A$  est l'un des  $\mathfrak q_i$. Nous obtenons  $Q(A) = \prod A_{\mathfrak q_i}$  par
Algèbre, Lemme  \ref{algebra-lemma-total-ring-fractions-no-embedded-points}  et  \ref{algebra-lemma-weakly-ass-zero-divisors}. En d'autres termes, déjà la valeur du
préfaisceau  $U \mapsto \mathcal{S}(U)^{-1}\mathcal{O}_X(U)$  coïncide avec  $j_*\mathcal{O}_Y(U)$  sur notre ouvert affine
 $U$. Par conséquent, l'application affichée est un isomorphisme, ce qui
prouve la première égalité affichée dans l'énoncé du lemme.

\medskip\noindent
Enfin, nous prouvons (1), (3), (4) et (5). Pour la partie (5), nous avons vu au
cours de la démonstration que  $\mathcal{K}_X = j_*\mathcal{O}_Y$. Le morphisme  $j$  est
quasi-compact par notre hypothèse que l'ensemble des composantes irréductibles de
 $X$  est localement fini. Ainsi,  $j$  est quasi-compact et
quasi-séparé (car  $Y$  est séparé). D'après Schémas, Lemme  \ref{schemes-lemma-push-forward-quasi-coherent} ,
 $j_*\mathcal{O}_Y$  est quasi-cohérent. Cela prouve (1). Soit  $x \in X$. Nous pouvons
choisir un voisinage affine ouvert  $U = \Spec(A)$  de  $x$  dont toutes les
composantes irréductibles passent par  $x$. Alors  $A \subset A_\mathfrak p$  car chaque
idéal idéal premier faiblement associé de  $A$  est contenu dans
 $\mathfrak p$ , donc les éléments de  $A \setminus \mathfrak p$  sont des non-diviseurs de zéro
d'après Algèbre, Lemme  \ref{algebra-lemma-weakly-ass-zero-divisors}. Il s'ensuit facilement que tout non-diviseur
de zéro de  $A_\mathfrak p$  est l'image d'un non-diviseur de zéro sur un voisinage
ouvert affine (possiblement plus petit) de  $x$. Cela prouve (3). La
partie (4) découle de la partie (3) en calculant les germes.
\end{proof}

\begin{definition}
\label{definition-regular-meromorphic-section}
Soit  $X$  un espace localement annelé. Soit  $\mathcal{L}$  un
 $\mathcal{O}_X$-module inversible. Une section méromorphe  $s$  de
 $\mathcal{L}$  est dite  {\it  régulière}  si
l'application induite  $\mathcal{K}_X \to \mathcal{K}_X(\mathcal{L})$  est injective. En d'autres termes,  $s$ 
est une section régulière du  $\mathcal{K}_X$-module inversible  $\mathcal{K}_X(\mathcal{L})$, voir
Définition  \ref{definition-regular-section}.
\end{definition}

\noindent
Détaillons quand les sections méromorphes (régulières) peuvent être tirées en
arrière.

\begin{lemma}
\label{lemma-meromorphic-sections-pullback}
Soit  $f : X \to Y$  un morphisme d'espaces localement annelés. Supposons que les
images réciproques de fonctions méromorphes soient définies pour  $f$ 
(voir Définition  \ref{definition-pullback-meromorphic-sections}).
\begin{enumerate}
\item Soit  $\mathcal{F}$  un faisceau de  $\mathcal{O}_Y$-modules. Il existe une application
canonique de tiré en arrière  $f^* : \Gamma(Y, \mathcal{K}_Y(\mathcal{F})) \to
\Gamma(X, \mathcal{K}_X(f^*\mathcal{F}))$  pour les sections méromorphes de
 $\mathcal{F}$.
\item Soit  $\mathcal{L}$  un  $\mathcal{O}_Y$-module inversible. Une section méromorphe
régulière  $s$  de  $\mathcal{L}$  se tire en arrière en une section
méromorphe régulière  $f^*s$  de  $f^*\mathcal{L}$.
\end{enumerate}
\end{lemma}

\begin{proof}
Omis.
\end{proof}

\begin{lemma}
\label{lemma-regular-meromorphic-ideal-denominators}
Soit  $X$  un schéma. Soit  $\mathcal{L}$  un  $\mathcal{O}_X$-module inversible.
Soit  $s$  une section méromorphe régulière de  $\mathcal{L}$. Notons
 $\mathcal{I} \subset \mathcal{O}_X$  le faisceau d'idéaux défini par la règle
$$
\mathcal{I}(V)
=
\{f \in \mathcal{O}_X(V) \mid fs \in \mathcal{L}(V)\}.
$$
La formule a du sens puisque  $\mathcal{L}(V) \subset \mathcal{K}_X(\mathcal{L})(V)$. Ensuite,  $\mathcal{I}$  est un faisceau
quasi-cohérent d'idéaux et nous avons des applications injectives
$$
1 : \mathcal{I} \longrightarrow \mathcal{O}_X,
\quad
s : \mathcal{I} \longrightarrow \mathcal{L}
$$
dont les conoyaux sont supportés sur des sous-ensembles fermés et de nulle densité
de  $X$.
\end{lemma}

\begin{proof}
La question est locale sur  $X$. Par conséquent, nous pouvons supposer que
 $X = \Spec(A)$ et  $\mathcal{L} = \mathcal{O}_X$. Après un rétrécissement supplémentaire, nous pouvons
supposer que  $s = a/b$  avec  $a, b \in A$   {\it  tous deux
}  non-diviseurs de zéro dans  $A$. Posons  $I = \{x \in A \mid x(a/b) \in A\}$.

\medskip\noindent
Pour montrer que  $\mathcal{I}$  est quasi-cohérent, nous devons montrer que
 $I_f = \{x \in A_f \mid x(a/b) \in A_f\}$  pour chaque  $f \in A$. Si  $c/f^n \in A_f$,  $(c/f^n)(a/b) \in A_f$, alors nous
voyons que  $f^mc(a/b) \in A$  pour certains  $m$, donc  $c/f^n \in I_f$.
Inversement, il est facile de voir que  $I_f$  est contenu dans  $\{x \in A_f \mid x(a/b) \in A_f\}$.
Cela prouve la quasi-cohérence.

\medskip\noindent
Prouvons la déclaration finale. Il est clair que  $(b) \subset I$. Par conséquent,
 $V(I) \subset V(b)$  est un sous-ensemble nulle part dense puisque  $b$  est un
non-diviseur de zéro. Ainsi, le conoyau de  $1$  est supporté dans un
ensemble fermé nulle part dense. Le même argument fonctionne pour le conoyau de
 $s$  puisque  $s(b) = (a) \subset sI \subset A$.
\end{proof}

\begin{definition}
\label{definition-regular-meromorphic-ideal-denominators}
Soit  $X$  un schéma. Soit  $\mathcal{L}$  un  $\mathcal{O}_X$-module inversible.
Soit  $s$  une section méromorphe régulière de  $\mathcal{L}$. Le faisceau
d'idéaux  $\mathcal{I}$  construit dans le Lemme  \ref{lemma-regular-meromorphic-ideal-denominators}  est appelé le faisceau
d'idéaux de dénominateurs  {\it  de  $s$}.
\end{definition}




\section{Fonctions et sections méromorphes ; cas noethérien}
\label{section-meromorphic-noetherian}

\noindent
Pour les schémas localement noethériens, nous pouvons prouver certains résultats
concernant le faisceau de fonctions méromorphes. Cependant, il existe un exemple
dans  \cite{misconceptions}  montrant que  $\mathcal{K}_X$  n'a pas besoin d'être quasi-cohérent
pour un schéma noethérien  $X$.

\begin{lemma}
\label{lemma-meromorphic-section-restricts-to-zero}
Soit  $X$  un schéma quasi-compact. Soient  $h \in \Gamma(X, \mathcal{O}_X)$  et  $f \in \Gamma(X, \mathcal{K}_X)$ 
tels que  $f$  se restreigne à zéro sur  $X_h$. Alors  $h^n f = 0$ 
pour un certain  $n \gg 0$.
\end{lemma}

\begin{proof}
Nous pouvons trouver un recouvrement de  $X$  par des ouverts affines
 $U$  tels que  $f|_U = s^{-1}a$  avec  $a \in \mathcal{O}_X(U)$  et  $s \in \mathcal{S}(U)$. Comme
 $X$  est quasi-compact, nous pouvons le recouvrir par un nombre fini
d'ouverts affines de cette forme. Ainsi, il suffit de prouver le lemme lorsque
 $X = \Spec(A)$  et  $f = s^{-1}a$. Notez que  $s \in A$  est un non-diviseur de zéro,
il suffit donc de prouver le résultat lorsque  $f = a$. La condition
 $f|_{X_h} = 0$  implique que  $a$  se mappe à zéro dans  $A_h = \mathcal{O}_X(X_h)$  comme
 $\mathcal{O}_X \subset \mathcal{K}_X$. Ainsi  $h^na = 0$  pour un certain  $n > 0$  comme désiré.
\end{proof}

\begin{lemma}
\label{lemma-locally-Noetherian-K}
Soit  $X$  un schéma localement noethérien.
\begin{enumerate}
\item Pour tout  $x \in X$  nous avons que  $\mathcal{S}_x \subset \mathcal{O}_{X, x}$  est l'ensemble des non-zéros
diviseurs, et donc  $\mathcal{K}_{X, x}$  est l'anneau de quotients total de  $\mathcal{O}_{X, x}$.
\item Pour tout ouvert affine  $U \subset X$ , l'anneau  $\mathcal{K}_X(U)$  est égal à l'anneau de
quotients total de  $\mathcal{O}_X(U)$.
\end{enumerate}
\end{lemma}

\begin{proof}
Pour démontrer ce lemme, nous pouvons supposer que  $X$  est le spectre
d'un anneau noethérien  $A$. Disons que  $x \in X$  correspond à
 $\mathfrak p \subset A$.

\medskip\noindent
Preuve de (1). Il est clair que  $\mathcal{S}_x$  est contenu dans l'ensemble des
non-diviseurs de zéro de  $\mathcal{O}_{X, x} = A_\mathfrak p$. Pour la réciproque, soit  $f, g \in A$,
 $g \not \in \mathfrak p$  et supposons que  $f/g$  soit un non-diviseur de zéro dans
 $A_{\mathfrak p}$. Soit  $I = \{a \in A \mid af = 0\}$. Alors nous voyons que  $I_{\mathfrak p} = 0$  par
l'exactitude de la localisation. Puisque  $A$  est noethérien, nous voyons
que  $I$  est de type fini et donc que  $g'I = 0$  pour certains
 $g' \in A$,  $g' \not \in \mathfrak p$. Ainsi  $f$  est un non-diviseur de zéro dans
 $A_{g'}$, c'est-à-dire dans un voisinage ouvert de Zariski de  $\mathfrak p$.
Ainsi  $f/g$  est un élément de  $\mathcal{S}_x$.

\medskip\noindent
Preuve de (2). Soit  $f \in \Gamma(X, \mathcal{K}_X)$  une fonction méromorphe. Posons  $I = \{a \in A \mid af \in A\}$.
Fixons un premier  $\mathfrak p \subset A$  correspondant au point  $x \in X$. D'après (1),
nous pouvons écrire l'image de  $f$  dans la germe en  $\mathfrak p$  comme
 $a/b$,  $a, b \in A_{\mathfrak p}$  avec  $b \in A_{\mathfrak p}$  n'étant pas un diviseur de zéro.
Écrivons  $b = c/d$  avec  $c, d \in A$,  $d \not \in \mathfrak p$. Alors  $ad - cf$  est une
section de  $\mathcal{K}_X$  qui s'annule dans un voisinage ouvert de  $x$.
Disons qu'elle s'annule sur  $D(e)$  avec  $e \in A$,  $e \not \in \mathfrak p$. Alors
 $e^n(ad - cf) = 0$  pour un certain  $n \gg 0$  d'après le Lemme  \ref{lemma-meromorphic-section-restricts-to-zero}. Ainsi,
 $e^nc \in I$  et  $e^nc$  se projettent sur un diviseur de zéro non nul dans
 $A_{\mathfrak p}$. Soit  $\text{Ass}(A) = \{\mathfrak q_1, \ldots, \mathfrak q_t\}$  les premiers associés de  $A$. En
regardant  $IA_{\mathfrak q_i}$  et en utilisant Algèbre, Lemme  \ref{algebra-lemma-associated-primes-localize} , ce qui précède
dit que  $I \not \subset \mathfrak q_i$  pour chaque  $i$. Par Algèbre, Lemme  \ref{algebra-lemma-silly} , il
existe un élément  $z \in I$,  $z \not \in \bigcup \mathfrak q_i$. Par Algèbre, Lemme  \ref{algebra-lemma-ass-zero-divisors} ,
nous voyons que  $z$  n'est pas un diviseur de zéro sur  $A$. Par
conséquent,  $f = (zf)/z$  est un élément de l'anneau total des fractions de
 $A$. Cela prouve (2).
\end{proof}

\begin{lemma}
\label{lemma-quasi-coherent-K}
Soit  $X$  un schéma localement noethérien n'ayant pas de points immergés.
Soit  $X^0$  l'ensemble des points génériques des composantes irréductibles
de  $X$. Alors nous avons
$$
\mathcal{K}_X =
\bigoplus\nolimits_{\eta \in X^0} j_{\eta, *}\mathcal{O}_{X, \eta} =
\prod\nolimits_{\eta \in X^0} j_{\eta, *}\mathcal{O}_{X, \eta}
$$
où  $j_\eta : \Spec(\mathcal{O}_{X, \eta}) \to X$  est l'application canonique des Schémas, Section  \ref{schemes-section-points}. De
plus,
\begin{enumerate}
\item $\mathcal{K}_X$  est un faisceau quasi-cohérent d'algèbres  $\mathcal{O}_X$,
\item pour tout  $\mathcal{O}_X$-module quasi-cohérent  $\mathcal{F}$  le faisceau
$$
\mathcal{K}_X(\mathcal{F}) =
\bigoplus\nolimits_{\eta \in X^0} j_{\eta, *}\mathcal{F}_\eta =
\prod\nolimits_{\eta \in X^0} j_{\eta, *}\mathcal{F}_\eta
$$
des sections méromorphes de  $\mathcal{F}$  est quasi-cohérente, et
\item l'anneau des fonctions rationnelles de  $X$  est l'anneau des fonctions
méromorphes sur  $X$, sous une formule:  $R(X) = \Gamma(X, \mathcal{K}_X)$.
\end{enumerate}
\end{lemma}

\begin{proof}
Ce lemme est un cas particulier du Lemme  \ref{lemma-meromorphic-weakass-finite}  car, dans le cas localement
noethérien, les points faiblement associés sont la même chose que les points
associés par le Lemme  \ref{lemma-ass-weakly-ass}.
\end{proof}

\begin{lemma}
\label{lemma-regular-meromorphic-section-exists-noetherian}
Soit  $X$  un schéma localement noethérien n'ayant pas de points immergés.
Soit  $\mathcal{L}$  un  $\mathcal{O}_X$-module inversible. Alors  $\mathcal{L}$  possède
une section méromorphe régulière.
\end{lemma}

\begin{proof}
Pour chaque point générique  $\eta$  de  $X$ , choisissez un générateur
 $s_\eta$  du module libre de rang  $1$   $\mathcal{L}_\eta$  sur l'anneau local
artinien  $\mathcal{O}_{X, \eta}$. Il découle immédiatement de la description de  $\mathcal{K}_X$ 
et  $\mathcal{K}_X(\mathcal{L})$  dans le Lemme  \ref{lemma-quasi-coherent-K}  que  $s = \prod s_\eta$  est une section
méromorphe régulière de  $\mathcal{L}$.
\end{proof}

\begin{lemma}
\label{lemma-make-maps-regular-section}
Supposons données
\begin{enumerate}
\item $X$  un schéma localement noethérien,
\item $\mathcal{L}$  un  $\mathcal{O}_X$-module inversible,
\item $s$  une section méromorphe régulière de  $\mathcal{L}$, et
\item $\mathcal{F}$  cohérent sur  $X$  sans points associés immergés et
 $\text{Supp}(\mathcal{F}) = X$.
\end{enumerate}
Soit  $\mathcal{I} \subset \mathcal{O}_X$  l'idéal des dénominateurs de  $s$. Soit  $T \subset X$ 
l'union des supports de  $\mathcal{O}_X/\mathcal{I}$  et  $\mathcal{L}/s(\mathcal{I})$ , qui est un sous-ensemble
fermé et rare  $T \subset X$  selon le Lemme  \ref{lemma-regular-meromorphic-ideal-denominators}. Alors, il existe des
applications injectives canoniques
$$
1 : \mathcal{I}\mathcal{F} \to \mathcal{F}, \quad
s : \mathcal{I}\mathcal{F} \to \mathcal{F} \otimes_{\mathcal{O}_X}\mathcal{L}
$$
dont les conoyaux sont supportés sur  $T$.
\end{lemma}

\begin{proof}
Réduire au cas affine avec  $\mathcal{L} \cong \mathcal{O}_X$, et  $s = a/b$  avec  $a, b \in A$  tous
deux non diviseurs de zéro. La preuve de l'étape de réduction est omise. Écrire
 $\mathcal{F} = \widetilde{M}$. Soit  $I = \{x \in A \mid x(a/b) \in A\}$  de sorte que  $\mathcal{I} = \widetilde{I}$  (voir la preuve du Lemme
 \ref{lemma-regular-meromorphic-ideal-denominators}). Notons que  $T = V(I) \cup V((a/b)I)$. Pour tout  $A$-module
 $M$ , considérons la carte  $1 : IM \to M$; c'est la carte qui donne
naissance à la carte  $1$  du lemme. Considérons d'autre part la carte
 $\sigma : IM \to M_b, x \mapsto ax/b$. Puisque  $b$  n'est pas un diviseur de zéro dans
 $A$, et puisque  $M$  a pour support  $\Spec(A)$  et aucun idéal
premier associé, nous voyons que  $b$  est également un non-diviseur de
zéro sur  $M$ . Par conséquent  $M \subset M_b$. Par définition de
 $I$ , nous avons  $\sigma(IM) \subset M$  comme sous-modules de  $M_b$. Ainsi,
nous obtenons une application de  $A$-modules  $s : IM \to M$  (à savoir
l'application unique telle que  $s(z)/1 = \sigma(z)$  dans  $M_b$  pour tout
 $z \in IM$). Elle est injective parce que  $a$  est aussi un non-diviseur
de zéro (à la fois sur  $A$  et  $M$). Il est clair que
 $M/IM$  est annihilé par  $I$  et que  $M/s(IM)$  est annihilé par
 $(a/b)I$. Ainsi, le lemme s'ensuit.
\end{proof}





\section{Fonctions et sections méromorphes ; cas réduit}
\label{section-meromorphic-reduced}

\noindent
Pour un schéma qui est réduit et qui a localement un nombre fini de composantes
irréductibles, le faisceau des fonctions méromorphes est quasi-cohérent.

\begin{lemma}
\label{lemma-reduced-finite-irreducible}
Soit  $X$  un schéma réduit tel que tout ouvert quasi-compact possède un
nombre fini de composantes irréductibles. Soit  $X^0$  l'ensemble des points
génériques des composantes irréductibles de  $X$. Alors nous avons
$$
\mathcal{K}_X =
\bigoplus\nolimits_{\eta \in X^0} j_{\eta, *}\kappa(\eta) =
\prod\nolimits_{\eta \in X^0} j_{\eta, *}\kappa(\eta)
$$
où  $j_\eta : \Spec(\kappa(\eta)) \to X$  est l'application canonique des Schémas, Section  \ref{schemes-section-points}. De
plus,
\begin{enumerate}
\item $\mathcal{K}_X$  est un faisceau quasi-cohérent d'algèbres  $\mathcal{O}_X$,
\item pour tout  $\mathcal{O}_X$-module quasi-cohérent  $\mathcal{F}$  le faisceau
$$
\mathcal{K}_X(\mathcal{F}) =
\bigoplus\nolimits_{\eta \in X^0} j_{\eta, *}\mathcal{F}_\eta =
\prod\nolimits_{\eta \in X^0} j_{\eta, *}\mathcal{F}_\eta
$$
de sections méromorphes de  $\mathcal{F}$  est quasi-cohérent,
\item $\mathcal{S}_x \subset \mathcal{O}_{X, x}$  est l'ensemble des non-diviseurs de zéro pour tout  $x \in X$,
\item $\mathcal{K}_{X, x}$  est l'anneau total des fractions de  $\mathcal{O}_{X, x}$  pour tout
 $x \in X$,
\item $\mathcal{K}_X(U)$  est égal à l'anneau total des fractions de  $\mathcal{O}_X(U)$  pour tout
ouvert affine  $U \subset X$,
\item l'anneau des fonctions rationnelles de  $X$  est l'anneau des fonctions
méromorphes sur  $X$, sous une formule:  $R(X) = \Gamma(X, \mathcal{K}_X)$.
\end{enumerate}
\end{lemma}

\begin{proof}
Ce lemme est un cas particulier du Lemme  \ref{lemma-meromorphic-weakass-finite}  car sur un schéma réduit,
les points faiblement associés sont les points génériques selon le Lemme
 \ref{lemma-weakass-reduced}.
\end{proof}

\begin{lemma}
\label{lemma-reduced-normalization}
Soit  $X$  un schéma. Supposons que  $X$  est réduit et que tout
ouvert quasi-compact  $U \subset X$  ait un nombre fini de composantes irréductibles.
Alors le morphisme de normalisation  $\nu : X^\nu \to X$  est le morphisme
$$
\underline{\Spec}_X(\mathcal{O}') \longrightarrow X
$$
où  $\mathcal{O}' \subset \mathcal{K}_X$  est la fermeture intégrale de  $\mathcal{O}_X$  dans le faisceau des
fonctions méromorphes.
\end{lemma}

\begin{proof}
Comparez la définition du morphisme de normalisation  $\nu : X^\nu \to X$  (voir
Morphismes, Définition  \ref{morphisms-definition-normalization}) avec la description de  $\mathcal{K}_X$  dans le
Lemme  \ref{lemma-reduced-finite-irreducible}  ci-dessus.
\end{proof}

\begin{lemma}
\label{lemma-meromorphic-functions-integral-scheme}
Soit  $X$  un schéma intègre avec point générique  $\eta$. Nous avons
\begin{enumerate}
\item le faisceau des fonctions méromorphes est isomorphe au faisceau constant de valeur
le corps des fonctions (voir Morphismes, Définition  \ref{morphisms-definition-function-field}) de  $X$.
\item Pour tout faisceau quasi-cohérent  $\mathcal{F}$  sur  $X$ , le faisceau
 $\mathcal{K}_X(\mathcal{F})$  est isomorphe au faisceau constant de valeur  $\mathcal{F}_\eta$.
\end{enumerate}
\end{lemma}

\begin{proof}
Omis.
\end{proof}

\noindent
Dans certains cas, nous pouvons montrer que des sections méromorphes régulières
existent.

\begin{lemma}
\label{lemma-regular-meromorphic-section-exists}
Soit  $X$  un schéma. Soit  $\mathcal{L}$  un  $\mathcal{O}_X$-module inversible.
Dans chacun des cas suivants,  $\mathcal{L}$  possède une section méromorphe
régulière:
\begin{enumerate}
\item $X$  est intègre,
\item $X$  est réduit et tout ouvert quasi-compact a un nombre fini de
composantes irréductibles,
\item $X$  est localement noethérien et n’a pas de points embarqués.
\end{enumerate}
\end{lemma}

\begin{proof}
Dans le cas (1), soit  $\eta \in X$  le point générique. Nous avons vu dans le Lemme
 \ref{lemma-meromorphic-functions-integral-scheme}  que  $\mathcal{K}_X$, respectivement \ $\mathcal{K}_X(\mathcal{L})$ , est le
faisceau constant de valeur  $\kappa(\eta)$, respectivement \ $\mathcal{L}_\eta$.
Puisque  $\dim_{\kappa(\eta)} \mathcal{L}_\eta = 1$ , nous pouvons choisir un élément non nul  $s \in \mathcal{L}_\eta$. Il est
clair que  $s$  est une section méromorphe régulière de  $\mathcal{L}$. Dans
le cas (2), choisissez  $s_\eta \in \mathcal{L}_\eta$  non nul pour tous les points génériques
 $\eta$  de  $X$; cela est possible car  $\mathcal{L}_\eta$  est un espace
vectoriel de dimension  $1$ sur  $\kappa(\eta)$. Il découle immédiatement de
la description de  $\mathcal{K}_X$  et  $\mathcal{K}_X(\mathcal{L})$  dans le Lemme  \ref{lemma-reduced-finite-irreducible}  que
 $s = \prod s_\eta$  est une section méromorphe régulière de  $\mathcal{L}$. Le cas (3) est
le Lemme  \ref{lemma-regular-meromorphic-section-exists-noetherian}.
\end{proof}






\section{Diviseurs de Weil}
\label{section-Weil-divisors}

\noindent
Nous allons introduire les diviseurs de Weil et l'équivalence rationnelle des
diviseurs de Weil pour les schémas intègres localement noethériens. Comme nous ne
supposons pas que nos schémas sont quasi-compacts, nous devons être un peu
prudents lors de la définition des diviseurs de Weil. Nous devons permettre des
sommes infinies de diviseurs premiers car une fonction rationnelle peut avoir, par
exemple, un nombre infini de pôles. Pour les schémas quasi-compacts, nos diviseurs
de Weil sont des sommes finies comme d'habitude. Voici un lemme de base que nous
utiliserons souvent pour prouver que des collections de sous-schémas fermés sont
localement finies. Soit

\begin{lemma}
\label{lemma-components-locally-finite}
. Soit  $X$  un schéma localement noethérien. Soit  $Z \subset X$  un
sous-schéma fermé. La collection des composantes irréductibles de  $Z$  est
localement finie dans  $X$.
\end{lemma}

\begin{proof}
Soit  $U \subset X$  un sous-schéma ouvert quasi-compact. Alors  $U$  est un
schéma noethérien, et a donc un espace topologique sous-jacent noethérien
(Propriétés, Lemme  \ref{properties-lemma-Noetherian-topology}). Par conséquent, chaque sous-espace est
noethérien et possède un nombre fini de composantes irréductibles (voir Topologie,
Lemme  \ref{topology-lemma-Noetherian}).
\end{proof}

\noindent
Rappelons que si  $Z$  est un sous-ensemble fermé irréductible d'un schéma
 $X$, alors la codimension de  $Z$  dans  $X$  est égale à
la dimension de l'anneau local  $\mathcal{O}_{X, \xi}$, où  $\xi \in Z$  est le point
générique. Voir Propriétés, Lemme  \ref{properties-lemma-codimension-local-ring}.

\begin{definition}
\label{definition-Weil-divisor}
Soit  $X$  un schéma intègre localement noethérien.
\begin{enumerate}
\item Un  {\it  diviseur premier}  est un sous-schéma fermé
intègre  $Z \subset X$  de codimension  $1$.
\item une {\it Diviseur de Weil}  est une somme formelle
 $D = \sum n_Z Z$ où la somme porte sur les diviseurs premiers de $X$  et la
collection  $\{Z \mid n_Z \not = 0\}$ est localement fini (Topologie, Définition \ref{topology-definition-locally-finite}).
\end{enumerate}
Le groupe de tous les diviseurs de Weil sur  $X$  est désigné  $\text{Div}(X)$.
\end{definition}

\noindent
Notre prochaine tâche est de définir le diviseur de Weil associé à une fonction
rationnelle. Pour ce faire, nous utilisons l'ordre d'annulation d'une fonction
rationnelle le long d'un diviseur premier, qui est défini comme suit.

\begin{definition}
\label{definition-order-vanishing}
Laisse $X$ soit un schéma intègre localement noethérien. Soit $f \in R(X)^*$.
Pour chaque diviseur premier $Z \subset X$ nous définissons
le {\it ordre de disparition de $f$ le
long $Z$} comme l'entier
$$
\text{ord}_Z(f) = \text{ord}_{\mathcal{O}_{X, \xi}}(f)
$$
où le côté droit est la notion d'Algèbre, Définition  \ref{algebra-definition-ord}  et  $\xi$ 
est le point générique de  $Z$.
\end{definition}

\noindent
Notez que pour  $f, g \in R(X)^*$  nous avons
$$
\text{ord}_Z(fg) = \text{ord}_Z(f) + \text{ord}_Z(g).
$$
Bien sûr, il peut arriver que  $\text{ord}_Z(f) < 0$. Dans ce cas, nous disons que
 $f$  a un  {\it  pôle}  le long de
 $Z$  et que  $-\text{ord}_Z(f) > 0$  est le  {\it  ordre du pôle de
 $f$  le long de  $Z$}. Il est important de noter que la
condition  $\text{ord}_Z(f) \geq 0$  est  {\bf  pas}  équivalente à
la condition  $f \in \mathcal{O}_{X, \xi}$  sauf si l’anneau local  $\mathcal{O}_{X, \xi}$  est un anneau de
valuation discrète.

\begin{lemma}
\label{lemma-divisor-locally-finite}
Soit  $X$  un schéma intégral localement noethérien. Soit  $f \in R(X)^*$.
Alors les collections
$$
\{Z \subset X \mid Z\text{ a prime divisor with generic point }\xi
\text{ and }f\text{ not in }\mathcal{O}_{X, \xi}\}
$$
et
$$
\{Z \subset X \mid Z \text{ a prime divisor and }\text{ord}_Z(f) \not = 0\}
$$
sont localement finis dans  $X$.
\end{lemma}

\begin{proof}
Il existe un sous-schéma ouvert non vide  $U \subset X$  tel que  $f$ 
corresponde à une section de  $\Gamma(U, \mathcal{O}_X^*)$. Ainsi, les diviseurs premiers qui
peuvent apparaître dans les ensembles du lemme sont tous des composantes
irréductibles de  $X \setminus U$. Ainsi, le Lemme  \ref{lemma-components-locally-finite}  donne le résultat
souhaité.
\end{proof}

\noindent
Ce lemme nous permet de faire la définition suivante.

\begin{definition}
\label{definition-principal-divisor}
Soit  $X$  un schéma intègre localement noethérien. Soit  $f \in R(X)^*$. Le
diviseur de Weil principal  {\it  associé à
 $f$}  est le diviseur de Weil
$$
\text{div}(f) = \text{div}_X(f) = \sum \text{ord}_Z(f) [Z]
$$
où la somme est sur les diviseurs premiers et  $\text{ord}_Z(f)$  est comme dans la
Définition  \ref{definition-order-vanishing}. Cela a un sens selon le Lemme  \ref{lemma-divisor-locally-finite}.
\end{definition}

\begin{lemma}
\label{lemma-div-additive}
Soit  $X$  un schéma intègre localement noethérien. Soit  $f, g \in R(X)^*$.
Alors
$$
\text{div}_X(fg) = \text{div}_X(f) + \text{div}_X(g)
$$
comme diviseurs de Weil sur  $X$.
\end{lemma}

\begin{proof}
Cela est clair à partir de l'additivité des fonctions  $\text{ord}$ .
\end{proof}

\noindent
Nous voyons d'après le lemme ci-dessus que l'ensemble des diviseurs de Weil
principaux forme un sous-groupe du groupe de tous les diviseurs de Weil. Cela
conduit à la définition suivante.

\begin{definition}
\label{definition-class-group}
Soit  $X$  un schéma intégral localement noethérien. Le
 {\it  groupe de classes de diviseurs de Weil }  de
 $X$  est le quotient du groupe des diviseurs de Weil par le sous-groupe
des diviseurs de Weil principaux. Notation:  $\text{Cl}(X)$.
\end{definition}

\noindent
Par construction, nous obtenons un complexe exact
\begin{equation}
\label{equation-Weil-divisor-class}
R(X)^* \xrightarrow{\text{div}} \text{Div}(X) \to \text{Cl}(X) \to 0
\end{equation}
que nous pouvons considérer comme une présentation de  $\text{Cl}(X)$. Notre
prochaine tâche est de relier le groupe de classes de diviseurs de Weil au groupe
de Picard.





\section{La classe de diviseur de Weil associée à un module inversible}
\label{section-c1}

\noindent
Dans cette section, nous suivons exactement la même progression que dans la
Section  \ref{section-Weil-divisors}  pour définir une application canonique  $\Pic(X) \to \text{Cl}(X)$  sur un
schéma intégral localement noethérien.

\medskip\noindent
Soit  $X$  un schéma. Soit  $\mathcal{L}$  un  $\mathcal{O}_X$-module inversible.
Soit  $\xi \in X$  un point. Si  $s_\xi, s'_\xi \in \mathcal{L}_\xi$  génère  $\mathcal{L}_\xi$  comme
 $\mathcal{O}_{X, \xi}$-module, alors il existe une unité  $u \in \mathcal{O}_{X, \xi}^*$  telle que
 $s_\xi = u s'_\xi$. La germe du faisceau des sections méromorphes  $\mathcal{K}_X(\mathcal{L})$  de
 $\mathcal{L}$  en  $x$  est égale à  $\mathcal{K}_{X, x} \otimes_{\mathcal{O}_{X, x}} \mathcal{L}_x$. Ainsi, l'image de toute
section méromorphe  $s$  de  $\mathcal{L}$  dans le germe en  $x$  peut
s'écrire comme  $s = fs_\xi$  avec  $f \in \mathcal{K}_{X, x}$. Ci-dessous, nous abrégerons cela en
disant  $f = s/s_\xi$. De plus, si  $X$  est intégral, nous avons que
 $\mathcal{K}_{X, x} = R(X)$  est égal au corps des fonctions de  $X$, donc  $s/s_\xi \in R(X)$.
Si  $s$  est une section méromorphe régulière, alors en fait  $s/s_\xi \in R(X)^*$.
Sur un schéma intégral, une section méromorphe régulière est la même chose qu'une
section méromorphe non nulle. Enfin, nous voyons que  $s/s_\xi$  est indépendant
du choix de  $s_\xi$  à un facteur près par multiplication par une unité de
l'anneau local  $\mathcal{O}_{X, x}$. En regroupant le tout, nous voyons que la définition
suivante a du sens.

\begin{definition}
\label{definition-order-vanishing-meromorphic}
Soit  $X$  un schéma intégral localement noethérien. Soit  $\mathcal{L}$  un
 $\mathcal{O}_X$-module inversible. Soit  $s \in \Gamma(X, \mathcal{K}_X(\mathcal{L}))$  une section méromorphe régulière
de  $\mathcal{L}$. Pour chaque diviseur premier  $Z \subset X$ , nous définissons
l'ordre de vanishing  {\it  de  $s$  le long de
 $Z$}  comme l'entier
$$
\text{ord}_{Z, \mathcal{L}}(s)
= \text{ord}_{\mathcal{O}_{X, \xi}}(s/s_\xi)
$$
où le côté droit est la notion d'Algèbre, Définition  \ref{algebra-definition-ord},  $\xi \in Z$ 
est le point générique, et  $s_\xi \in \mathcal{L}_\xi$  est un générateur.
\end{definition}

\noindent
Comme dans le cas des diviseurs principaux, nous avons le lemme suivant.

\begin{lemma}
\label{lemma-divisor-meromorphic-locally-finite}
Soit  $X$  un schéma intègre localement noethérien. Soit  $\mathcal{L}$  un
 $\mathcal{O}_X$-module inversible. Soit  $s \in \mathcal{K}_X(\mathcal{L})$  une section méromorphe régulière
(c’est-à-dire non nulle) de  $\mathcal{L}$. Alors les ensembles
$$
\{Z \subset X \mid Z \text{ a prime divisor with generic point }\xi
\text{ and }s\text{ not in }\mathcal{L}_\xi\}
$$
et
$$
\{Z \subset X \mid Z \text{ is a prime divisor and }
\text{ord}_{Z, \mathcal{L}}(s) \not = 0\}
$$
sont localement finis dans  $X$.
\end{lemma}

\begin{proof}
Il existe un sous-schéma ouvert non vide  $U \subset X$  tel que  $s$ 
corresponde à une section de  $\Gamma(U, \mathcal{L})$  qui génère  $\mathcal{L}$  sur
 $U$. Ainsi, les diviseurs premiers qui peuvent apparaître dans les
ensembles du lemme sont tous des composantes irréductibles de  $X \setminus U$. Ainsi,
le Lemme  \ref{lemma-components-locally-finite} donne le résultat souhaité.
\end{proof}

\begin{lemma}
\label{lemma-divisor-meromorphic-well-defined}
Soit  $X$  un schéma intègre localement noethérien. Soit  $\mathcal{L}$  un
 $\mathcal{O}_X$-module inversible. Soit  $s, s' \in \mathcal{K}_X(\mathcal{L})$  une section méromorphe non nulle
de  $\mathcal{L}$. Alors  $f = s/s'$  est un élément de  $R(X)^*$  et nous avons
$$
\sum \text{ord}_{Z, \mathcal{L}}(s)[Z]
=
\sum \text{ord}_{Z, \mathcal{L}}(s')[Z]
+
\text{div}(f)
$$
en tant que diviseurs de Weil.
\end{lemma}

\begin{proof}
Cela est clair d'après les définitions. Il faut noter que le Lemme  \ref{lemma-divisor-meromorphic-locally-finite} 
garantit que les sommes sont bien des diviseurs de Weil.
\end{proof}

\begin{definition}
\label{definition-divisor-invertible-sheaf}
Soit  $X$  un schéma intègre localement noethérien. Soit  $\mathcal{L}$  un
 $\mathcal{O}_X$-module inversible.
\begin{enumerate}
\item Pour toute section méromorphe non nulle  $s$  de  $\mathcal{L}$ , nous
définissons le diviseur de Weil  {\it  associé à
 $s$}  comme
$$
\text{div}_\mathcal{L}(s) =
\sum \text{ord}_{Z, \mathcal{L}}(s) [Z] \in \text{Div}(X)
$$
où la somme est effectuée sur les diviseurs premiers.
\item Nous définissons  {\it Classe de diviseur de Weil associée
à $\mathcal{L}$} comme l'image
de $\text{div}_\mathcal{L}(s)$ dans $\text{Cl}(X)$ où $s$ est toute section méromorphe non nulle
de $\mathcal{L}$ au-dessus $X$. Cela est bien défini par le
Lemme \ref{lemma-divisor-meromorphic-well-defined}.
\end{enumerate}
\end{definition}

\noindent
Comme prévu, cette construction est additive dans le module inversible.

\begin{lemma}
\label{lemma-c1-additive}
Soit  $X$  un schéma intègre localement noethérien. Soient  $\mathcal{L}$,
 $\mathcal{N}$  des  $\mathcal{O}_X$-modules inversibles. Soient  $s$, resp.
\ $t$  une section méromorphe non nulle de  $\mathcal{L}$, resp.
\ $\mathcal{N}$. Alors  $st$  est une section méromorphe non nulle
de  $\mathcal{L} \otimes \mathcal{N}$, et
$$
\text{div}_{\mathcal{L} \otimes \mathcal{N}}(st)
=
\text{div}_\mathcal{L}(s) + \text{div}_\mathcal{N}(t)
$$
dans  $\text{Div}(X)$. En particulier, la classe de diviseur de Weil de  $\mathcal{L} \otimes_{\mathcal{O}_X} \mathcal{N}$ 
est la somme des classes de diviseurs de Weil de  $\mathcal{L}$  et  $\mathcal{N}$.
\end{lemma}

\begin{proof}
Soit  $s$, resp. \ $t$ , une section méromorphe non
nulle de  $\mathcal{L}$, resp. \ $\mathcal{N}$. Alors  $st$  est une
section méromorphe non nulle de  $\mathcal{L} \otimes \mathcal{N}$. Soit  $Z \subset X$  un diviseur
premier. Soit  $\xi \in Z$  son point générique. Choisissez les générateurs
 $s_\xi \in \mathcal{L}_\xi$ et  $t_\xi \in \mathcal{N}_\xi$. Alors  $s_\xi t_\xi$  est un générateur de
 $(\mathcal{L} \otimes \mathcal{N})_\xi$. Donc  $st/(s_\xi t_\xi) = (s/s_\xi)(t/t_\xi)$. Ainsi, nous voyons que
$$
\text{ord}_{Z, \mathcal{L} \otimes \mathcal{N}}(st)
=
\text{ord}_{Z, \mathcal{L}}(s) + \text{ord}_{Z, \mathcal{N}}(t)
$$
par l'additivité de la fonction  $\text{ord}_Z$ .
\end{proof}

\noindent
De cette manière, nous obtenons un homomorphisme de groupes abéliens
\begin{equation}
\label{equation-c1}
\Pic(X) \longrightarrow \text{Cl}(X)
\end{equation}
qui assigne à un module inversible sa classe de diviseur de Weil.

\begin{lemma}
\label{lemma-normal-c1-injective}
Soit  $X$  un schéma intègre localement noethérien. Si  $X$  est
normal, alors l'application (\ref{equation-c1})  $\Pic(X) \to \text{Cl}(X)$  est injective.
\end{lemma}

\begin{proof}
Soit  $\mathcal{L}$  un  $\mathcal{O}_X$-module inversible dont la classe de diviseur de
Weil associée est triviale. Soit  $s$  une section méromorphe régulière de
 $\mathcal{L}$. L'hypothèse signifie que  $\text{div}_\mathcal{L}(s) = \text{div}(f)$  pour un certain  $f \in R(X)^*$.
Alors nous voyons que  $t = f^{-1}s$  est une section méromorphe régulière de
 $\mathcal{L}$  avec  $\text{div}_\mathcal{L}(t) = 0$, voir le Lemme  \ref{lemma-divisor-meromorphic-well-defined}. Nous montrerons que
 $t$  définit une trivialisation de  $\mathcal{L}$ , ce qui termine la preuve
du lemme. Pour prouver cela, nous pouvons travailler localement sur  $X$.
Par conséquent, nous pouvons supposer que  $X = \Spec(A)$  est affine et que
 $\mathcal{L}$  est triviale. Alors  $A$  est un domaine normal noethérien et
 $t$  est un élément de son corps de fractions tel que  $\text{ord}_{A_\mathfrak p}(t) = 0$  pour
tous les  $1$  premiers de hauteur  $\mathfrak p$  de  $A$. Notre
objectif est de montrer que  $t$  est une unité de  $A$. Puisque
 $A_\mathfrak p$  est un anneau de valuation discrète pour les premiers de hauteur un
de  $A$  (Algèbre, Lemme  \ref{algebra-lemma-criterion-normal}), la condition signifie que
 $t \in A_\mathfrak p^*$  pour tous les premiers  $\mathfrak p$  de hauteur  $1$. Cela
implique  $t \in A$  et  $t^{-1} \in A$  selon Algèbre, Lemme  \ref{algebra-lemma-normal-domain-intersection-localizations-height-1}  et la
preuve est complète.
\end{proof}

\begin{lemma}
\label{lemma-local-rings-UFD-c1-bijective}
Soit  $X$  un schéma intégral localement noethérien. Considérons la carte
(\ref{equation-c1})  $\Pic(X) \to \text{Cl}(X)$. Les conditions suivantes sont équivalentes
\begin{enumerate}
\item les anneaux locaux de  $X$  sont des UFDs, et
\item $X$  est normal et  $\Pic(X) \to \text{Cl}(X)$  est surjective.
\end{enumerate}
Dans ce cas,  $\Pic(X) \to \text{Cl}(X)$  est un isomorphisme.
\end{lemma}

\begin{proof}
Si (1) est vrai, alors  $X$  est normal selon Algèbre, Lemme  \ref{algebra-lemma-UFD-normal}.
Par conséquent, l'application (\ref{equation-c1}) est injective d'après le Lemme
 \ref{lemma-normal-c1-injective}. De plus, tout diviseur premier  $D \subset X$  est un diviseur de
Cartier effectif selon le Lemme  \ref{lemma-weil-divisor-is-cartier-UFD}. Dans ce cas, la section canonique
 $1_D$  de  $\mathcal{O}_X(D)$  (Définition  \ref{definition-invertible-sheaf-effective-Cartier-divisor}) s'annule exactement le long
de  $D$  et nous voyons que la classe de  $D$  est l'image de
 $\mathcal{O}_X(D)$  sous l'application (\ref{equation-c1}). Ainsi, l'application est également
surjective.

\medskip\noindent
Supposons que (2) soit vrai. Choisissez un diviseur premier  $D \subset X$. Puisque
(\ref{equation-c1}) est surjectif, il existe un faisceau inversible  $\mathcal{L}$, une
section méromorphe régulière  $s$, et  $f \in R(X)^*$  tels que  $\text{div}_\mathcal{L}(s) + \text{div}(f) = [D]$.
En d'autres termes,  $\text{div}_\mathcal{L}(fs) = [D]$. Soit  $x \in X$  et soit  $A = \mathcal{O}_{X, x}$. Ainsi
 $A$  est un domaine local normal noethérien avec le corps des fractions
 $K = R(X)$. Chaque idéal premier de hauteur  $1$  de  $A$ 
correspond à un diviseur premier sur  $X$  et chaque  $\mathcal{O}_X$-module
inversible se restreint au module inversible trivial sur  $\Spec(A)$. Il s'ensuit
que pour chaque idéal premier de hauteur  $1$   $\mathfrak p \subset A$ , il existe un
élément  $f \in K$  tel que  $\text{ord}_{A_\mathfrak p}(f) = 1$  et  $\text{ord}_{A_{\mathfrak p'}}(f) = 0$  pour chaque autre idéal
premier de hauteur un  $\mathfrak p'$. Ensuite,  $f \in A$  par Algèbre, Lemme
 \ref{algebra-lemma-normal-domain-intersection-localizations-height-1}. En procédant de la même manière, nous voyons que chaque élément
 $g \in \mathfrak p$  est de la forme  $g = af$  pour un certain  $a \in A$. Ainsi,
nous voyons que chaque idéal premier de hauteur un de  $A$  est principal
et  $A$  est un UFD par Algèbre, Lemme  \ref{algebra-lemma-characterize-UFD-height-1}.
\end{proof}







\section{Compléments sur les modules inversibles}
\label{section-invertible-modules}

\noindent
Dans cette section, nous discutons de certaines propriétés des modules
inversibles.

\begin{lemma}
\label{lemma-in-image-pullback}
Soit  $\varphi : X \to Y$  un morphisme de schémas. Soit  $\mathcal{L}$  un
 $\mathcal{O}_X$-module inversible. Supposons que
\begin{enumerate}
\item $X$  soit localement noethérien,
\item $Y$  soit localement noethérien, intègre et normal,
\item $\varphi$  soit plat avec des fibres intègres (donc non vides)
\item $\varphi$  est soit quasi-compact, soit localement de type fini,
\item $\mathcal{L}$  est trivial lorsqu'on le restreint à la fibre générique de
 $\varphi$.
\end{enumerate}
Alors  $\mathcal{L} \cong \varphi^*\mathcal{N}$  pour un certain  $\mathcal{O}_Y$-module inversible  $\mathcal{N}$.
\end{lemma}

\begin{proof}
Soit  $\xi \in Y$  le point générique. Soit  $X_\xi$  la fibre schématique de
 $\varphi$  au-dessus de  $\xi$. On note  $\mathcal{L}_\xi$  le tiré en arrière de
 $\mathcal{L}$  sur  $X_\xi$. L'hypothèse (5) signifie que  $\mathcal{L}_\xi$  est
trivial. Choisir une section trivialisante  $s \in \Gamma(X_\xi, \mathcal{L}_\xi)$. On observe que
 $X$  est intègre d'après le Lemme  \ref{lemma-flat-over-integral-integral-fibre}. Ainsi, nous pouvons
considérer  $s$  comme une section méromorphe régulière de  $\mathcal{L}$.
Les tirés en arrière des fonctions méromorphes sont définis pour  $\varphi$ 
selon le Lemme  \ref{lemma-pullback-meromorphic-sections-defined}. Soit  $\mathcal{N} \subset \mathcal{K}_Y$  le  $\mathcal{O}_Y$-module dont les
sections sur un ouvert  $V \subset Y$  sont ces fonctions méromorphes  $g \in \mathcal{K}_Y(V)$ 
telles que  $\varphi^*(g)s \in \mathcal{L}(\varphi^{-1}V)$. A priori  $\varphi^*(g)s$  est une section de  $\mathcal{K}_X(\mathcal{L})$  sur
 $\varphi^{-1}V$. Nous affirmons que  $\mathcal{N}$  est un  $\mathcal{O}_Y$-module
inversible et que l'application
$$
\varphi^*\mathcal{N} \longrightarrow \mathcal{L},\quad
g \longmapsto gs
$$
est un isomorphisme.

\medskip\noindent
Nous prouvons d'abord l'affirmation dans la situation suivante:  $X$  et
 $Y$  sont affines et  $\mathcal{L}$  est trivial. Disons que  $Y = \Spec(R)$,
 $X = \Spec(A)$  et  $s$  sont donnés par l'élément  $s \in A \otimes_R K$  où
 $K$  est le corps des fractions de  $R$. Nous pouvons écrire
 $s = a/r$  pour un certain  $r \in R$  non nul et  $a \in A$. Puisque
 $s$  engendre  $\mathcal{L}$  sur la fibre générique, nous voyons qu'il
existe un  $s' \in A \otimes_R K$  tel que  $ss' = 1$. Ainsi, nous voyons que  $s = r'/a'$ 
pour un certain  $r' \in R$  et  $a' \in A$ non nuls. Soit  $\mathfrak p_1, \ldots, \mathfrak p_n \subset R$  les
idéaux premiers minimaux au-dessus de  $rr'$. Chaque  $R_{\mathfrak p_i}$  est un
anneau de valuation discrète (Algèbre, Lemmata  \ref{algebra-lemma-minimal-over-1}  et  \ref{algebra-lemma-criterion-normal}). Par
hypothèse,  $\mathfrak q_i = \mathfrak p_i A$  est un nombre premier. Donc  $\mathfrak q_i A_{\mathfrak q_i}$  est engendré par
un seul élément et nous trouvons que  $A_{\mathfrak q_i}$  est aussi un anneau de valuation
discrète (Algèbre, Lemme  \ref{algebra-lemma-characterize-dvr}). Bien sûr,  $R_{\mathfrak p_i} \to A_{\mathfrak q_i}$  a un indice de
ramification  $1$. Soit  $e_i, e'_i \geq 0$  la valuation de  $a, a'$  dans
 $A_{\mathfrak q_i}$. Alors  $e_i + e'_i$  est la valuation de  $rr'$  dans
 $R_{\mathfrak p_i}$. Notez que
$$
\mathfrak p_1^{(e_1 + e'_1)} \cap \ldots \cap \mathfrak p_i^{(e_n + e'_n)} =
(rr')
$$
dans  $R$  selon Algèbre, Lemme  \ref{algebra-lemma-normal-domain-intersection-localizations-height-1}. Posons
$$
I = \mathfrak p_1^{(e_1)} \cap \ldots \cap \mathfrak p_i^{(e_n)}
\quad\text{and}\quad
I' = \mathfrak p_1^{(e'_1)} \cap \ldots \cap \mathfrak p_i^{(e'_n)}
$$
de sorte que  $II' \subset (rr')$. Observez que
$$
IA =
(\mathfrak p_1^{(e_1)} \cap \ldots \cap \mathfrak p_i^{(e_n)})A =
(\mathfrak p_1A)^{(e_1)} \cap \ldots \cap (\mathfrak p_i A)^{(e_n)}
$$
par l'Algèbre, lemmes  \ref{algebra-lemma-symbolic-power-flat-extension}  et  \ref{algebra-lemma-flat-intersect-ideals}. De même pour  $I'A$.
D'où  $a \in IA$  et  $a' \in I'A$. Nous concluons que  $IA \otimes_A I'A \to rr'A$  est surjectif.
Par la platitude fidèle de  $R \to A$ , nous trouvons que  $I \otimes_R I' \to (rr')$  est
également surjectif. Il s'ensuit que  $II' = (rr')$ ,  $I$  et  $I'$ 
sont localement libres de rang fini  $1$, voir Algèbre, lemme
 \ref{algebra-lemma-product-ideals-principal}. Ainsi, localement pour la topologie de Zariski sur  $R$ ,
nous pouvons écrire  $I = (g)$  et  $I' = (g')$  avec  $gg' = rr'$. Ensuite,
 $a = ug$  et  $a' = u'g'$  pour certains  $u, u' \in A$. Nous concluons que
 $u, u'$  sont des unités. Ainsi, localement pour la topologie de Zariski sur
 $R$ , nous avons  $s = ug/r$  et l'affirmation s'ensuit dans ce cas.

\medskip\noindent
Soit  $y \in Y$  un point. Choisissons  $x \in X$  se mappant sur  $y$.
Nous pouvons appliquer le résultat du paragraphe précédent à  $\Spec(\mathcal{O}_{X, x}) \to \Spec(\mathcal{O}_{Y, y})$. Nous en
concluons qu'il existe un élément  $g \in R(Y)^*$  bien défini à un facteur près d'un
élément de  $\mathcal{O}_{Y, y}^*$  tel que  $\varphi^*(g)s$  génère  $\mathcal{L}_x$. Ainsi,
 $\varphi^*(g)s$  génère  $\mathcal{L}$  dans un voisinage  $U$  de  $x$.
Supposons que  $x'$  soit un second point au-dessus de  $y$  et que
 $g' \in R(Y)^*$  soit tel que  $\varphi^*(g')s$  génère  $\mathcal{L}$  dans un voisinage
ouvert  $U'$  de  $x'$. Alors nous pouvons choisir un point
 $x''$  dans  $U \cap U' \cap \varphi^{-1}(\{y\})$  car la fibre est irréductible. Par l'unicité pour
l'application d'anneaux  $\mathcal{O}_{Y, y} \to \mathcal{O}_{X, x''}$ , nous trouvons que  $g$  et
 $g'$  diffèrent (multiplicativement) par un élément de  $\mathcal{O}_{Y, y}^*$. Ainsi,
nous voyons que  $\varphi^*(g)s$  est un générateur pour  $\mathcal{L}$  sur un voisinage
ouvert de  $\varphi^{-1}(y)$. Soit  $Z \subset X$  l'ensemble des points  $z \in X$  tels
que  $\varphi^*(g)s$  ne génère pas  $\mathcal{L}_z$. Les arguments ci-dessus montrent que
 $Z$  est fermé et que  $Z = \varphi^{-1}(T)$  pour un certain sous-ensemble
 $T \subset Y$  avec  $y \not \in T$. Si nous pouvons montrer que  $T$  est
fermé, alors  $g$  sera un générateur pour  $\mathcal{N}$  en tant que
 $\mathcal{O}_Y$-module dans le voisinage ouvert  $Y \setminus T$  de  $y$ ,
terminant ainsi la preuve (certains détails omis).

\medskip\noindent
Si  $\varphi$  est quasi-compact, alors  $T$  est fermé par les
morphismes, Lemme  \ref{morphisms-lemma-fpqc-quotient-topology}. Si  $\varphi$  est localement de type fini, alors
 $\varphi$  est ouvert par les morphismes, Lemme  \ref{morphisms-lemma-fppf-open}. Alors  $Y \setminus T$ 
est ouvert en tant qu'image de l'ouvert  $X \setminus Z$.
\end{proof}

\begin{lemma}
\label{lemma-closure-effective-cartier-divisor}
Soit  $X$  un schéma localement noethérien. Soit  $U \subset X$  un ouvert et
soit  $D \subset U$  un diviseur de Cartier effectif. Si  $\mathcal{O}_{X, x}$  est un UFD pour
tout  $x \in X \setminus U$, alors il existe un diviseur de Cartier effectif  $D' \subset X$ 
avec  $D = U \cap D'$.
\end{lemma}

\begin{proof}
Soit  $D' \subset X$  l'image schématique du morphisme  $D \to X$. Puisque
 $X$  est localement noethérien, le morphisme  $D \to X$  est
quasi-compact, voir Propriétés, Lemme  \ref{properties-lemma-immersion-into-noetherian}. Par conséquent, la formation
de  $D'$  commute avec le passage aux ouverts dans  $X$  selon
Morphismes, Lemme  \ref{morphisms-lemma-quasi-compact-scheme-theoretic-image}. Ainsi, nous pouvons supposer que  $X = \Spec(A)$  est
affine. Soit  $I \subset A$  l'idéal correspondant à  $D'$. Soit  $\mathfrak p \subset A$ 
un idéal premier correspondant à un point de  $X \setminus U$. Pour terminer la
démonstration, il suffit de montrer que  $I_\mathfrak p$  est engendré par un élément,
voir Lemme  \ref{lemma-effective-Cartier-in-points}. Ainsi, nous pouvons remplacer  $X$  par
 $\Spec(A_\mathfrak p)$, voir Morphismes, Lemme  \ref{morphisms-lemma-flat-base-change-scheme-theoretic-image}. En d'autres termes, nous
pouvons supposer que  $X$  est le spectre d'un UFD local  $A$.
Alors tous les anneaux locaux de  $A$  sont des UFD. Il s'ensuit que
 $D = \sum a_i D_i$  avec  $D_i \subset U$  un diviseur de Cartier effectif et intégral, voir
Lemme  \ref{lemma-effective-Cartier-divisor-is-a-sum}. Les points génériques  $\xi_i$  de  $D_i$ 
correspondent aux idéaux premiers  $\mathfrak p_i \subset A$  de hauteur  $1$, voir Lemme
 \ref{lemma-effective-Cartier-codimension-1}. Alors  $\mathfrak p_i = (f_i)$  pour un certain élément premier  $f_i \in A$  et
nous concluons que  $D'$  est défini par  $\prod f_i^{a_i}$  comme souhaité.
\end{proof}

\begin{lemma}
\label{lemma-extend-invertible-module}
Soit  $X$  un schéma localement noethérien. Soit  $U \subset X$  un ouvert et
soit  $\mathcal{L}$  un  $\mathcal{O}_U$-module inversible. Si  $\mathcal{O}_{X, x}$  est un UFD
pour tout  $x \in X \setminus U$, alors il existe un  $\mathcal{O}_X$-module inversible
 $\mathcal{L}'$  avec  $\mathcal{L} \cong \mathcal{L}'|_U$.
\end{lemma}

\begin{proof}
Choisissez  $x \in X$,  $x \not \in U$. Nous montrerons qu'il existe un voisinage
ouvert affine  $W \subset X$, tel que  $\mathcal{L}|_{W \cap U}$  s'étend à un faisceau inversible
sur  $W$. Cela implique, par collage de faisceaux (Faisceaux, Section
 \ref{sheaves-section-glueing-sheaves}), que nous pouvons étendre  $\mathcal{L}$  à l'ouvert strictement plus
grand  $U \cup W$. Soit  $W = \Spec(A)$  un voisinage ouvert affine. Puisque
 $U \cap W$  est quasi-affine, nous voyons que nous pouvons écrire  $\mathcal{L}|_{W \cap U}$ 
comme  $\mathcal{O}(D_1) \otimes \mathcal{O}(D_2)^{\otimes -1}$  pour certains diviseurs de Cartier effectifs  $D_1, D_2 \subset W \cap U$, voir
le Lemme  \ref{lemma-quasi-projective-Noetherian-pic-effective-Cartier}. Alors  $D_1$  et  $D_2$  s'étendent à des
diviseurs de Cartier effectifs de  $W$  par le Lemme  \ref{lemma-closure-effective-cartier-divisor} , ce qui
nous donne l'extension du faisceau inversible.

\medskip\noindent
Si  $X$  est noethérien (ce qui est le cas le plus utilisé en pratique), ce
qui précède, combiné à l'induction noethérienne, termine la preuve. Dans le cas
général, nous procédons comme suit. Tout d'abord, parce que chaque anneau local
d'un point à l'extérieur de  $U$  est un domaine et que  $X$  est
localement noethérien, nous voyons que l'adhérence de  $U$  dans
 $X$  est ouverte. Ainsi, nous pouvons supposer que  $U \subset X$  est dense
et schématiquement dense. Maintenant, nous considérons l'ensemble  $T$  des
triplets  $(U', \mathcal{L}', \alpha)$  où  $U \subset U' \subset X$  est un sous-schéma ouvert,  $\mathcal{L}'$  est
un  $\mathcal{O}_{U'}$-module inversible, et  $\alpha : \mathcal{L}'|_U \to \mathcal{L}$  est un isomorphisme. Nous
dotons  $T$  d'un ordre partiel  $\leq$  défini par la règle
 $(U', \mathcal{L}', \alpha) \leq (U'', \mathcal{L}'', \alpha')$  si et seulement si  $U' \subset U''$  et qu'il existe un isomorphisme
 $\beta : \mathcal{L}''|_{U'} \to \mathcal{L}'$  compatible avec  $\alpha$  et  $\alpha'$. Observez que
 $\beta$  est unique (si elle existe) car  $U \subset X$  est dense. La première
partie de la preuve montre que pour tout élément  $t = (U', \mathcal{L}', \alpha)$  de  $T$  avec
 $U' \not = X$ , il existe un  $t' \in T$  avec  $t' > t$. Ainsi, pour terminer la
preuve, il suffit de montrer que le lemme de Zorn s'applique. Considérons donc un
sous-ensemble totalement ordonné  $I \subset T$. Si  $i \in I$  correspond au
triple  $(U_i, \mathcal{L}_i, \alpha_i)$, alors nous pouvons construire un module inversible
 $\mathcal{L}'$  sur  $U' = \bigcup U_i$  de la manière suivante. Pour  $W \subset U'$  ouvert et
quasi-compact, nous voyons que  $W \subset U_i$  pour un certain  $i$  et nous
posons
$$
\mathcal{L}'(W) = \mathcal{L}_i(W)
$$
Pour les cartes de transition, nous utilisons le  $\beta$'s (qui sont uniques
et donc se composent correctement). Cela définit un inversible
 $\mathcal{O}$-module $\mathcal{L}'$ sur la base des ouverts quasi-compacts
de $U'$ ce qui est suffisant pour définir un module inversible (Faisceaux,
Section  \ref{sheaves-section-bases}). Nous omettons les détails.
\end{proof}

\begin{lemma}
\label{lemma-open-subscheme-UFD}
Soit  $R$  un UFD. Les groupes de Picard des suivants sont triviaux.
\begin{enumerate}
\item $\Spec(R)$  et tout sous-schéma ouvert de celui-ci.
\item $\mathbf{A}^n_R = \Spec(R[x_1, \ldots, x_n])$  et tout sous-schéma ouvert de celui-ci.
\end{enumerate}
En particulier, le groupe de Picard de tout sous-schéma ouvert de l'espace affine
 $n$  $\mathbf{A}^n_k$  sur un corps  $k$  est trivial.
\end{lemma}

\begin{proof}
Tout anneau de polynômes sur  $R$  est un UFD selon Algèbre, Lemme
 \ref{algebra-lemma-polynomial-ring-UFD}. Ainsi, la question est de montrer que  $\Pic(U)$  est trivial si
 $U \subset \Spec(R)$  est un ouvert (non vide).

\medskip\noindent
Dans le cas noethérien, nous pouvons utiliser le Lemme  \ref{lemma-extend-invertible-module}  pour réduire à
montrer que  $\Pic(\Spec(R))$  est trivial. Cela est à son tour prouvé dans Plus sur
Algèbre, Lemme  \ref{more-algebra-lemma-UFD-Pic-trivial}. Nous invitons le lecteur à sauter le reste de la
preuve.

\medskip\noindent
Soit  $U \subset \Spec(R)$  un ouvert non vide. Choisissez un  $f \in R$  non nul avec
 $D(f) \subset U$. Puisque  $f$  est un produit d'éléments premiers de
 $R$ , nous concluons qu'il existe un nombre fini d'éléments premiers
 $a_1, \ldots, a_n \in R$  tels que  $(a_i) \not \in U$  (ou équivalemment  $U \cap V(a_i) = \emptyset$). Après avoir
remplacé  $R$  par  $R_{a_1 \ldots a_n}$ , nous pouvons supposer que  $U$ 
contient tous les idéaux premiers principaux. (Notez qu'une localisation d'un UFD
est un UFD, par exemple selon Algèbre, Lemme  \ref{algebra-lemma-invert-prime-elements}.)

\medskip\noindent
Supposons que  $U$  contienne tous les idéaux premiers principaux. Écrivons
 $U = \Spec(R) \setminus V(I)$  pour un certain idéal  $I \subset R$. Prenons  $f \in I$  non nul et
écrivons  $f = a_1 \ldots a_n$  comme un produit d’éléments premiers. Comme  $(a_i) \not \in V(I)$ ,
nous pouvons choisir  $g \in I$  avec  $g \not \in (a_i)$  pour tous les  $i$ 
(évitement des premiers). Imaginez
$$
U' = D(f) \cup D(g) \subset U \subset \Spec(R)
$$
D'après Plus sur Algèbre, Lemme  \ref{more-algebra-lemma-UFD-Pic-trivial} , les groupes de Picard de
 $D(f)$  et  $D(g)$  sont triviaux et, par un argument élémentaire sur les
UFD, l'application  $R_f^* \times R_g^* \to R_{fg}^*$,  $(u, v) \mapsto uv^{-1}$  est surjective. Nous concluons que
 $\Pic(U')$  est trivial.

\medskip\noindent
Ainsi, il suffit de montrer qu'un module inversible  $\mathcal{L}$  sur  $U$ 
dont le pull-back par  $j : U' \to U'$  est trivial, est trivial. Pour ce faire, il
suffit de montrer que l'application  $\mathcal{L} \to j_*j^*\mathcal{L}$  est un isomorphisme. Pour cela à
son tour, soit  $u \in U \setminus U'$. Alors  $f, g \in \mathfrak m_u \subset \mathcal{O}_{U, u} = A$  sont des éléments d'un UFD local
qui n'ont pas de facteur premier en commun. Le lecteur montre facilement que cela
signifie que
$$
0 \to A \to A_f \times A_g \to A_{fg}
$$
est exact en regardant les facteurs premiers. Puisque  $A = \mathcal{L}_u$  parce que
 $\mathcal{L}$  est inversible et puisque  $(j_*j^*\mathcal{L})_u$  est le noyau de  $A_f \times A_g \to A_{fg}$ ,
nous concluons.
\end{proof}

\begin{lemma}
\label{lemma-Pic-projective-space-UFD}
Soit  $R$  un UFD. Le groupe de Picard de  $\mathbf{P}^n_R$  est  $\mathbf{Z}$.
Plus précisément, il existe un isomorphisme
$$
\mathbf{Z} \longrightarrow \Pic(\mathbf{P}^n_R),\quad
m \longmapsto \mathcal{O}_{\mathbf{P}^n_R}(m)
$$
. En particulier, le groupe de Picard de  $\mathbf{P}^n_k$  de l'espace projectif sur un
corps  $k$  est  $\mathbf{Z}$.
\end{lemma}

\begin{proof}
D'après le Lemme  \ref{lemma-open-subscheme-UFD} , les groupes de Picard des ouverts  $D_+(T_i) \cong \mathbf{A}^n_R$  sont
triviaux. Ainsi, si  $\mathcal{L}$  est un module inversible sur  $\mathbf{P}^n_R$, alors
il est donné par un cocycle  $1$ à valeurs dans le faisceau des foncteurs
inversibles  $\mathcal{O}^*$  pour le recouvrement ouvert
$$
\mathbf{P}^n_R = \bigcup\nolimits_{i = 0, \ldots, n} D_+(T_i)
$$
. Remarquons que pour  $i \not = j$  nous avons
$$
\mathcal{O}^*(D_+(T_i) \cap D_+(T_j)) =
\mathcal{O}^*(D_+(T_iT_j)) = \left(R[T_0, \ldots, T_n]_{(T_iT_j)}\right)^* =
R^* \times (T_i/T_j)^\mathbf{Z}
$$
. Ainsi, un tel cocycle  $(g_{ij})$  est donné par des unités
$$
g_{ij} = u_{ij} (T_i/T_j)^{e_{ij}}
$$
avec  $u_{ij} \in R^*$  et  $e_{ij} \in \mathbf{Z}$  satisfaisant à la condition de cocycle. La
condition de cocycle sur  $D_+(T_iT_jT_k)$  pour  $\#\{i, j, k\} = 3$  nous indique que
$$
u_{ik} = u_{ij} u_{jk} \quad\text{and}\quad
e_{ik} = e_{ij} = e_{jk}
$$
. D'où  $u_{ij} = u_{i1} u_{j1}^{-1}$  est une frontière. Ainsi, toutes les classes d'isomorphisme de
modules inversibles sont obtenues en prenant le cocycle avec
$$
g_{ij} = (T_i/T_j)^e
$$
pour un certain  $e \in \mathbf{Z}$. Puisque  $\mathcal{O}(n)$  possède une section
trivialisation  $T_i^n$  sur  $D_+(T_i)$ , nous voyons que le cocycle
correspondant de  $\mathcal{O}(n)$  est  $(T_i/T_j)^n$  et la preuve est complète.
\end{proof}






\section{Diviseurs de Weil sur les schémas normaux}
\label{section-weil-divisors-normal}

\noindent
Tout d'abord, nous discutons des propriétés des modules réflexifs.

\begin{lemma}
\label{lemma-reflexive-normal}
Soit  $X$  un schéma normal intègre et localement noethérien. Pour
 $\mathcal{F}$  et  $\mathcal{G}$  des  $\mathcal{O}_X$-modules réflexifs cohérents,
l'application
$$
(\SheafHom_{\mathcal{O}_X}(\mathcal{F}, \mathcal{O}_X)
\otimes_{\mathcal{O}_X} \mathcal{G})^{**} \to
\SheafHom_{\mathcal{O}_X}(\mathcal{F}, \mathcal{G})
$$
est un isomorphisme. La règle  $\mathcal{F}, \mathcal{G} \mapsto
(\mathcal{F} \otimes_{\mathcal{O}_X} \mathcal{G})^{**}$  définit une loi de groupe abélien sur
l'ensemble des classes d'isomorphisme des  $1$   $\mathcal{O}_X$-modules
réflexifs cohérents.
\end{lemma}

\begin{proof}
Bien que cela ne soit pas strictement nécessaire, nous recommandons de lire la
Remarque  \ref{remark-tensor}  avant de procéder à la démonstration. Choisissez un
sous-schéma ouvert  $j : U \to X$  de sorte que chaque composante irréductible de
 $X \setminus U$  ait codimension  $\geq 2$  dans  $X$  et de sorte que
 $j^*\mathcal{F}$  et  $j^*\mathcal{G}$  soient localement libres de rang fini, voir le Lemme
 \ref{lemma-reflexive-over-normal}. L'application
$$
\SheafHom_{\mathcal{O}_U}(j^*\mathcal{F}, \mathcal{O}_U)
\otimes_{\mathcal{O}_U} j^*\mathcal{G} \to
\SheafHom_{\mathcal{O}_U}(j^*\mathcal{F}, j^*\mathcal{G})
$$
est un isomorphisme, car nous pouvons le vérifier localement et cela est clair
lorsque les modules sont libres de type fini. Observez que  $j^*$  appliqué à
la flèche affichée du lemme donne la flèche que nous venons de montrer comme étant
un isomorphisme (petit détail omis). Puisque  $j^*$  définit une équivalence
entre les modules réflexifs cohérents sur  $U$  et les modules réflexifs
cohérents sur  $X$  (d'après le Lemme  \ref{lemma-reflexive-S2-extend}  et le critère de Serre,
Lemme  \ref{properties-lemma-criterion-normal}), nous concluons que la flèche du lemme est aussi un
isomorphisme. Si  $\mathcal{F}$  a un rang  $1$, alors  $j^*\mathcal{F}$  est un
 $\mathcal{O}_U$-module inversible et le module réflexif  $\mathcal{F}^\vee = \SheafHom(\mathcal{F}, \mathcal{O}_X)$  se restreint à
son inverse. Il s'ensuit de la même manière qu'auparavant que  $(\mathcal{F} \otimes_{\mathcal{O}_X} \mathcal{F}^\vee)^{**} = \mathcal{O}_X$. De
cette façon, nous voyons que nous avons des inverses pour la loi de groupe donnée
dans l'énoncé du lemme.
\end{proof}

\begin{lemma}
\label{lemma-normal-class-group}
Laisse $X$ soit un schéma normal localement noethérien et intègre. Le
groupe de rang $1$ réflexif cohérent $\mathcal{O}_X$-modules est isomorphe au
groupe des classes de diviseurs de Weil $\text{Cl}(X)$  de  $X$.
\end{lemma}

\begin{proof}
Laisse $\mathcal{F}$ être un rang $1$ réflexif cohérent $\mathcal{O}_X$-module.
Choisissez un ouvert $U \subset X$  de telle sorte que chaque composante irréductible
de  $X \setminus U$ a une codimension $\geq 2$ dans $X$  et tel que
 $\mathcal{F}|_U$  est inversible, voir le Lemme  \ref{lemma-reflexive-over-normal}. Remarquez que
 $\text{Cl}(U) = \text{Cl}(X)$ en tant que groupe de classes de diviseurs de Weil de $X$ ne
dépend que de son corps de fonctions rationnelles et des points de
codimension $1$  et leurs anneaux locaux. Ainsi, nous pouvons définir la
classe de diviseur de Weil de  $\mathcal{F}$ être la classe de diviseur de Weil
de $\mathcal{F}|_U$ dans $\text{Cl}(U)$. Nous omettons la vérification que ceci est
indépendant du choix de  $U$.

\medskip\noindent
Désignons par  $\text{Cl}'(X)$  l'ensemble des classes d'isomorphisme de  $1$ 
modules cohérents réflexifs de rang  $\mathcal{O}_X$. La construction ci-dessus donne
un homomorphisme de groupes
$$
\text{Cl}'(X) \longrightarrow \text{Cl}(X)
$$
car pour toute paire  $\mathcal{F}, \mathcal{G}$  d'éléments de  $\text{Cl}'(X)$ , nous pouvons choisir
un  $U$  qui convient pour les deux et l'attribution (\ref{equation-c1})
envoyant un module inversible à sa classe de diviseur de Weil est un
homomorphisme. Si  $\mathcal{F}$  est dans le noyau de cette application, alors nous
trouvons que  $\mathcal{F}|_U$  est trivial (Lemme  \ref{lemma-normal-c1-injective}) et donc  $\mathcal{F}$ 
est aussi trivial d'après le Lemme  \ref{lemma-reflexive-S2-extend}  et le critère de Serre Propriétés,
Lemme  \ref{properties-lemma-criterion-normal}. Pour terminer la preuve, il suffit de vérifier que
l'application est surjective.

\medskip\noindent
Soit  $D = \sum n_Z Z$  un diviseur de Weil sur  $X$. Nous affirmons qu'il
existe un ouvert  $U \subset X$  tel que chaque composante irréductible de
 $X \setminus U$  ait codimension  $\geq 2$  dans  $X$  et tel que
 $Z|_U$  soit un diviseur de Cartier effectif pour  $n_Z \not = 0$. Pour prouver
cette affirmation, nous pouvons supposer que  $X$  est affine. Ensuite,
nous pouvons supposer que  $D = n_1 Z_1 + \ldots + n_r Z_r$  est une somme finie avec  $Z_1, \ldots, Z_r$  deux
à deux distincts. Après avoir éliminé  $Z_i \cap Z_j$  pour  $i \not = j$ , nous pouvons
supposer que  $Z_1, \ldots, Z_r$  sont deux à deux disjoints. Cela nous ramène au cas d'un
seul diviseur premier  $Z$  sur  $X$. Comme  $X$  est
 $(R_1)$  par les propriétés, le lemme  \ref{properties-lemma-criterion-normal}  montre que l'anneau local
 $\mathcal{O}_{X, \xi}$  au point générique  $\xi$  de  $Z$  est un anneau de
valuation discrète. Soit  $f \in \mathcal{O}_{X, \xi}$  un uniformisant. Soit  $V \subset X$  un
voisinage ouvert de  $\xi$  tel que  $f$  soit l'image d'un élément
 $f \in \mathcal{O}_X(V)$. Après avoir réduit  $V$ , nous pouvons supposer que
 $Z \cap V = V(f)$  est schématiquement, puisque c'est vrai dans l'anneau local en
 $\xi$. Dans ce cas, prendre
$$
U = X \setminus (Z \setminus V) = (X \setminus Z) \cup V
$$
donne l'ouvert souhaité, prouvant ainsi l'affirmation.

\medskip\noindent
Afin de montrer que la classe de diviseur de  $D$  est dans l'image, nous
pouvons écrire  $D = \sum_{n_Z < 0} n_Z Z - \sum_{n_Z > 0} (-n_Z) Z$. Par additivité de l'application construite ci-dessus,
nous pouvons et supposons  $n_Z \leq 0$  pour tous les diviseurs premiers
 $Z$  (cette étape peut être évitée si le lecteur le souhaite). Soit
 $U \subset X$  comme dans l'assertion ci-dessus. Si  $U$  est quasi-compact,
alors nous écrivons  $D|_U = -n_1 Z_1 - \ldots - n_r Z_r$  pour des diviseurs premiers deux à deux distincts
 $Z_i$  et  $n_i > 0$  et nous considérons le  $\mathcal{O}_U$-module inversible
$$
\mathcal{L} =
\mathcal{I}_1^{n_1} \ldots \mathcal{I}_r^{n_r} \subset \mathcal{O}_U
$$
où  $\mathcal{I}_i$  est le faisceau d'idéaux de  $Z_i$. C'est inversible par
notre choix de  $U$  et le Lemme  \ref{lemma-sum-effective-Cartier-divisors}. De plus,  $\text{div}_\mathcal{L}(1) = D|_U$.
Puisque  $\mathcal{L} = \mathcal{F}|_U$  pour un certain rang  $1$  module  $\mathcal{O}_X$
réflexif cohérent  $\mathcal{F}$  par le Lemme  \ref{lemma-reflexive-S2-extend} , nous constatons que
 $D$  est dans l'image de notre application.

\medskip\noindent
Si  $U$  n'est pas quasi-compact, alors nous définissons  $\mathcal{L} \subset \mathcal{O}_U$ 
localement par la formule affichée ci-dessus. Le lecteur montre que la
construction se recolle et termine la preuve exactement comme auparavant. Détails
omis.
\end{proof}

\begin{lemma}
\label{lemma-structure-sheaf-Xs}
Soit  $X$  un schéma normal intègre localement noethérien. Soit
 $\mathcal{F}$  un  $\mathcal{O}_X$-module cohérent réflexif de rang 1. Soit
 $s \in \Gamma(X, \mathcal{F})$.
$$
U = \{x \in X \mid s : \mathcal{O}_{X, x} \to \mathcal{F}_x
\text{ is an isomorphism}\}
$$
Alors  $j : U \to X$  est un sous-schéma ouvert de  $X$  et
$$
j_*\mathcal{O}_U =
\colim (\mathcal{O}_X \xrightarrow{s} \mathcal{F}
\xrightarrow{s} \mathcal{F}^{[2]}
\xrightarrow{s} \mathcal{F}^{[3]}
\xrightarrow{s} \ldots)
$$
où  $\mathcal{F}^{[1]} = \mathcal{F}$  et par induction  $\mathcal{F}^{[n + 1]} =
(\mathcal{F} \otimes_{\mathcal{O}_X} \mathcal{F}^{[n]})^{**}$.
\end{lemma}

\begin{proof}
L'ensemble  $U$  est ouvert selon Modules, Lemme  \ref{modules-lemma-finite-type-surjective-on-stalk}  et
 \ref{modules-lemma-finite-type-to-coherent-injective-on-stalk}. On observe que  $j$  est quasi-compact selon Propriétés,
Lemme  \ref{properties-lemma-immersion-into-noetherian}. Pour prouver l'affirmation finale, il suffit de montrer que
pour tout ouvert quasi-compact  $W \subset X$ , il existe un isomorphisme
$$
\colim \Gamma(W, \mathcal{F}^{[n]})
\longrightarrow
\Gamma(U \cap W, \mathcal{O}_U)
$$
de  $\mathcal{O}_X(W)$-modules compatibles avec les applications de restriction. Nous
omettrons la vérification des compatibilités. Après avoir remplacé  $X$ 
par  $W$  et réécrit ce qui précède en termes de homs, nous voyons qu'il
suffit de construire un isomorphisme
$$
\colim \Hom_{\mathcal{O}_X}(\mathcal{O}_X, \mathcal{F}^{[n]})
\longrightarrow
\Hom_{\mathcal{O}_U}(\mathcal{O}_U, \mathcal{O}_U)
$$
Choisissez un ouvert  $V \subset X$  tel que chaque composante irréductible de
 $X \setminus V$  ait une codimension  $\geq 2$  dans  $X$  et tel que
 $\mathcal{F}|_V$  soit inversible, voir le Lemme  \ref{lemma-reflexive-over-normal}. Alors la restriction
définit une équivalence de catégories entre les  $1$ -modules cohérents
réflexifs de rang sur  $X$  et  $V$  et entre les
 $1$ -modules cohérents réflexifs de rang sur  $U$  et  $V \cap U$.
Voir le Lemme  \ref{lemma-reflexive-S2-extend}  et le critère de Serre Propriétés, Lemme  \ref{properties-lemma-criterion-normal}.
Ainsi il suffit de construire un isomorphisme
$$
\colim \Gamma(V, (\mathcal{F}|_V)^{\otimes n}) \longrightarrow
\Gamma(V \cap U, \mathcal{O}_U)
$$
Puisque  $\mathcal{F}|_V$  est inversible et puisque  $U \cap V$  est égal à l'ensemble
des points où  $s|_V$  génère ce module inversible, il s'agit d'un cas
particulier des Propriétés, Lemme  \ref{properties-lemma-invert-s-sections}  (il existe également une formule
explicite pour l'application).
\end{proof}

\begin{lemma}
\label{lemma-Xs-codim-complement}
Hypothèses et notations comme dans le Lemme  \ref{lemma-structure-sheaf-Xs}. Si  $s$  est non
nul, alors chaque composante irréductible de  $X \setminus U$  a une codimension
 $1$  dans  $X$.
\end{lemma}

\begin{proof}
Soit  $\xi \in X$  un point générique d'une composante irréductible  $Z$  de
 $X \setminus U$. Après avoir remplacé  $X$  par un voisinage ouvert de
 $\xi$ , on peut supposer que  $Z = X \setminus U$  est irréductible. Puisque
 $s : \mathcal{O}_U \to \mathcal{F}|_U$  est un isomorphisme, si la codimension de  $Z$  dans
 $X$  est  $\geq 2$, alors  $s : \mathcal{O}_X \to \mathcal{F}$  est un isomorphisme d'après le
Lemme  \ref{lemma-reflexive-S2-extend}  et le critère de Serre Propriétés, Lemme  \ref{properties-lemma-criterion-normal}. Cela
signifierait que  $Z = \emptyset$, ce qui constitue une contradiction.
\end{proof}

\begin{remark}
\label{remark-structure-sheaf-Xs}
Soit  $A$  un domaine normal noethérien. Soit  $M$  un  $1$ 
de rang  $A$-module réflexif fini. Soit  $s \in M$  non nul. Soit
 $\mathfrak p_1, \ldots, \mathfrak p_r$  les idéaux premiers de hauteur  $1$  de  $A$  dans le
support de  $M/As$. Alors l'ouvert  $U$  du Lemme  \ref{lemma-structure-sheaf-Xs}  est
$$
U = \Spec(A) \setminus
\left(V(\mathfrak p_1) \cup \ldots \cup V(\mathfrak p_r)\right)
$$
selon le Lemme  \ref{lemma-Xs-codim-complement}. De plus, si  $M^{[n]}$  désigne l'enveloppe
réflexive de  $M \otimes_A \ldots \otimes_A M$  (facteurs $n$), alors
$$
\Gamma(U, \mathcal{O}_U) = \colim M^{[n]}
$$
selon le Lemme  \ref{lemma-structure-sheaf-Xs}.
\end{remark}

\begin{lemma}
\label{lemma-affine-Xs}
Hypothèses et notation comme dans le Lemme  \ref{lemma-structure-sheaf-Xs}. Les éléments suivants
sont équivalents
\begin{enumerate}
\item le morphisme d'inclusion  $j : U \to X$  est affine, et
\item pour chaque  $x \in X \setminus U$  il existe un  $n > 0$  tel que  $s^n \in \mathfrak m_x \mathcal{F}^{[n]}_x$.
\end{enumerate}
\end{lemma}

\begin{proof}
Supposons (1). Alors pour  $x \in X \setminus U$  l'image inverse  $U_x$  de
 $U$  sous le morphisme canonique  $f_x : \Spec(\mathcal{O}_{X, x}) \to X$  est affine et ne contient
pas  $x$. Ainsi,  $\mathfrak m_x \Gamma(U_x, \mathcal{O}_{U_x})$  est l'idéal unité. En particulier, nous
voyons que nous pouvons écrire
$$
1 = \sum f_i g_i
$$
avec  $f_i \in \mathfrak m_x$  et  $g_i \in \Gamma(U_x, \mathcal{O}_{U_x})$. Par le Lemme  \ref{lemma-structure-sheaf-Xs}  nous avons
 $\Gamma(U_x, \mathcal{O}_{U_x}) = \colim \mathcal{F}^{[n]}_x$  avec des applications de transition données par la multiplication par
 $s$. Donc pour certains  $n > 0$  nous avons
$$
s^n = \sum f_i t_i
$$
pour certains  $t_i = s^ng_i \in \mathcal{F}^{[n]}_x$. Ainsi (2) est vrai.

\medskip\noindent
Inversement, supposons que (2) est vrai. Pour prouver  $j$  est affine est
local sur  $X$, voir Morphismes, Lemme  \ref{morphisms-lemma-characterize-affine}. Ainsi, on peut
supposer que  $X$  est affine. Notre objectif est de montrer que
 $U$  est affine. Par cohomologie des schémas, Lemme  \ref{coherent-lemma-affine-if-quasi-affine}  il suffit
de montrer que  $H^p(U, \mathcal{O}_U) = 0$  pour  $p > 0$. Puisque  $H^p(U, \mathcal{O}_U) = H^0(X, R^pj_*\mathcal{O}_U)$  (Cohomologie
des schémas, Lemme  \ref{coherent-lemma-quasi-coherence-higher-direct-images-application}) et puisque  $R^pj_*\mathcal{O}_U$  est quasi-cohérente
(Cohomologie des schémas, Lemme  \ref{coherent-lemma-quasi-coherence-higher-direct-images}), il suffit de montrer que la fibre
 $(R^pj_*\mathcal{O}_U)_x$  en un point  $x \in X$  est nulle. Considérons le diagramme de
changement de base
$$
\xymatrix{
U_x \ar[d]_{j_x} \ar[r] & U \ar[d]^j \\
\Spec(\mathcal{O}_{X, x}) \ar[r] & X
}
$$
Par la Cohomologie des Schémas, Lemme  \ref{coherent-lemma-flat-base-change-cohomology}  nous avons  $(R^pj_*\mathcal{O}_U)_x = R^pj_{x, *}\mathcal{O}_{U_x}$. Par
conséquent, nous pouvons supposer que  $X$  est local avec un point fermé
 $x$  et nous devons montrer que  $U$  est affine (car cela équivaut
à l'annulation souhaitée selon la référence donnée ci-dessus). En particulier,
 $d = \dim(X)$  est fini (Algèbre, Proposition  \ref{algebra-proposition-dimension}). Si  $x \in U$, alors
 $U = X$  et le résultat est clair. Si  $d = 0$  et  $x \not \in U$, alors
 $U = \emptyset$  et le résultat est clair. Supposons maintenant  $d > 0$  et
 $x \not \in U$. Comme  $j_*\mathcal{O}_U = \colim \mathcal{F}^{[n]}$ , notre hypothèse signifie que nous pouvons écrire
$$
1 = \sum f_i g_i
$$
pour certains  $n > 0$,  $f_i \in \mathfrak m_x$, et  $g_i \in \mathcal{O}(U)$. Par induction sur
 $d$ , nous savons que  $D(f_i) \cap U$  est affine pour tout  $i$: en
parcourant tout l'argument donné avec  $X$  remplacé par  $D(f_i)$ , nous
obtenons des anneaux locaux noethériens dont la dimension est strictement
inférieure à  $d$. Par conséquent,  $U$  est affine d'après
Propriétés, Lemme  \ref{properties-lemma-characterize-affine}  comme désiré.
\end{proof}





\section{Proj relatif}
\label{section-relative-proj}

\noindent
Quelques résultats sur Proj relatif. Tout d'abord, quelques résultats très
basiques. Rappelons qu'un Proj relatif est toujours séparé sur la base, voir
Constructions, Lemme  \ref{constructions-lemma-relative-proj-separated}.

\begin{lemma}
\label{lemma-relative-proj-quasi-compact}
Soit  $S$  un schéma. Soit  $\mathcal{A}$  une  $\mathcal{O}_S$-algèbre graduée
quasi-cohérente. Soit  $p : X = \underline{\text{Proj}}_S(\mathcal{A}) \to S$  le Proj relatif de  $\mathcal{A}$. Si l'une des
conditions suivantes est satisfaite
\begin{enumerate}
\item $\mathcal{A}$  est de type fini comme faisceau d'algèbres de  $\mathcal{A}_0$,
\item $\mathcal{A}$  est généré par  $\mathcal{A}_1$  comme un  $\mathcal{A}_0$-algèbre et
 $\mathcal{A}_1$  est un  $\mathcal{A}_0$-module de type fini,
\item il existe un  $\mathcal{A}_0$-sous-module quasi-cohérent de type fini  $\mathcal{F} \subset \mathcal{A}_{+}$  tel
que  $\mathcal{A}_{+}/\mathcal{F}\mathcal{A}$  soit un faisceau d'idéaux localement nilpotent de  $\mathcal{A}/\mathcal{F}\mathcal{A}$,
\end{enumerate}
alors  $p$  est quasi-compact.
\end{lemma}

\begin{proof}
La question est locale sur la base, voir Schémas, Lemme  \ref{schemes-lemma-quasi-compact-affine}. Ainsi, nous
pouvons supposer que  $S$  est affine. Disons que  $S = \Spec(R)$  et
 $\mathcal{A}$  correspondent à l'algèbre  $R$ graduée  $A$. Alors
 $X = \text{Proj}(A)$, voir Constructions, Section  \ref{constructions-section-relative-proj-via-glueing}. Dans le cas (1), nous
pouvons, après éventuellement avoir localisé davantage, supposer que  $A$ 
est généré par des éléments homogènes  $f_1, \ldots, f_n \in A_{+}$  sur  $A_0$. Alors
 $A_{+} = (f_1, \ldots, f_n)$  par Algèbre, Lemme  \ref{algebra-lemma-S-plus-generated}. Dans le cas (3), nous voyons que
 $\mathcal{F} = \widetilde{M}$  pour un certain module  $A_0$ de type fini  $M \subset A_{+}$. Disons
 $M = \sum A_0f_i$. Disons que  $f_i = \sum f_{i, j}$  est la décomposition en parties homogènes.
La condition du cas (3) signifie que  $A_{+} \subset \sqrt{(f_{i, j})}$. Ainsi, dans les deux cas, nous
concluons que  $\text{Proj}(A)$  est quasi-compact par Constructions, Lemme
 \ref{constructions-lemma-proj-quasi-compact}. Enfin, (2) découle de (1).
\end{proof}

\begin{lemma}
\label{lemma-relative-proj-finite-type}
Soit  $S$  un schéma. Soit  $\mathcal{A}$  une  $\mathcal{O}_S$-algèbre graduée
quasi-cohérente. Soit  $p : X = \underline{\text{Proj}}_S(\mathcal{A}) \to S$  le Proj relatif de  $\mathcal{A}$. Si
 $\mathcal{A}$  est de type fini en tant que faisceau d'algèbres  $\mathcal{O}_S$, alors
 $p$  est de type fini et  $\mathcal{O}_X(d)$  est un  $\mathcal{O}_X$-module de type
fini.
\end{lemma}

\begin{proof}
L'hypothèse implique que  $p$  est quasi-compact, voir le Lemme
 \ref{lemma-relative-proj-quasi-compact}. Il suffit donc de montrer que  $p$  est localement de type
fini. Ainsi, la question est locale sur la base et la cible, voir Morphismes,
Lemme  \ref{morphisms-lemma-locally-finite-type-characterize}. Supposons que  $S = \Spec(R)$  et  $\mathcal{A}$  correspondent à
l'algèbre  $R$ graduée  $A$. Après une nouvelle localisation sur
 $S$ , nous pouvons supposer que  $A$  est une  $R$-algèbre
de type fini. Le schéma  $X$  est construit en collant les spectres des
anneaux  $A_{(f)}$  pour  $f \in A_{+}$  homogène. Chacun d'eux est de type fini sur
 $R$  d'après Algèbre, Lemme  \ref{algebra-lemma-dehomogenize-finite-type}  partie (1). Ainsi  $\text{Proj}(A)$ 
est de type fini sur  $R$. Pour voir l'énoncé sur  $\mathcal{O}_X(d)$ , utilisez
la partie (2) de Algèbre, Lemme  \ref{algebra-lemma-dehomogenize-finite-type}.
\end{proof}

\begin{lemma}
\label{lemma-relative-proj-universally-closed}
Soit  $S$  un schéma. Soit  $\mathcal{A}$  une  $\mathcal{O}_S$-algèbre graduée
quasi-cohérente. Soit  $p : X = \underline{\text{Proj}}_S(\mathcal{A}) \to S$  le Proj relatif de  $\mathcal{A}$. Si
 $\mathcal{O}_S \to \mathcal{A}_0$  est une application d'algèbres intègre \footnote{ En
d'autres termes, la clôture intégrale de  $\mathcal{O}_S$  dans  $\mathcal{A}_0$, voir
Morphismes, Définition  \ref{morphisms-definition-integral-closure}, est égale à  $\mathcal{A}_0$. }  et
 $\mathcal{A}$  est de type fini comme  $\mathcal{A}_0$-algèbre, alors  $p$  est
universellement fermé.
\end{lemma}

\begin{proof}
La question est locale sur la base. Ainsi, nous pouvons supposer que  $X = \Spec(R)$ 
est affine. Soit  $\mathcal{A}$  l'algèbre  $\mathcal{O}_X$ quasi-cohérente associée à
l'algèbre  $R$ graduée  $A$. L'hypothèse est que  $R \to A_0$  est
intègre et que  $A$  est de type fini sur  $A_0$. Écrivons
 $X \to \Spec(R)$  comme la composition  $X \to \Spec(A_0) \to \Spec(R)$. Comme  $R \to A_0$  est un
morphisme d'anneaux intègre, nous voyons que  $\Spec(A_0) \to \Spec(R)$  est universellement
fermé, voir Morphismes, Lemme  \ref{morphisms-lemma-integral-universally-closed}. Le morphisme quasi-compact (voir
Constructions, Lemme  \ref{constructions-lemma-proj-quasi-compact})
$$
X = \text{Proj}(A) \to \Spec(A_0)
$$
satisfait la partie existence du critère valuatif par Constructions, Lemme
 \ref{constructions-lemma-proj-valuative-criterion}  et donc il est universellement fermé par Schémas, Proposition
 \ref{schemes-proposition-characterize-universally-closed}. Ainsi,  $X \to \Spec(R)$  est universellement fermé comme composition de
morphismes universellement fermés.
\end{proof}

\begin{lemma}
\label{lemma-relative-proj-proper}
Soit  $S$  un schéma. Soit  $\mathcal{A}$  une  $\mathcal{O}_S$-algèbre graduée
quasi-cohérente. Soit  $p : X = \underline{\text{Proj}}_S(\mathcal{A}) \to S$  le Proj relatif de  $\mathcal{A}$. Les conditions
suivantes sont équivalentes
\begin{enumerate}
\item $\mathcal{A}_0$  est un  $\mathcal{O}_S$-module de type fini et  $\mathcal{A}$  est de type
fini comme  $\mathcal{A}_0$-algèbre,
\item $\mathcal{A}_0$  est un  $\mathcal{O}_S$-module de type fini et  $\mathcal{A}$  est de type
fini comme  $\mathcal{O}_S$-algèbre
\end{enumerate}
Si ces conditions sont vérifiées, alors  $p$  est localement projectif et
en particulier propre.
\end{lemma}

\begin{proof}
Supposons que  $\mathcal{A}_0$  soit un  $\mathcal{O}_S$-module de type fini. Choisissez un
ouvert affine  $U = \Spec(R) \subset X$  tel que  $\mathcal{A}$  corresponde à une
 $R$-algèbre graduée  $A$  avec  $A_0$  un  $R$-module
fini. La condition (1) signifie que (après éventuellement une localisation
supplémentaire sur  $S$)  $A$  est une  $A_0$-algèbre de type
fini et la condition (2) signifie que (après éventuellement une localisation
supplémentaire sur  $S$)  $A$  est une  $R$-algèbre de type
fini. Ainsi, ces conditions s'impliquent mutuellement d'après le Lemme d'Algèbre
 \ref{algebra-lemma-compose-finite-type}.

\medskip\noindent
Un morphisme localement projectif est propre, voir Morphismes, Lemme  \ref{morphisms-lemma-locally-projective-proper}.
Ainsi, nous pouvons maintenant supposer que  $S = \Spec(R)$  et  $X = \text{Proj}(A)$  et que
 $A_0$  est fini sur  $R$  et  $A$  de type fini sur
 $R$. Nous allons montrer que  $X = \text{Proj}(A) \to \Spec(R)$  est projectif. Nous
encourageons le lecteur à le prouver lui-même, en construisant directement une
immersion fermée de  $X$  dans un espace projectif sur  $R$, plutôt
qu'en lisant l'argument que nous donnons ci-dessous.

\medskip\noindent
D'après le Lemme  \ref{lemma-relative-proj-finite-type} , nous voyons que  $X$  est de type fini sur
 $\Spec(R)$. Constructions, le Lemme  \ref{constructions-lemma-ample-on-proj}  nous dit que  $\mathcal{O}_X(d)$  est
ample sur  $X$  pour un certain  $d \geq 1$  (voir Propriétés, Section
 \ref{properties-section-ample}). Par conséquent,  $X \to \Spec(R)$  est quasi-projectif (d'après
Morphismes, Définition  \ref{morphisms-definition-quasi-projective}). D'après Morphismes, le Lemme  \ref{morphisms-lemma-quasi-projective-open-projective} ,
nous concluons que  $X$  est isomorphe à un sous-schéma ouvert d'un schéma
projectif sur  $\Spec(R)$. Par conséquent, pour terminer la démonstration, il
suffit de montrer que  $X \to \Spec(R)$  est universellement fermé (utiliser Morphismes,
Lemme  \ref{morphisms-lemma-image-proper-scheme-closed}). Cela découle du Lemme  \ref{lemma-relative-proj-universally-closed}.
\end{proof}

\begin{lemma}
\label{lemma-relative-proj-projective}
Soit  $S$  un schéma. Soit  $\mathcal{A}$  une  $\mathcal{O}_S$-algèbre graduée
quasi-cohérente. Soit  $p : X = \underline{\text{Proj}}_S(\mathcal{A}) \to S$  le Proj relatif de  $\mathcal{A}$. Si
 $\mathcal{A}$  est engendré par  $\mathcal{A}_1$  sur  $\mathcal{A}_0$  et  $\mathcal{A}_1$  est un
 $\mathcal{O}_S$-module de type fini, alors  $p$  est projectif.
\end{lemma}

\begin{proof}
À savoir, le morphisme associé à l'application d'algèbres graduées  $\mathcal{O}_S$
$$
\text{Sym}_{\mathcal{O}_X}^*(\mathcal{A}_1)
\longrightarrow
\mathcal{A}
$$
est une immersion fermée  $X \to \mathbf{P}(\mathcal{A}_1)$, voir Constructions, Lemme  \ref{constructions-lemma-surjective-generated-degree-1-map-relative-proj}.
\end{proof}

\begin{lemma}
\label{lemma-relative-proj-flat}
Soit  $S$  un schéma. Soit  $\mathcal{A}$  une  $\mathcal{O}_S$-algèbre graduée
quasi-cohérente. Soit  $p : X = \underline{\text{Proj}}_S(\mathcal{A}) \to S$  le Proj relatif de  $\mathcal{A}$. Si
 $\mathcal{A}_d$  est un  $\mathcal{O}_S$-module plat pour  $d \gg 0$, alors  $p$ 
est plat et  $\mathcal{O}_X(d)$  est plat sur  $S$.
\end{lemma}

\begin{proof}
La platitude localement affine de  $X$  sur  $S$  se réduit à
l'énoncé suivant: Soit  $R$  un anneau, soit  $A$  une
 $R$-algèbre graduée avec  $A_d$  plate sur  $R$  pour
 $d \gg 0$, soit  $f \in A_d$  pour un certain  $d > 0$, alors  $A_{(f)}$ 
est plat sur  $R$. Comme  $A_{(f)} = \colim A_{nd}$  où les cartes de transition sont
données par multiplication par  $f$, ceci découle du Lemme  \ref{algebra-lemma-colimit-flat}
d'Algèbre. Argumentez de manière similaire pour obtenir la platitude de
 $\mathcal{O}_X(d)$  sur  $S$.
\end{proof}

\begin{lemma}
\label{lemma-relative-proj-finite-presentation}
Soit  $S$  un schéma. Soit  $\mathcal{A}$  une  $\mathcal{O}_S$-algèbre graduée
quasi-cohérente. Soit  $p : X = \underline{\text{Proj}}_S(\mathcal{A}) \to S$  le Proj relatif de  $\mathcal{A}$. Si
 $\mathcal{A}$  est une  $\mathcal{O}_S$-algèbre de présentation finie, alors
 $p$  est de présentation finie et  $\mathcal{O}_X(d)$  est un  $\mathcal{O}_X$-module
de présentation finie.
\end{lemma}

\begin{proof}
Localement affine, cela se réduit à l'énoncé suivant: Soit  $R$  un anneau
et soit  $A$  une  $R$-algèbre graduée de présentation finie. Alors
 $\text{Proj}(A) \to \Spec(R)$  est de présentation finie et  $\mathcal{O}_{\text{Proj}(A)}(d)$  est un  $\mathcal{O}_{\text{Proj}(A)}$-module
de présentation finie. La condition de présentation finie implique que nous
pouvons choisir une présentation
$$
A = R[X_1, \ldots, X_n]/(F_1, \ldots, F_m)
$$
où  $R[X_1, \ldots, X_n]$  est un anneau polynomial gradué en donnant des poids  $d_i$ 
à  $X_i$  et  $F_1, \ldots, F_m$  sont des polynômes homogènes de degré  $e_j$.
Soit  $R_0 \subset R$  le sous-anneau généré par les coefficients des polynômes
 $F_1, \ldots, F_m$. Ensuite, nous posons  $A_0 = R_0[X_1, \ldots, X_n]/(F_1, \ldots, F_m)$. Par construction,  $A = A_0 \otimes_{R_0} R$.
Ainsi, d'après Constructions, Lemme  \ref{constructions-lemma-base-change-map-proj} , il suffit de prouver le résultat
pour  $X_0 = \text{Proj}(A_0)$  sur  $R_0$. D'après le Lemme  \ref{lemma-relative-proj-finite-type} , nous savons que
 $X_0$  est de type fini sur  $R_0$  et que  $\mathcal{O}_{X_0}(d)$  est un
 $\mathcal{O}_{X_0}$-module quasi-cohérent de type fini. Puisque  $R_0$  est
noethérien (en tant qu'algèbre  $\mathbf{Z}$ de type fini), nous voyons que
 $X_0$  est de présentation finie sur  $R_0$  (Morphismes, Lemme
 \ref{morphisms-lemma-noetherian-finite-type-finite-presentation}) et que  $\mathcal{O}_{X_0}(d)$  est de présentation finie d'après Cohomology of
Schémas, Lemme  \ref{coherent-lemma-coherent-Noetherian}. Cela termine la démonstration.
\end{proof}










\section{Sous-schémas fermés du Proj relatif}
\label{section-closed-in-relative-proj}

\noindent
Quelques lemmes auxiliaires sur les sous-schémas fermés de Proj relatif.

\begin{lemma}
\label{lemma-closed-subscheme-proj}
Soit  $S$  un schéma. Soit  $\mathcal{A}$  une  $\mathcal{O}_S$-algèbre graduée
quasi-cohérente. Soit  $p : X = \underline{\text{Proj}}_S(\mathcal{A}) \to S$  le Proj relatif de  $\mathcal{A}$. Soit
 $i : Z \to X$  un sous-schéma fermé. On note  $\mathcal{I} \subset \mathcal{A}$  le noyau de l'application
canonique
$$
\mathcal{A}
\longrightarrow
\bigoplus\nolimits_{d \geq 0} p_*\left((i_*\mathcal{O}_Z)(d)\right).
$$
Si  $p$  est quasi-compact, alors il existe un isomorphisme  $Z = \underline{\text{Proj}}_S(\mathcal{A}/\mathcal{I})$.
\end{lemma}

\begin{proof}
Le morphisme  $p$  est séparé par Constructions, Lemme  \ref{constructions-lemma-relative-proj-separated}. Comme
 $p$  est quasi-compact,  $p_*$  transforme les modules
quasi-cohérents en modules quasi-cohérents, voir Schémas, Lemme  \ref{schemes-lemma-push-forward-quasi-coherent}. Par
conséquent,  $\mathcal{I}$  est un  $\mathcal{O}_S$-module quasi-cohérent. En
particulier,  $\mathcal{B} = \mathcal{A}/\mathcal{I}$  est une  $\mathcal{O}_S$-algèbre graduée quasi-cohérente. Le
morphisme de fonctorialité  $Z' = \underline{\text{Proj}}_S(\mathcal{B}) \to
\underline{\text{Proj}}_S(\mathcal{A})$  est défini partout et est une immersion
fermée, voir Constructions, Lemme  \ref{constructions-lemma-surjective-graded-rings-map-relative-proj}. Il suffit donc de prouver
 $Z = Z'$  en tant que sous-schémas fermés de  $X$.

\medskip\noindent
Ceci étant dit, la question est locale sur la base et nous pouvons supposer que
 $S = \Spec(R)$  et que  $X = \text{Proj}(A)$  pour une certaine  $R$-algèbre graduée
 $A$. Supposons  $\mathcal{I} = \widetilde{I}$  pour  $I \subset A$  un idéal gradué. D'après la
Constructions, Lemme  \ref{constructions-lemma-proj-quasi-compact} , il existe  $f_0, \ldots, f_n \in A_{+}$  tel que  $A_{+} \subset \sqrt{(f_0, \ldots, f_n)}$ , en
d'autres termes  $X = \bigcup D_{+}(f_i)$. Par conséquent, il suffit de vérifier que
 $Z \cap D_{+}(f_i) = Z' \cap D_{+}(f_i)$  pour chaque  $i$. En renumérotant, nous pouvons supposer
 $i = 0$. Disons que  $Z \cap D_{+}(f_0)$, resp. \ $Z' \cap D_{+}(f_0)$  est défini
par l'idéal  $J$, resp. \ $J'$  de  $A_{(f_0)}$.

\medskip\noindent
L'inclusion  $J' \subset J$. Soit  $d$  le plus petit commun multiple de
 $\deg(f_0), \ldots, \deg(f_n)$. Notez que chacun des torsions  $\mathcal{O}_X(nd)$  est inversible,
trivialise par  $f_i^{nd/\deg(f_i)}$  sur  $D_{+}(f_i)$, et que pour tout module
quasi-cohérent  $\mathcal{F}$  sur  $X$ , les applications de multiplication
 $\mathcal{O}_X(nd) \otimes_{\mathcal{O}_X} \mathcal{F}(m)
\to \mathcal{F}(nd + m)$  sont des isomorphismes, voir Constructions, Lemme  \ref{constructions-lemma-when-invertible}.
Observez que  $J'$  est l'idéal engendré par les éléments  $g/f_0^e$  où
 $g \in I$  est homogène de degré  $e\deg(f_0)$  (voir la preuve de Constructions,
Lemme  \ref{constructions-lemma-surjective-graded-rings-map-proj}). Bien sûr, en remplaçant  $g$  par  $f_0^lg$  pour un
 $l$  convenable, nous pouvons toujours supposer que  $d | e$. Ensuite,
puisque  $g$  s'annule en tant que section de  $\mathcal{O}_X(e\deg(f_0))$  restreinte à
 $Z$ , nous voyons que  $g/f_0^d$  est un élément de  $J$. Ainsi
 $J' \subset J$.

\medskip\noindent
Inversement, supposons que  $g/f_0^e \in J$. Encore une fois, nous pouvons supposer
 $d | e$. Choisissez  $i \in \{1, \ldots, n\}$. Alors  $Z \cap D_{+}(f_i)$  est coupé par un certain
idéal  $J_i \subset A_{(f_i)}$. De plus,
$$
J \cdot A_{(f_0f_i)} = J_i \cdot A_{(f_0f_i)}.
$$
. Le côté droit est la localisation de  $J_i$  par rapport à  $f_0^{\deg(f_i)}/f_i^{\deg(f_0)}$. Il
s'ensuit que
$$
f_0^{e_i}g/f_i^{(e_i + e)\deg(f_0)/\deg(f_i)} \in J_i
$$
pour un certain  $e_i \gg 0$  suffisamment divisible. Cela prouve que  $f_0^{\max(e_i)}g$ 
est un élément de  $I$, car sa restriction à chaque ouvert affine
 $D_{+}(f_i)$  s'annule sur le sous-schéma fermé  $Z \cap D_{+}(f_i)$. Ainsi  $g/f_0^e \in J'$  et
nous concluons  $J \subset J'$  comme souhaité.
\end{proof}

\begin{example}
\label{example-closed-subscheme-of-proj}
Soit  $A$  un anneau gradué. Soient  $X = \text{Proj}(A)$  et  $S = \Spec(A_0)$. Étant
donné un idéal gradué  $I \subset A$ , nous obtenons un sous-schéma fermé
 $V_+(I) = \text{Proj}(A/I) \to X$  par le Lemme de Constructions  \ref{constructions-lemma-surjective-graded-rings-map-proj}. En traduisant le résultat
du Lemme  \ref{lemma-closed-subscheme-proj} , nous voyons que si  $X$  est quasi-compact, alors
tout sous-schéma fermé  $Z$  est de la forme  $V_+(I(Z))$  où l'idéal gradué
 $I(Z) \subset A$  est donné par la règle
$$
I(Z) = \Ker(A \longrightarrow
\bigoplus\nolimits_{n \geq 0} \Gamma(Z, \mathcal{O}_Z(n)))
$$
. Nous pouvons alors poser les deux questions naturelles suivantes:
\begin{enumerate}
\item Quels idéaux  $I$  sont de la forme  $I(Z)$?
\item Pouvons-nous décrire l'opération  $I \mapsto I(V_+(I))$?
\end{enumerate}
Nous y répondrons lorsque  $A$  est noethérien.

\medskip\noindent
Tout d'abord, supposons que  $A$  est généré par  $A_1$  sur
 $A_0$. Dans ce cas, pour tout idéal  $I \subset A$ , le noyau de l'application
 $A/I \to \bigoplus \Gamma(\text{Proj}(A/I), \mathcal{O})$  est l'ensemble des éléments de torsion de  $A/I$, voir
Cohomologie des Schémas, Proposition  \ref{coherent-proposition-coherent-modules-on-proj}. Par conséquent, nous concluons
que
$$
I(V_+(I)) = \{x \in A \mid A_n x \subset I\text{ for some }n \geq 0\}
$$
L'idéal à droite est parfois appelé la saturation de  $I$. Cela répond à
(2) et la réponse à (1) est qu'un idéal est de la forme  $I(Z)$  si et
seulement s'il est saturé, c'est-à-dire égal à sa propre saturation.

\medskip\noindent
Si  $A$  est un anneau gradué noethérien général, alors nous utilisons
Cohomologie des Schémas, Proposition  \ref{coherent-proposition-coherent-modules-on-proj-general}. Ainsi, nous voyons que pour
 $d$  égal au ppcm des degrés des générateurs de  $A$  sur
 $A_0$  nous obtenons
$$
I(V_+(I)) = \{x \in A \mid (Ax)_{nd} \subset I\text{ for all }n \gg 0\}
$$
Cela peut être différent de la saturation de  $I$  si  $d \not = 1$. Par
exemple, supposons que  $A = \mathbf{Q}[x, y]$  avec  $\deg(x) = 2$  et  $\deg(y) = 3$. Alors
 $d = 6$. Soit  $I = (y^2)$. Alors nous voyons  $y \in I(V_+(I))$  parce que pour
tout  $f \in A$  homogène tel que  $6 | \deg(fy)$  nous avons  $y | f$, donc
 $fy \in I$. Il s'ensuit que  $I(V_+(I)) = (y)$  mais  $x^n y \not \in I$  pour tout
 $n$  donc  $I(V_+(I))$  n'est pas égal à la saturation.
\end{example}

\begin{lemma}
\label{lemma-equation-codim-1-in-projective-space}
Soit  $R$  un UFD noethérien. Soit  $Z \subset \mathbf{P}^n_R$  un sous-schéma fermé qui
n'a pas de points incorporés tel que chaque composante irréductible de
 $Z$  ait une codimension  $1$  dans  $\mathbf{P}^n_R$. Alors l'idéal
 $I(Z) \subset R[T_0, \ldots, T_n]$  correspondant à  $Z$  est principal.
\end{lemma}

\begin{proof}
Remarquez que les anneaux locaux de  $X = \mathbf{P}^n_R$  sont des UFD car  $X$  est
recouvert par des ouverts affines isomorphes à  $\mathbf{A}^n_R$  et  $R[x_1, \ldots, x_n]$  est un
UFD (Algèbre, Lemme  \ref{algebra-lemma-polynomial-ring-UFD}). Ainsi  $Z$  est un diviseur de Cartier
effectif d'après le Lemme  \ref{lemma-codimension-1-is-effective-Cartier}. Soit  $\mathcal{I} \subset \mathcal{O}_X$  le faisceau
quasi-cohérent d'idéaux correspondant à  $Z$. Choisissez un isomorphisme
 $\mathcal{O}(m) \to \mathcal{I}$  pour un certain  $m \in \mathbf{Z}$, voir le Lemme  \ref{lemma-Pic-projective-space-UFD}. Alors la
composition
$$
\mathcal{O}_X(m) \to \mathcal{I} \to \mathcal{O}_X
$$
est non nulle. Nous concluons que  $m \leq 0$  et que la section correspondante de
 $\mathcal{O}_X(m)^{\otimes -1} = \mathcal{O}_X(-m)$  est donnée par un certain  $F \in R[T_0, \ldots, T_n]$  de degré  $-m$, voir
Cohomologie des Schémas, Lemme  \ref{coherent-lemma-cohomology-projective-space-over-ring}. Ainsi, sur l'ouvert standard
 $i$ème  $U_i = D_+(T_i)$ , le sous-schéma fermé  $Z \cap U_i$  est découpé par
l'idéal
$$
(F(T_0/T_i, \ldots, T_n/T_i)) \subset R[T_0/T_i, \ldots, T_n/T_i]
$$
. Ainsi, les éléments homogènes de l'idéal gradué  $I(Z) = \Ker(R[T_0, \ldots, T_n] \to \bigoplus \Gamma(\mathcal{O}_Z(m)))$  sont l'ensemble des
polynômes homogènes  $G$  tels que
$$
G(T_0/T_i, \ldots, T_n/T_i) \in (F(T_0/T_i, \ldots, T_n/T_i))
$$
pour  $i = 0, \ldots, n$. En éliminant les dénominateurs, nous voyons qu'il existe
 $e_i \geq 0$  tel que
$$
T_i^{e_i}G \in (F)
$$
pour  $i = 0, \ldots, n$. Comme  $R$  est un UFD, il en est de même pour
 $R[T_0, \ldots, T_n]$. Alors  $F | T_0^{e_0}G$  et  $F | T_1^{e_1}G$  impliquent  $F | G$  car
 $T_0^{e_0}$  et  $T_1^{e_1}$  n'ont pas de facteur commun. Ainsi  $I(Z) = (F)$.
\end{proof}

\noindent
Dans le cas où le sous-schéma fermé est localement défini par un nombre fini
d'équations, nous pouvons le définir par un faisceau d'idéaux de type fini de
 $\mathcal{A}$.

\begin{lemma}
\label{lemma-closed-subscheme-proj-finite}
Soit  $S$  un schéma quasi-compact et quasi-séparé. Soit  $\mathcal{A}$  une
 $\mathcal{O}_S$-algèbre graduée quasi-cohérente. Soit  $p : X = \underline{\text{Proj}}_S(\mathcal{A}) \to S$  le Proj relatif de
 $\mathcal{A}$. Soit  $i : Z \to X$  un sous-schéma fermé. Si  $p$  est
quasi-compact et  $i$  de présentation finie, alors il existe un
 $d > 0$  et un  $\mathcal{O}_S$-submodule quasi-cohérent de type fini  $\mathcal{F} \subset \mathcal{A}_d$ 
tel que  $Z = \underline{\text{Proj}}_S(\mathcal{A}/\mathcal{F}\mathcal{A})$.
\end{lemma}

\begin{proof}
D'après le lemme \ref{lemma-closed-subscheme-proj} nous savons qu'il existe un faisceau de quasi-cohérent
en degrés d'idéaux $\mathcal{I} \subset \mathcal{A}$ tel que $Z = \underline{\text{Proj}}(\mathcal{A}/\mathcal{I})$. Depuis  $S$ est
quasi-compact, nous pouvons choisir un recouvrement ouvert affine
fini $S = U_1 \cup \ldots \cup U_n$. Dis  $U_i = \Spec(R_i)$. Laisser  $\mathcal{A}|_{U_i}$ correspondre au
gradué $R_i$-algèbre $A_i$  et  $\mathcal{I}|_{U_i}$ à l'idéal
gradué $I_i \subset A_i$. Notez que  $p^{-1}(U_i) = \text{Proj}(A_i)$ comme des plans sur $R_i$. Depuis
 $p$ est quasi-compact, nous pouvons choisir un nombre fini d'éléments
homogènes $f_{i, j} \in A_{i, +}$ tel que $p^{-1}(U_i) = \bigcup_j D_{+}(f_{i, j})$. La condition sur  $Z \to X$  signifie
que le faisceau idéal de  $Z$  dans  $\mathcal{O}_X$  est de type fini, voir
Morphismes, Lemme  \ref{morphisms-lemma-closed-immersion-finite-presentation}. Par conséquent, nous pouvons trouver un nombre
fini d'éléments homogènes  $h_{i, j, k} \in I_i \cap A_{i, +}$  tels que l'idéal de  $Z \cap D_{+}(f_{i, j})$  soit
engendré par les éléments  $h_{i, j, k}/f_{i, j}^{e_{i, j, k}}$. Choisissez  $d > 0$  comme un multiple
commun de tous les entiers  $\deg(f_{i, j})$  et  $\deg(h_{i, j, k})$. D'après Propriétés, Lemme
 \ref{properties-lemma-quasi-coherent-colimit-finite-type} , il existe un faisceau quasi-cohérent de type fini  $\mathcal{F} \subset \mathcal{I}_d$  tel
que toutes les sections locales
$$
h_{i, j, k}f_{i, j}^{(d - \deg(h_{i, j, k}))/\deg(f_{i, j})}
$$
soient des sections de  $\mathcal{F}$. Par construction,  $\mathcal{F}$  est une
solution.
\end{proof}

\noindent
La version suivante du Lemme  \ref{lemma-closed-subscheme-proj-finite}  sera utilisée dans la démonstration du
Lemme  \ref{lemma-composition-admissible-blowups}.

\begin{lemma}
\label{lemma-closed-subscheme-proj-finite-type}
Soit  $S$  un schéma quasi-compact et quasi-séparé. Soit  $\mathcal{A}$  une
 $\mathcal{O}_S$-algèbre graduée quasi-cohérente. Soit  $p : X = \underline{\text{Proj}}_S(\mathcal{A}) \to S$  le Proj relatif de
 $\mathcal{A}$. Soit  $i : Z \to X$  un sous-schéma fermé. Soit  $U \subset S$  un ouvert.
Supposons que
\begin{enumerate}
\item $p$  soit quasi-compact,
\item $i$  de présentation finie,
\item $U \cap p(i(Z)) = \emptyset$,
\item $U$  soit quasi-compact,
\item $\mathcal{A}_n$  soit un  $\mathcal{O}_S$-module de type fini pour tout  $n$.
\end{enumerate}
Alors il existe un  $d > 0$  et un  $\mathcal{O}_S$-sous-module quasi-cohérent de
type fini  $\mathcal{F} \subset \mathcal{A}_d$  avec (a)  $Z = \underline{\text{Proj}}_S(\mathcal{A}/\mathcal{F}\mathcal{A})$  et (b) le support de  $\mathcal{A}_d/\mathcal{F}$  est
disjoint de  $U$.
\end{lemma}

\begin{proof}
Soit  $\mathcal{I} \subset \mathcal{A}$  le faisceau d'idéaux gradués quasi-cohérents construit dans le
Lemme  \ref{lemma-closed-subscheme-proj}. Soient  $U_i$,  $R_i$,  $A_i$,  $I_i$,
 $f_{i, j}$,  $h_{i, j, k}$ et  $d$  tels que construits dans la
démonstration du Lemme  \ref{lemma-closed-subscheme-proj-finite}. Puisque  $U \cap p(i(Z)) = \emptyset$ , nous voyons que
 $\mathcal{I}_d|_U = \mathcal{A}_d|_U$  (par notre construction de  $\mathcal{I}$  en tant que noyau). Comme
 $U$  est quasi-compact, nous pouvons choisir un recouvrement ouvert affine
fini  $U = W_1 \cup \ldots \cup W_m$. Puisque  $\mathcal{A}_d$  est de type fini, nous pouvons trouver un
nombre fini de sections  $g_{t, s} \in \mathcal{A}_d(W_t)$  qui engendrent  $\mathcal{A}_d|_{W_t} = \mathcal{I}_d|_{W_t}$  en tant que
 $\mathcal{O}_{W_t}$-module. Pour terminer la preuve, notez que d'après les Propriétés,
Lemme  \ref{properties-lemma-quasi-coherent-colimit-finite-type} , il existe un  $\mathcal{F} \subset \mathcal{I}_d$  de type fini tel que toutes les
sections locales
$$
h_{i, j, k}f_{i, j}^{(d - \deg(h_{i, j, k}))/\deg(f_{i, j})}
\quad\text{and}\quad
g_{t, s}
$$
soient des sections de  $\mathcal{F}$. Par construction,  $\mathcal{F}$  est une
solution.
\end{proof}

\begin{lemma}
\label{lemma-conormal-sheaf-section-projective-bundle}
Soit  $X$  un schéma. Soit  $\mathcal{E}$  un  $\mathcal{O}_X$-module
quasi-cohérent. Il existe une bijection
$$
\left\{
\begin{matrix}
\text{sections }\sigma\text{ of the } \\
\text{morphism } \mathbf{P}(\mathcal{E}) \to X
\end{matrix}
\right\}
\leftrightarrow
\left\{
\begin{matrix}
\text{surjections }\mathcal{E} \to \mathcal{L}\text{ where} \\
\mathcal{L}\text{ is an invertible }\mathcal{O}_X\text{-module}
\end{matrix}
\right\}
$$
. Dans ce cas,  $\sigma$  est une immersion fermée et il existe un isomorphisme
canonique
$$
\Ker(\mathcal{E} \to \mathcal{L})
\otimes_{\mathcal{O}_X} \mathcal{L}^{\otimes -1}
\longrightarrow
\mathcal{C}_{\sigma(X)/\mathbf{P}(\mathcal{E})}
$$
. La bijection et l'isomorphisme sont compatibles avec le changement de base.
\end{lemma}

\begin{proof}
Rappelez-vous que  $\pi : \mathbf{P}(\mathcal{E}) \to X$  est le proj relatif de l'algèbre symétrique sur
 $\mathcal{E}$, voir Constructions, Définition  \ref{constructions-definition-projective-bundle}. Par conséquent, les
descriptions des sections  $\sigma$  suivent immédiatement de la description du
foncteur des points de  $\mathbf{P}(\mathcal{E})$  dans Constructions, Lemme  \ref{constructions-lemma-apply-relative}.
Puisque  $\pi$  est séparé, toute section est une immersion fermée
(Constructions, Lemme  \ref{constructions-lemma-relative-proj-separated}  et Schémas, Lemme  \ref{schemes-lemma-section-immersion}). Soit
 $U \subset X$  un ouvert affine et  $k \in \mathcal{E}(U)$  et  $s \in \mathcal{E}(U)$  des sections locales
telles que  $k$  s'envoie sur zéro dans  $\mathcal{L}$  et  $s$ 
s'envoie sur un générateur  $\overline{s}$  de  $\mathcal{L}$. Alors  $f = k/s$  est
une section de  $\mathcal{O}_{\mathbf{P}(\mathcal{E})}$  définie dans un voisinage ouvert  $D_+(s)$  de
 $s(U)$  dans  $\pi^{-1}(U)$. De plus, puisque  $k$  se mape sur zéro
dans  $\mathcal{L}$ , nous voyons que  $f$  est une section du faisceau
d'idéaux de  $s(U)$  dans  $\pi^{-1}(U)$. Ainsi, nous pouvons prendre l'image
 $\overline{f}$  de  $f$  dans  $\mathcal{C}_{\sigma(X)/\mathbf{P}(\mathcal{E})}(U)$. Nous affirmons (1) que l'image
 $\overline{f}$  dépend seulement des sections  $k$  et  $\overline{s}$  et non du
choix de  $s$  et (2) que nous obtenons un isomorphisme sur  $U$  de
cette manière (voir ci-dessous). Cependant, une fois (1) et (2) établis, nous
voyons que la construction est compatible avec le changement de base par
 $U' \to U$  où  $U'$  est affine, ce qui prouve que ces applications
locales se recollent et sont compatibles avec un changement de base arbitraire.

\medskip\noindent
Pour prouver (1) et (2), nous explicitions ce qui se passe. À savoir, supposons
que  $U = \Spec(A)$  et que  $\mathcal{E} \to \mathcal{L}$  correspondent à l'application des
 $A$-modules  $M \to N$. Alors  $k \in K = \Ker(M \to N)$  et  $s \in M$  se mappent
sur un générateur  $\overline{s}$  de  $N$. D'où  $M = K \oplus A s$. Ainsi
$$
\text{Sym}(M) = \text{Sym}(K)[s]
$$
. Considérons l'identification  $\text{Sym}(K) \to \text{Sym}(M)_{(s)}$  via la règle  $g \mapsto g/s^n$  pour
 $g \in \text{Sym}^n(K)$. Cela donne un isomorphisme  $D_+(s) = \Spec(\text{Sym}(K))$  tel que  $\sigma$ 
correspond à l'application d'anneau  $\text{Sym}(K) \to A$  envoyant  $K$  sur zéro.
Via cet isomorphisme, nous voyons que le module quasi-cohérent correspondant à
 $K$  est identifié avec  $\mathcal{C}_{\sigma(U)/D_+(s)}$ , prouvant (2). Enfin, supposons que
 $s' = k' + s$  pour un certain  $k' \in K$. Ensuite
$$
k/s' = (k/s) (s/s') = (k/s) (s'/s)^{-1} = (k/s) (1 + k'/s)^{-1}
$$
dans un voisinage ouvert de  $\sigma(U)$  dans  $D_+(s)$. Ainsi, nous voyons que
 $s'/s$  se restreint à  $1$  sur  $\sigma(U)$  et nous voyons que
 $k/s'$  se mappe vers le même élément du faisceau conormal que  $k/s$ ,
prouvant ainsi (1).
\end{proof}





\section{Éclatements}
\label{section-blowing-up}

\noindent
L'éclatement est un outil important en géométrie algébrique.

\begin{definition}
\label{definition-blow-up}
Soit  $X$  un schéma. Soit  $\mathcal{I} \subset \mathcal{O}_X$  un faisceau quasi-cohérent
d'idéaux, et soit  $Z \subset X$  le sous-schéma fermé correspondant à  $\mathcal{I}$,
voir Schémas, Définition  \ref{schemes-definition-immersion}. L'éclatement  {\it  de
 $X$  le long de  $Z$}, ou l'éclatement
 {\it  de  $X$  dans le faisceau d'idéaux
 $\mathcal{I}$}  est le morphisme
$$
b :
\underline{\text{Proj}}_X
\left(\bigoplus\nolimits_{n \geq 0} \mathcal{I}^n\right)
\longrightarrow
X
$$
Le  {\it diviseur exceptionnel} de l'éclatement est
l'image inverse $b^{-1}(Z)$. Parfois  $Z$  est appelé le
 {\it centre}  de l'éclatement.
\end{definition}

\noindent
Nous verrons plus tard que le diviseur exceptionnel est un diviseur de Cartier
effectif. De plus, l'éclatement se caractérise comme le schéma « le plus petit »
au-dessus de $X$ tel que l'image réciproque de $Z$ est un diviseur
effectif de Cartier.

\medskip\noindent
Si  $b : X' \to X$  est l’éclatement de  $X$  dans  $Z$, alors nous
notons souvent  $\mathcal{O}_{X'}(n)$  les torsions du faisceau structural. Notez que ce sont
des  $\mathcal{O}_{X'}$-modules inversibles et que  $\mathcal{O}_{X'}(n) = \mathcal{O}_{X'}(1)^{\otimes n}$  car  $X'$  est le
Proj relatif d’une  $\mathcal{O}_X$-algèbre quasi-cohérente graduée qui est engendrée
en degré  $1$, voir Constructions, Lemme  \ref{constructions-lemma-apply-relative}. Notez que
 $\mathcal{O}_{X'}(1)$  est relativement très ample sur  $b$, même si  $b$ 
n’a pas besoin d’être de type fini ou même quasi-compact, car  $X'$  est
muni d’une immersion fermée dans  $\mathbf{P}(\mathcal{I})$, voir Morphismes, Exemple
 \ref{morphisms-example-very-ample}.

\begin{lemma}
\label{lemma-blowing-up-affine}
Soit  $X$  un schéma. Soit  $\mathcal{I} \subset \mathcal{O}_X$  un faisceau quasi-cohérent
d'idéaux. Soit  $U = \Spec(A)$  un sous-schéma ouvert affine de  $X$  et soit
 $I \subset A$  l'idéal correspondant à  $\mathcal{I}|_U$. Si  $b : X' \to X$  est
l'éclatement de  $X$  dans  $\mathcal{I}$, alors il existe un isomorphisme
canonique
$$
b^{-1}(U) = \text{Proj}(\bigoplus\nolimits_{d \geq 0} I^d)
$$
de  $b^{-1}(U)$  avec le spectre homogène de l'algèbre de Rees de  $I$  dans
 $A$. De plus,  $b^{-1}(U)$  possède un recouvrement ouvert affine par des
spectres des algèbres d'éclatement affines  $A[\frac{I}{a}]$.
\end{lemma}

\begin{proof}
La première affirmation découle de la construction du relatif Proj par collage,
voir Constructions, Section  \ref{constructions-section-relative-proj-via-glueing}. Pour  $a \in I$ , notons  $a^{(1)}$ 
l'élément  $a$  considéré comme un élément de degré  $1$  dans
l'algèbre de Rees  $\bigoplus_{n \geq 0} I^n$. Puisque ces éléments génèrent l'algèbre de Rees
sur  $A$ , nous voyons que  $\text{Proj}(\bigoplus_{d \geq 0} I^d)$  est couvert par les ouverts affines
 $D_{+}(a^{(1)})$. Le schéma affine  $D_{+}(a^{(1)})$  est le spectre de l'algèbre d'éclatement
affine  $A' = A[\frac{I}{a}]$, voir Algèbre, Définition  \ref{algebra-definition-blow-up}. Cela termine la
preuve.
\end{proof}

\begin{lemma}
\label{lemma-flat-base-change-blowing-up}
\begin{slogan}
L'éclatement commute avec le changement de base plat.
\end{slogan}
Soit  $X_1 \to X_2$  un morphisme plat de schémas. Soit  $Z_2 \subset X_2$  un sous-schéma
fermé. Soit  $Z_1$  l'image réciproque de  $Z_2$  dans  $X_1$.
Soit  $X'_i$  l'éclatement de  $Z_i$  dans  $X_i$. Alors il existe un
diagramme cartésien
$$
\xymatrix{
X_1' \ar[r] \ar[d] & X_2' \ar[d] \\
X_1 \ar[r] & X_2
}
$$
de schémas.
\end{lemma}

\begin{proof}
Soit  $\mathcal{I}_2$  le faisceau d'idéaux de  $Z_2$  dans  $X_2$. On note
 $g : X_1 \to X_2$  le morphisme donné. Alors le faisceau d'idéaux  $\mathcal{I}_1$  de
 $Z_1$  est l'image de  $g^*\mathcal{I}_2 \to \mathcal{O}_{X_1}$  (par définition de l'image réciproque,
voir Schémas, Définition  \ref{schemes-definition-inverse-image-closed-subscheme}). D'après Constructions, Lemme  \ref{constructions-lemma-relative-proj-base-change} ,
nous voyons que  $X_1 \times_{X_2} X_2'$  est le Proj relatif de  $\bigoplus_{n \geq 0} g^*\mathcal{I}_2^n$. Comme
 $g$  est plat, l'application  $g^*\mathcal{I}_2^n \to \mathcal{O}_{X_1}$  est injective avec image
 $\mathcal{I}_1^n$. Ainsi nous voyons que  $X_1 \times_{X_2} X_2' = X_1'$.
\end{proof}

\begin{lemma}
\label{lemma-blowing-up-gives-effective-Cartier-divisor}
Soit  $X$  un schéma. Soit  $Z \subset X$  un sous-schéma fermé. Le blow-up
 $b : X' \to X$  de  $Z$  dans  $X$  possède les propriétés suivantes:
\begin{enumerate}
\item $b|_{b^{-1}(X \setminus Z)} : b^{-1}(X \setminus Z) \to X \setminus Z$  est un isomorphisme,
\item le diviseur exceptionnel  $E = b^{-1}(Z)$  est un diviseur de Cartier effectif sur
 $X'$,
\item il existe un isomorphisme canonique  $\mathcal{O}_{X'}(-1) = \mathcal{O}_{X'}(E)$
\end{enumerate}
\end{lemma}

\begin{proof}
Comme l'éclatement commute avec les restrictions aux sous-schémas ouverts (Lemme
 \ref{lemma-flat-base-change-blowing-up}), la première assertion signifie simplement que  $X' = X$  si
 $Z = \emptyset$. Dans ce cas, nous éclatons dans le faisceau d'idéaux  $\mathcal{I} = \mathcal{O}_X$  et
le résultat découle de Constructions, Exemple  \ref{constructions-example-trivial-proj}.

\medskip\noindent
La deuxième assertion est locale sur  $X$, donc nous pouvons supposer
 $X$  affine. Disons  $X = \Spec(A)$  et  $Z = \Spec(A/I)$. Par le Lemme
 \ref{lemma-blowing-up-affine} , nous voyons que  $X'$  est couvert par les spectres des
algèbres d'éclatement affines  $A' = A[\frac{I}{a}]$. Alors  $IA' = aA'$  et  $a$  se
projettent sur un élément non-diviseur de zéro dans  $A'$  selon Algèbre,
Lemme  \ref{algebra-lemma-affine-blowup}. Cela prouve le lemme car l'image réciproque de  $Z$ 
dans  $\Spec(A')$  correspond à  $\Spec(A'/IA') \subset \Spec(A')$.

\medskip\noindent
Considérons l'application canonique  $\psi_{univ, 1} : b^*\mathcal{I} \to \mathcal{O}_{X'}(1)$, voir la discussion suivant
Constructions, Définition  \ref{constructions-definition-relative-proj}. Nous affirmons que cela se factorise à
travers un isomorphisme  $\mathcal{I}_E \to \mathcal{O}_{X'}(1)$  (ce qui prouve l'affirmation finale). À
savoir, sur l'ouvert affine correspondant à l'algèbre de l'éclatement  $A' = A[\frac{I}{a}]$ 
mentionnée ci-dessus,  $\psi_{univ, 1}$  correspond à l'application de
 $A'$-module
$$
I \otimes_A A'
\longrightarrow
\left(\Big(\bigoplus\nolimits_{d \geq 0} I^d\Big)_{a^{(1)}}\right)_1
$$
où  $a^{(1)}$  est comme dans Algèbre, Définition  \ref{algebra-definition-blow-up}. Nous omettons la
vérification que ceci est l'application  $I \otimes_A A' \to IA' = aA'$.
\end{proof}

\begin{lemma}[Propriété universelle de l'éclatement]
\label{lemma-universal-property-blowing-up}
Soit  $X$  un schéma. Soit  $Z \subset X$  un sous-schéma fermé. Soit
 $\mathcal{C}$  la sous-catégorie pleine de  $(\Sch/X)$  consistant en  $Y \to X$ 
tels que l'image inverse de  $Z$  est un diviseur de Cartier effectif sur
 $Y$. Alors l'éclatement  $b : X' \to X$  de  $Z$  dans  $X$ 
est un objet final de  $\mathcal{C}$.
\end{lemma}

\begin{proof}
On voit que  $b : X' \to X$  est un objet de  $\mathcal{C}$  selon le Lemme  \ref{lemma-blowing-up-gives-effective-Cartier-divisor}.
Soit  $f : Y \to X$  un objet de  $\mathcal{C}$. Nous devons montrer qu'il existe un
unique morphisme  $Y \to X'$  sur  $X$. Soit  $D = f^{-1}(Z)$. Soit
 $\mathcal{I} \subset \mathcal{O}_X$  le faisceau d'idéaux de  $Z$  et soit  $\mathcal{I}_D$  le faisceau
d'idéaux de  $D$. Alors  $f^*\mathcal{I} \to \mathcal{I}_D$  est une surjection vers un module
 $\mathcal{O}_Y$ inversible. Cela s'étend à une application  $\psi : \bigoplus f^*\mathcal{I}^d \to \bigoplus \mathcal{I}_D^d$  d'algèbres
graduées  $\mathcal{O}_Y$. (On observe que  $\mathcal{I}_D^d = \mathcal{I}_D^{\otimes d}$  comme  $D$  est un
diviseur de Cartier effectif.) Selon le matériel dans Constructions, Section
 \ref{constructions-section-relative-proj} , le triple  $(1, f : Y \to X, \psi)$  définit un morphisme  $Y \to X'$  sur
 $X$. La restriction
$$
Y \setminus D \longrightarrow X' \setminus b^{-1}(Z) = X \setminus Z
$$
est unique. L'ouvert  $Y \setminus D$  est schématiquement dense dans  $Y$ 
selon le Lemme  \ref{lemma-complement-effective-Cartier-divisor}. Ainsi, le morphisme  $Y \to X'$  est unique d'après
Morphismes, Lemme  \ref{morphisms-lemma-equality-of-morphisms}  (de même  $b$  est séparé par Constructions,
Lemme  \ref{constructions-lemma-relative-proj-separated}).
\end{proof}

\begin{lemma}
\label{lemma-characterize-affine-blowup}
Soit  $b : X' \to X$  le blow-up du schéma  $X$  le long d'un sous-schéma fermé
 $Z$. Soit  $U = \Spec(A)$  un ouvert affine de  $X$  et soit
 $I \subset A$  l'idéal correspondant à  $Z \cap U$. Soit  $a \in I$  et soit
 $x' \in X'$  un point se projetant sur un point de  $U$. Alors
 $x'$  est un point de l'ouvert affine  $U' = \Spec(A[\frac{I}{a}])$  si et seulement si
l'image de  $a$  dans  $\mathcal{O}_{X', x'}$  définit le diviseur exceptionnel.
\end{lemma}

\begin{proof}
Puisque le diviseur exceptionnel au-dessus de  $U'$  est coupé par l'image
de  $a$  dans  $A' = A[\frac{I}{a}]$ , une direction est claire. Inversement,
supposons que l'image de  $a$  dans  $\mathcal{O}_{X', x'}$  coupe  $E$.
Puisque chaque élément de  $I$  se mappe à un élément de l'idéal
définissant  $E$  sur  $b^{-1}(U)$ , nous voyons que les éléments de
 $I$  deviennent divisibles par  $a$  dans  $\mathcal{O}_{X', x'}$. Ainsi, pour
 $f \in I^n$ , nous pouvons écrire  $f = \psi(f) a^n$  pour un certain  $\psi(f) \in \mathcal{O}_{X', x'}$.
Remarquez que puisque  $a$  se mappe à un non-zéro diviseur de
 $\mathcal{O}_{X', x'}$ , l'élément  $\psi(f)$  est caractérisé de manière unique par ceci.
Ensuite, nous définissons
$$
A' \longrightarrow \mathcal{O}_{X', x'},\quad
f/a^n \longmapsto \psi(f)
$$
Ici, nous utilisons la description des algèbres de soufflage donnée suivant
Algèbre, Définition  \ref{definition-blow-up}. L'unicité mentionnée ci-dessus montre qu'il
s'agit d'un homomorphisme d'algèbres  $A$. Cela donne un morphisme
 $\Spec(\mathcal{O}_{X', x"}) \to \Spec(A') = U'$. Par la propriété universelle du soufflage (Lemme  \ref{lemma-universal-property-blowing-up}), il
s'agit d'un morphisme au-dessus de  $X'$, ce qui implique bien sûr que
 $x' \in U'$.
\end{proof}

\begin{lemma}
\label{lemma-blow-up-effective-Cartier-divisor}
Soit  $X$  un schéma. Soit  $Z \subset X$  un diviseur de Cartier effectif. Le
soufflage de  $X$  dans  $Z$  est le morphisme identité de
 $X$.
\end{lemma}

\begin{proof}
Immédiat d'après la propriété universelle des soufflages (Lemme  \ref{lemma-universal-property-blowing-up}).
\end{proof}

\begin{lemma}
\label{lemma-blow-up-reduced-scheme}
Soit  $X$  un schéma. Soit  $\mathcal{I} \subset \mathcal{O}_X$  un faisceau quasi-cohérent
d'idéaux. Si  $X$  est réduit, alors le soufflage  $X'$  de
 $X$  dans  $\mathcal{I}$  est réduit.
\end{lemma}

\begin{proof}
Combiner le Lemme  \ref{lemma-blowing-up-affine}  avec Algèbre, Lemme  \ref{algebra-lemma-blowup-reduced}.
\end{proof}

\begin{lemma}
\label{lemma-blow-up-integral-scheme}
Soit  $X$  un schéma. Soit  $\mathcal{I} \subset \mathcal{O}_X$  un faisceau quasi-cohérent d'idéaux
non nul. Si  $X$  est intègre, alors l'éclatement  $X'$  de  $X$ 
en  $\mathcal{I}$  est intègre.
\end{lemma}

\begin{proof}
Combiner le Lemme  \ref{lemma-blowing-up-affine}  avec Algèbre, Lemme  \ref{algebra-lemma-blowup-domain}.
\end{proof}

\begin{lemma}
\label{lemma-blow-up-and-irreducible-components}
Soit  $X$  un schéma. Soit  $Z \subset X$  un sous-schéma fermé. Soit
 $b : X' \to X$  l'éclatement de  $X$  le long de  $Z$. Alors  $b$ 
induit une application bijective de l'ensemble des points génériques des
composantes irréductibles de  $X'$  vers l'ensemble des points génériques
des composantes irréductibles de  $X$  qui ne sont pas dans  $Z$.
\end{lemma}

\begin{proof}
Le diviseur exceptionnel  $E \subset X'$  est un diviseur de Cartier effectif et
 $X' \setminus E \to X \setminus Z$  est un isomorphisme, voir le Lemme  \ref{lemma-blowing-up-gives-effective-Cartier-divisor}. Il suffit donc de
montrer ce qui suit: étant donné un diviseur de Cartier effectif  $D \subset S$  d'un
schéma  $S$ , aucun des points génériques des composantes irréductibles de
 $S$  n'est contenu dans  $D$. Pour voir ceci, nous pouvons
remplacer  $S$  par les membres d'un recouvrement ouvert affine. Ainsi,
d'après le Lemme  \ref{lemma-characterize-effective-Cartier-divisor} , nous pouvons supposer  $S = \Spec(A)$  et  $D = V(f)$ 
où  $f \in A$  est un non-diviseur de zéro. Ensuite, nous devons montrer que
 $f$  n'est contenu dans aucun idéal premier minimal  $\mathfrak p \subset A$. Si
c'était le cas, alors  $f$  se mapperait sur un non-diviseur de zéro
contenu dans l'idéal maximal de  $R_\mathfrak p$ , ce qui est en contradiction avec
Algèbre, Lemme  \ref{algebra-lemma-minimal-prime-reduced-ring}.
\end{proof}

\begin{lemma}
\label{lemma-blow-up-pullback-effective-Cartier}
Soit  $X$  un schéma. Soit  $b : X' \to X$  un éclatement de  $X$  dans
un sous-schéma fermé. Le tiré en arrière  $b^{-1}D$  est défini pour tous les
diviseurs de Cartier effectifs  $D \subset X$  et les tirés en arrière de fonctions
méromorphes sont définis pour  $b$  (Définitions  \ref{definition-pullback-effective-Cartier-divisor}  et
 \ref{definition-pullback-meromorphic-sections}).
\end{lemma}

\begin{proof}
D'après les lemmes  \ref{lemma-blowing-up-affine}  et  \ref{lemma-characterize-effective-Cartier-divisor} , cela se réduit au fait algébrique
suivant: Soit  $A$  un anneau,  $I \subset A$  un idéal,  $a \in I$, et
 $x \in A$  un non-diviseur de zéro. Alors l'image de  $x$  dans
 $A[\frac{I}{a}]$  est un non-diviseur de zéro. À savoir, supposons que  $x (y/a^n) = 0$ 
soit dans  $A[\frac{I}{a}]$. Alors  $a^mxy = 0$  est dans  $A$  pour un certain
 $m$. Par conséquent  $a^my = 0$  car  $x$  est un non-diviseur de
zéro. D'où  $y/a^n$  est nul dans  $A[\frac{I}{a}]$  comme souhaité.
\end{proof}

\begin{lemma}
\label{lemma-blowing-up-two-ideals}
Soit  $X$  un schéma. Soient  $\mathcal{I}, \mathcal{J} \subset \mathcal{O}_X$  des faisceaux quasi-cohérents
d'idéaux. Soit  $b : X' \to X$  l'éclatement de  $X$  dans  $\mathcal{I}$. Soit
 $b' : X'' \to X'$  l'éclatement de  $X'$  dans  $b^{-1}\mathcal{J} \mathcal{O}_{X'}$. Alors  $X'' \to X$ 
est canoniquement isomorphe à l'éclatement de  $X$  dans  $\mathcal{I}\mathcal{J}$.
\end{lemma}

\begin{proof}
Soit  $E \subset X'$  le diviseur exceptionnel de  $b$  qui est un diviseur de
Cartier effectif d'après le Lemme  \ref{lemma-blowing-up-gives-effective-Cartier-divisor}. Alors  $(b')^{-1}E$  est un diviseur
de Cartier effectif sur  $X''$  d'après le Lemme  \ref{lemma-blow-up-pullback-effective-Cartier}. Soit
 $E' \subset X''$  le diviseur exceptionnel de  $b'$  (également un diviseur de
Cartier effectif). Considérons le diviseur de Cartier effectif  $E'' = E' + (b')^{-1}E$. Par
construction, l'idéal de  $E''$  est  $(b \circ b')^{-1}\mathcal{I} (b \circ b')^{-1}\mathcal{J} \mathcal{O}_{X''}$. Par conséquent, selon le
Lemme  \ref{lemma-universal-property-blowing-up} , il existe un morphisme canonique de  $X''$  vers
l'éclatement  $c : Y \to X$  de  $X$  dans  $\mathcal{I}\mathcal{J}$. Inversement, lorsque
 $\mathcal{I}\mathcal{J}$  est ramené à un idéal inversible, nous voyons que  $c^{-1}\mathcal{I}\mathcal{O}_Y$ 
définit un diviseur de Cartier effectif, voir le Lemme  \ref{lemma-sum-closed-subschemes-effective-Cartier}. Ainsi, un
morphisme  $c' : Y \to X'$  sur  $X$  par le Lemme  \ref{lemma-universal-property-blowing-up}. Puis
 $(c')^{-1}b^{-1}\mathcal{J}\mathcal{O}_Y = c^{-1}\mathcal{J}\mathcal{O}_Y$  qui définit également un diviseur de Cartier effectif. Ainsi, un
morphisme  $c'' : Y \to X''$  sur  $X'$. Nous omettons la vérification que ce
morphisme est l'inverse du morphisme  $X'' \to Y$  construit précédemment.
\end{proof}

\begin{lemma}
\label{lemma-blowing-up-projective}
Soit  $X$  un schéma. Soit  $\mathcal{I} \subset \mathcal{O}_X$  un faisceau quasi-cohérent
d'idéaux. Soit  $b : X' \to X$  le blow-up de  $X$  dans le faisceau d'idéaux
 $\mathcal{I}$. Si  $\mathcal{I}$  est de type fini, alors
\begin{enumerate}
\item $b : X' \to X$  est un morphisme projectif, et
\item $\mathcal{O}_{X'}(1)$  est un faisceau inversible  $b$-relativement ample.
\end{enumerate}
\end{lemma}

\begin{proof}
La surjection des  $\mathcal{O}_X$-algèbres graduées
$$
\text{Sym}_{\mathcal{O}_X}^*(\mathcal{I})
\longrightarrow
\bigoplus\nolimits_{d \geq 0} \mathcal{I}^d
$$
définit via Constructions, Lemme  \ref{constructions-lemma-surjective-generated-degree-1-map-relative-proj}  une immersion fermée
$$
X' = \underline{\text{Proj}}_X (\bigoplus\nolimits_{d \geq 0} \mathcal{I}^d)
\longrightarrow
\mathbf{P}(\mathcal{I}).
$$
Par conséquent,  $b$  est projectif, voir Morphismes, Définition
 \ref{morphisms-definition-projective}. La deuxième affirmation découle par exemple de la caractérisation
des faisceaux inversibles relativement amples dans Morphismes, Lemme  \ref{morphisms-lemma-characterize-relatively-ample}.
Certains détails omis.
\end{proof}

\begin{lemma}
\label{lemma-composition-finite-type-blowups}
\begin{slogan}
La composition des éclatements est un éclatement
\end{slogan}
Soit  $X$  un schéma quasi-compact et quasi-séparé. Soit  $Z \subset X$  un
sous-schéma fermé de présentation finie. Soit  $b : X' \to X$  l’éclatement de centre
 $Z$. Soit  $Z' \subset X'$  un sous-schéma fermé de présentation finie. Soit
 $X'' \to X'$  l’éclatement de centre  $Z'$. Il existe un sous-schéma fermé
 $Y \subset X$  de présentation finie, tel que
\begin{enumerate}
\item $Y = Z \cup b(Z')$  théoriquement, et
\item la composition  $X'' \to X$  est isomorphe au soufflage de  $X$  dans
 $Y$.
\end{enumerate}
\end{lemma}

\begin{proof}
La condition que  $Z \to X$  soit de présentation finie signifie que  $Z$ 
est découpé par un faisceau quasi-cohérent d'idéaux de type fini  $\mathcal{I} \subset \mathcal{O}_X$,
voir Morphismes, Lemme  \ref{morphisms-lemma-closed-immersion-finite-presentation}. Écrivons  $\mathcal{A} = \bigoplus_{n \geq 0} \mathcal{I}^n$  de sorte que
 $X' = \underline{\text{Proj}}(\mathcal{A})$. Remarquons que  $X \setminus Z$  est un ouvert quasi-compact de
 $X$  d'après Propriétés, Lemme  \ref{properties-lemma-quasi-coherent-finite-type-ideals}. Puisque  $b^{-1}(X \setminus Z) \to X \setminus Z$  est un
isomorphisme (Lemme  \ref{lemma-blowing-up-gives-effective-Cartier-divisor}), le même résultat montre que  $b^{-1}(X \setminus Z) \setminus Z'$  est un
ouvert quasi-compact dans  $X'$. Ainsi,  $U = X \setminus (Z \cup b(Z'))$  est un ouvert
quasi-compact dans  $X$. D'après le Lemme  \ref{lemma-closed-subscheme-proj-finite-type} , il existe un
 $d > 0$  et un sous-module  $\mathcal{O}_X$ de type fini de  $\mathcal{F} \subset \mathcal{I}^d$  tel que
 $Z' = \underline{\text{Proj}}(\mathcal{A}/\mathcal{F}\mathcal{A})$  et tel que le support de  $\mathcal{I}^d/\mathcal{F}$  soit contenu dans  $X \setminus U$.

\medskip\noindent
Puisque  $\mathcal{F} \subset \mathcal{I}^d$  est un sous-module  $\mathcal{O}_X$, nous pouvons considérer
 $\mathcal{F} \subset \mathcal{I}^d \subset \mathcal{O}_X$  comme un faisceau quasi-cohérent d'idéaux de type fini sur
 $X$. Désignons ceci par  $\mathcal{J} \subset \mathcal{O}_X$  pour éviter toute confusion. Puisque
 $\mathcal{I}^d / \mathcal{J}$  et  $\mathcal{O}/\mathcal{I}^d$  sont supportés sur  $X \setminus U$ , nous voyons que
 $V(\mathcal{J})$  est contenu dans  $X \setminus U$. Inversement, en tant que  $\mathcal{J} \subset \mathcal{I}^d$ ,
nous voyons que  $Z \subset V(\mathcal{J})$. Sur  $X \setminus Z \cong X' \setminus b^{-1}(Z)$ , le faisceau d'idéaux  $\mathcal{J}$ 
définit  $Z'$  (voir la formule affichée ci-dessous). Par conséquent,
 $V(\mathcal{J})$  est égal à  $Z \cup b(Z')$. Il s'ensuit que  $V(\mathcal{I}\mathcal{J}) = Z \cup b(Z')$  également, d'un
point de vue des ensembles. De plus,  $\mathcal{I}\mathcal{J}$  est un idéal de type fini en
tant que produit de deux tels. Nous affirmons que  $X'' \to X$  est isomorphe à
l'éclatement de  $X$  en  $\mathcal{I}\mathcal{J}$ , ce qui termine la démonstration du
lemme en posant  $Y = V(\mathcal{I}\mathcal{J})$.

\medskip\noindent
Tout d'abord, rappelons que l'éclatement de  $X$  dans  $\mathcal{I}\mathcal{J}$  est le
même que l'éclatement de  $X'$  dans  $b^{-1}\mathcal{J} \mathcal{O}_{X'}$, voir le Lemme
 \ref{lemma-blowing-up-two-ideals}. Il suffit donc de montrer que l'éclatement de  $X'$  dans
 $b^{-1}\mathcal{J} \mathcal{O}_{X'}$  coïncide avec l'éclatement de  $X'$  dans  $Z'$. Nous
allons montrer que
$$
b^{-1}\mathcal{J} \mathcal{O}_{X'} = \mathcal{I}_E^d \mathcal{I}_{Z'}
$$
en tant que faisceaux d'idéaux sur  $X''$. Cela prouvera ce que nous voulons
car  $\mathcal{I}_E^d$  découpe le diviseur de Cartier effectif  $dE$  et nous
pouvons utiliser les Lemmata  \ref{lemma-blow-up-effective-Cartier-divisor}  et  \ref{lemma-blowing-up-two-ideals}.

\medskip\noindent
Pour voir l'égalité affichée des idéaux, nous pouvons travailler localement. Avec
la notation  $A$,  $I$,  $a \in I$  comme dans le Lemme
 \ref{lemma-blowing-up-affine} , nous voyons que  $\mathcal{F}$  correspond à un
 $R$-sous-module  $M \subset I^d$  se mappant isomorphiquement à un idéal
 $J \subset R$. La condition  $Z' = \underline{\text{Proj}}(\mathcal{A}/\mathcal{F}\mathcal{A})$  signifie que  $Z' \cap \Spec(A[\frac{I}{a}])$  est découpé par
l'idéal généré par les éléments  $m/a^d$,  $m \in M$. Disons que l'élément
 $m \in M$  correspond à la fonction  $f \in J$. Alors, dans l'algèbre affine
de l'éclatement  $A' = A[\frac{I}{a}]$ , nous voyons que  $f = (a^dm)/a^d = a^d (m/a^d)$. Ainsi, l'égalité
tient.
\end{proof}






\section{Transformée stricte}
\label{section-strict-transform}

\noindent
Dans cette section, nous discutons brièvement de la transformée stricte lors d'un
éclatement. Soit  $S$  un schéma et soit  $Z \subset S$  un sous-schéma fermé.
Soit  $b : S' \to S$  l'éclatement de  $S$  dans  $Z$  et notons
 $E \subset S'$  le diviseur exceptionnel  $E = b^{-1}Z$. Dans ce qui suit, nous
considérerons souvent un schéma  $X$  sur  $S$  et formerons le
diagramme cartésien
$$
\xymatrix{
\text{pr}_{S'}^{-1}E \ar[r] \ar[d] &
X \times_S S' \ar[r]_-{\text{pr}_X} \ar[d]_{\text{pr}_{S'}} &
X \ar[d]^f \\
E \ar[r] & S' \ar[r] & S
}
$$
Puisque  $E$  est un diviseur de Cartier effectif (Lemme  \ref{lemma-blowing-up-gives-effective-Cartier-divisor}),
nous voyons que  $\text{pr}_{S'}^{-1}E \subset X \times_S S'$  est localement principal (Lemme  \ref{lemma-pullback-locally-principal}). Ainsi,
le complément de  $\text{pr}_{S'}^{-1}E$  dans  $X \times_S S'$  est rétrocaptif (Lemme
 \ref{lemma-complement-locally-principal-closed-subscheme}). Par conséquent, pour un  $\mathcal{O}_{X \times_S S'}$-module quasi-cohérent
 $\mathcal{G}$ , le sous-faisceau de sections supportées sur  $\text{pr}_{S'}^{-1}E$  est un
sous-module quasi-cohérent, voir Propriétés, Lemme  \ref{properties-lemma-sections-supported-on-closed-subset}. Si  $\mathcal{G}$ 
est un faisceau quasi-cohérent d'algèbres, par exemple  $\mathcal{G} = \mathcal{O}_{X \times_S S'}$, alors ce
sous-faisceau est un idéal de  $\mathcal{G}$.

\begin{definition}
\label{definition-strict-transform}
Avec  $Z \subset S$  et  $f : X \to S$  comme ci-dessus.
\begin{enumerate}
\item Étant donné un  $\mathcal{O}_X$-module quasi-cohérent  $\mathcal{F}$ , la
 {\it  transformée stricte}  de  $\mathcal{F}$  par
rapport au blow-up de  $S$  dans  $Z$  est le quotient  $\mathcal{F}'$ 
de  $\text{pr}_X^*\mathcal{F}$  par le sous-module de sections supportées sur  $\text{pr}_{S'}^{-1}E$.
\item Le  {\it transformation stricte}  de  $X$ est
le sous-schéma fermé $X' \subset X \times_S S'$  découpé par l'idéal quasi-cohérent de sections de
 $\mathcal{O}_{X \times_S S'}$ pris en charge sur $\text{pr}_{S'}^{-1}E$.
\end{enumerate}
\end{definition}

\noindent
Notez que prendre la transformée stricte le long d'un éclatement dépend du
sous-schéma fermé utilisé pour l'éclatement (et pas seulement du
morphisme $S' \to S$). Cette notion est souvent utilisée pour les sous-schémas
fermés de  $S$. Il s'avère que la transformée stricte de  $X$  est
un agrandissement de  $X$.

\begin{lemma}
\label{lemma-strict-transform}
Dans la situation de la Définition  \ref{definition-strict-transform}.
\begin{enumerate}
\item La transformation stricte $X'$  de  $X$  est l'éclatement de
 $X$ dans le sous-schéma fermé $f^{-1}Z$  de  $X$.
\item Pour un  $\mathcal{O}_X$-module quasi-cohérent  $\mathcal{F}$ , la transformée stricte
 $\mathcal{F}'$  est canoniquement isomorphe à l'image directe le long de  $X' \to X \times_S S'$ 
de la transformée stricte de  $\mathcal{F}$  relative à l'éclatement  $X' \to X$.
\end{enumerate}
\end{lemma}

\begin{proof}
Soit  $X'' \to X$  l'éclatement de  $X$  dans  $f^{-1}Z$. Par la propriété
universelle de l'éclatement (Lemme  \ref{lemma-universal-property-blowing-up}) il existe un diagramme commutatif
$$
\xymatrix{
X'' \ar[r] \ar[d] & X \ar[d] \\
S' \ar[r] & S
}
$$
d'où un morphisme  $X'' \to X \times_S S'$. Ainsi, la première affirmation est que ce
morphisme est une immersion fermée avec image  $X'$. La question est locale
sur  $X$. Ainsi, nous pouvons supposer que  $X$  et  $S$ 
sont affines. Disons que  $S = \Spec(A)$,  $X = \Spec(B)$ et  $Z$  sont définis
par l'idéal  $I \subset A$. Posons  $J = IB$. L'application  $B \otimes_A \bigoplus_{n \geq 0} I^n \to \bigoplus_{n \geq 0} J^n$  définit
une immersion fermée  $X'' \to X \times_S S'$, voir Constructions, lemmes  \ref{constructions-lemma-base-change-map-proj}  et
 \ref{constructions-lemma-surjective-graded-rings-generated-degree-1-map-proj}. Nous omettons la vérification que ce morphisme est le même que celui
construit ci-dessus à partir de la propriété universelle. Choisissons  $a \in I$ 
correspondant à l'ouvert affine  $\Spec(A[\frac{I}{a}]) \subset S'$, voir le Lemme  \ref{lemma-blowing-up-affine}. L'image
réciproque de  $\Spec(A[\frac{I}{a}])$  dans la transformée stricte  $X'$  de
 $X$  est le spectre de
$$
B' = (B \otimes_A A[\textstyle{\frac{I}{a}}])/a\text{-power-torsion}
$$
voir Propriétés, Lemme  \ref{properties-lemma-sections-supported-on-closed-subset}. D'autre part, en laissant  $b \in J$  être
l'image de  $a$ , nous voyons que  $\Spec(B[\frac{J}{b}])$  est l'image réciproque de
 $\Spec(A[\frac{I}{a}])$  dans  $X''$. Par Algèbre, Lemme  \ref{algebra-lemma-blowup-base-change} , l'ouvert
 $\Spec(B[\frac{J}{b}])$  se mappe isomorphiquement sur le sous-schéma ouvert  $\text{pr}_{S'}^{-1}(\Spec(A[\frac{I}{a}]))$  de
 $X'$. Ainsi  $X'' \to X'$  est un isomorphisme.

\medskip\noindent
Dans la notation ci-dessus, que  $\mathcal{F}$  corresponde au  $B$-module
 $N$. La transformée stricte de  $\mathcal{F}$  correspond au
 $B \otimes_A A[\frac{I}{a}]$-module
$$
N' = (N \otimes_A A[\textstyle{\frac{I}{a}}])/a\text{-power-torsion}
$$
voir Propriétés, Lemme  \ref{properties-lemma-sections-supported-on-closed-subset}. La transformée stricte de  $\mathcal{F}$ 
relativement à l'éclatement de  $X$  dans  $f^{-1}Z$  correspond au
 $B[\frac{J}{b}]$-module  $N \otimes_B B[\frac{J}{b}]/b\text{-power-torsion}$. De manière exactement identique à ci-dessus, on
démontre que ces deux modules sont isomorphes. Détails omis.
\end{proof}

\begin{lemma}
\label{lemma-strict-transform-flat}
Dans la situation de la Définition  \ref{definition-strict-transform}.
\begin{enumerate}
\item Si  $X$  est plat sur  $S$  en tous les points au-dessus de
 $Z$, alors la transformée stricte de  $X$  est égale au changement
de base  $X \times_S S'$.
\item Soit  $\mathcal{F}$  un  $\mathcal{O}_X$-module quasi-cohérent. Si  $\mathcal{F}$  est plat
sur  $S$  en tous les points au-dessus de  $Z$, alors la
transformée stricte  $\mathcal{F}'$  de  $\mathcal{F}$  est égale au tiré en arrière
 $\text{pr}_X^*\mathcal{F}$.
\end{enumerate}
\end{lemma}

\begin{proof}
Nous prouverons la partie (2) car elle implique la partie (1) par la définition de
la transformée stricte d'un schéma sur  $S$. La question est locale sur
 $X$. Ainsi, nous pouvons supposer que  $S = \Spec(A)$,  $X = \Spec(B)$, et que
 $\mathcal{F}$  correspond au  $B$-module  $N$. Ensuite,  $\mathcal{F}'$ 
sur l'ouvert  $\Spec(B \otimes_A A[\frac{I}{a}])$  de  $X \times_S S'$  correspond au module
$$
N' = (N \otimes_A A[\textstyle{\frac{I}{a}}])/a\text{-power-torsion}
$$
voir Propriétés, Lemme  \ref{properties-lemma-sections-supported-on-closed-subset}. Ainsi, nous devons montrer que la
 $a$-puissance-torsion de  $N \otimes_A A[\frac{I}{a}]$  est nulle. Soit  $y \in N \otimes_A A[\frac{I}{a}]$  avec
 $a^n y = 0$. Si  $\mathfrak q \subset B$  est un nombre premier et  $a \not \in \mathfrak q$, alors
 $y$  se mappe à zéro dans  $(N \otimes_A A[\frac{I}{a}])_\mathfrak q$. d'autre part, si  $a \in \mathfrak q$,
alors  $N_\mathfrak q$  est un  $A$-module plat et nous voyons que  $N_\mathfrak q \otimes_A A[\frac{I}{a}]
=(N \otimes_A A[\frac{I}{a}])_\mathfrak q$ 
n'a pas de torsion de puissance  $a$ (comme  $A[\frac{I}{a}]$  n'en a pas). Par
conséquent,  $y$  se mappe également à zéro dans cette localisation. Nous
concluons que  $y$  est nul d'après Algèbre, Lemme  \ref{algebra-lemma-characterize-zero-local}.
\end{proof}

\begin{lemma}
\label{lemma-strict-transform-affine}
Soit  $S$  un schéma. Soit  $Z \subset S$  un sous-schéma fermé. Soit
 $b : S' \to S$  le blow-up de  $Z$  dans  $S$. Soit  $g : X \to Y$  un
morphisme affine de schémas sur  $S$. Soit  $\mathcal{F}$  un faisceau
quasi-cohérent sur  $X$. Soit  $g' : X \times_S S' \to Y \times_S S'$  le changement de base de
 $g$. Soit  $\mathcal{F}'$  la transformée stricte de  $\mathcal{F}$  par rapport
à  $b$. Alors  $g'_*\mathcal{F}'$  est la transformée stricte de  $g_*\mathcal{F}$.
\end{lemma}

\begin{proof}
Remarquez que  $g'_*\text{pr}_X^*\mathcal{F} = \text{pr}_Y^*g_*\mathcal{F}$  par Cohomologie des Schémas, Lemme  \ref{coherent-lemma-affine-base-change}. Soit
 $\mathcal{K} \subset \text{pr}_X^*\mathcal{F}$  le sous-faisceau des sections supportées dans l'image inverse de
 $Z$  dans  $X \times_S S'$. D'après Propriétés, Lemme  \ref{properties-lemma-push-sections-supported-on-closed-subset} , le
faisceau direct  $g'_*\mathcal{K}$  est le sous-faisceau des sections de  $\text{pr}_Y^*g_*\mathcal{F}$ 
supportées dans l'image inverse de  $Z$  dans  $Y \times_S S'$. Comme
 $g'$  est affine (Morphismes, Lemme  \ref{morphisms-lemma-base-change-affine}), nous voyons que
 $g'_*$  est exact, donc nous concluons.
\end{proof}

\begin{lemma}
\label{lemma-strict-transform-different-centers}
Soit  $S$  un schéma. Soit  $Z \subset S$  un sous-schéma fermé. Soit
 $D \subset S$  un diviseur de Cartier effectif. Soit  $Z' \subset S$  le sous-schéma
fermé défini par le produit des faisceaux d'idéaux de  $Z$  et de
 $D$. Soit  $S' \to S$  l'éclatement de  $S$  dans  $Z$.
\begin{enumerate}
\item L'éclatement de  $S$  dans  $Z'$  est isomorphe à  $S' \to S$.
\item Soit  $f : X \to S$  un morphisme de schémas et soit  $\mathcal{F}$  un
 $\mathcal{O}_X$-module quasi-cohérent. Si  $\mathcal{F}$  n'a pas de sections locales
non nulles supportées dans  $f^{-1}D$, alors la transformée stricte de
 $\mathcal{F}$  par rapport au soufflage dans  $Z$  coïncide avec la
transformée stricte de  $\mathcal{F}$  par rapport au soufflage de  $S$  dans
 $Z'$.
\end{enumerate}
\end{lemma}

\begin{proof}
La première affirmation découle de la combinaison des lemmes  \ref{lemma-blowing-up-two-ideals}  et
 \ref{lemma-blow-up-effective-Cartier-divisor}. En utilisant le lemme  \ref{lemma-blowing-up-affine} , la deuxième affirmation se
traduit en le problème algébrique suivant. Soit  $A$  un anneau,
 $I \subset A$  un idéal,  $x \in A$  un élément non diviseur de zéro, et
 $a \in I$. Soit  $M$  un  $A$-module dont la torsion
 $x$ est nulle. À montrer: la torsion de puissance  $a$ dans
 $M \otimes_A A[\frac{I}{a}]$  est égale à la torsion de puissance  $xa$. La raison en est
que le noyau et le conoyau de l'application  $A \to A[\frac{I}{a}]$  sont des torsions de
puissance  $a$, donc cette application devient un isomorphisme après
inversion de  $a$. Ainsi, le noyau et le conoyau de  $M \to M \otimes_A A[\frac{I}{a}]$  sont
également des torsions de puissance  $a$. Cela implique le résultat.
\end{proof}

\begin{lemma}
\label{lemma-strict-transform-composition-blowups}
Soit  $S$  un schéma. Soit  $Z \subset S$  un sous-schéma fermé. Soit
 $b : S' \to S$  le plongement éclaté de centre  $Z$. Soit  $Z' \subset S'$  un
sous-schéma fermé. Soit  $S'' \to S'$  le plongement éclaté de centre  $Z'$.
Soit  $Y \subset S$  un sous-schéma fermé tel que  $Y = Z \cup b(Z')$  soit théoriquement le
même ensemble et que la composition  $S'' \to S$  soit isomorphe au plongement
éclaté de  $S$  en  $Y$. Dans cette situation, pour tout schéma
 $X$  sur  $S$  et  $\mathcal{F} \in \QCoh(\mathcal{O}_X)$ , nous avons
\begin{enumerate}
\item le transformé strict de  $\mathcal{F}$  par rapport au plongement éclaté de
 $S$  en  $Y$  est égal au transformé strict par rapport au
plongement éclaté  $S'' \to S'$  dans  $Z'$  du transformé strict de
 $\mathcal{F}$  par rapport au plongement éclaté  $S' \to S$  de  $S$  en
 $Z$, et
\item La transformée stricte de  $X$  par rapport au blow-up de  $S$  dans
 $Y$  est égale à la transformée stricte par rapport au blow-up
 $S'' \to S'$  dans  $Z'$  de la transformée stricte de  $X$  par
rapport au blow-up  $S' \to S$  de  $S$  dans  $Z$.
\end{enumerate}
\end{lemma}

\begin{proof}
Soit  $\mathcal{F}'$  la transformée stricte de  $\mathcal{F}$  par rapport au blow-up
 $S' \to S$  de  $S$  dans  $Z$. Soit  $\mathcal{F}''$  la transformée
stricte de  $\mathcal{F}'$  par rapport au blow-up  $S'' \to S'$  de  $S'$  dans
 $Z'$. Soit  $\mathcal{G}$  la transformée stricte de  $\mathcal{F}$  par rapport
au blow-up  $S'' \to S$  de  $S$  dans  $Y$. Nous étiquetons
également les morphismes
$$
\xymatrix{
X \times_S S'' \ar[r]_q \ar[d]^{f''} &
X \times_S S' \ar[r]_p \ar[d]^{f'} &
X \ar[d]^f \\
S'' \ar[r] & S' \ar[r] & S
}
$$
Par définition, il existe une surjection  $p^*\mathcal{F} \to \mathcal{F}'$  et une surjection
 $q^*\mathcal{F}' \to \mathcal{F}''$  qui se combinent par exactitude à droite de  $q^*$  en une
surjection  $(p \circ q)^*\mathcal{F} \to \mathcal{F}''$. De plus, nous avons la surjection  $(p \circ q)^*\mathcal{F} \to \mathcal{G}$. Il suffit
donc de prouver que ces deux surjections ont le même noyau.

\medskip\noindent
Le noyau de la surjection  $p^*\mathcal{F} \to \mathcal{F}'$  est supporté sur  $(f \circ p)^{-1}Z$, donc cette
application est un isomorphisme aux points dans le complément. Par conséquent, le
noyau de  $q^*p^*\mathcal{F} \to q^*\mathcal{F}'$  est supporté sur  $(f \circ p \circ q)^{-1}Z$. Le noyau de  $q^*\mathcal{F}' \to \mathcal{F}''$  est
supporté sur  $(f' \circ q)^{-1}Z'$. Combiné, nous voyons que le noyau de  $(p \circ q)^*\mathcal{F} \to \mathcal{F}''$  est
supporté sur  $(f \circ p \circ q)^{-1}Z \cup (f' \circ q)^{-1}Z' =
(f \circ p \circ q)^{-1}Y$. Par construction de  $\mathcal{G}$ , nous voyons que nous
obtenons une factorisation  $(p \circ q)^*\mathcal{F} \to \mathcal{F}'' \to \mathcal{G}$. Pour terminer la preuve, il suffit de
montrer que  $\mathcal{F}''$  n’a pas de sections (locales) non nulles supportées sur
 $(f \circ p \circ q)^{-1}(Y) =
(f \circ p \circ q)^{-1}Z \cup (f' \circ q)^{-1}Z'$. Cela découle du Lemme  \ref{lemma-strict-transform-different-centers}  appliqué à  $\mathcal{F}'$  sur
 $X \times_S S'$  au-dessus de  $S'$, le sous-schéma fermé  $Z'$  et le
diviseur de Cartier effectif  $b^{-1}Z$.
\end{proof}

\begin{lemma}
\label{lemma-strict-transform-universally-injective}
Dans la situation de la Définition  \ref{definition-strict-transform}. Supposons que
$$
0 \to \mathcal{F}_1 \to \mathcal{F}_2 \to \mathcal{F}_3 \to 0
$$
soit une suite exacte de faisceaux quasi-cohérents sur  $X$  qui reste
exacte après tout changement de base  $T \to S$. Alors les transformés stricts
de  $\mathcal{F}_i'$  relativement à tout éclatement  $S' \to S$  forment également une
suite exacte courte  $0 \to \mathcal{F}'_1 \to \mathcal{F}'_2 \to \mathcal{F}'_3 \to 0$ .
\end{lemma}

\begin{proof}
Nous pouvons nous localiser sur  $S$  et  $X$  et supposer que les
deux sont affines. Ensuite, nous pouvons pousser  $\mathcal{F}_i$  vers  $S$,
voir le Lemme  \ref{lemma-strict-transform-affine}. Nous pouvons supposer que notre éclatement est le
morphisme  $1 : S \to S$  associé à un diviseur de Cartier effectif  $D \subset S$.
Alors la traduction en algèbre est la suivante: Supposons que  $A$  soit un
anneau et  $0 \to M_1 \to M_2 \to M_3 \to 0$  soit une suite universellement exacte de
 $A$-modules. Soit  $a\in A$. Alors la suite
$$
0 \to
M_1/a\text{-power torsion} \to
M_2/a\text{-power torsion} \to
M_3/a\text{-power torsion} \to 0
$$
est également exact. À savoir, la surjectivité de la dernière application et
l'injectivité de la première application sont immédiates. Le problème est
l'exactitude au milieu. Supposons que  $x \in M_2$  se mappe à zéro dans
 $M_3/a\text{-power torsion}$. Alors  $y = a^n x \in M_1$  pour certains  $n$. Ensuite,  $y$ 
se mappe à zéro dans  $M_2/a^nM_2$. Comme  $M_1 \to M_2$  est universellement
injectif, nous voyons que  $y$  se mappe à zéro dans  $M_1/a^nM_1$. Ainsi,
 $y = a^n z$  pour certains  $z \in M_1$. Ainsi,  $a^n(x - y) = 0$. Par conséquent,
 $y$  se mappe à la classe de  $x$  dans  $M_2/a\text{-power torsion}$  comme
souhaité.
\end{proof}







\section{Éclatements admissibles}
\label{section-admissible-blowups}

\noindent
Pour avoir un peu plus de contrôle sur nos éclatements, nous introduisons la
terminologie standard suivante.

\begin{definition}
\label{definition-admissible-blowup}
Soit  $X$  un schéma. Soit  $U \subset X$  un sous-schéma ouvert. Un morphisme
 $X' \to X$  est appelé un  {\it   $U$-éclatement
admissible}  s'il existe une immersion fermée  $Z \to X$  de
présentation finie avec  $Z$  disjointe de  $U$  telle que
 $X'$  est isomorphe à l'éclatement de  $X$  dans  $Z$.
\end{definition}

\noindent
Nous rappelons que  $Z \to X$  est de présentation finie si et seulement si le
faisceau d'idéaux  $\mathcal{I}_Z \subset \mathcal{O}_X$  est de type fini, voir Morphismes, Lemme
 \ref{morphisms-lemma-closed-immersion-finite-presentation}. En particulier, un éclatement  $U$-admissible est un
morphisme projectif, voir Lemme  \ref{lemma-blowing-up-projective}. Notez qu'il peut y avoir plusieurs
centres donnant lieu au même morphisme. Par conséquent, l'exigence est simplement
l'existence d'un certain centre disjoint de  $U$  qui produit  $X'$.
Enfin, comme le morphisme  $b : X' \to X$  est un isomorphisme au-dessus de
 $U$  (voir Lemme  \ref{lemma-blowing-up-gives-effective-Cartier-divisor}), nous abuserons souvent de la notation et
considérerons  $U$  comme un sous-schéma ouvert de  $X'$  également.

\begin{lemma}
\label{lemma-composition-admissible-blowups}
\begin{slogan}
Les éclatements admissibles sont stables par composition.
\end{slogan}
Soit  $X$  un schéma quasi-compact et quasi-séparé. Soit  $U \subset X$  un
sous-schéma ouvert quasi-compact. Soit  $b : X' \to X$  un éclatement
 $U$-admissible. Soit  $X'' \to X'$  un éclatement  $U$-admissible.
Alors la composition  $X'' \to X$  est un éclatement  $U$-admissible.
\end{lemma}

\begin{proof}
Immédiat à partir du lemme plus précis  \ref{lemma-composition-finite-type-blowups}.
\end{proof}

\begin{lemma}
\label{lemma-extend-admissible-blowups}
Soit  $X$  un schéma quasi-compact et quasi-séparé. Soit  $U, V \subset X$  des
sous-schémas ouverts quasi-compacts. Soit  $b : V' \to V$  un éclatement
 $U \cap V$-admissible. Alors il existe un éclatement  $U$-admissible
 $X' \to X$  dont la restriction à  $V$  est  $V'$.
\end{lemma}

\begin{proof}
Soit  $\mathcal{I} \subset \mathcal{O}_V$  le faisceau quasi-cohérent d'idéaux de type fini tel que
 $V(\mathcal{I})$  soit disjoint de  $U \cap V$  et tel que  $V'$  soit isomorphe
à l'éclatement de  $V$  dans  $\mathcal{I}$. Soit  $\mathcal{I}' \subset \mathcal{O}_{U \cup V}$  le faisceau
quasi-cohérent d'idéaux dont la restriction à  $U$  est  $\mathcal{O}_U$  et
dont la restriction à  $V$  est  $\mathcal{I}$  (voir Faisceaux, Section
 \ref{sheaves-section-glueing-sheaves}). D'après Propriétés, Lemme  \ref{properties-lemma-extend} , il existe un faisceau
quasi-cohérent d'idéaux de type fini  $\mathcal{J} \subset \mathcal{O}_X$  dont la restriction à
 $U \cup V$  est  $\mathcal{I}'$. Le lemme s'ensuit.
\end{proof}

\begin{lemma}
\label{lemma-dominate-admissible-blowups}
Soit  $X$  un schéma quasi-compact et quasi-séparé. Soit  $U \subset X$  un
sous-schéma ouvert quasi-compact. Soient  $b_i : X_i \to X$,  $i = 1, \ldots, n$  des
éclatements admissibles de  $U$. Il existe un éclatement admissible
 $U$,  $b : X' \to X$  tel que (a)  $b$  se factorise comme
 $X' \to X_i \to X$  pour  $i = 1, \ldots, n$  et (b) chacun des morphismes  $X' \to X_i$  est un
éclatement admissible  $U$.
\end{lemma}

\begin{proof}
Soit  $\mathcal{I}_i \subset \mathcal{O}_X$  le faisceau cohérent quasi-type fini d'idéaux tel que
 $V(\mathcal{I}_i)$  soit disjoint de  $U$  et tel que  $X_i$  soit isomorphe
à l'éclatement de  $X$  dans  $\mathcal{I}_i$. Posons  $\mathcal{I} = \mathcal{I}_1 \cdot \ldots \cdot \mathcal{I}_n$  et soit
 $X'$  l'éclatement de  $X$  dans  $\mathcal{I}$. Alors  $X' \to X$  se
factorise à travers  $b_i$  par le Lemme  \ref{lemma-blowing-up-two-ideals}.
\end{proof}

\begin{lemma}
\label{lemma-separate-disjoint-opens-by-blowing-up}
\begin{slogan}
Séparer les composantes irréductibles en effectuant un éclatement.
\end{slogan}
Soit  $X$  un schéma quasi-compact et quasi-séparé. Soient  $U, V$  des
sous-schémas ouverts disjoints quasi-compacts de  $X$. Alors il existe un
éclatement  $U \cup V$-admissible  $b : X' \to X$  tel que  $X'$  soit une
union disjointe de sous-schémas ouverts  $X' = X'_1 \amalg X'_2$  avec  $b^{-1}(U) \subset X'_1$  et
 $b^{-1}(V) \subset X'_2$.
\end{lemma}

\begin{proof}
Choisissez un faisceau d'idéaux quasi-cohérent de type fini  $\mathcal{I}$,
respectivement \ $\mathcal{J}$  tel que  $X \setminus U = V(\mathcal{I})$, respectivement
\ $X \setminus V = V(\mathcal{J})$, voir Propriétés, Lemme  \ref{properties-lemma-quasi-coherent-finite-type-ideals}. Alors  $V(\mathcal{I}\mathcal{J}) = X$ 
au sens des ensembles, donc  $\mathcal{I}\mathcal{J}$  est un faisceau d'idéaux localement
nilpotent. Puisque  $\mathcal{I}$  et  $\mathcal{J}$  sont de type fini et que
 $X$  est quasi-compact, il existe un  $n > 0$  tel que  $\mathcal{I}^n \mathcal{J}^n = 0$.
Nous pouvons et nous remplaçons  $\mathcal{I}$  par  $\mathcal{I}^n$  et  $\mathcal{J}$  par
 $\mathcal{J}^n$. D'où  $\mathcal{I} \mathcal{J} = 0$. Soit  $b : X' \to X$  l'éclatement dans  $\mathcal{I} + \mathcal{J}$.
Ceci est  $U \cup V$-admissible car  $V(\mathcal{I} + \mathcal{J}) = X \setminus U \cup V$. Nous montrerons que
 $X'$  est une réunion disjointe de sous-schémas ouverts  $X' = X'_1 \amalg X'_2$  tels
que  $b^{-1}\mathcal{I}|_{X'_2} = 0$  et  $b^{-1}\mathcal{J}|_{X'_1} = 0$ , ce qui prouvera le lemme.

\medskip\noindent
Nous utiliserons la description de l'éclatement dans le Lemme  \ref{lemma-blowing-up-affine}.
Supposons que  $U = \Spec(A) \subset X$  soit un ouvert affine tel que  $\mathcal{I}|_U$, resp.
\ $\mathcal{J}|_U$  correspond à l'idéal de type fini  $I \subset A$, resp.
\ $J \subset A$. Alors
$$
b^{-1}(U) = \text{Proj}(A \oplus (I + J) \oplus (I + J)^2 \oplus \ldots)
$$
Ceci est recouvert par les sous-ensembles ouverts affines  $A[\frac{I + J}{x}]$  et
 $A[\frac{I + J}{y}]$  avec  $x \in I$  et  $y \in J$. Puisque  $x \in I$  est un
non-diviseur de zéro dans  $A[\frac{I + J}{x}]$  et  $IJ = 0$ , nous voyons que
 $J A[\frac{I + J}{x}] = 0$. Puisque  $y \in J$  est un non-diviseur de zéro dans  $A[\frac{I + J}{y}]$ 
et  $IJ = 0$ , nous voyons que  $I A[\frac{I + J}{y}] = 0$. De plus,
$$
\Spec(A[\textstyle{\frac{I + J}{x}}]) \cap
\Spec(A[\textstyle{\frac{I + J}{y}}]) =
\Spec(A[\textstyle{\frac{I + J}{xy}}]) = \emptyset
$$
parce que  $xy$  est à la fois un non-diviseur du zéro et zéro. Ainsi,
 $b^{-1}(U)$  est l'union disjointe du sous-schéma ouvert  $U_1$  défini comme
l'union des ouverts standards  $\Spec(A[\frac{I + J}{x}])$  pour  $x \in I$  et du sous-schéma
ouvert  $U_2$  qui est l'union des ouverts affines  $\Spec(A[\frac{I + J}{y}])$  pour
 $y \in J$. Nous avons vu que  $b^{-1}\mathcal{I}\mathcal{O}_{X'}$  se restreint à zéro sur  $U_2$ 
et  $b^{-1}\mathcal{I}\mathcal{O}_{X'}$  se restreint à zéro sur  $U_1$. Nous omettons la
vérification que ces sous-schémas ouverts se recollent en sous-schémas ouverts
globaux  $X'_1$  et  $X'_2$.
\end{proof}

\begin{lemma}
\label{lemma-blowing-up-denominators}
Soit  $X$  un schéma localement noethérien. Soit  $\mathcal{L}$  un
 $\mathcal{O}_X$-module inversible. Soit  $s$  une section méromorphe régulière
de  $\mathcal{L}$. Soit  $U \subset X$  le plus grand sous-schéma ouvert tel que
 $s$  corresponde à une section de  $\mathcal{L}$  sur  $U$.
L'éclatement  $b : X' \to X$  dans l'idéal des dénominateurs de  $s$  est
 $U$-admissible. Il existe un diviseur de Cartier effectif  $D \subset X'$  et
un isomorphisme
$$
b^*\mathcal{L} = \mathcal{O}_{X'}(D - E),
$$
où  $E \subset X'$  est le diviseur exceptionnel de sorte que la section méromorphe
 $b^*s$  corresponde, via l'isomorphisme, à la section méromorphe
 $1_D \otimes (1_E)^{-1}$.
\end{lemma}

\begin{proof}
D'après la définition de l'idéal des dénominateurs dans la Définition
 \ref{definition-regular-meromorphic-ideal-denominators} , nous voyons immédiatement que  $b$  est un éclatement
admissible  $U$. Pour la notation  $1_D$,  $1_E$ et
 $\mathcal{O}_{X'}(D - E)$ , veuillez vous référer à la Définition  \ref{definition-invertible-sheaf-effective-Cartier-divisor}. Le tiré en
arrière  $b^*s$  est défini par les lemmes  \ref{lemma-blow-up-pullback-effective-Cartier}  et  \ref{lemma-meromorphic-sections-pullback}.
Ainsi, l'énoncé du lemme a un sens. Nous pouvons réinterpréter la dernière
assertion en disant que  $b^*s$  est une section régulière globale de
 $b^*\mathcal{L}(E)$  dont le schéma des zéros est  $D$. Cela définit de manière
unique  $D$ , donc pour prouver le lemme, nous pouvons travailler
localement de manière affine sur  $X$  et  $X'$. Supposons que
 $X = \Spec(A)$  soit affine et  $\mathcal{L} = \mathcal{O}_X$. Alors  $s$  est une fonction
méromorphe régulière et en rétrécissant davantage nous pouvons supposer
 $s = a'/a$  avec  $a', a \in A$  des non-diviseurs de zéro. Alors l'idéal des
dénominateurs de  $s$  correspond à l'idéal  $I = \{x \in A \mid xa' \in aA\}$. Rappelons que
 $X'$  est recouvert par les spectres des algèbres d'éclatement affines
 $A' = A[\frac{I}{x}]$  avec  $x \in I$  (Lemme  \ref{lemma-blowing-up-affine}). Fixons  $x \in I$  et
écrivons  $xa' = a a''$  pour un certain  $a'' \in A$. Le diviseur  $E \subset X'$  est
défini par  $x \in A'$  sur le spectre de  $A'$  et donc  $1/x$  est un
générateur de  $\mathcal{O}_{X'}(E)$  sur  $\Spec(A')$. Enfin, dans l'anneau total de
fractions de  $A'$  nous avons  $a'/a = a''/x$. Par conséquent,  $b^*s = a'/a$  se
restreint à une section régulière de  $\mathcal{O}_{X'}(E)$  qui est au-dessus de
 $\Spec(A')$  donnée par  $a''/x$. Cela termine la preuve. (Le diviseur
 $D \cap \Spec(A')$  est découpé par l'image de  $a''$  dans  $A'$.)
\end{proof}





\section{Éclatements et platitude}
\label{section-blowup-flat}

\noindent
Nous continuons la discussion commencée dans Plus sur l'algèbre, Section
 \ref{more-algebra-section-blowup-flat}. Nous prouverons des résultats supplémentaires dans Plus sur la
platitude, Section  \ref{flat-section-blowup-flat}.


\begin{lemma}
\label{lemma-strict-transform-blowup-fitting-ideal}
Soit  $S$  un schéma. Soit  $\mathcal{F}$  un  $\mathcal{O}_S$-module
quasi-cohérent de type fini. Soit  $Z_k \subset S$  le sous-schéma fermé découpé par
 $\text{Fit}_k(\mathcal{F})$, voir Section  \ref{section-fitting-ideals}. Soit  $S' \to S$  le blow-up de
 $S$  dans  $Z_k$  et soit  $\mathcal{F}'$  la transformée stricte de
 $\mathcal{F}$. Alors  $\mathcal{F}'$  peut localement être généré par  $\leq k$ 
sections.
\end{lemma}

\begin{proof}
Rappelons que  $\mathcal{F}'$  peut être localement engendré par des sections
 $\leq k$  si et seulement si  $\text{Fit}_k(\mathcal{F}') = \mathcal{O}_{S'}$, voir le Lemme  \ref{lemma-fitting-ideal-generate-locally}. Ainsi ce
lemme est une traduction de Plus sur l’Algèbre, Lemme  \ref{more-algebra-lemma-blowup-fitting-ideal}.
\end{proof}

\begin{lemma}
\label{lemma-strict-transform-blowup-fitting-ideal-locally-free}
Soit  $S$  un schéma. Soit  $\mathcal{F}$  un  $\mathcal{O}_S$-module
quasi-cohérent de type fini. Soit  $Z_k \subset S$  le sous-schéma fermé découpé par
 $\text{Fit}_k(\mathcal{F})$, voir la Section  \ref{section-fitting-ideals}. Supposons que  $\mathcal{F}$  soit
localement libre de rang  $k$  sur  $S \setminus Z_k$. Soit  $S' \to S$ 
l’éclatement de  $S$  dans  $Z_k$  et soit  $\mathcal{F}'$  la transformée
stricte de  $\mathcal{F}$. Alors  $\mathcal{F}'$  est localement libre de rang
 $k$.
\end{lemma}

\begin{proof}
Traduction de Plus sur l’Algèbre, Lemme  \ref{more-algebra-lemma-blowup-fitting-ideal-locally-free}.
\end{proof}

\begin{lemma}
\label{lemma-blowup-fitting-ideal}
Soit  $X$  un schéma. Soit  $\mathcal{F}$  un  $\mathcal{O}_X$-module de
présentation finie. Soit  $U \subset X$  un ouvert schématiquement dense tel que
 $\mathcal{F}|_U$  soit localement libre de rang constant  $r$. Alors
\begin{enumerate}
\item l'éclatement  $b : X' \to X$  de  $X$  dans le  $r$-ième idéal de
Fitting de  $\mathcal{F}$  est  $U$-admissible,
\item la transformée stricte  $\mathcal{F}'$  de  $\mathcal{F}$  par rapport à  $b$  est
localement libre de rang  $r$,
\item le noyau  $\mathcal{K}$  de la surjection  $b^*\mathcal{F} \to \mathcal{F}'$  est de présentation finie et
 $\mathcal{K}|_U = 0$,
\item $b^*\mathcal{F}$  et  $\mathcal{K}$  sont des  $\mathcal{O}_{X'}$-modules parfaits de dimension
tor  $\leq 1$.
\end{enumerate}
\end{lemma}

\begin{proof}
L'idéal  $\text{Fit}_r(\mathcal{F})$  est de type fini d'après le Lemme  \ref{lemma-fitting-ideal-of-finitely-presented}  et sa
restriction à  $U$  est égale à  $\mathcal{O}_U$  d'après le Lemme  \ref{lemma-fitting-ideal-finite-locally-free}.
Par conséquent,  $b : X' \to X$  est  $U$-admissible, voir la Définition
 \ref{definition-admissible-blowup}.

\medskip\noindent
D'après le Lemme  \ref{lemma-fitting-ideal-finite-locally-free} , la restriction de  $\text{Fit}_{r - 1}(\mathcal{F})$  à  $U$  est
nulle, et puisque  $U$  est schématiquement dense, nous concluons que
 $\text{Fit}_{r - 1}(\mathcal{F}) = 0$  sur tout  $X$. Ainsi il découle du Lemme  \ref{lemma-fitting-ideal-finite-locally-free}  que
 $\mathcal{F}$  est localement libre de rang  $r$  sur le complément du
sous-schéma défini par le  $r$-ième idéal de Fitting de  $\mathcal{F}$  (ce
complément peut être plus grand que  $U$ , ce qui explique pourquoi nous
avons dû effectuer cette étape dans l'argument). Par conséquent, d'après le Lemme
 \ref{lemma-strict-transform-blowup-fitting-ideal-locally-free} , la transformée stricte
$$
b^*\mathcal{F} \longrightarrow \mathcal{F}'
$$
est localement libre de rang  $r$. Le noyau  $\mathcal{K}$  de cette
application est supporté sur le diviseur exceptionnel de l'éclatement  $b$ 
et donc  $\mathcal{K}|_U = 0$. Enfin, puisque  $\mathcal{F}'$  est localement libre de type
fini et que la flèche affichée est surjective, nous pouvons localement sur
 $X'$  écrire  $b^*\mathcal{F}$  comme la somme directe de  $\mathcal{K}$  et
 $\mathcal{F}'$. Puisque  $b^*\mathcal{F}'$  est de présentation finie (Modules, Lemme
 \ref{modules-lemma-pullback-finite-presentation}), il en est de même pour  $\mathcal{K}$.

\medskip\noindent
L'énoncé sur la dimension tor découle de More on Algebra, Lemme  \ref{more-algebra-lemma-fitting-ideals-and-pd1}.
\end{proof}





\section{Modifications}
\label{section-modifications}

\noindent
Dans cette section, nous rassemblerons des résultats du type: après une
modification, telle et telle propriété est vraie. Nous verrons plus tard qu'une
modification peut être dominée par un éclatement (More on Flatness, Lemme
 \ref{flat-lemma-dominate-modification-by-blowup}).

\begin{lemma}
\label{lemma-filter-after-modification}
Soit  $X$  un schéma intègre. Soit  $\mathcal{E}$  un  $\mathcal{O}_X$-module
localement libre de type fini. Il existe une modification  $f : X' \to X$  telle que
 $f^*\mathcal{E}$  possède une filtration dont les quotients successifs sont des
 $\mathcal{O}_{X'}$-modules inversibles.
\end{lemma}

\begin{proof}
Nous prouvons cela par induction sur le rang  $r$  de  $\mathcal{E}$. Si
 $r = 1$  ou  $r = 0$ , le lemme est évident. Supposons  $r > 1$. Soit
 $P = \mathbf{P}(\mathcal{E})$  avec un morphisme de structure  $\pi : P \to X$, voir Constructions,
Section  \ref{constructions-section-projective-bundle}. Alors  $\pi$  est propre (Lemme  \ref{lemma-relative-proj-proper}). Il
existe une surjection canonique
$$
\pi^*\mathcal{E} \to \mathcal{O}_P(1)
$$
dont le noyau est localement libre de rang  $r - 1$. Choisissez un sous-schéma
ouvert non vide  $U \subset X$  tel que  $\mathcal{E}|_U \cong \mathcal{O}_U^{\oplus r}$. Alors  $P_U = \pi^{-1}(U)$  est
isomorphe à  $\mathbf{P}^{r - 1}_U$. En particulier, il existe une section  $s : U \to P_U$  de
 $\pi$. Soit  $X' \subset P$  l'image schématique du morphisme  $U \to P_U \to P$.
Alors  $X'$  est intègre (Morphismes, Lemme  \ref{morphisms-lemma-scheme-theoretic-image-reduced}), le morphisme
 $f = \pi|_{X'} : X' \to X$  est propre (Morphismes, Lemme  \ref{morphisms-lemma-closed-immersion-proper}  et  \ref{morphisms-lemma-composition-proper}), et
 $f^{-1}(U) \to U$  est un isomorphisme. Par conséquent,  $f$  est une
modification (Morphismes, Définition  \ref{morphisms-definition-modification}). Par construction, le produit
fibré  $f^*\mathcal{E}$  possède une filtration en deux étapes dont le quotient est
inversible car il est égal à  $\mathcal{O}_P(1)|_{X'}$  et dont le sous  $\mathcal{E}'$  est
localement libre de rang  $r - 1$. Par récurrence, nous pouvons trouver une
modification  $g : X'' \to X'$  telle que  $g^*\mathcal{E}'$  possède une filtration comme dans
l'énoncé du lemme. Ainsi,  $f \circ g : X'' \to X$  est la modification requise.
\end{proof}

\begin{lemma}
\label{lemma-extend-rational-map-after-modification}
Soit  $S$  un schéma. Soient  $X$,  $Y$  des schémas sur
 $S$. Supposons que  $X$  soit noethérien et que  $Y$  soit
propre sur  $S$. Étant donné une application rationnelle  $S$
 $f : U \to Y$  de  $X$  vers  $Y$ , il existe un morphisme
 $p : X' \to X$  et un  $S$-morphisme  $f' : X' \to Y$  tels que
\begin{enumerate}
\item $p$  soit propre et  $p^{-1}(U) \to U$  soit un isomorphisme,
\item $f'|_{p^{-1}(U)}$  soit égal à  $f \circ p|_{p^{-1}(U)}$.
\end{enumerate}
\end{lemma}

\begin{proof}
On note  $j : U \to X$  le morphisme d'inclusion. Soit  $X' \subset Y \times_S X$  l'image
schématique de  $(f, j) : U \to Y \times_S X$  (Morphismes, Définition  \ref{morphisms-definition-scheme-theoretic-image}). La projection
 $g : Y \times_S X \to X$  est propre (Morphismes, Lemme  \ref{morphisms-lemma-base-change-proper}). La composition
 $p : X' \to X$  de  $X' \to Y \times_S X$  et  $g$  est propre (Morphismes, Lemmata
 \ref{morphisms-lemma-closed-immersion-proper}  et  \ref{morphisms-lemma-composition-proper}). Comme  $g$  est séparé et  $U \subset X$  est
rétrocompact (puisque  $X$  est noethérien), nous concluons que
 $p^{-1}(U) \to U$  est un isomorphisme d'après Morphismes, Lemme  \ref{morphisms-lemma-scheme-theoretic-image-of-partial-section}. D'autre
part, la composition  $f' : X' \to Y$  de  $X' \to Y \times_S X$  et de la projection  $Y \times_S X \to Y$ 
coïncide avec  $f$  sur  $p^{-1}(U)$.
\end{proof}







\begin{multicols}{2}[\section{Autres chapitres}]
\noindent
Préliminaires
\begin{enumerate}
\item \hyperref[introduction-section-phantom]{Introduction}
\item \hyperref[conventions-section-phantom]{Conventions}
\item \hyperref[sets-section-phantom]{Théorie des ensembles}
\item \hyperref[categories-section-phantom]{Catégories}
\item \hyperref[topology-section-phantom]{Topologie}
\item \hyperref[sheaves-section-phantom]{Faisceaux sur les espaces}
\item \hyperref[sites-section-phantom]{Sites et faisceaux}
\item \hyperref[stacks-section-phantom]{Champs}
\item \hyperref[fields-section-phantom]{Corps}
\item \hyperref[algebra-section-phantom]{Algèbre commutative}
\item \hyperref[brauer-section-phantom]{Groupes de Brauer}
\item \hyperref[homology-section-phantom]{Algèbre homologique}
\item \hyperref[derived-section-phantom]{Catégories dérivées}
\item \hyperref[simplicial-section-phantom]{Méthodes simpliciales}
\item \hyperref[more-algebra-section-phantom]{Compléments d'algèbre}
\item \hyperref[smoothing-section-phantom]{Lissage des morphismes d'anneaux}
\item \hyperref[modules-section-phantom]{Faisceaux de modules}
\item \hyperref[sites-modules-section-phantom]{Modules sur les sites}
\item \hyperref[injectives-section-phantom]{Objets injectifs}
\item \hyperref[cohomology-section-phantom]{Cohomologie des faisceaux}
\item \hyperref[sites-cohomology-section-phantom]{Cohomologie sur les sites}
\item \hyperref[dga-section-phantom]{Algèbre différentielle graduée}
\item \hyperref[dpa-section-phantom]{Algèbre à puissances divisées}
\item \hyperref[sdga-section-phantom]{Faisceaux différentiels gradués}
\item \hyperref[hypercovering-section-phantom]{Hyperrecouvrements}
\end{enumerate}
Schémas
\begin{enumerate}
\setcounter{enumi}{25}
\item \hyperref[schemes-section-phantom]{Schémas}
\item \hyperref[constructions-section-phantom]{Constructions de schémas}
\item \hyperref[properties-section-phantom]{Propriétés des schémas}
\item \hyperref[morphisms-section-phantom]{Morphismes de schémas}
\item \hyperref[coherent-section-phantom]{Cohomologie des schémas}
\item \hyperref[divisors-section-phantom]{Diviseurs}
\item \hyperref[limits-section-phantom]{Limites de schémas}
\item \hyperref[varieties-section-phantom]{Variétés}
\item \hyperref[topologies-section-phantom]{Topologies sur les schémas}
\item \hyperref[descent-section-phantom]{Descente}
\item \hyperref[perfect-section-phantom]{Catégories dérivées de schémas}
\item \hyperref[more-morphisms-section-phantom]{Compléments sur les morphismes}
\item \hyperref[flat-section-phantom]{Compléments sur la platitude}
\item \hyperref[groupoids-section-phantom]{Schémas en groupoïdes}
\item \hyperref[more-groupoids-section-phantom]{Compléments sur les schémas en groupoïdes}
\item \hyperref[etale-section-phantom]{Morphismes étales de schémas}
\end{enumerate}
Thèmes de théorie des schémas
\begin{enumerate}
\setcounter{enumi}{41}
\item \hyperref[chow-section-phantom]{Homologie de Chow}
\item \hyperref[intersection-section-phantom]{Théorie de l'intersection}
\item \hyperref[pic-section-phantom]{Schémas de Picard des courbes}
\item \hyperref[weil-section-phantom]{Théories cohomologiques de Weil}
\item \hyperref[adequate-section-phantom]{Modules adéquats}
\item \hyperref[dualizing-section-phantom]{Complexes dualisants}
\item \hyperref[duality-section-phantom]{Dualité pour les schémas}
\item \hyperref[discriminant-section-phantom]{Discriminants et différentes}
\item \hyperref[derham-section-phantom]{Cohomologie de de Rham}
\item \hyperref[local-cohomology-section-phantom]{Cohomologie locale}
\item \hyperref[algebraization-section-phantom]{Géométrie algébrique et formelle}
\item \hyperref[curves-section-phantom]{Courbes algébriques}
\item \hyperref[resolve-section-phantom]{Résolution des surfaces}
\item \hyperref[models-section-phantom]{Réduction semi-stable}
\item \hyperref[functors-section-phantom]{Foncteurs et morphismes}
\item \hyperref[equiv-section-phantom]{Catégories dérivées de variétés}
\item \hyperref[pione-section-phantom]{Groupes fondamentaux de schémas}
\item \hyperref[etale-cohomology-section-phantom]{Cohomologie étale}
\item \hyperref[crystalline-section-phantom]{Cohomologie cristalline}
\item \hyperref[proetale-section-phantom]{Cohomologie pro-étale}
\item \hyperref[relative-cycles-section-phantom]{Cycles relatifs}
\item \hyperref[more-etale-section-phantom]{Compléments de cohomologie étale}
\item \hyperref[trace-section-phantom]{La formule des traces}
\end{enumerate}
Espaces algébriques
\begin{enumerate}
\setcounter{enumi}{64}
\item \hyperref[spaces-section-phantom]{Espaces algébriques}
\item \hyperref[spaces-properties-section-phantom]{Propriétés des espaces algébriques}
\item \hyperref[spaces-morphisms-section-phantom]{Morphismes d'espaces algébriques}
\item \hyperref[decent-spaces-section-phantom]{Espaces algébriques décents}
\item \hyperref[spaces-cohomology-section-phantom]{Cohomologie des espaces algébriques}
\item \hyperref[spaces-limits-section-phantom]{Limites d'espaces algébriques}
\item \hyperref[spaces-divisors-section-phantom]{Diviseurs sur les espaces algébriques}
\item \hyperref[spaces-over-fields-section-phantom]{Espaces algébriques sur des corps}
\item \hyperref[spaces-topologies-section-phantom]{Topologies sur les espaces algébriques}
\item \hyperref[spaces-descent-section-phantom]{Descente et espaces algébriques}
\item \hyperref[spaces-perfect-section-phantom]{Catégories dérivées d'espaces}
\item \hyperref[spaces-more-morphisms-section-phantom]{Compléments sur les morphismes d'espaces}
\item \hyperref[spaces-flat-section-phantom]{Platitude sur les espaces algébriques}
\item \hyperref[spaces-groupoids-section-phantom]{Groupoïdes dans les espaces algébriques}
\item \hyperref[spaces-more-groupoids-section-phantom]{Compléments sur les groupoïdes dans les espaces}
\item \hyperref[bootstrap-section-phantom]{Amorçage}
\item \hyperref[spaces-pushouts-section-phantom]{Sommes amalgamées d'espaces algébriques}
\end{enumerate}
Thèmes de géométrie
\begin{enumerate}
\setcounter{enumi}{81}
\item \hyperref[spaces-chow-section-phantom]{Groupes de Chow des espaces}
\item \hyperref[groupoids-quotients-section-phantom]{Quotients de groupoïdes}
\item \hyperref[spaces-more-cohomology-section-phantom]{Compléments sur la cohomologie des espaces}
\item \hyperref[spaces-simplicial-section-phantom]{Espaces simpliciaux}
\item \hyperref[spaces-duality-section-phantom]{Dualité pour les espaces}
\item \hyperref[formal-spaces-section-phantom]{Espaces algébriques formels}
\item \hyperref[restricted-section-phantom]{Algébrisation des espaces formels}
\item \hyperref[spaces-resolve-section-phantom]{Résolution des surfaces revisitée}
\end{enumerate}
Théorie des déformations
\begin{enumerate}
\setcounter{enumi}{89}
\item \hyperref[formal-defos-section-phantom]{Théorie formelle des déformations}
\item \hyperref[defos-section-phantom]{Théorie des déformations}
\item \hyperref[cotangent-section-phantom]{Le complexe cotangent}
\item \hyperref[examples-defos-section-phantom]{Problèmes de déformation}
\end{enumerate}
Champs algébriques
\begin{enumerate}
\setcounter{enumi}{93}
\item \hyperref[algebraic-section-phantom]{Champs algébriques}
\item \hyperref[examples-stacks-section-phantom]{Exemples de champs}
\item \hyperref[stacks-sheaves-section-phantom]{Faisceaux sur les champs algébriques}
\item \hyperref[criteria-section-phantom]{Critères de représentabilité}
\item \hyperref[artin-section-phantom]{Axiomes d'Artin}
\item \hyperref[quot-section-phantom]{Espaces de Quot et de Hilbert}
\item \hyperref[stacks-properties-section-phantom]{Propriétés des champs algébriques}
\item \hyperref[stacks-morphisms-section-phantom]{Morphismes de champs algébriques}
\item \hyperref[stacks-limits-section-phantom]{Limites de champs algébriques}
\item \hyperref[stacks-cohomology-section-phantom]{Cohomologie des champs algébriques}
\item \hyperref[stacks-perfect-section-phantom]{Catégories dérivées de champs}
\item \hyperref[stacks-introduction-section-phantom]{Introduction aux champs algébriques}
\item \hyperref[stacks-more-morphisms-section-phantom]{Compléments sur les morphismes de champs}
\item \hyperref[stacks-geometry-section-phantom]{La géométrie des champs}
\end{enumerate}
Thèmes de théorie des espaces de modules
\begin{enumerate}
\setcounter{enumi}{107}
\item \hyperref[moduli-section-phantom]{Champs de modules}
\item \hyperref[moduli-curves-section-phantom]{Espaces de modules de courbes}
\end{enumerate}
Divers
\begin{enumerate}
\setcounter{enumi}{109}
\item \hyperref[examples-section-phantom]{Exemples}
\item \hyperref[exercises-section-phantom]{Exercices}
\item \hyperref[guide-section-phantom]{Guide bibliographique}
\item \hyperref[desirables-section-phantom]{Desiderata}
\item \hyperref[coding-section-phantom]{Style de programmation}
\item \hyperref[obsolete-section-phantom]{Obsolète}
\item \hyperref[fdl-section-phantom]{Licence de documentation libre GNU}
\item \hyperref[index-section-phantom]{Index généré automatiquement}
\end{enumerate}
\end{multicols}


\bibliography{my}
\bibliographystyle{amsalpha}

\end{document}
